% =============================================================
% Bölüm 6 - İşlemci Tasarımı
% Bilgisayar Mimarisi - Oğuz Ergin
% =============================================================

\chapter{İşlemci Tasarımı}
\label{chap:islemci-tasarimi}
\bolumfoto[Mimar Sinan, Selimiye'yi ``ustalık eserim'' diye tanımlar. İşlemci tasarımı da bir mimarın ustalık sınavıdır. Yazmaçlar, AMB, denetim birimi ve veri yolları bir araya getirilerek çalışan bir bütün oluşturulur.]{gorseller/bolum06-selimiye.jpg}{Selimiye Camisi -- Edirne}

\begin{tcolorbox}[
    colframe=sinyalyesil!60!black,
    title={\textbf{Öğrenme Hedefleri}},
    fonttitle=\bfseries,
    rounded corners,
    boxrule=0.8pt
]
Bu bölümü tamamladığınızda aşağıdakileri yapabileceksiniz:
\begin{itemize}[nosep, leftmargin=*]
    \item Bir buyruğun yürütümünü oluşturan beş temel aşamayı (getir, çöz, yürüt, bellek, sonucu yaz) ve her aşamada veri yolunda neler olduğunu açıklamak.
    \item RISC-V buyruklarını yürütebilen tek vuruşluk bir veri yolu tasarlamak. Ele alınan buyruk türleri (R, I, S, B, J) için etkin yolu çıkarmak.
    \item Tek vuruşluk denetim birimini doğruluk tablosu ve Karno haritasıyla gerçekleştirmek.
    \item Tek vuruşluk işlemcide kritik yolu hesaplamak ve saat çevrim zamanını belirlemek.
    \item Çok vuruşluk işlemcinin nasıl çalıştığını anlamak. Her buyruğun yalnızca ihtiyacı kadar çevrim kullanarak donanım tekrar kullanımının nasıl sağlandığını göstermek.
    \item Çok vuruşluk denetim birimini sonlu durum makinesi (SDM) ile tasarlamak.
    \item Mikroprogramlamanın temel ilkesini kavramak. Yatay ile dikey mikrokod arasındaki farkları değerlendirmek.
    \item Tek vuruşluk, çok vuruşluk ve mikroprogramlamalı yaklaşımları başarım, alan ve esneklik açısından karşılaştırmak.
\end{itemize}
\end{tcolorbox}

\section*{Giriş}

Buraya kadar bilgisayarın temel yapı taşlarını ayrı ayrı inceledik. Bölüm~\ref{chap:aritmetik}'te toplama, çarpma ve kayan nokta aritmetiği gerçekleştiren birimleri, özellikle de aritmetik mantık birimini (AMB) gördük. Ek~\ref{ch:sayisal-tasarim}'de ise mantık kapılarını, flip-flopları, yazmaçları ve durum makinelerini öğrendik. Bu parçalar tek başlarına işlemci değildir. Bir araya getirilip eşgüdüm içinde çalıştırıldıklarında \emph{tam} bir işlemci ortaya çıkar.

Bir işlemcinin temel görevi, bellekte tutulan makine kodu buyruklarını sırayla yürütmektir. Her buyruk için işlemci beş temel aşamadan geçer: program sayacının gösterdiği adresten buyruğu \emph{getirir}\index{getirme}\index{fetch|see{getirme}}, getirilen buyruğun türünü ve işlenenlerini \emph{çözer}\index{buyruk çözme}\index{instruction decode|see{buyruk çözme}}, AMB üzerinde gerekli işlemi \emph{yürütür}\index{yürütme}, gerekiyorsa belleğe \emph{erişir} ve sonucu yazmaca \emph{yazar}. Bu adımları gerçekleştirebilmek için işlemcinin iki temel birime ihtiyacı vardır:

\begin{itemize}
\item \textbf{Veri Yolu}\index{veri yolu}\index{datapath|see{veri yolu}}: Verilerin bellekten ya da yazmaç öbeğinden alınıp AMB üzerinde işlenmesinden sonucun yazılmasına kadar süren fiziksel yoldur. Yazmaç öbeği\index{yazmaç öbeği}\index{register file|see{yazmaç öbeği}}, AMB\index{AMB}\index{ALU|see{AMB}}, bellek arabirimleri ve bunları birbirine bağlayan yollar veri yolunu oluşturur.
\item \textbf{Denetim birimi}\index{denetim birimi}\index{control unit|see{denetim birimi}}: Veri yolundaki çoklayıcıları, yazmaçları ve AMB'yi yönlendiren sinyalleri üreten beyindir. Yürütülen buyruğa göre doğru denetim sinyallerini doğru anda üretir. Bu sinyaller veri yoluna ne yapacağını söyler.
\end{itemize}

İşlemci tasarımının iki ana ekseni vardır. İlki \emph{veri yolunun yapısı}dır. \textbf{Tek vuruşluk} tasarımda her buyruk tek bir saat çevriminde başlar ve biter. Bu en yalın yaklaşımdır ama saat çevrim zamanını \emph{en yavaş} buyruğa göre belirlemek gerekir, dolayısıyla hızlı buyruklar bile uzun çevrim zamanına mahkum olur. \textbf{Çok vuruşluk} tasarımın amacı tek vuruşluğun bu darboğazını aşarak başarımı artırmaktır. Her buyruk birkaç kısa saat çevrimine bölünür ve her buyruk türü yalnızca ihtiyacı kadar çevrim kullanır. Böylece hızlı buyruklar artık yavaş buyrukları beklemek zorunda kalmaz. AMB ve bellek gibi alanca büyük birimler farklı çevrimlerde yeniden kullanılır, böylece hem yongada yer kazanılır hem de saat sıklığı yükseltilebilir.

İkinci eksen \emph{denetim biriminin gerçekleştirilmesidir}: \textbf{donanım tabanlı denetim}, sinyalleri doğrudan mantık kapılarıyla ürettiği için ek bir bellek erişimi gerektirmez ve hızlıdır; ancak buyruk kümesi büyüdükçe kapılarla elle sadeleştirilen devre karmaşıklaşır ve buyruk kümesi her değiştiğinde yeniden tasarlanması gerekir. \textbf{Mikroprogramlama}\index{mikroprogramlama} ise her buyruğu, mikrokod belleğinde tutulan daha küçük adımlar (mikrobuyruklar) dizisine indirger. Denetim sinyalleri bu mikrokoddan üretilir. Mikroprogramlama, denetim mantığını yazılım gibi güncellenebilir kılarak karmaşık buyrukların eklenmesini ve hatadan dönülmesini kolaylaştırır. CISC mimarileri büyük ölçüde bu sayede gerçeklenebilir hale gelmiştir.

Bu bölümde önce tek vuruşluk bir RISC-V işlemcisini sıfırdan tasarlayacak, ardından çok vuruşluk sürümüne geçerek AMB ve bellek gibi alanca büyük birimleri yeniden kullanarak aynı buyruk kümesini nasıl yürüteceğimizi göreceğiz. Çok vuruşluk işlemcinin denetim birimini önce sonlu durum makinesiyle (donanım tabanlı) kuracak, sonra aynı işlevi mikroprogramlama yöntemiyle gerçekleştireceğiz. Bölümün sonunda bu tasarım kararlarını başarım ve alan ölçütleriyle karşılaştıracak, gerçek dünya işlemcilerinde hangi tercihlerin neden yapıldığını tartışacağız.

Bu bölümde öğreneceklerimiz, Bölüm~\ref{chap:boruhatti}'de göreceğimiz \emph{boru hattı} tasarımının ön koşuludur. Günümüzün masaüstü ve sunucu işlemcileri (Intel Core, AMD Ryzen, Apple M serisi, ARM Cortex, RISC-V tabanlı yongalar) çok derin boru hatlarına ve sırasız yürütüme sahip olsa da, hepsi özünde burada öğreneceğimiz veri yolu ve denetim birimi ilkeleri üzerine kuruludur.

\section{Veri Yolu}
\label{sec:veriyolu}

Veri yolu, işlemcinin asıl iş gören yarısıdır: buyruğun bellekten getirilmesi, işlenecek verilerin yazmaç öbeğinden okunması, AMB üzerinde aritmetik ya da mantık işleminin gerçekleştirilmesi ve sonucun doğru yere yazılması. Bütün bu fiziksel adımlar veri yolu üzerinde olur. Bir RISC-V veri yolu beş temel birimden kurulur: \textbf{program sayacı} (PS), \textbf{buyruk belleği}, \textbf{yazmaç öbeği}, \textbf{aritmetik mantık birimi} (AMB) ve \textbf{veri belleği}. Bu birimler çoklayıcılar, anlık değer üreteçleri ve denetim sinyalleri aracılığıyla birbirine bağlanarak buyruk yürütümünü tek bir koordineli zincire dönüştürür.

Bir buyruğun yürütümünü bu kitap boyunca klasik \emph{beş aşamalı} modelle inceleyeceğiz. Bu model RISC-V ve MIPS gibi yalın buyruk kümeleri için doğal bir bölünmedir; ancak çağdaş işlemciler buyruk yürütümünü çok daha fazla aşamaya bölebilir. Sırasız yürütüm yapan tasarımlarda adlandırma (\emph{rename}), planlama (\emph{schedule}) ve tamamlama (\emph{retire}) gibi ek aşamalar bulunur. Intel Core ailesinin boru hattı 14--19 aşamalıdır. Bu bölümde tasarlayacağımız tek vuruşluk ve çok vuruşluk işlemciler bu klasik beş aşamayı temel alır:
\begin{itemize}
    \item \textbf{Getirme}\index{getirme}\index{fetch|see{getirme}} (\emph{fetch}): Program sayacının gösterdiği adresteki buyruk, buyruk belleğinden okunur.
    \item \textbf{Çözme}\index{buyruk çözme}\index{instruction decode|see{buyruk çözme}} (\emph{decode}): Buyruk parçalarına ayrılır. İşlem kodu, kaynak ve hedef yazmaç numaraları ile varsa anlık değer çıkarılır. Yazmaç öbeğinden kaynak yazmaçlar okunur.
    \item \textbf{Yürütme}\index{yürütme} (\emph{execute}): Okunan değerler AMB'ye verilir ve buyruğun gerektirdiği işlem gerçekleştirilir. Dallanma buyruklarında dallanma koşulu da bu aşamada değerlendirilir.
    \item \textbf{Bellek erişimi}\index{bellek erişimi} (\emph{memory access}): Yükle ve sakla buyruklarında veri belleğine erişilir. Yükle buyruğunda bellekten değer okunur, sakla buyruğunda yazmaçtan gelen değer belleğe yazılır. Diğer buyruklarda bu aşamada bir iş yapılmaz.
    \item \textbf{Geri yazma}\index{geri yazma}\index{write-back|see{geri yazma}} (\emph{write-back}): Buyruğun ürettiği sonuç yazmaç öbeğine yazılır. Yükle buyruğunda bellekten okunan değer, R/I/U/J türü buyruklarda AMB sonucu, atlama buyruklarında ise dönüş adresi (PS+4) yazılır. S ve B türü buyruklarda bu aşama atlanır.
\end{itemize}

Bu zincirin merkezinde \textbf{program sayacı} bulunur. PS\index{program sayacı}\index{program counter|see{program sayacı}}, Bölüm~\ref{chap:buyruk-kumesi}'te tanıtıldığı üzere bir sonraki yürütülecek buyruğun bellek adresini tutar ve sabit 4 baytlık RISC-V buyrukları nedeniyle her yürütümden sonra dört artırılır. Dallanma ya da atlama buyruklarında ise hedef adresle güncellenir. Veri yolu açısından PS, getirme aşamasında buyruk belleğine adres olarak gönderilen ve yürütme sonunda PS+4 ya da dallanma hedefiyle yeniden yazılan bir yazmaçtır.

İlerleyen alt bölümlerde önce veri yolunun saat beslemesini ve gerçekleştireceğimiz RISC-V buyruk kümesini tanıtacak, ardından veri yolunun her bir bileşenini ayrı ayrı ele alacağız. Bölümün sonunda bu bileşenleri bir araya getirerek tek vuruşluk bir RISC-V veri yolu inşa etmiş olacağız.

\subsection{Saat Beslemesi}
\label{subsec:saat-beslemesi}

Veri yolundaki işlemlerin doğru sırayla gerçekleşmesi için tüm yazmaçların değerlerini \emph{aynı anda} ve \emph{öngörülebilir aralıklarla} güncellemesi gerekir. Bu eşgüdümü sağlayan sinyale \textbf{saat}\index{saat}\index{clock|see{saat}} (\emph{clock}) denir. Saat sinyali iki gerilim düzeyi (\emph{yüksek} ve \emph{düşük}) arasında düzenli olarak salınır. Ardışık iki yükseliş arasındaki süreye bir \textbf{çevrim}\index{saat çevrimi}\index{clock cycle|see{saat çevrimi}} (\emph{clock cycle}) denir (Şekil~\ref{fig:saat-vurusu}).

Saatin eşgüdüm dışında bir başka temel görevi daha vardır. İşlemciye bir \emph{zaman ölçeği} kazandırır. Donanım kendi başına ``ne kadar zaman geçti?'' sorusunu yanıtlayamaz; ancak kaç saat çevriminin geçtiğini sayabilir. Bu nedenle Bölüm~\ref{chap:basarim}'de gördüğümüz başarım hesabında zaman doğrudan saat çevrimi cinsinden ifade edilir. Saat olmadan işlemcide ``zaman'' tanımlı değildir.

Saatin neden gerekli olduğunu görmek için veri yolundaki iki tür devreyi ayırmak gerekir. AMB, çoklayıcılar ve kapılar gibi \textbf{birleşik}\index{birleşik devre}\index{combinational|see{birleşik devre}} devreler, girişleri değiştikçe çıkışlarını sürekli yeniden hesaplar; ancak bu çıkışlar yalnızca tellerde ilerler ve teller bilgiyi saklayamaz. Saklama görevi yazmaçlara ve belleğe aittir (Ek~\ref{ch:sayisal-tasarim}). Saat sinyali, birleşik devrelerin ürettiği değerlerin doğru anda yazmaçlara ``kilitlenmesini'' sağlar. Bu kilitleme yapılmadıkça bir buyruğun sonucu bir sonraki buyruğun girdisi olamaz.

\begin{figure}[htbp]
\centering
\begin{tikzpicture}[zamansinyal/.style={very thick, sinyalmavi},
                    zamanok/.style={<->, >=Stealth, thin},
                    zamanetiket/.style={font=\footnotesize}]
  % Saat sinyali dalga formu
  \draw[zamansinyal]
    (0,0) -- (0,1.5) -- (1.5,1.5) -- (1.5,0) -- (3,0) -- (3,1.5) -- (4.5,1.5) --
    (4.5,0) -- (6,0) -- (6,1.5) -- (7.5,1.5) -- (7.5,0) -- (9,0) -- (9,1.5) -- (10.5,1.5) -- (10.5,0) -- (11,0);
  % Yüksek/Düşük etiketleri
  \node[zamanetiket, left] at (-0.2, 1.5) {Yüksek};
  \node[zamanetiket, left] at (-0.2, 0) {Düşük};
  % 1 çevrim oku
  \draw[zamanok] (0,-0.7) -- (3,-0.7);
  \node[zamanetiket, below] at (1.5,-0.7) {1 çevrim};
  \draw[zamanok] (3,-0.7) -- (6,-0.7);
  \node[zamanetiket, below] at (4.5,-0.7) {1 çevrim};
  % Yükselen kenar
  \draw[->, >=Stealth, thick, sinyalkirmizi] (3.0, 2.3) -- (3.0, 1.7);
  \node[zamanetiket, above, text=sinyalkirmizi, align=center] at (3.0, 2.3) {Yükselen\\kenar};
  % Düşen kenar
  \draw[->, >=Stealth, thick, sinyalyesil] (4.5, 2.3) -- (4.5, 1.7);
  \node[zamanetiket, above, text=sinyalyesil, align=center] at (4.5, 2.3) {Düşen\\kenar};
  % Zaman ekseni
  \draw[->, >=Stealth, thin] (-0.5,-1.5) -- (11.5,-1.5);
  \node[zamanetiket, right] at (11.5,-1.5) {Zaman};
\end{tikzpicture}
\caption[Saat sinyali ve çevrim]{Saat sinyali ve çevrim. Yükselen ve düşen kenarlar işaretlenmiştir. Tipik tasarımda yazmaçlar yükselen kenarda tetiklenir.}
\label{fig:saat-vurusu}
\end{figure}

Yazmaçlar, saat sinyalinin yalnızca bir kenarında tetiklenir. Çağdaş tasarımların büyük çoğunluğunda bu \emph{yükselen kenar}\index{yükselen kenar}\index{kenar tetiklemeli}\index{edge-triggered|see{kenar tetiklemeli}}'dır. Bir çevrim böyle ilerler: yükselen kenarda yazmaçlar yeni değerlerini yakalar. Ardından gelen tüm çevrim boyunca birleşik devreler bu değerler üzerinde işlem yapar. Bir sonraki yükselen kenara kadar üretilen sonuç hedef yazmacın girişinde kararlı hâle gelmiş olmalıdır, yoksa yanlış değer yakalanır. Bu kısıt, çevrim süresinin alt sınırını belirler.

İşlemcinin \textbf{saat sıklığı}\index{saat sıklığı}\index{clock frequency|see{saat sıklığı}}, çevrim süresinin tersidir. Tek vuruşluk tasarımda çevrim, en yavaş buyruğun veri yolundaki en uzun birleşik mantık yolunu (kritik yol\index{kritik yol}\index{critical path|see{kritik yol}}, \emph{critical path}) tamamlamasına yetecek kadar uzun olmalıdır. Bu yüzden saat sıklığı doğrudan bu yol gecikmesiyle sınırlanır. Çok vuruşluk tasarımda ise yalnızca aşamalardan birini tamamlamak yeterli olduğu için aynı donanımda saat sıklığı daha yükseğe çıkarılabilir.

\subsection{Gerçekleştirilecek RISC-V Buyruk Kümesi}
\label{subsec:riscv-buyruk-kumesi}

Bir işlemci tasarımına başlamadan önce iki düzeyde karar verilmesi gerekir: önce \emph{hangi buyruk kümesi mimarisinin} (BKM) gerçekleştirileceği (RISC-V, x86, ARM gibi), sonra o mimarinin \emph{hangi alt kümesinin} işlemcide yer alacağı. Birinci karar bu kitapta zaten verilmiştir. Açık kaynak, yalın ve modüler bir tasarımı olan RISC-V'i kullanıyoruz (Bölüm~\ref{chap:buyruk-kumesi}). İkinci karar olan alt küme seçimi ise bu bölümün konusudur. RISC-V'in temel RV32I kümesi kırktan fazla buyruk içerir; ancak bir veri yolunun nasıl tasarlandığını öğrenmek için bunların tümünü ele almak gerekmez. Tasarımın özünü bozmadan üzerinde çalışılabilir bir alt küme seçmek hem öğretici hem de pratiktir. Gerçek işlemcilerin tasarımı da çoğu zaman çekirdek bir altkümeyi destekleyen veri yolundan başlar, geri kalan buyruklar bunun üzerine eklenir.

Seçtiğimiz alt küme, Bölüm~\ref{chap:buyruk-kumesi}'te tanıtılan biçimlerden R, I, S, B ve J'yi en az bir buyrukla temsil eden sekiz buyruktan oluşur (U biçimi, aşağıda değinildiği gibi bu alt kümenin dışında bırakılmıştır) ve bir programın gerek duyduğu temel davranışları (aritmetik--mantık işlemleri, bellek erişimi, koşullu dallanma, koşulsuz atlama) kapsar:

\begin{itemize}
    \item \textbf{Aritmetik ve mantık} (R türü): \texttt{add} (topla), \texttt{sub} (çıkar).
    \item \textbf{Anlık değerle mantık} (I türü): \texttt{ori} (mantıksal VEYA anlık).
    \item \textbf{Bellekten yükleme} (I türü): \texttt{lw} (yükle).
    \item \textbf{Belleğe yazma} (S türü): \texttt{sw} (sakla).
    \item \textbf{Koşullu dallanma} (B türü): \texttt{beq} (eşitse dallan), \texttt{bne} (eşit değilse dallan).
    \item \textbf{Koşulsuz atlama} (J türü): \texttt{jal} (atla ve bağla).
\end{itemize}

Bu sekiz buyruğun her biri veri yolunda farklı bir davranış ortaya çıkarır ve veri yolunun farklı bir bölümünü ``kullanır''. R ve I türü aritmetik buyruklar yazmaç öbeğinden okur, AMB'den geçer ve sonucu yine yazmaç öbeğine yazar. \texttt{lw} ek olarak veri belleğinden okur ve okunan değeri yazmaca yazar. \texttt{sw} yazmaç öbeğinden okuduğu değeri veri belleğine yazar. \texttt{beq} ve \texttt{bne} AMB'yi karşılaştırma için kullanır ve sonuca göre PS'i günceller. \texttt{jal} ise dönüş adresini yazmaca yazarken PS'i koşulsuz olarak hedefe taşır. Bir veri yolunun bu beş farklı davranışı destekleyebilmesi için PS, buyruk belleği, yazmaç öbeği, AMB ve veri belleğini doğru noktalarda birbirine bağlaması ve her buyruğun türüne göre uygun yolu açıp diğerlerini kapatması gerekir. Bu açma--kapama işini denetim birimi yapar.

RISC-V'in geri kalan buyrukları (R türünün \texttt{and}, \texttt{or}, \texttt{slt} gibi diğer üyeleri; \texttt{lui}, \texttt{auipc}, kaydırma buyrukları, \texttt{jalr} vb.) bu alt kümenin dışında bırakılmıştır. Örneğin \texttt{and}/\texttt{or}/\texttt{slt}, \texttt{add} ile tıpatıp aynı veri yolunu izler. Eklenmeleri yalnızca \texttt{AMBİşlem} kodlamasına yeni değerler katmayı gerektirir ve bölüm sonu alıştırmalarına bırakılmıştır.

Buyruk kümesini belirlemekle ilk adımı atmış olduk. Bir işlemci tasarımı genel olarak altı adımlı bir süreç olarak ele alınabilir:
\begin{enumerate}
    \item Gerçekleştirilecek buyruk kümesinin belirlenmesi.
    \item Gereken donanım birimlerinin ve saat stratejisinin seçilmesi.
    \item Veri yolunun birleştirilmesi.
    \item Denetim işaretlerinin belirlenmesi.
    \item Denetim biriminin tasarlanması.
    \item Denetim biriminin ve veri yolunun birleştirilmesi.
\end{enumerate}
İlerleyen alt bölümlerde bu adımları sırayla izleyerek tek vuruşluk bir RISC-V işlemcisi inşa edeceğiz.

\section{Tek Vuruşluk Veri Yolu}
\label{sec:tek-vurusluk}

Tek vuruşluk veri yolu, her buyruğun tam olarak \emph{bir} saat çevriminde başlayıp bittiği işlemci tasarımıdır. İşlemcide herhangi bir anda yalnızca bir buyruk bulunur. Bir buyruk tamamlanmadan ikincisi başlamaz. Her çevrimde tek bir buyruk tamamlanan bu yaklaşımda buyruk başına çevrim sayısı (BBÇ) tam olarak 1'dir. Bu, her saat çevriminde en çok bir buyruk alabilen işlemciler için ulaşılabilecek en iyi değerdir. Bölüm~\ref{chap:boruhatti}'de göreceğimiz boru hatlı tasarım bu sınıra yaklaşır. İleride değineceğimiz çoklu buyruk başlatımlı (\emph{superscalar}) işlemciler ise her çevrimde birden çok buyruk alarak BBÇ'yi 1'in altına indirebilir.

Ancak bu yalın görünüşün arkasında önemli bir ödünleşme vardır. Bölüm~\ref{chap:buyruk-kumesi}'te gördüğümüz gibi RISC-V buyrukları farklı işler yapar ve farklı donanım birimlerine ihtiyaç duyar. \texttt{add} yalnızca yazmaç öbeği ve AMB'ye uğrar; \texttt{lw} bunlara ek olarak veri belleğine de uğrar; \texttt{beq} ise yazmaca yazma yapmaz. Bu farklı yol uzunluklarına rağmen tek vuruşluk işlemcide saat çevrim zamanı en yavaş buyruğun veri yolunda izleyeceği en uzun yolu (kritik yolu) tamamlamasına yetecek kadar uzun seçilmek zorundadır. Aksi halde en yavaş buyruk tek çevrimde bitemez ve yanlış sonuç üretir. Bu kısıt nedeniyle, kısa yola sahip hızlı buyruklar (örneğin \texttt{add}, \texttt{beq}) bile en yavaş buyruğun (RISC-V'de \texttt{lw}) belirlediği uzun çevrim zamanına mahkum olur ve veri yolunda kullanmadıkları zamanı boş geçirir. Şekil~\ref{fig:tek-vurusluk-zaman}'de bu durum görülmektedir. Aynı saat çevrimi içinde \texttt{lw} bütün aşamaları doldururken \texttt{beq} çevrimin neredeyse yarısını boş geçirir.

\begin{figure}[htbp]
\centering
% --- TASLAK: Tek vuruşluk işlemcide buyruk türlerine göre aşama süreleri ve boşa geçen zaman ---
\begin{tikzpicture}[
    asama/.style={draw, minimum height=0.6cm, font=\footnotesize, anchor=west},
    bos/.style={draw, fill=gray!20, minimum height=0.6cm, font=\footnotesize, anchor=west, text=gray!70!black}
]
  \def\sira{0}
  % lw: BB(200) + YO(100) + AMB(200) + VB(200) + GY(100) = 800 ps, boşta 0 (en yavaş)
  \node[font=\footnotesize, anchor=east] at (0,0) {\texttt{lw}:};
  \node[asama, fill=blokmavi!30, minimum width=2cm] at (0.2,0) {BB};
  \node[asama, fill=sinyalyesil!30, minimum width=1cm] at (2.2,0) {YO};
  \node[asama, fill=sinyalkirmizi!30, minimum width=2cm] at (3.2,0) {AMB};
  \node[asama, fill=blokmavi!30, minimum width=2cm] at (5.2,0) {VB};
  \node[asama, fill=sinyalyesil!30, minimum width=1cm] at (7.2,0) {GY};
  % sw: BB + YO + AMB + VB = 700 ps, boşta 100
  \node[font=\footnotesize, anchor=east] at (0,-0.8) {\texttt{sw}:};
  \node[asama, fill=blokmavi!30, minimum width=2cm] at (0.2,-0.8) {BB};
  \node[asama, fill=sinyalyesil!30, minimum width=1cm] at (2.2,-0.8) {YO};
  \node[asama, fill=sinyalkirmizi!30, minimum width=2cm] at (3.2,-0.8) {AMB};
  \node[asama, fill=blokmavi!30, minimum width=2cm] at (5.2,-0.8) {VB};
  \node[bos, minimum width=1cm] at (7.2,-0.8) {boş};
  % R türü (add): BB(200) + YO(100) + AMB(200) + GY(100) = 600 ps, boşta 200
  \node[font=\footnotesize, anchor=east] at (0,-1.6) {R türü (\texttt{add}):};
  \node[asama, fill=blokmavi!30, minimum width=2cm] at (0.2,-1.6) {BB};
  \node[asama, fill=sinyalyesil!30, minimum width=1cm] at (2.2,-1.6) {YO};
  \node[asama, fill=sinyalkirmizi!30, minimum width=2cm] at (3.2,-1.6) {AMB};
  \node[asama, fill=sinyalyesil!30, minimum width=1cm] at (5.2,-1.6) {GY};
  \node[bos, minimum width=2cm] at (6.2,-1.6) {boşta};
  % ori: aynı R türü = 600 ps, boşta 200
  \node[font=\footnotesize, anchor=east] at (0,-2.4) {\texttt{ori}:};
  \node[asama, fill=blokmavi!30, minimum width=2cm] at (0.2,-2.4) {BB};
  \node[asama, fill=sinyalyesil!30, minimum width=1cm] at (2.2,-2.4) {YO};
  \node[asama, fill=sinyalkirmizi!30, minimum width=2cm] at (3.2,-2.4) {AMB};
  \node[asama, fill=sinyalyesil!30, minimum width=1cm] at (5.2,-2.4) {GY};
  \node[bos, minimum width=2cm] at (6.2,-2.4) {boşta};
  % beq: BB + YO + AMB = 500 ps, boşta 300 (en hızlı)
  \node[font=\footnotesize, anchor=east] at (0,-3.2) {\texttt{beq}:};
  \node[asama, fill=blokmavi!30, minimum width=2cm] at (0.2,-3.2) {BB};
  \node[asama, fill=sinyalyesil!30, minimum width=1cm] at (2.2,-3.2) {YO};
  \node[asama, fill=sinyalkirmizi!30, minimum width=2cm] at (3.2,-3.2) {AMB};
  \node[bos, minimum width=3cm] at (5.2,-3.2) {boşta};
  % --- Eksen ---
  \draw[->, thick] (0.2,-3.9) -- (8.5,-3.9) node[right, font=\footnotesize] {Zaman};
  \foreach \x/\t in {0.2/0, 2.2/200, 3.2/300, 5.2/500, 7.2/700, 8.2/800} {
    \draw (\x,-3.85) -- (\x,-3.95);
    \node[font=\tiny, below] at (\x,-3.95) {\t};
  }
  \draw[dashed, sinyalkirmizi] (8.2,0.35) -- (8.2,-3.85);
  \node[font=\footnotesize, sinyalkirmizi, anchor=south] at (8.2,0.35) {Çevrim sonu};
\end{tikzpicture}
\caption[Tek vuruşluk işlemcide buyruk türlerine göre aşama süreleri]{Tek vuruşluk işlemcide buyruk türlerine göre aşama süreleri (en yavaştan en hızlıya sıralı). BB: buyruk belleği, YO: yazmaç okuma, AMB: aritmetik mantık birimi, VB: veri belleği, GY: geri yazma. Saat çevrimi en yavaş buyruğun (\texttt{lw}) süresine göre belirlendiğinden, \texttt{beq} ve \texttt{add} gibi kısa buyruklar çevrimin önemli bir bölümünü boş geçirir.}
\label{fig:tek-vurusluk-zaman}
\end{figure}

Bu ödünleşmenin başarım üzerindeki etkisini görmek için sayısal bir örnek yapalım.

\begin{ornek}[6.1 -- Tek vuruşluk işlemcide kritik yol ve başarım]
Bir tek vuruşluk RISC-V işlemcisinde her birimin gecikme süresi şöyle verilsin:
\begin{center}
\small
\begin{tabular}{lc}
\toprule
\textbf{Birim} & \textbf{Gecikme (ps)} \\
\midrule
Buyruk belleği erişimi (BB)   & 200 \\
Yazmaç öbeği okuma (YO)       & 100 \\
AMB                            & 200 \\
Veri belleği erişimi (VB)     & 200 \\
Yazmaç öbeği yazma (GY)       & 100 \\
\bottomrule
\end{tabular}
\end{center}
Bir programda buyrukların \%25'i \texttt{lw}, \%10'u \texttt{sw}, \%45'i R/I türü aritmetik, \%20'si dallanma (\texttt{beq}/\texttt{bne}) olsun. Tek vuruşluk işlemcinin ortalama buyruk başına süresini hesaplayın ve bu değeri her buyruğun yalnızca \emph{kendi gereksinimi kadar} süre alabileceği varsayımsal bir tasarımla karşılaştırın.

\textbf{Çözüm.} Her buyruk türünün kritik yol gecikmesi:
\begin{itemize}[nosep]
    \item \texttt{lw}\,: BB + YO + AMB + VB + GY $= 200 + 100 + 200 + 200 + 100 = \mathbf{800}$\,ps
    \item \texttt{sw}\,: BB + YO + AMB + VB $= 200 + 100 + 200 + 200 = 700$\,ps
    \item R/I aritmetik--mantık (\texttt{add}, \texttt{sub}, \texttt{ori} vb.): BB + YO + AMB + GY $= 200 + 100 + 200 + 100 = 600$\,ps
    \item \texttt{beq}\,: BB + YO + AMB $= 200 + 100 + 200 = 500$\,ps
\end{itemize}

Eğer her buyruk yalnızca kendi gecikmesi kadar süre alabilseydi, programın buyruk karışımına göre buyruk başına ortalama gecikme şu olurdu:
\[
\bar{T} = 0{,}25 \cdot 800 + 0{,}10 \cdot 700 + 0{,}45 \cdot 600 + 0{,}20 \cdot 500 = 640\,\text{ps}
\]

Tek vuruşluk işlemcide ise saat çevrimi en yavaş buyruğa göre belirlendiği için tüm buyruklar 800\,ps süre alır. Dolayısıyla buyruk başına ortalama gecikme de \emph{800\,ps}'ye çıkar. Tek vuruşluk işlemci ortalama gecikmeyi buyruk başına $\frac{800}{640} \approx 1{,}25$ kat, başka bir deyişle \%25 büyütmüş olur. Bu uzama tamamen \texttt{beq} ve \texttt{add} gibi hızlı buyrukların uzun çevrimde boş beklemesinden kaynaklanır.

Bu sonucu Bölüm~\ref{chap:basarim}'de gördüğümüz başarım denklemiyle de yorumlayabiliriz:
\[
T_{\text{yürütme}} = N \times \text{BBÇ} \times T_{\text{çevrim}}
\]
Tek vuruşluk işlemcide BBÇ tam olarak 1'dir. Bu, denklemdeki BBÇ teriminin (her çevrimde bir buyruk alan tasarımlar için) ulaşabileceği en iyi değerdir. Ancak bu kazancın bedeli $T_{\text{çevrim}}$'in en yavaş buyruğa eşitlenmesiyle ödenir.

\begin{center}
\small
\begin{tabular}{lccc}
\toprule
\textbf{Tasarım} & \textbf{BBÇ} & \boldmath$T_{\text{çevrim}}$ & \textbf{Buyruk başına süre} \\
\midrule
İdeal (varsayımsal) & --- & --- & 640\,ps \\
Tek vuruşluk         & 1 (en iyi) & 800\,ps (en yavaş buyruk) & 800\,ps \\
\bottomrule
\end{tabular}
\end{center}

Tek vuruşluk işlemcinin başarımı artırmak için tek aracı $T_{\text{çevrim}}$'i kısaltmaktır; ancak çevrim zamanı en yavaş buyrukla sınırlı olduğundan bu kısaltma da sınırlıdır. Sonraki alt bölümde göreceğimiz çok vuruşluk tasarım, BBÇ'yi 1'in üstüne çıkarma pahasına $T_{\text{çevrim}}$'i küçültür. Toplam yürütme zamanı bu iki etkenin çarpımıyla belirlenir.
\end{ornek}

Bu ödünleşmeyi anladıktan sonra artık tek vuruşluk veri yolunu inşa etmeye başlayabiliriz. Tasarımı önümüze yığarak değil, \emph{buyruk buyruk inşa ederek} ele alacağız. Önce en yalın RISC-V buyruğu olan \texttt{add}'i yürütebilecek olabildiğince yalın bir veri yolu çizecek. Sonra her yeni buyruk için ``bunu yapabilmek için veri yoluna ne eklemek gerek?'' sorusunu yanıtlayarak donanımı büyüteceğiz. Her aşamada yeni eklenen bileşeni vurgulayacağız. Böylece her çoklayıcının, her bellek bağlantısının ve her denetim sinyalinin neden var olduğu açık biçimde görülecek. İzleyeceğimiz buyruk sırası şudur: \texttt{add} $\rightarrow$ \texttt{sub} $\rightarrow$ \texttt{ori} $\rightarrow$ \texttt{lw} $\rightarrow$ \texttt{sw} $\rightarrow$ \texttt{beq}/\texttt{bne} $\rightarrow$ \texttt{jal}.

\subsection{İlk Adım: \texttt{add} Buyruğu için Yalın Veri Yolu}
\label{subsec:vy-add}

İlk olarak yalnızca \texttt{add x5, x6, x7} (\texttt{x5 = x6 + x7}) buyruğunu yürütebilecek bir işlemci tasarlayalım. Bu buyruğu yürütmek için işlemcinin sırasıyla şunları yapması gerekir:
\begin{enumerate}
    \item Program sayacının gösterdiği adresteki buyruğu bellekten getir.
    \item Buyruğun bit alanlarından kaynak yazmaç numaralarını (\texttt{rs1}, \texttt{rs2}) ve hedef yazmaç numarasını (\texttt{rd}) çıkar.
    \item Yazmaç öbeğinden \texttt{rs1} ve \texttt{rs2} değerlerini oku.
    \item Bir toplayıcıda bu iki değeri topla.
    \item Sonucu yazmaç öbeğindeki \texttt{rd}'ye yaz.
    \item Bir sonraki buyruğa geçmek için PS'i 4 artır.
\end{enumerate}

Bu adımları gerçekleştirebilmek için veri yolunda beş temel birim yeterlidir: \textbf{program sayacı (PS)}, \textbf{+4 toplayıcı}, \textbf{buyruk belleği}, \textbf{yazmaç öbeği} ve bir \textbf{toplayıcı}. Henüz çıkarma, yükleme veya dallanma yapmadığımız için tam donanımlı bir AMB'ye değil, yalnızca toplama yapan bir birime ihtiyacımız var. Bir sonraki adımda \texttt{sub} buyruğunu eklediğimizde bu toplayıcının yerini AMB alacak. Şekil~\ref{fig:vy-add}'te bu beş birimin nasıl bağlandığı gösterilmektedir.

\begin{figure}[htbp]
\centering
\begin{tikzpicture}[>=Stealth, semithick, font=\footnotesize,
    yzm/.style={draw, fill=yellow!25, minimum width=1cm, minimum height=1cm, align=center},
    bllk/.style={draw, fill=green!20, minimum width=1.8cm, minimum height=2.2cm, align=center},
    yobk/.style={draw, fill=blue!10, minimum width=2.6cm, minimum height=3cm},
    top/.style={draw, fill=red!15, minimum width=1.6cm, minimum height=2cm, align=center},
    artop/.style={draw, fill=gray!15, minimum width=1.1cm, minimum height=0.9cm, align=center},
    bitayar/.style={fill=white, inner sep=1pt, font=\tiny}
]
  % --- PS ve +4 toplayıcı (PS bloğu 2 cm yukarı: Top. üst kenarı buyruk şeridi üst kenarıyla aynı hizada) ---
  \node[yzm] (ps) at (0,2) {PS};
  \node[artop] (top4) at (0, 3.9) {Top.};
  \node[bllk] (bb) at (2.6,2) {Buyruk\\Belleği};
  % PS -> BB adres yolu (etiket telin altında); bit etiketi BB tarafına kaydırıldı (nokta üstüyle çakışmasın)
  \draw[->] (ps.east) -- (bb.west) node[midway, below, font=\scriptsize] {adres};
  % Adres yolunda dallanma noktası
  \coordinate (adresJ) at (1.1, 2);
  \node[circle, fill=black, inner sep=1pt] at (adresJ) {};

  % 4 sabiti girişi - Top.'un sağ kenarı üst tarafına yatay (sağdan sola); ok uzatıldı, "32 bit" tel üstünde
  % "4"un altında 32 bitlik ikili gösterim parantez içinde (sayının 32 bit olarak nasıl göründüğünü anlatır)
  \node[font=\scriptsize, align=center] at ([xshift=1.55cm, yshift=0.22cm]top4.east) {$4$\\[-2pt]\tiny$(0{\ldots}0100)_2$};
  \draw[->] ([xshift=1.15cm, yshift=0.22cm]top4.east) -- ([yshift=0.22cm]top4.east)
            node[bitayar, pos=0.5] {32 bit};
  % PS dallanması - adres yolundan yukarı, Top.'un sağ kenar alt tarafına
  \draw[->] (adresJ) -- ++(0,1.68) coordinate (psup)
            -- ([yshift=-0.22cm]top4.east);
  \node[font=\scriptsize, anchor=east] at ([xshift=-0.08cm, yshift=0.84cm]adresJ) {PS};
  % Top. çıkışı - sol kenardan PS+4 olarak çıkar, PS'in çok solunda dirsek yaparak PS sol kenarına girer
  % "32 bit" etiketi son yatay segmentin ortasında (PS yazısının solunda boşluk olacak şekilde)
  \draw[->] (top4.west) -- ++(-1.4,0) node[midway, above, font=\scriptsize] {PS$+$4}
            |- (ps.west) node[bitayar, pos=0.78] {32 bit};

  % --- Buyruk yolu (BB'den çıkan ana tel) ---
  \coordinate (parse) at (4.18, 2);
  \draw[->] (bb.east) -- (parse) node[midway, below, font=\scriptsize] {buyruk} node[bitayar, pos=0.55] {32 bit};

  % --- R türü buyruk biçimi şeridi (parse'dan büyüteç gibi yukarı) ---
  % Şerit yerleşim: parse'ın üstünde, toplam genişlik 6cm (her bit = 0.1875 cm)
  % Sola kaydırıldı: rs1 alanı A_adr portuyla aynı x'te (7.4) hizalı, rs1 teli düz dikey iner
  \coordinate (seritsol) at (4.68, 3.8);
  % Büyüteç çizgileri (parse'tan şerit köşelerine kesik)
  \draw[gray, dashed, semithick] (parse) -- (seritsol);
  \draw[gray, dashed, semithick] (parse) -- ([xshift=6cm]seritsol);
  % Şerit hücreleri (R türü: funct7 | rs2 | rs1 | funct3 | rd | opcode)
  % bit genişlikleri: 7, 5, 5, 3, 5, 7 → cm: 1.3125, 0.9375, 0.9375, 0.5625, 0.9375, 1.3125
  \begin{scope}[shift={(seritsol)}]
    \draw[fill=gray!10] (0,0) rectangle (1.3125,0.55) node[pos=.5, font=\tiny] {funct7};
    \draw[fill=green!18] (1.3125,0) rectangle (2.25,0.55) node[pos=.5, font=\tiny] {rs2};
    \draw[fill=green!18] (2.25,0) rectangle (3.1875,0.55) node[pos=.5, font=\tiny] {rs1};
    \draw[fill=gray!10] (3.1875,0) rectangle (3.75,0.55) node[pos=.5, font=\tiny] {f3};
    \draw[fill=blue!18] (3.75,0) rectangle (4.6875,0.55) node[pos=.5, font=\tiny] {rd};
    \draw[fill=gray!10] (4.6875,0) rectangle (6,0.55) node[pos=.5, font=\tiny] {opcode};
    % Bit numaraları altında
    \node[font=\tiny, anchor=north] at (0,0) {31};
    \node[font=\tiny, anchor=north] at (1.3125,0) {25};
    \node[font=\tiny, anchor=north] at (2.25,0) {20};
    \node[font=\tiny, anchor=north] at (3.1875,0) {15};
    \node[font=\tiny, anchor=north] at (3.75,0) {12};
    \node[font=\tiny, anchor=north] at (4.6875,0) {7};
    \node[font=\tiny, anchor=north] at (6,0) {0};
    % Üstte tip etiketi
    \node[font=\scriptsize, anchor=south] at (3,0.55) {R türü buyruk};
  \end{scope}

  % --- Yazmaç Öbeği ---
  \node[yobk] (yo) at (8, 0) {};
  \node[anchor=center, font=\footnotesize\bfseries, align=center] at (yo.center) {Yazmaç\\Öbeği};

  % Şeritten rs1, rs2, rd → Yazmaç Öbeği üst portlarına (A, B, C)
  % RISC-V kuralı: rs1 -> A_adr (1. okuma), rs2 -> B_adr (2. okuma); AMB hesabı A-B = rs1-rs2
  % Şerit içi merkezler: rs2 = seritsol + (1.78, 0); rs1 = (2.72, 0); rd = (4.22, 0)
  % Port x'leri: A_adr = yo.x - 0.6 = 7.40; B_adr = yo.x = 8.00; C_adr = yo.x + 0.6 = 8.60
  % Üç tel farklı y'lerde yatay segmentlerle iner; çapraz olsa da çakışmaz
  \draw[->] ([xshift=2.72cm]seritsol) -- ++(0,-1.2) -| ([xshift=-0.6cm]yo.north)
            node[bitayar, pos=0.75] {5 bit};  % rs1 -> A_adr (y=2.6'da sola dirsek)
  \draw[->] ([xshift=1.78cm]seritsol) -- ++(0,-0.8) -| ([xshift=0cm]yo.north)
            node[bitayar, pos=0.75] {5 bit};  % rs2 -> B_adr (y=3.0'da sağa dirsek)
  \draw[->] ([xshift=4.22cm]seritsol) -- ++(0,-0.4) -| ([xshift=0.6cm]yo.north)
            node[bitayar, pos=0.75] {5 bit};  % rd  -> C_adr (y=3.4'te sola dirsek)

  % Üst port etiketleri (yazmaç adresi/numarası) - kutu üst kenarına yakın
  \node[font=\scriptsize, anchor=north] at ([xshift=-0.6cm,yshift=-0.05cm]yo.north) {$A_{\text{adr}}$};
  \node[font=\scriptsize, anchor=north] at ([xshift=0cm,yshift=-0.05cm]yo.north) {$B_{\text{adr}}$};
  \node[font=\scriptsize, anchor=north] at ([xshift=0.6cm,yshift=-0.05cm]yo.north) {$C_{\text{adr}}$};

  % Sağ port etiketleri (A ve B portlarının verileri) - birbirine yakınlaştırıldı
  \node[font=\scriptsize, anchor=east] at ([xshift=-0.1cm,yshift=0.5cm]yo.east) {$A_{\text{veri}}$};
  \node[font=\scriptsize, anchor=east] at ([xshift=-0.1cm,yshift=-0.5cm]yo.east) {$B_{\text{veri}}$};

  % Sol port: C portuna yazılacak veri (alt sol)
  \node[font=\scriptsize, anchor=west] at ([xshift=0.1cm,yshift=-1.0cm]yo.west) {$C_{\text{veri}}$};

  % Yazmaca Yaz §6.2.5 sw'de tanıtılacak (henüz her zaman 1 — şimdilik gizli)

  % --- Toplayıcı (add buyruğu için AMB yerine, aşağı kaydırıldı) ---
  \node[top, font=\scriptsize] (tpl) at (11.6, 0) {Toplayıcı};

  % Yazmaç Öbeği -> Toplayıcı: Averi/Bveri TAM YATAY (Toplayıcı giriş portları aynı y'de)
  \draw[->] ([yshift=0.5cm]yo.east) -- ([yshift=0.5cm]tpl.west) node[bitayar, pos=0.5] {32 bit};
  \draw[->] ([yshift=-0.5cm]yo.east) -- ([yshift=-0.5cm]tpl.west) node[bitayar, pos=0.5] {32 bit};

  % Toplayıcı çıkışı -> sağa düz, sonra dirsek; "toplam" yatay segmentin ortasında üstte, "32 bit" diğer tellerdeki gibi tel üzerinde (fill=white)
  % Alttaki dirsek Yazmaca Yaz yazısının hemen altında olacak şekilde yukarı çekildi (-3.3 → -2.5)
  \draw[->] (tpl.east) -- ++(1.6,0)
            node[pos=0.5, above, font=\scriptsize] {toplam}
            node[bitayar, pos=0.5] {32 bit}
            |- ++(0,-2.5)
            -| ([xshift=-1.0cm,yshift=-1.0cm]yo.west)
            -- ([yshift=-1.0cm]yo.west);
  % Sol alt köşe: bu şekilde uygulanmış buyrukların listesi
  \node[draw, rounded corners, fill=blue!5, font=\tiny, align=center, inner sep=4pt, anchor=north east] at (14.0, 4.6) {Buyruk durumu (1/7)\\[2pt]\begin{tabular}{@{}l@{\quad}l@{}}
1.~{\color{green!60!black}\checkmark}~\texttt{add} & 5.~{\color{gray}$\circ$}~\texttt{sw}\\
2.~{\color{gray}$\circ$}~\texttt{sub} & 6.~{\color{gray}$\circ$}~\texttt{beq}/\texttt{bne}\\
3.~{\color{gray}$\circ$}~\texttt{ori} & 7.~{\color{gray}$\circ$}~\texttt{jal}\\
4.~{\color{gray}$\circ$}~\texttt{lw} & \\
\end{tabular}};
\end{tikzpicture}
\caption[\texttt{add} buyruğunu yürütebilen yalın veri yolu]{\texttt{add} buyruğunu yürütebilen yalın veri yolu. Yazmaç öbeğinin üç \emph{adres girişi} vardır: $A_{\text{adr}}$ ve $B_{\text{adr}}$ okunacak iki yazmacın numarasını, $C_{\text{adr}}$ ise sonucun yazılacağı yazmacın numarasını alır. Buyruktan çıkan \texttt{rs1} ve \texttt{rs2} alanları $A_{\text{adr}}$ ve $B_{\text{adr}}$'a, \texttt{rd} alanı $C_{\text{adr}}$'a bağlanır. Karşılığında veri çıkışları $A_{\text{veri}}$ ve $B_{\text{veri}}$ toplayıcıya, toplayıcı sonucu da $C_{\text{veri}}$ üzerinden yazmaç öbeğine geri yazılır. PS ayrı bir toplayıcıyla 4 artırılır. Bu basit veri yolunda henüz hiçbir denetim sinyali yoktur: bütün yollar her saat çevriminde aktiftir.}
\label{fig:vy-add}
\end{figure}

Veri yolundaki her bileşenin rolü şudur:
\begin{itemize}
    \item \textbf{PS} sıradaki buyruğun adresini tutar. Her saat çevriminde bu adres buyruk belleğine adres olarak verilir ve çevrim sonunda PS+4 değeriyle güncellenir.
    \item \textbf{Buyruk belleği}, adres olarak PS'i alır ve o adresteki 32 bitlik buyruğu çıkışına verir.
    \item \textbf{Yazmaç öbeği} aynı anda iki yazmaç okuyup bir yazmaç yazabilen, 32 yazmaçlı bir saklama birimidir. Buyruktan gelen 5 bitlik yazmaç numaraları (\texttt{rs1}, \texttt{rs2}, \texttt{rd}) hangi yazmaçların okunup yazılacağını belirler.
    \item \textbf{Toplayıcı}, yazmaç öbeğinden gelen iki 32 bitlik veriyi toplar ve sonucu yazmaç öbeğine geri yazılmak üzere üretir. Bir sonraki adımda bu birimin yerini, hem toplama hem çıkarma yapabilen \textbf{AMB} (Aritmetik Mantık Birimi, Bölüm~\ref{chap:aritmetik}) alacak.
    \item \textbf{+4 toplayıcı}, PS'i 4 artırarak bir sonraki buyruğun adresini hesaplar.
\end{itemize}

Bu veri yolu \texttt{add} dışında hiçbir RISC-V buyruğunu yürütemez, ama önemli bir başlangıç noktasıdır. Şimdi yeni buyruklar ekleyerek bu temel iskeleti büyüteceğiz.

\subsection{\texttt{sub} Buyruğunu Eklemek: İlk Denetim Sinyali}
\label{subsec:vy-sub}

\texttt{sub x5, x6, x7} buyruğu \texttt{add}'in neredeyse aynısıdır. Yazmaç öbeğinden iki değer okunur, AMB'den geçirilir, sonuç hedef yazmaca yazılır. Tek fark, AMB'nin toplama yerine çıkarma yapmasıdır.

Donanım açısından AMB zaten hem toplama hem çıkarma yapabilen bir devredir (Bölüm~\ref{chap:aritmetik}). Eksik olan tek şey AMB'ye ``\emph{şu işlemi yap}'' diyen bir bağlantıdır. Örneğin tek bitlik bir hatta 0 gelirse toplama, 1 gelirse çıkarma yapılır. Bu bağlantıya \textbf{denetim sinyali} adı verilir ve bu özel sinyali \texttt{AMBİşlem} olarak adlandıralım. Şekil~\ref{fig:vy-sub}'te toplayıcının yerine AMB konmuş ve bu yeni denetim sinyali kırmızı çerçeveyle vurgulanmış veri yolu görülmektedir.

\begin{figure}[htbp]
\centering
\begin{tikzpicture}[>=Stealth, semithick, font=\footnotesize,
    yzm/.style={draw, fill=yellow!25, minimum width=1cm, minimum height=1cm, align=center},
    bllk/.style={draw, fill=green!20, minimum width=1.8cm, minimum height=2.2cm, align=center},
    yobk/.style={draw, fill=blue!10, minimum width=2.6cm, minimum height=3cm},
    amb/.style={draw, fill=red!20, minimum width=1.6cm, minimum height=2cm, align=center},
    artop/.style={draw, fill=gray!15, minimum width=1.1cm, minimum height=0.9cm, align=center},
    bitayar/.style={fill=white, inner sep=1pt, font=\tiny},
    yeni/.style={draw=red!70!black, very thick}
]
  % PS bloğu (Şekil 6.3 ile aynı)
  \node[yzm] (ps) at (0,2) {PS};
  \node[artop] (top4) at (0, 3.9) {Top.};
  \node[bllk] (bb) at (2.6,2) {Buyruk\\Belleği};
  \draw[->] (ps.east) -- (bb.west) node[midway, below, font=\scriptsize] {adres};
  \coordinate (adresJ) at (1.1, 2);
  \node[circle, fill=black, inner sep=1pt] at (adresJ) {};
  % 4 sabiti
  \node[font=\scriptsize, align=center] at ([xshift=1.55cm, yshift=0.22cm]top4.east) {$4$\\[-2pt]\tiny$(0{\ldots}0100)_2$};
  \draw[->] ([xshift=1.15cm, yshift=0.22cm]top4.east) -- ([yshift=0.22cm]top4.east)
            node[bitayar, pos=0.5] {32 bit};
  % PS dallanması
  \draw[->] (adresJ) -- ++(0,1.68) coordinate (psup)
            -- ([yshift=-0.22cm]top4.east);
  \node[font=\scriptsize, anchor=east] at ([xshift=-0.08cm, yshift=0.84cm]adresJ) {PS};
  % Top. çıkışı
  \draw[->] (top4.west) -- ++(-1.4,0) node[midway, above, font=\scriptsize] {PS$+$4}
            |- (ps.west) node[bitayar, pos=0.78] {32 bit};

  % --- Buyruk yolu ---
  \coordinate (parse) at (4.18, 2);
  \draw[->] (bb.east) -- (parse) node[midway, below, font=\scriptsize] {buyruk} node[bitayar, pos=0.55] {32 bit};

  % --- R türü buyruk şeridi (sola kaydı: rs1 -> A_adr düz dikey) ---
  \coordinate (seritsol) at (4.68, 3.8);
  \draw[gray, dashed, semithick] (parse) -- (seritsol);
  \draw[gray, dashed, semithick] (parse) -- ([xshift=6cm]seritsol);
  \begin{scope}[shift={(seritsol)}]
    \draw[fill=gray!10] (0,0) rectangle (1.3125,0.55) node[pos=.5, font=\tiny] {funct7};
    \draw[fill=green!18] (1.3125,0) rectangle (2.25,0.55) node[pos=.5, font=\tiny] {rs2};
    \draw[fill=green!18] (2.25,0) rectangle (3.1875,0.55) node[pos=.5, font=\tiny] {rs1};
    \draw[fill=gray!10] (3.1875,0) rectangle (3.75,0.55) node[pos=.5, font=\tiny] {f3};
    \draw[fill=blue!18] (3.75,0) rectangle (4.6875,0.55) node[pos=.5, font=\tiny] {rd};
    \draw[fill=gray!10] (4.6875,0) rectangle (6,0.55) node[pos=.5, font=\tiny] {opcode};
    \node[font=\tiny, anchor=north] at (0,0) {31};
    \node[font=\tiny, anchor=north] at (1.3125,0) {25};
    \node[font=\tiny, anchor=north] at (2.25,0) {20};
    \node[font=\tiny, anchor=north] at (3.1875,0) {15};
    \node[font=\tiny, anchor=north] at (3.75,0) {12};
    \node[font=\tiny, anchor=north] at (4.6875,0) {7};
    \node[font=\tiny, anchor=north] at (6,0) {0};
    \node[font=\scriptsize, anchor=south] at (3,0.55) {R türü buyruk};
  \end{scope}

  % --- Yazmaç Öbeği ---
  \node[yobk] (yo) at (8, 0) {};
  \node[anchor=center, font=\footnotesize\bfseries, align=center] at (yo.center) {Yazmaç\\Öbeği};

  % RISC-V kuralı: rs1 -> A_adr, rs2 -> B_adr (AMB hesabı A-B = rs1-rs2)
  \draw[->] ([xshift=2.72cm]seritsol) -- ++(0,-1.2) -| ([xshift=-0.6cm]yo.north) node[bitayar, pos=0.75] {5 bit};  % rs1 -> A_adr
  \draw[->] ([xshift=1.78cm]seritsol) -- ++(0,-0.8) -| ([xshift=0cm]yo.north) node[bitayar, pos=0.75] {5 bit};      % rs2 -> B_adr
  \draw[->] ([xshift=4.22cm]seritsol) -- ++(0,-0.4) -| ([xshift=0.6cm]yo.north) node[bitayar, pos=0.75] {5 bit};   % rd -> C_adr

  % Port etiketleri
  \node[font=\scriptsize, anchor=north] at ([xshift=-0.6cm,yshift=-0.05cm]yo.north) {$A_{\text{adr}}$};
  \node[font=\scriptsize, anchor=north] at ([xshift=0cm,yshift=-0.05cm]yo.north) {$B_{\text{adr}}$};
  \node[font=\scriptsize, anchor=north] at ([xshift=0.6cm,yshift=-0.05cm]yo.north) {$C_{\text{adr}}$};
  \node[font=\scriptsize, anchor=east] at ([xshift=-0.1cm,yshift=0.5cm]yo.east) {$A_{\text{veri}}$};
  \node[font=\scriptsize, anchor=east] at ([xshift=-0.1cm,yshift=-0.5cm]yo.east) {$B_{\text{veri}}$};
  \node[font=\scriptsize, anchor=west] at ([xshift=0.1cm,yshift=-1.0cm]yo.west) {$C_{\text{veri}}$};
  % Yazmaca Yaz §6.2.5 sw'de tanıtılacak (şu ana kadar her zaman 1 — gizli)

  % --- AMB (Toplayıcı yerine — YENİ, kırmızı çerçeve) ---
  \node[amb, yeni, font=\scriptsize] (ab) at (11.6, 0) {AMB};

  % Yazmaç Öbeği -> AMB
  \draw[->] ([yshift=0.5cm]yo.east) -- ([yshift=0.5cm]ab.west) node[bitayar, pos=0.5] {32 bit};
  \draw[->] ([yshift=-0.5cm]yo.east) -- ([yshift=-0.5cm]ab.west) node[bitayar, pos=0.5] {32 bit};

  % AMB çıkışı -> Cveri (toplam stil: yatay segmentin ortasında, dirsek Yaz alt boşluğuna)
  \draw[->] (ab.east) -- ++(1.6,0)
            node[pos=0.5, above, font=\scriptsize] {sonuç}
            node[bitayar, pos=0.5] {32 bit}
            |- ++(0,-2.5)
            -| ([xshift=-1.0cm,yshift=-1.0cm]yo.west)
            -- ([yshift=-1.0cm]yo.west);

  % AMBİşlem denetim sinyali (kırmızı, üstten — YENİ ve İLK denetim sinyali, kırmızı çerçeve)
  \node[font=\scriptsize, text=sinyalkirmizi, anchor=south, align=center] at ([yshift=0.5cm]ab.north) {AMBİşlem\\[-1pt]\tiny 0\,$=$\,Topla,\ 1\,$=$\,Çıkar};
  \draw[->, sinyalkirmizi, thick] ([yshift=0.4cm]ab.north) -- (ab.north);
  % Sol alt köşe: bu şekilde uygulanmış buyrukların listesi
  \node[draw, rounded corners, fill=blue!5, font=\tiny, align=center, inner sep=4pt, anchor=north east] at (14.0, 4.6) {Buyruk durumu (2/7)\\[2pt]\begin{tabular}{@{}l@{\quad}l@{}}
1.~{\color{green!60!black}\checkmark}~\texttt{add} & 5.~{\color{gray}$\circ$}~\texttt{sw}\\
2.~{\color{green!60!black}\checkmark}~\texttt{sub} & 6.~{\color{gray}$\circ$}~\texttt{beq}/\texttt{bne}\\
3.~{\color{gray}$\circ$}~\texttt{ori} & 7.~{\color{gray}$\circ$}~\texttt{jal}\\
4.~{\color{gray}$\circ$}~\texttt{lw} & \\
\end{tabular}};
\end{tikzpicture}
\caption[\texttt{add} ve \texttt{sub} buyruklarını yürüten veri yolu]{\texttt{add} ve \texttt{sub} buyruklarını yürüten veri yolu. Önceki aşamadan tek farklar: yalnızca toplama yapan \emph{toplayıcı} yerine hem toplama hem çıkarma yapabilen \textbf{AMB} (kırmızı çerçeve) konuldu. Bu birime \emph{yapacağı işlemi bildiren ilk denetim sinyalimiz} olan \texttt{AMBİşlem} eklendi. \texttt{AMBİşlem} sinyali buyruğun \texttt{funct7} ve \texttt{funct3} alanlarından üretilir. \texttt{add} için toplama, \texttt{sub} için çıkarma seçer.}
\label{fig:vy-sub}
\end{figure}

\texttt{AMBİşlem} sinyali buyruğun bit alanlarından (\texttt{funct7} ve \texttt{funct3}) üretilir. Şu an için bu sinyalin ``buyruktan geldiğini'' varsayalım. Denetim biriminin nasıl tasarlandığını Bölüm~\ref{sec:denetim-birimi}'te göreceğiz.

Bu adımda toplayıcının yerini hem toplama hem çıkarma yapabilen AMB aldı. Veri yolunun temel yapısı aynı kalsa da donanım karmaşıklığı arttı. Şimdilik AMB'nin yalnızca iki işlemi (toplama ve çıkarma) gerçekleştirmesi yeterli oluyor; ancak sonraki adımda \texttt{ori} buyruğunu eklediğimizde AMB'ye yeni bir işlem (VEYA) daha katılacak, dolayısıyla \texttt{AMBİşlem} sinyali tek bitten genişleyecek. Buyruk kümesi büyüdükçe AMB'nin desteklemesi gereken işlem sayısı, iç karmaşıklığı ve kapladığı alan da büyür. Bunun yanında veri yoluna ilk kez \emph{denetim} kavramı dahil edildi. Bundan sonra her yeni buyruk hem yeni donanım hem de yeni denetim sinyalleri getirecek.

\subsection{\texttt{ori} Buyruğunu Eklemek: Anlık Değer, İlk Çoklayıcı ve AMB'ye Yeni İşlem}
\label{subsec:vy-ori}

\texttt{ori} buyruğunu eklemek için veri yoluna yeni bileşenler katmamız ve AMB'nin işlem repertuarını genişletmemiz gerekir. Şekil~\ref{fig:vy-ori}'te bu adımdan sonra elde edilen veri yolu gösterilmektedir. Aşağıda her bir eklemeyi sırayla ele alıyoruz.

\begin{figure}[htbp]
\centering
\begin{tikzpicture}[>=Stealth, semithick, font=\footnotesize,
    yzm/.style={draw, fill=yellow!25, minimum width=1cm, minimum height=1cm, align=center},
    bllk/.style={draw, fill=green!20, minimum width=1.8cm, minimum height=1.6cm, align=center},
    yobk/.style={draw, fill=blue!10, minimum width=2.6cm, minimum height=3cm},
    amb/.style={draw, fill=red!20, minimum width=1.6cm, minimum height=2cm, align=center},
    artop/.style={draw, fill=gray!15, minimum width=1.1cm, minimum height=0.7cm, align=center},
    anlikb/.style={draw, fill=orange!25, minimum width=1.4cm, minimum height=0.9cm, align=center},
    muxdik/.style={draw, fill=purple!15, minimum width=0.7cm, minimum height=1.0cm, align=center},
    bitayar/.style={fill=white, inner sep=1pt, font=\tiny},
    yeni/.style={draw=red!70!black, very thick}
]
  % --- Sol dikey istif: Top., PS, BB (yukarı kaydı +1.5) ---
  \node[artop] (top4) at (0, 1.0) {Top.};
  \node[yzm] (ps) at (0, 2.4) {PS};
  \node[bllk] (bb) at (0, 4.6) {Buyruk\\Belleği};

  % 4 sabiti Top.east'ten sağdan (tek satır "4", tel hizasında; ok uzunluğu önceki ölçü)
  \node[font=\scriptsize, anchor=west] at ([xshift=1.05cm]top4.east) {$4$};
  \draw[->] ([xshift=0.95cm]top4.east) -- (top4.east) node[bitayar, pos=0.5] {32 bit};

  % PS değeri Top'a (PS+4 hesabı için)
  \draw[->] (ps.south) -- (top4.north) node[midway, right=1pt, font=\scriptsize] {PS};

  % Top.west çıkışı -> sola dirsekli yukarı -> PS.west (PS+4 döngüsü); PS+4 etiketi telin ALTINDA
  \draw[->] (top4.west) -- ++(-0.8,0) |- (ps.west) node[bitayar, pos=0.6] {32 bit};
  % PS+4 etiketi yatay segmentin BELİRGİN altına (yatay tel y=1.0'ın 0.35cm altında)
  \node[font=\scriptsize] at (-0.95, 0.65) {PS$+$4};

  % PS.north -> BB.south (adres dikey yukarı)
  \draw[->] (ps.north) -- (bb.south) node[midway, right=1pt, font=\scriptsize] {adres};

  % --- Buyruk yolu ve şerit (parse BB hizasında y=4.6, şerit yukarıda) ---
  \coordinate (parse) at (2.5, 4.6);
  \draw[->] (bb.east) -- (parse) node[midway, below, font=\scriptsize] {buyruk} node[bitayar, pos=0.55] {32 bit};

  \coordinate (seritsol) at (3.0, 6.4);
  \draw[gray, dashed, semithick] (parse) -- (seritsol);
  \draw[gray, dashed, semithick] (parse) -- ([xshift=6cm]seritsol);
  \begin{scope}[shift={(seritsol)}]
    \draw[fill=orange!25, yeni] (0,0) rectangle (2.25,0.55) node[pos=.5, font=\tiny] {imm[11:0]};
    \draw[fill=green!18] (2.25,0) rectangle (3.1875,0.55) node[pos=.5, font=\tiny] {rs1};
    \draw[fill=gray!10] (3.1875,0) rectangle (3.75,0.55) node[pos=.5, font=\tiny] {f3};
    \draw[fill=blue!18] (3.75,0) rectangle (4.6875,0.55) node[pos=.5, font=\tiny] {rd};
    \draw[fill=gray!10] (4.6875,0) rectangle (6,0.55) node[pos=.5, font=\tiny] {opcode};
    \node[font=\tiny, anchor=north] at (0,0) {31};
    \node[font=\tiny, anchor=north] at (2.25,0) {20};
    \node[font=\tiny, anchor=north] at (3.1875,0) {15};
    \node[font=\tiny, anchor=north] at (3.75,0) {12};
    \node[font=\tiny, anchor=north] at (4.6875,0) {7};
    \node[font=\tiny, anchor=north] at (6,0) {0};
    \node[font=\scriptsize, anchor=south] at (3,0.55) {I türü buyruk};
  \end{scope}

  % --- Yazmaç Öbeği (sağ blok da yukarı kaydı, şekil dikey kısalsın) ---
  \node[yobk] (yo) at (5.4, 2.6) {};
  \node[anchor=center, font=\footnotesize\bfseries, align=center] at (yo.center) {Yazmaç\\Öbeği};

  % I türü: rs2 yok; sadece rs1 -> A_adr ve rd -> C_adr
  \draw[->] ([xshift=2.72cm]seritsol) -- ++(0,-1.2) -| ([xshift=-0.6cm]yo.north) node[bitayar, pos=0.75] {5 bit};  % rs1 -> A_adr
  \draw[->] ([xshift=4.22cm]seritsol) -- ++(0,-0.4) -| ([xshift=0.6cm]yo.north) node[bitayar, pos=0.75] {5 bit};   % rd -> C_adr

  % Port etiketleri
  \node[font=\scriptsize, anchor=north] at ([xshift=-0.6cm,yshift=-0.05cm]yo.north) {$A_{\text{adr}}$};
  \node[font=\scriptsize, anchor=north] at ([xshift=0cm,yshift=-0.05cm]yo.north) {$B_{\text{adr}}$};
  \node[font=\scriptsize, anchor=north] at ([xshift=0.6cm,yshift=-0.05cm]yo.north) {$C_{\text{adr}}$};
  \node[font=\scriptsize, anchor=east] at ([xshift=-0.1cm,yshift=0.5cm]yo.east) {$A_{\text{veri}}$};
  \node[font=\scriptsize, anchor=east] at ([xshift=-0.1cm,yshift=-0.25cm]yo.east) {$B_{\text{veri}}$};
  \node[font=\scriptsize, anchor=west] at ([xshift=0.1cm,yshift=-1.0cm]yo.west) {$C_{\text{veri}}$};

  % --- AMBKaynak çoklayıcısı (yo ile birlikte yukarı kaydı) ---
  \node[muxdik, yeni] (mux1) at (8.0, 2.1) {};
  \node[font=\tiny] at (mux1.center) {Çok.};
  \node[font=\tiny, anchor=west] at ([xshift=0.02cm, yshift=0.25cm]mux1.west) {0};
  \node[font=\tiny, anchor=west] at ([xshift=0.02cm, yshift=-0.25cm]mux1.west) {1};

  % --- AMB (yo ile birlikte yukarı kaydı) ---
  \node[amb, yeni, font=\scriptsize] (ab) at (9.6, 2.6) {AMB};

  % A_veri -> AMB üst giriş (düz yatay)
  \draw[->] ([yshift=0.5cm]yo.east) -- ([yshift=0.5cm]ab.west) node[bitayar, pos=0.5] {32 bit};

  % B_veri -> mux1 üst giriş (B_veri y=1.85, mux üst giriş y=1.85, düz yatay)
  \draw[->] ([yshift=-0.25cm]yo.east) -- ([yshift=0.25cm]mux1.west) node[bitayar, pos=0.45] {32 bit};

  % mux1 çıkışı -> AMB B alt giriş (DÜZ YATAY, 32 bit yazısı kaldırıldı)
  \draw[->] (mux1.east) -- ([yshift=-0.5cm]ab.west);

  % --- Anlık değer üreteci (yo ile birlikte yukarı kaydı, yo.x ile aynı) ---
  \node[anlikb, yeni] (anlik) at (5.4, -0.2) {Anlık\\Üreteci};

  % imm[11:0] -> Yazmaç Öbeği'nin SOLUNDAN inerek anlik.west'e
  \draw[->] ([xshift=1.125cm]seritsol) -- ++(0,-1.0) -- ++(-0.6,0) |- (anlik.west) node[bitayar, pos=0.75] {12 bit};

  % anlik çıkışı -> mux1 alt giriş; "32 bit" sağa kaydı
  \draw[->] (anlik.east) -- ++(0.75,0) node[bitayar, pos=0.55] {32 bit} |- ([yshift=-0.25cm]mux1.west);

  % AMB çıkışı -> yazmaç öbeğine geri (Cveri); anlik altından geçsin diye aşağı kaydı
  \draw[->] (ab.east) -- ++(1.6,0)
            node[pos=0.5, above, font=\scriptsize] {sonuç}
            node[bitayar, pos=0.5] {32 bit}
            |- ++(0,-3.7)
            -| ([xshift=-1.0cm,yshift=-1.0cm]yo.west)
            -- ([yshift=-1.0cm]yo.west);

  % AMBİşlem denetim sinyali (üstten, 2 bit) — DEĞİŞTİ: 1 bit -> 2 bit (VEYA eklendi), kırmızı çerçeve
  \node[font=\scriptsize, text=sinyalkirmizi, anchor=south, align=center] at ([yshift=0.5cm]ab.north) {AMBİşlem\\[-1pt]\tiny 00\,$=$\,Topla,\ 01\,$=$\,Çıkar,\ 10\,$=$\,VEYA};
  \draw[->, sinyalkirmizi, thick] ([yshift=0.4cm]ab.north) -- (ab.north);

  % AMBKaynak denetim sinyali (alttan) — YENİ (ilk tanıtım, kırmızı çerçeve)
  \node[font=\scriptsize, text=sinyalkirmizi, anchor=north] at ([yshift=-0.4cm]mux1.south) {AMBKaynak};
  \draw[->, sinyalkirmizi, thick] ([yshift=-0.3cm]mux1.south) -- (mux1.south);

  % Sağ üst köşe: uygulanmış buyruklar (şerit sağında, sonuç teli üstündeki boş alanda)
  \node[draw, rounded corners, fill=blue!5, font=\tiny, align=center, inner sep=4pt, anchor=north east] at (12.0, 7.2) {Buyruk durumu (3/7)\\[2pt]\begin{tabular}{@{}l@{\quad}l@{}}
1.~{\color{green!60!black}\checkmark}~\texttt{add} & 5.~{\color{gray}$\circ$}~\texttt{sw}\\
2.~{\color{green!60!black}\checkmark}~\texttt{sub} & 6.~{\color{gray}$\circ$}~\texttt{beq}/\texttt{bne}\\
3.~{\color{green!60!black}\checkmark}~\texttt{ori} & 7.~{\color{gray}$\circ$}~\texttt{jal}\\
4.~{\color{gray}$\circ$}~\texttt{lw} & \\
\end{tabular}};
\end{tikzpicture}
\caption[\texttt{add}, \texttt{sub} ve \texttt{ori} buyruklarını yürüten veri yolu]{\texttt{add}, \texttt{sub} ve \texttt{ori} buyruklarını yürüten veri yolu. Önceki aşamadan üç değişiklik: (1) buyruk şeridi I türünü gösterir; \texttt{imm[11:0]} alanı \texttt{funct7}+\texttt{rs2}'nin yerini alır. (2) AMB'nin B yolu önüne \textbf{2$\times$1 AMBKaynak çoklayıcısı} eklendi; yazmaç öbeğinden gelen $B_{\text{veri}}$ ile yeni eklenen \textbf{Anlık değer üreteci} çıkışı arasında seçim yapar. (3) AMB içine bit düzeyinde VEYA kapıları eklendi; \texttt{AMBİşlem} sinyali artık 2 bit.}
\label{fig:vy-ori}
\end{figure}

\texttt{ori x5, x6, 100} buyruğu \texttt{x6}'nın değerini 100 anlık değeriyle bit düzeyinde VEYA'lar ve sonucu \texttt{x5}'e yazar. Bu buyrukta AMB'nin ikinci girişi yazmaç öbeğinden değil, buyruğun içindeki anlık değer alanından gelir. RISC-V buyruk biçiminde bu alana ancak 12 bit (bazı buyruklar için 20 bit) sığabilmektedir. Oysa yazmaçlar ve AMB 32 bit veriyle çalışır. Bu nedenle anlık değer AMB'ye verilmeden önce 32 bite genişletilmelidir. Bunun için veri yoluna \textbf{Anlık değer üreteci} adlı bir bileşen eklenir (Ek~\ref{ch:sayi-sistemleri}). Üreteç çıkışı AMB'nin ikinci girişine bağlanır.

AMB'nin ikinci girişi artık iki kaynaktan birinden gelmelidir: \texttt{add}/\texttt{sub} buyruklarında yazmaç öbeğinden, \texttt{ori}'de ise anlık değer üretecinden. Bu seçimi yapacak bileşen \textbf{\texttt{AMBKaynak} çoklayıcısıdır}; iki girişinden birini denetim sinyaliyle belirlenen seçime göre çıkışına aktarır.

İkinci temel fark, AMB'nin yapacağı işlemin toplama/çıkarma değil VEYA olmasıdır. AMB içine bit düzeyinde VEYA gerçekleştiren mantık kapıları eklenir. Artık AMB üç ayrı işlemi destekler (toplama, çıkarma, VEYA); bunlar arasında seçim yapacak \texttt{AMBİşlem} denetim sinyali tek bitten 2 bite çıkar.

İşaretle genişletme: RISC-V'de tüm I türü buyrukların 12 bitlik anlık değeri 32 bite işaretle genişletilir. \texttt{ori} için bunun bariz bir sonucu vardır. \texttt{ori x5, x6, -1} (anlık değer 0xFFF, ikiye tümleyende $-1$) yazıldığında üreteç çıkışı 0xFFFFFFFF olur. AMB'de \texttt{x6 | 0xFFFFFFFF} hesaplanır ve \texttt{x5}'in 32 bitinin tümü birlenir. Aynı kural \texttt{addi}'deki gibi aritmetik buyruklarda da negatif anlık değerleri doğru işlemek için zorunludur.

\begin{karsilastirma}[Anlık Değer Genişletme: Farklı Yaklaşımlar]
RISC-V tüm I türü buyruklarda işaretle genişletme uygular. Tek bir anlık değer üreteci bütün buyruk türlerini destekler ve denetim mantığı sadeleşir. MIPS farklı bir yol izler. Aritmetik buyruklar (\texttt{addi}, \texttt{slti}) işaretle, mantıksal buyruklar (\texttt{andi}, \texttt{ori}, \texttt{xori}) ise sıfırla genişletilir. Gerekçe: mantıksal işlemde negatif anlık değer doğal değildir. Üst bitleri 0 yapmak ``maske'' düşüncesiyle tutarlıdır. Bedeli, anlık değer üretecinin buyruk türüne göre farklı davranması ve denetim biriminin ek karar üretmesidir. ARM (A64) ve x86-64'te anlık değer kodlaması daha karmaşıktır. ARM 12 bitlik anlık değeri bir döngüsel kaydırma alanıyla birlikte saklar, x86 ise anlık değer uzunluğunu buyruğa göre değiştirir (8/16/32 bit). Az bit ile çok değer ifade etmek mümkün olur ama anlık değer çözücü büyür. Anlık değer politikası, bir ISA'nın ifade gücü ile donanım basitliği arasındaki ödünleşimi doğrudan yansıtır. RISC-V basitlik tarafını seçmiştir.
\end{karsilastirma}

\subsection{\texttt{lw} Buyruğunu Eklemek: Veri Belleği ve Geri Yazma Çoklayıcısı}
\label{subsec:vy-lw}

\texttt{lw} buyruğu, bir bellek sözcüğünü yazmaç öbeğine yükler. Buna kadar tasarladığımız veri yolunda buyruk belleği vardı ama \emph{veri belleği} yoktu; çünkü önceki buyrukların hiçbiri yazmaç öbeği dışında bir saklama birimine erişmiyordu. Yükleme/saklama buyrukları, yazmaç öbeğinin çok küçük olması (32 yazmaç) nedeniyle veri için ayrı bir saklama alanı kullanırlar. Bu alan \textbf{veri belleğidir}.

Veri belleği büyük bir saklama birimidir (gerçek sistemlerde DRAM kullanılır) ve üç bağlantısı vardır: \textbf{adres} girişi (32 bit), \textbf{okunan veri} çıkışı (32 bit) ve \textbf{yazılacak veri} girişi (32 bit). Bu bölümde tek çevrimlik bir işlemci tasarladığımız için bellek erişiminin de tek çevrimde tamamlandığını varsayıyoruz. Adres girişine verilen konumun içeriği okunan veri çıkışında belirir. Bunun için ayrı bir denetim sinyali gerekmez. Yazma yönü ise yalnızca \texttt{BelleğeYaz} denetim sinyali etkin olduğunda gerçekleşir ve bu sinyal \texttt{sw} buyruğu için bir sonraki adımda eklenecektir; \texttt{lw} için yalnızca okuma yolu kullanılır. Bellek hiyerarşisi ve DRAM erişim hızları Bölüm~\ref{chap:bellek}'de ele alınacaktır.

\texttt{lw x5, 100(x6)} buyruğu, \texttt{x6 + 100} adresindeki bellek sözcüğünü \texttt{x5}'e yükler:
\begin{enumerate}
    \item \texttt{x6} değerini yazmaç öbeğinden oku.
    \item 100 anlık değerini işaretle genişlet.
    \item AMB'de toplayarak bellek adresini hesapla.
    \item Bu adresi veri belleğine ver. Bellekten dönen değeri al.
    \item Dönen değeri yazmaç öbeğindeki \texttt{x5}'e yaz.
\end{enumerate}

İlk üç adım \texttt{ori} için kurduğumuz yapıda zaten mevcuttur (adres hesabı, anlık + yazmaç toplaması). Bu adımda veri yoluna iki yeni bileşen katılır: yukarıda tanıtılan \textbf{veri belleği} ile \textbf{\texttt{GeriYazmaKaynağı} çoklayıcısı}. Aritmetik buyruklarda (\texttt{add}, \texttt{sub}, \texttt{ori}) yazmaç öbeğine yazılan değer AMB sonucundan gelirken \texttt{lw}'de bu değer bellekten okunandır. \texttt{GeriYazmaKaynağı} çoklayıcısı bu iki kaynak arasından birini denetim sinyaline göre seçer ve yazmaç öbeğinin yazma girişine yönlendirir. Yeni yapı Şekil~\ref{fig:vy-lw}'da gösterilmektedir.

\begin{figure}[htbp]
\centering
\begin{tikzpicture}[>=Stealth, semithick, font=\footnotesize,
    yzm/.style={draw, fill=yellow!25, minimum width=1cm, minimum height=1cm, align=center},
    bllk/.style={draw, fill=green!20, minimum width=1.8cm, minimum height=1.6cm, align=center},
    yobk/.style={draw, fill=blue!10, minimum width=2.6cm, minimum height=3cm},
    amb/.style={draw, fill=red!20, minimum width=1.6cm, minimum height=2cm, align=center},
    artop/.style={draw, fill=gray!15, minimum width=1.1cm, minimum height=0.7cm, align=center},
    anlikb/.style={draw, fill=orange!25, minimum width=1.4cm, minimum height=0.9cm, align=center},
    muxdik/.style={draw, fill=purple!15, minimum width=0.7cm, minimum height=1.0cm, align=center},
    muxyatay/.style={draw, fill=purple!15, minimum width=1.0cm, minimum height=0.7cm, align=center},
    vbll/.style={draw, fill=green!20, minimum width=1.8cm, minimum height=2.2cm, align=center},
    bitayar/.style={fill=white, inner sep=1pt, font=\tiny},
    yeni/.style={draw=red!70!black, very thick}
]
  % --- Sol dikey istif (Şekil 6.5 ile aynı) ---
  \node[artop] (top4) at (0, 1.0) {Top.};
  \node[yzm] (ps) at (0, 2.4) {PS};
  \node[bllk] (bb) at (0, 4.6) {Buyruk\\Belleği};

  \node[font=\scriptsize, anchor=west] at ([xshift=1.05cm]top4.east) {$4$};
  \draw[->] ([xshift=0.95cm]top4.east) -- (top4.east) node[bitayar, pos=0.5] {32 bit};

  \draw[->] (ps.south) -- (top4.north) node[midway, right=1pt, font=\scriptsize] {PS};

  \draw[->] (top4.west) -- ++(-0.8,0) |- (ps.west) node[bitayar, pos=0.6] {32 bit};
  \node[font=\scriptsize] at (-0.95, 0.65) {PS$+$4};

  \draw[->] (ps.north) -- (bb.south) node[midway, right=1pt, font=\scriptsize] {adres};

  \coordinate (parse) at (2.5, 4.6);
  \draw[->] (bb.east) -- (parse) node[midway, below, font=\scriptsize] {buyruk} node[bitayar, pos=0.55] {32 bit};

  \coordinate (seritsol) at (3.0, 6.4);
  \draw[gray, dashed, semithick] (parse) -- (seritsol);
  \draw[gray, dashed, semithick] (parse) -- ([xshift=6cm]seritsol);
  \begin{scope}[shift={(seritsol)}]
    \draw[fill=orange!25] (0,0) rectangle (2.25,0.55) node[pos=.5, font=\tiny] {imm[11:0]};
    \draw[fill=green!18] (2.25,0) rectangle (3.1875,0.55) node[pos=.5, font=\tiny] {rs1};
    \draw[fill=gray!10] (3.1875,0) rectangle (3.75,0.55) node[pos=.5, font=\tiny] {f3};
    \draw[fill=blue!18] (3.75,0) rectangle (4.6875,0.55) node[pos=.5, font=\tiny] {rd};
    \draw[fill=gray!10] (4.6875,0) rectangle (6,0.55) node[pos=.5, font=\tiny] {opcode};
    \node[font=\tiny, anchor=north] at (0,0) {31};
    \node[font=\tiny, anchor=north] at (2.25,0) {20};
    \node[font=\tiny, anchor=north] at (3.1875,0) {15};
    \node[font=\tiny, anchor=north] at (3.75,0) {12};
    \node[font=\tiny, anchor=north] at (4.6875,0) {7};
    \node[font=\tiny, anchor=north] at (6,0) {0};
    \node[font=\scriptsize, anchor=south] at (3,0.55) {I türü buyruk};
  \end{scope}

  \node[yobk] (yo) at (5.4, 2.6) {};
  \node[anchor=center, font=\footnotesize\bfseries, align=center] at (yo.center) {Yazmaç\\Öbeği};

  \draw[->] ([xshift=2.72cm]seritsol) -- ++(0,-1.2) -| ([xshift=-0.6cm]yo.north) node[bitayar, pos=0.75] {5 bit};
  \draw[->] ([xshift=4.22cm]seritsol) -- ++(0,-0.4) -| ([xshift=0.6cm]yo.north) node[bitayar, pos=0.75] {5 bit};

  \node[font=\scriptsize, anchor=north] at ([xshift=-0.6cm,yshift=-0.05cm]yo.north) {$A_{\text{adr}}$};
  \node[font=\scriptsize, anchor=north] at ([xshift=0cm,yshift=-0.05cm]yo.north) {$B_{\text{adr}}$};
  \node[font=\scriptsize, anchor=north] at ([xshift=0.6cm,yshift=-0.05cm]yo.north) {$C_{\text{adr}}$};
  \node[font=\scriptsize, anchor=east] at ([xshift=-0.1cm,yshift=0.5cm]yo.east) {$A_{\text{veri}}$};
  \node[font=\scriptsize, anchor=east] at ([xshift=-0.1cm,yshift=-0.25cm]yo.east) {$B_{\text{veri}}$};
  \node[font=\scriptsize, anchor=west] at ([xshift=0.1cm,yshift=-1.0cm]yo.west) {$C_{\text{veri}}$};

  % --- AMBKaynak çoklayıcısı ---
  \node[muxdik] (mux1) at (8.0, 2.1) {};
  \node[font=\tiny] at (mux1.center) {Çok.};
  \node[font=\tiny, anchor=west] at ([xshift=0.02cm, yshift=0.25cm]mux1.west) {0};
  \node[font=\tiny, anchor=west] at ([xshift=0.02cm, yshift=-0.25cm]mux1.west) {1};

  % --- AMB ---
  \node[amb] (ab) at (9.6, 2.6) {AMB};

  \draw[->] ([yshift=0.5cm]yo.east) -- ([yshift=0.5cm]ab.west) node[bitayar, pos=0.5] {32 bit};
  \draw[->] ([yshift=-0.25cm]yo.east) -- ([yshift=0.25cm]mux1.west) node[bitayar, pos=0.45] {32 bit};
  \draw[->] (mux1.east) -- ([yshift=-0.5cm]ab.west);

  % --- Veri Belleği (YENİ, kırmızı çerçeve; 0.75 birim sola kaydı) ---
  \node[vbll, yeni] (vb) at (11.65, 0.5) {Veri\\Belleği};
  % Port etiketleri italic (yazmaç öbeği konvansiyonu): A_adr ve V_oku üst kenarda simetrik, V_yaz sol-alt
  \node[font=\scriptsize, anchor=north] at ([xshift=-0.5cm, yshift=-0.05cm]vb.north) {$A_{\text{adr}}$};
  \node[font=\scriptsize, anchor=north] at ([xshift=0.5cm, yshift=-0.05cm]vb.north) {$V_{\text{oku}}$};
  \node[font=\scriptsize, anchor=south west] at ([xshift=0.05cm, yshift=0.05cm]vb.south west) {$V_{\text{yaz}}$};
  % AMB sonucu dallanma noktası: A_adr port'un TAM ÜSTÜNDE, AMB hizası y=2.6
  \coordinate (ambsplit) at (11.15, 2.6);
  \node[circle, fill=black, inner sep=1pt] at (ambsplit) {};
  % Dallanma: ambsplit -> A_adr port (dikey aşağı)
  \draw[->] (ambsplit) -- ([xshift=-0.5cm]vb.north);

  % --- GeriYazmaKaynağı çoklayıcısı (DİKEY; 0.75 birim sola, üst giriş y=2.6) ---
  \node[muxdik, yeni] (mux2) at (13.25, 2.35) {};
  \node[font=\tiny] at (mux2.center) {Çok.};
  % Etiketler: 0 üst (vb, lw), 1 alt (AMB sonucu)
  \node[font=\tiny, anchor=west] at ([xshift=0.02cm, yshift=0.25cm]mux2.west) {0};
  \node[font=\tiny, anchor=west] at ([xshift=0.02cm, yshift=-0.25cm]mux2.west) {1};

  % V_oku (vb üst kenar sağ tarafı) -> mux2 ALT giriş (1=V_oku): dikey yukarı, yatay sağa
  \draw[->] ([xshift=0.5cm]vb.north) |- ([yshift=-0.25cm]mux2.west) node[bitayar, pos=0.6] {32 bit};

  % AMB çıkışı tek yatay tel: ab.east -> mux2 üst giriş (y=2.6 düz yatay); "32 bit" AMB-ambsplit arası orta
  \draw[->] (ab.east) -- ([yshift=0.25cm]mux2.west) node[bitayar, pos=0.15] {32 bit};

  % mux2 çıkışı (sağdan) -> Cveri yolu
  \draw[->] (mux2.east) -- ++(0.5,0)
            |- ++(0,-3.7)
            -| ([xshift=-1.0cm,yshift=-1.0cm]yo.west)
            -- ([yshift=-1.0cm]yo.west);
  % "32 bit" mux2 çıkışının sağında (ayrı node, kesin konum)
  \node[bitayar, anchor=west] at ([xshift=0.15cm]mux2.east) {32 bit};

  % --- Anlık değer üreteci ---
  \node[anlikb] (anlik) at (5.4, -0.2) {Anlık\\Üreteci};

  \draw[->] ([xshift=1.125cm]seritsol) -- ++(0,-1.0) -- ++(-0.6,0) |- (anlik.west) node[bitayar, pos=0.75] {12 bit};

  \draw[->] (anlik.east) -- ++(0.75,0) node[bitayar, pos=0.55] {32 bit} |- ([yshift=-0.25cm]mux1.west);

  % AMBİşlem denetim sinyali (üstten, 2 bit)
  \node[font=\scriptsize, text=sinyallacivert, anchor=south, align=center] at ([yshift=0.5cm]ab.north) {AMBİşlem\\[-1pt]\tiny 00\,$=$\,Topla,\ 01\,$=$\,Çıkar,\ 10\,$=$\,VEYA};
  \draw[->, sinyallacivert, thick] ([yshift=0.4cm]ab.north) -- (ab.north);

  % AMBKaynak denetim sinyali (alttan)
  \node[font=\scriptsize, text=sinyallacivert, anchor=north] at ([yshift=-0.4cm]mux1.south) {AMBKaynak};
  \draw[->, sinyallacivert, thick] ([yshift=-0.3cm]mux1.south) -- (mux1.south);

  % GeriYazmaKaynağı denetim sinyali (alttan) — YENİ (ilk tanıtım, kırmızı çerçeve)
  \node[font=\scriptsize, text=sinyalkirmizi, anchor=north, align=center] at ([yshift=-0.4cm]mux2.south) {GeriYazma\\Kaynağı};
  \draw[->, sinyalkirmizi, thick] ([yshift=-0.3cm]mux2.south) -- (mux2.south);

  % Uygulanmış buyruklar kutusu (sağ üst köşe)
  \node[draw, rounded corners, fill=blue!5, font=\tiny, align=center, inner sep=4pt, anchor=north east] at (13.0, 7.2) {Buyruk durumu (4/7)\\[2pt]\begin{tabular}{@{}l@{\quad}l@{}}
1.~{\color{green!60!black}\checkmark}~\texttt{add} & 5.~{\color{gray}$\circ$}~\texttt{sw}\\
2.~{\color{green!60!black}\checkmark}~\texttt{sub} & 6.~{\color{gray}$\circ$}~\texttt{beq}/\texttt{bne}\\
3.~{\color{green!60!black}\checkmark}~\texttt{ori} & 7.~{\color{gray}$\circ$}~\texttt{jal}\\
4.~{\color{green!60!black}\checkmark}~\texttt{lw} & \\
\end{tabular}};
\end{tikzpicture}
\caption[\texttt{add}, \texttt{sub}, \texttt{ori} ve \texttt{lw} buyruklarını yürüten veri yolu]{\texttt{add}, \texttt{sub}, \texttt{ori} ve \texttt{lw} buyruklarını yürüten veri yolu. Şekil~\ref{fig:vy-ori}'e göre iki yeni bileşen eklendi (kırmızı çerçeve): \textbf{Veri belleği}, AMB'nin ürettiği adresi alıp 32 bitlik bir sözcüğü okur. \textbf{GeriYazmaKaynağı çoklayıcısı} yazmaç öbeğine yazılacak değerin AMB sonucundan (\texttt{add}, \texttt{sub}, \texttt{ori} için) mı yoksa veri belleğinden gelen değerden (\texttt{lw} için) mi seçileceğini \texttt{GeriYazmaKaynağı} denetim sinyaliyle belirler.}
\label{fig:vy-lw}
\end{figure}

Veri yolu artık dört buyruğu yürütebilir. \texttt{lw} eklemenin tek temel fiziksel maliyeti veri belleği ile bir çoklayıcıdır. Geri kalan her şey daha önceki buyruklardan miras kalmıştır.

\subsection{\texttt{sw} Buyruğunu Eklemek: Belleğe Yazma}
\label{subsec:vy-sw}

\texttt{sw x7, 100(x6)} buyruğu, \texttt{x7}'nin değerini \texttt{x6 + 100} adresine yazar. Adres hesabı \texttt{lw} ile özdeştir. Tek fark yön değişikliğidir:
\begin{enumerate}
    \item \texttt{x6} (baz adres) ve \texttt{x7} (yazılacak değer) değerlerini yazmaç öbeğinden oku.
    \item 100 anlık değerini işaretle genişlet.
    \item AMB'de toplayarak bellek adresini hesapla.
    \item \texttt{x7}'nin değerini veri belleğinin yazma kapısına ver. \texttt{BelleğeYaz} sinyalini etkinleştir.
\end{enumerate}

Şekil~\ref{fig:vy-sw}'de bu adımdan sonra elde edilen veri yolu gösterilmektedir. Aşağıda buyruk biçimini ve eklenen denetim sinyallerini sırayla ele alıyoruz.

\begin{figure}[htbp]
\centering
\begin{tikzpicture}[>=Stealth, semithick, font=\footnotesize,
    yzm/.style={draw, fill=yellow!25, minimum width=1cm, minimum height=1cm, align=center},
    bllk/.style={draw, fill=green!20, minimum width=1.8cm, minimum height=1.6cm, align=center},
    yobk/.style={draw, fill=blue!10, minimum width=2.6cm, minimum height=3cm},
    amb/.style={draw, fill=red!20, minimum width=1.6cm, minimum height=2cm, align=center},
    artop/.style={draw, fill=gray!15, minimum width=1.1cm, minimum height=0.7cm, align=center},
    anlikb/.style={draw, fill=orange!25, minimum width=1.4cm, minimum height=0.9cm, align=center},
    muxdik/.style={draw, fill=purple!15, minimum width=0.7cm, minimum height=1.0cm, align=center},
    muxyatay/.style={draw, fill=purple!15, minimum width=1.0cm, minimum height=0.7cm, align=center},
    vbll/.style={draw, fill=green!20, minimum width=1.8cm, minimum height=2.2cm, align=center},
    bitayar/.style={fill=white, inner sep=1pt, font=\tiny},
    yeni/.style={draw=red!70!black, very thick}
]
  % --- Sol dikey istif (Şekil 6.6 ile aynı) ---
  \node[artop] (top4) at (0, 1.0) {Top.};
  \node[yzm] (ps) at (0, 2.4) {PS};
  \node[bllk] (bb) at (0, 4.6) {Buyruk\\Belleği};

  \node[font=\scriptsize, anchor=west] at ([xshift=1.05cm]top4.east) {$4$};
  \draw[->] ([xshift=0.95cm]top4.east) -- (top4.east) node[bitayar, pos=0.5] {32 bit};

  \draw[->] (ps.south) -- (top4.north) node[midway, right=1pt, font=\scriptsize] {PS};

  \draw[->] (top4.west) -- ++(-0.8,0) |- (ps.west) node[bitayar, pos=0.6] {32 bit};
  \node[font=\scriptsize] at (-0.95, 0.65) {PS$+$4};

  \draw[->] (ps.north) -- (bb.south) node[midway, right=1pt, font=\scriptsize] {adres};

  \coordinate (parse) at (2.5, 4.6);
  \draw[->] (bb.east) -- (parse) node[midway, below, font=\scriptsize] {buyruk} node[bitayar, pos=0.55] {32 bit};

  % S türü buyruk şeridi (DEĞİŞTİ: anlık değer iki parçaya bölünmüş; kırmızı çerçeve)
  \coordinate (seritsol) at (3.0, 6.4);
  \draw[gray, dashed, semithick] (parse) -- (seritsol);
  \draw[gray, dashed, semithick] (parse) -- ([xshift=6cm]seritsol);
  \begin{scope}[shift={(seritsol)}]
    \draw[fill=orange!25, yeni] (0,0) rectangle (1.3125,0.55) node[pos=.5, font=\tiny] {imm[11:5]};
    \draw[fill=green!18] (1.3125,0) rectangle (2.25,0.55) node[pos=.5, font=\tiny] {rs2};
    \draw[fill=green!18] (2.25,0) rectangle (3.1875,0.55) node[pos=.5, font=\tiny] {rs1};
    \draw[fill=gray!10] (3.1875,0) rectangle (3.75,0.55) node[pos=.5, font=\tiny] {f3};
    \draw[fill=orange!25, yeni] (3.75,0) rectangle (4.6875,0.55) node[pos=.5, font=\tiny] {imm[4:0]};
    \draw[fill=gray!10] (4.6875,0) rectangle (6,0.55) node[pos=.5, font=\tiny] {opcode};
    \node[font=\tiny, anchor=north] at (0,0) {31};
    \node[font=\tiny, anchor=north] at (1.3125,0) {25};
    \node[font=\tiny, anchor=north] at (2.25,0) {20};
    \node[font=\tiny, anchor=north] at (3.1875,0) {15};
    \node[font=\tiny, anchor=north] at (3.75,0) {12};
    \node[font=\tiny, anchor=north] at (4.6875,0) {7};
    \node[font=\tiny, anchor=north] at (6,0) {0};
    \node[font=\scriptsize, anchor=south] at (3,0.55) {S türü buyruk};
  \end{scope}

  \node[yobk] (yo) at (5.4, 2.6) {};
  \node[anchor=center, font=\footnotesize\bfseries, align=center] at (yo.center) {Yazmaç\\Öbeği};

  % rs1 -> A_adr, rs2 -> B_adr (S türü için: rs2 yazılacak veri; rd -> C_adr yok)
  \draw[->] ([xshift=2.72cm]seritsol) -- ++(0,-1.2) -| ([xshift=-0.6cm]yo.north) node[bitayar, pos=0.75] {5 bit};
  \draw[->] ([xshift=1.78cm]seritsol) -- ++(0,-0.6) -| ([xshift=0cm]yo.north) node[bitayar, pos=0.75] {5 bit};

  \node[font=\scriptsize, anchor=north] at ([xshift=-0.6cm,yshift=-0.05cm]yo.north) {$A_{\text{adr}}$};
  \node[font=\scriptsize, anchor=north] at ([xshift=0cm,yshift=-0.05cm]yo.north) {$B_{\text{adr}}$};
  \node[font=\scriptsize, anchor=north] at ([xshift=0.6cm,yshift=-0.05cm]yo.north) {$C_{\text{adr}}$};
  \node[font=\scriptsize, anchor=east] at ([xshift=-0.1cm,yshift=0.5cm]yo.east) {$A_{\text{veri}}$};
  \node[font=\scriptsize, anchor=east] at ([xshift=-0.1cm,yshift=-0.25cm]yo.east) {$B_{\text{veri}}$};
  \node[font=\scriptsize, anchor=west] at ([xshift=0.1cm,yshift=-1.0cm]yo.west) {$C_{\text{veri}}$};

  % YENİ — YazmacaYaz denetim sinyali (sw için 0; kırmızı çerçeve)
  \node[font=\scriptsize, text=sinyalkirmizi, anchor=south] at ([xshift=1.05cm, yshift=0.5cm]yo.north) {YazmacaYaz};
  \draw[->, sinyalkirmizi, thick] ([xshift=1.05cm, yshift=0.4cm]yo.north) -- ([xshift=1.05cm]yo.north);

  % --- AMBKaynak çoklayıcısı ---
  \node[muxdik] (mux1) at (8.0, 2.1) {};
  \node[font=\tiny] at (mux1.center) {Çok.};
  \node[font=\tiny, anchor=west] at ([xshift=0.02cm, yshift=0.25cm]mux1.west) {0};
  \node[font=\tiny, anchor=west] at ([xshift=0.02cm, yshift=-0.25cm]mux1.west) {1};

  % --- AMB ---
  \node[amb] (ab) at (9.6, 2.6) {AMB};

  \draw[->] ([yshift=0.5cm]yo.east) -- ([yshift=0.5cm]ab.west) node[bitayar, pos=0.5] {32 bit};
  % B_veri -> mux1 (junction'lı: yeni tel veri belleğine de gidecek); "32 bit" etiketi kaldırıldı (B_veri çıkışı tek teldir)
  \draw[->] ([yshift=-0.25cm]yo.east) -- ([yshift=0.25cm]mux1.west);
  \draw[->] (mux1.east) -- ([yshift=-0.5cm]ab.west);

  % --- Veri Belleği ---
  \node[vbll] (vb) at (11.65, 0.5) {Veri\\Belleği};
  \node[font=\scriptsize, anchor=north] at ([xshift=-0.5cm, yshift=-0.05cm]vb.north) {$A_{\text{adr}}$};
  \node[font=\scriptsize, anchor=north] at ([xshift=0.5cm, yshift=-0.05cm]vb.north) {$V_{\text{oku}}$};
  % V_yaz: vb'nin sol kenarının alt kısmında (giriş sol taraftan)
  \node[font=\scriptsize, anchor=west] at ([xshift=0.1cm, yshift=-0.85cm]vb.west) {$V_{\text{yaz}}$};

  % AMB sonucu dallanma noktası: A_adr port'un TAM ÜSTÜNDE
  \coordinate (ambsplit) at (11.15, 2.6);
  \node[circle, fill=black, inner sep=1pt] at (ambsplit) {};
  \draw[->] (ambsplit) -- ([xshift=-0.5cm]vb.north);

  % YENİ — B_veri -> V_yaz teli (kırmızı); sol taraftan (vb.west) giriş
  \coordinate (bverijnct) at (7.0, 2.35);
  \node[circle, fill=black, inner sep=1pt] at (bverijnct) {};
  \draw[->, yeni] (bverijnct) -- (7.0, -0.35) -- ([yshift=-0.85cm]vb.west) node[bitayar, pos=0.55] {32 bit};

  % YENİ — BelleğeYaz denetim sinyali (sw için 1); vb'nin altından gelir; kırmızı çerçeve
  \node[font=\scriptsize, text=sinyalkirmizi, anchor=north] at ([yshift=-0.5cm]vb.south) {BelleğeYaz};
  \draw[->, sinyalkirmizi, thick] ([yshift=-0.4cm]vb.south) -- (vb.south);

  % --- GeriYazmaKaynağı çoklayıcısı (Şekil 6.6 ile aynı) ---
  \node[muxdik] (mux2) at (13.25, 2.35) {};
  \node[font=\tiny] at (mux2.center) {Çok.};
  \node[font=\tiny, anchor=west] at ([xshift=0.02cm, yshift=0.25cm]mux2.west) {0};
  \node[font=\tiny, anchor=west] at ([xshift=0.02cm, yshift=-0.25cm]mux2.west) {1};

  % V_oku -> mux2 ALT giriş (1=V_oku)
  \draw[->] ([xshift=0.5cm]vb.north) |- ([yshift=-0.25cm]mux2.west) node[bitayar, pos=0.6] {32 bit};

  % AMB çıkışı -> mux2 üst giriş (y=2.6 düz yatay)
  \draw[->] (ab.east) -- ([yshift=0.25cm]mux2.west) node[bitayar, pos=0.15] {32 bit};

  % mux2 çıkışı -> Cveri yolu (alt yatay segmenti aşağıya kaydırıldı: BelleğeYaz sinyaline yer açmak için)
  \draw[->] (mux2.east) -- ++(0.5,0)
            |- ++(0,-4.0)
            -| ([xshift=-1.0cm,yshift=-1.0cm]yo.west)
            -- ([yshift=-1.0cm]yo.west);
  \node[bitayar, anchor=west] at ([xshift=0.15cm]mux2.east) {32 bit};

  % --- Anlık değer üreteci ---
  \node[anlikb] (anlik) at (5.4, -0.2) {Anlık\\Üreteci};
  % S türü: imm[11:5] ile imm[4:0] anlık üretecinde birleşir; tek temsil teli gösterilmiştir
  \draw[->] ([xshift=1.125cm]seritsol) -- ++(0,-1.0) -- ++(-0.6,0) |- (anlik.west) node[bitayar, pos=0.75] {12 bit};
  \draw[->] (anlik.east) -- ++(0.75,0) node[bitayar, pos=0.55] {32 bit} |- ([yshift=-0.25cm]mux1.west);

  % AMBİşlem denetim sinyali (üstten, 2 bit)
  \node[font=\scriptsize, text=sinyallacivert, anchor=south, align=center] at ([yshift=0.5cm]ab.north) {AMBİşlem\\[-1pt]\tiny 00\,$=$\,Topla,\ 01\,$=$\,Çıkar,\ 10\,$=$\,VEYA};
  \draw[->, sinyallacivert, thick] ([yshift=0.4cm]ab.north) -- (ab.north);

  % AMBKaynak denetim sinyali (alttan)
  \node[font=\scriptsize, text=sinyallacivert, anchor=north] at ([yshift=-0.4cm]mux1.south) {AMBKaynak};
  \draw[->, sinyallacivert, thick] ([yshift=-0.3cm]mux1.south) -- (mux1.south);

  % GeriYazmaKaynağı denetim sinyali (alttan)
  \node[font=\scriptsize, text=sinyallacivert, anchor=north, align=center] at ([yshift=-0.4cm]mux2.south) {GeriYazma\\Kaynağı};
  \draw[->, sinyallacivert, thick] ([yshift=-0.3cm]mux2.south) -- (mux2.south);

  % Uygulanmış buyruklar kutusu (sağ üst köşe; 5/7)
  \node[draw, rounded corners, fill=blue!5, font=\tiny, align=center, inner sep=4pt, anchor=north east] at (13.0, 7.2) {Buyruk durumu (5/7)\\[2pt]\begin{tabular}{@{}l@{\quad}l@{}}
1.~{\color{green!60!black}\checkmark}~\texttt{add} & 5.~{\color{green!60!black}\checkmark}~\texttt{sw}\\
2.~{\color{green!60!black}\checkmark}~\texttt{sub} & 6.~{\color{gray}$\circ$}~\texttt{beq}/\texttt{bne}\\
3.~{\color{green!60!black}\checkmark}~\texttt{ori} & 7.~{\color{gray}$\circ$}~\texttt{jal}\\
4.~{\color{green!60!black}\checkmark}~\texttt{lw} & \\
\end{tabular}};
\end{tikzpicture}
\caption[\texttt{add}, \texttt{sub}, \texttt{ori}, \texttt{lw} ve \texttt{sw} buyruklarını yürüten veri yolu]{\texttt{add}, \texttt{sub}, \texttt{ori}, \texttt{lw} ve \texttt{sw} buyruklarını yürüten veri yolu. Şekil~\ref{fig:vy-lw}'ya göre üç şey eklendi (kırmızı): (1) yazmaç öbeğinin $B_{\text{veri}}$ çıkışından veri belleğinin $V_{\text{yaz}}$ kapısına giden yeni tel, (2) \texttt{BelleğeYaz} denetim sinyali (\texttt{sw} için 1; yazma yönünü etkinleştirir), (3) \texttt{YazmacaYaz} denetim sinyali (\texttt{sw} için 0; bu buyrukta yazmaç öbeğine yazma yapılmaz). Buyruk şeridi S türünü göstermektedir; anlık değerin \texttt{imm[11:5]} ve \texttt{imm[4:0]} olarak iki parçaya bölünmüş olması S türünün ayırt edici özelliğidir, iki parça anlık değer üretecinin içinde birleştirilir.}
\label{fig:vy-sw}
\end{figure}

\texttt{sw} S türü bir buyruktur. R ve I türü buyruklarda anlık değer alanı buyruğun yüksek bitlerinde tek parça olarak bulunurken, S türünde \texttt{imm[11:5]} (bit 31--25) ve \texttt{imm[4:0]} (bit 11--7) olarak ikiye bölünmüştür. Bu bölünme rastgele değildir. Amacı \texttt{rs1} (bit 15--19) ve \texttt{rs2} (bit 20--24) alanlarının buyruk türleri arasında aynı bit konumlarında kalmasını sağlamaktır. Böylece yazmaç öbeğinin okuma adres girişleri her buyruk türünde aynı dilimlerden beslenir. Çözücü buyruğun türünü öğrenmeden yazmaç adreslerini doğrudan ilgili bitlerden alabilir. Anlık değerin parçalarının yeniden birleştirilmesi sorumluluğu anlık değer üretecine kalır. Bu birim buyruk türüne göre doğru parçaları seçer, 12 biti kurar ve 32 bite işaretle genişletir. I türü ve S türü arasında geçişi yönetmek için anlık değer üretecine küçük bir denetim girişi (\texttt{AnlıkBiçim}) eklenebilir.

\texttt{sw} buyruğunun veri yoluna asıl katkısı iki yeni denetim sinyalidir.

\texttt{BelleğeYaz}, veri belleğinin yazma kanalını açar. Belleğin okuma yolu sürekli açıktır. Adres girişine bir değer verildiğinde okunan veri çıkışı kendiliğinden güncellenir. Yazma yolu ise yalnızca \texttt{BelleğeYaz}=1 olduğunda etkinleşir. Aksi halde belleğin içeriği saat kenarında değişmez. Bu asimetri kasıtlıdır. Rastgele bir adresin üzerine kazara yazmasın diye yazma her zaman açık bir izin sinyali gerektirir. \texttt{sw} için \texttt{BelleğeYaz}=1, diğer tüm buyruklar için 0'dır.

\texttt{YazmacaYaz}, yazmaç öbeğinin yazma yönünü etkinleştirir. \texttt{sw}'ye kadar tasarladığımız tüm buyruklar (\texttt{add}, \texttt{sub}, \texttt{ori}, \texttt{lw}) sonucu yazmaç öbeğine yazıyordu. Bu yüzden \texttt{YazmacaYaz} pratikte hep 1 değerini alıyor ve şekillerde gizli tutuluyordu. \texttt{sw} bu sinyalin 0 olmasını gerektiren ilk buyruktur, çünkü yazma hedefi yazmaç öbeği değil veri belleğidir. Sinyal fiziksel olarak yazmaç öbeğinin içindeki yazma yetki kapısına bağlıdır. Saat darbesinin yükselen kenarında \texttt{YazmacaYaz}=1 ise $C_{\text{veri}}$ girişindeki değer $C_{\text{adr}}$ ile gösterilen yazmaca yazılır, \texttt{YazmacaYaz}=0 ise yazma engellenir. \texttt{sw} sırasında $C_{\text{adr}}$ girişine hangi değer düşerse düşsün, $C_{\text{veri}}$ hattında ne olursa olsun yazmaç öbeği değişmez.

Donanım açısından bu adımda yeni bir bileşen eklemiyoruz. Veri belleğinin yazma kapısı zaten vardı, $B_{\text{veri}}$ çıkışı da öyle. Tek yeni fiziksel yol, $B_{\text{veri}}$'den veri belleğinin yazma kapısına giden teldir. Bütün yenilik denetim katmanındadır. İki yeni sinyal denetim biriminin üretmesi gereken karar uzayını genişletir. \texttt{BelleğeYaz} ile \texttt{YazmacaYaz} birbirinden bağımsız hareket eder. Bir sonraki bölümde ele alacağımız \texttt{beq}/\texttt{bne} buyrukları da \texttt{YazmacaYaz}=0 ile çalışacak fakat onlarda \texttt{BelleğeYaz}=0 olacaktır.

\subsection{Dallanma Buyrukları (\texttt{beq}, \texttt{bne}): PS Güncellemesi Genişler}
\label{subsec:vy-beq}

\texttt{beq x5, x6, etiket} buyruğu \texttt{x5} ile \texttt{x6} eşitse \texttt{etiket} adresine dallanır. Eşit değilse PS+4 ile devam eder. Bu, PS güncellemesinin ilk genişlemesidir. Şekil~\ref{fig:vy-beq}'de eklenen donanım ve denetim sinyalleri gösterilmektedir. Aşağıda adımları ve yenilikleri sırayla ele alıyoruz.

\begin{figure}[htbp]
\centering
\begin{tikzpicture}[>=Stealth, semithick, font=\footnotesize,
    yzm/.style={draw, fill=yellow!25, minimum width=1cm, minimum height=1cm, align=center},
    bllk/.style={draw, fill=green!20, minimum width=1.8cm, minimum height=1.6cm, align=center},
    yobk/.style={draw, fill=blue!10, minimum width=2.6cm, minimum height=3cm},
    amb/.style={draw, fill=red!20, minimum width=1.6cm, minimum height=2cm, align=center},
    artop/.style={draw, fill=gray!15, minimum width=1.1cm, minimum height=0.7cm, align=center},
    anlikb/.style={draw, fill=orange!25, minimum width=1.4cm, minimum height=0.9cm, align=center},
    muxdik/.style={draw, fill=purple!15, minimum width=0.7cm, minimum height=1.0cm, align=center},
    muxyatay/.style={draw, fill=purple!15, minimum width=1.0cm, minimum height=0.7cm, align=center},
    vbll/.style={draw, fill=green!20, minimum width=1.8cm, minimum height=2.2cm, align=center},
    bitayar/.style={fill=white, inner sep=1pt, font=\tiny},
    yeni/.style={draw=red!70!black, very thick}
]
  % --- Sol PS güncelleme bloğu (DİKEY TASARIM v3): mux PS'in ALTINDA, toplayıcılar mux'un ALTINDA.
  % Toplayıcı girişleri ÜSTTEN, çıkışları ALTTAN. Çıkış teli aşağıdan dirsekle mux.south'a gelir.
  % PS shared: PS.west ve PS.east (PS.south değil; PS.south yalnızca mux çıkışını alır).
  \node[bllk] (bb) at (0, 4.6) {Buyruk\\Belleği};
  \node[yzm] (ps) at (0, 2.7) {PS};
  \node[muxyatay, yeni] (psmux) at (0, 1.3) {};
  \node[font=\tiny] at (psmux.center) {Çok.};
  \node[font=\tiny, anchor=south] at ([xshift=-0.3cm, yshift=0.02cm]psmux.south) {0};
  \node[font=\tiny, anchor=south] at ([xshift=0.3cm, yshift=0.02cm]psmux.south) {1};

  \node[artop] (top4) at (-1.3, 0) {Top.};
  \node[artop, yeni, minimum width=1.2cm, minimum height=0.7cm] (daladd) at (1.3, 0) {Dal.\\Top.};

  % PS -> BB (adres)
  \draw[->] (ps.north) -- (bb.south) node[midway, right=1pt, font=\scriptsize] {adres};

  % psmux çıkışı (kuzey) -> PS güneyi (kısa dikey)
  \draw[->, yeni] (psmux.north) -- (ps.south);

  % PS shared: PS.west -> Top4.north_WEST (dış üst sol giriş; 4 içe alındığı için PS dışa kaydı)
  \draw[->] (ps.west) -- (-1.5, 2.7) -- ([xshift=-0.2cm]top4.north);
  % PS.east -> Dal.Top.north_west (iç üst sol giriş) — kırmızı (Dal.Top için yeni)
  \draw[->, yeni] (ps.east) -- (1.1, 2.7) -- ([xshift=-0.2cm]daladd.north);

  % 4 sabiti -> Top4.north_EAST (iç üst sağ giriş; sağa alındı, Dallan'a yer açmak için)
  \node[font=\scriptsize] at (-1.1, 0.85) {$4$};
  \draw[->] (-1.1, 0.7) -- ([xshift=0.2cm]top4.north);

  % B-anlık -> Dal.Top.north_east (dış üst sağ giriş, anlık.east teli üzerinden tap)
  \coordinate (anlikdal) at (6.85, 1.0);
  \node[circle, fill=black, inner sep=1pt, yeni] at (anlikdal) {};
  \draw[->, yeni] (anlikdal) -- (1.5, 1.0) -- ([xshift=0.2cm]daladd.north);

  % Top4 çıkışı (south=aşağı) -> dirsek aşağı-içe -> mux.south input 0 (alttan giriş)
  \draw[->] (top4.south) -- (-1.3, -0.7) -- (-0.3, -0.7) -- ([xshift=-0.3cm]psmux.south);
  \node[font=\scriptsize, anchor=north] at (-0.8, -0.75) {PS$+$4};
  % Dal.Top çıkışı (south=aşağı) -> dirsek aşağı-içe -> mux.south input 1 — kırmızı
  \draw[->, yeni] (daladd.south) -- (1.3, -0.7) -- (0.3, -0.7) -- ([xshift=0.3cm]psmux.south);
  \node[font=\scriptsize, anchor=north] at (0.8, -0.75) {Dal.hedef};

  % Dallan denetim sinyali (psmux'a SOLDAN, ok yatay; etiket okun üstünde)
  \draw[->, sinyalkirmizi, thick] ([xshift=-0.4cm]psmux.west) -- (psmux.west);
  \node[font=\scriptsize, text=sinyalkirmizi, anchor=south] at ([xshift=-0.5cm, yshift=0.05cm]psmux.west) {Dallan};

  \coordinate (parse) at (2.5, 4.6);
  \draw[->] (bb.east) -- (parse) node[midway, below, font=\scriptsize] {buyruk} node[bitayar, pos=0.55] {32 bit};

  % B türü buyruk şeridi (DEĞİŞTİ: anlık değer dört parçaya dağılmış; kırmızı çerçeve)
  \coordinate (seritsol) at (3.0, 6.4);
  \draw[gray, dashed, semithick] (parse) -- (seritsol);
  \draw[gray, dashed, semithick] (parse) -- ([xshift=6cm]seritsol);
  \begin{scope}[shift={(seritsol)}]
    \draw[fill=orange!25, yeni] (0,0) rectangle (1.3125,0.55) node[pos=.5, font=\tiny] {imm[12,10:5]};
    \draw[fill=green!18] (1.3125,0) rectangle (2.25,0.55) node[pos=.5, font=\tiny] {rs2};
    \draw[fill=green!18] (2.25,0) rectangle (3.1875,0.55) node[pos=.5, font=\tiny] {rs1};
    \draw[fill=gray!10] (3.1875,0) rectangle (3.75,0.55) node[pos=.5, font=\tiny] {f3};
    \draw[fill=orange!25, yeni] (3.75,0) rectangle (4.6875,0.55) node[pos=.5, font=\tiny] {imm[4:1,11]};
    \draw[fill=gray!10] (4.6875,0) rectangle (6,0.55) node[pos=.5, font=\tiny] {opcode};
    \node[font=\tiny, anchor=north] at (0,0) {31};
    \node[font=\tiny, anchor=north] at (1.3125,0) {25};
    \node[font=\tiny, anchor=north] at (2.25,0) {20};
    \node[font=\tiny, anchor=north] at (3.1875,0) {15};
    \node[font=\tiny, anchor=north] at (3.75,0) {12};
    \node[font=\tiny, anchor=north] at (4.6875,0) {7};
    \node[font=\tiny, anchor=north] at (6,0) {0};
    \node[font=\scriptsize, anchor=south] at (3,0.55) {B türü buyruk};
  \end{scope}

  \node[yobk] (yo) at (5.4, 2.6) {};
  \node[anchor=center, font=\footnotesize\bfseries, align=center] at (yo.center) {Yazmaç\\Öbeği};

  % rs1 -> A_adr, rs2 -> B_adr (B türü için her iki yazmaç da okunur)
  \draw[->] ([xshift=2.72cm]seritsol) -- ++(0,-1.2) -| ([xshift=-0.6cm]yo.north) node[bitayar, pos=0.75] {5 bit};
  \draw[->] ([xshift=1.78cm]seritsol) -- ++(0,-0.6) -| ([xshift=0cm]yo.north) node[bitayar, pos=0.75] {5 bit};

  \node[font=\scriptsize, anchor=north] at ([xshift=-0.6cm,yshift=-0.05cm]yo.north) {$A_{\text{adr}}$};
  \node[font=\scriptsize, anchor=north] at ([xshift=0cm,yshift=-0.05cm]yo.north) {$B_{\text{adr}}$};
  \node[font=\scriptsize, anchor=north] at ([xshift=0.6cm,yshift=-0.05cm]yo.north) {$C_{\text{adr}}$};
  \node[font=\scriptsize, anchor=east] at ([xshift=-0.1cm,yshift=0.5cm]yo.east) {$A_{\text{veri}}$};
  \node[font=\scriptsize, anchor=east] at ([xshift=-0.1cm,yshift=-0.25cm]yo.east) {$B_{\text{veri}}$};
  \node[font=\scriptsize, anchor=west] at ([xshift=0.1cm,yshift=-1.0cm]yo.west) {$C_{\text{veri}}$};

  % YazmacaYaz denetim sinyali (beq için 0)
  \node[font=\scriptsize, text=sinyallacivert, anchor=south] at ([xshift=1.05cm, yshift=0.5cm]yo.north) {YazmacaYaz};
  \draw[->, sinyallacivert, thick] ([xshift=1.05cm, yshift=0.4cm]yo.north) -- ([xshift=1.05cm]yo.north);

  % --- AMBKaynak çoklayıcısı ---
  \node[muxdik] (mux1) at (8.0, 2.1) {};
  \node[font=\tiny] at (mux1.center) {Çok.};
  \node[font=\tiny, anchor=west] at ([xshift=0.02cm, yshift=0.25cm]mux1.west) {0};
  \node[font=\tiny, anchor=west] at ([xshift=0.02cm, yshift=-0.25cm]mux1.west) {1};

  % --- AMB ---
  \node[amb] (ab) at (9.6, 2.6) {AMB};

  \draw[->] ([yshift=0.5cm]yo.east) -- ([yshift=0.5cm]ab.west) node[bitayar, pos=0.5] {32 bit};
  \draw[->] ([yshift=-0.25cm]yo.east) -- ([yshift=0.25cm]mux1.west);
  \draw[->] (mux1.east) -- ([yshift=-0.5cm]ab.west);

  % --- Veri Belleği ---
  \node[vbll] (vb) at (11.65, 0.5) {Veri\\Belleği};
  \node[font=\scriptsize, anchor=north] at ([xshift=-0.5cm, yshift=-0.05cm]vb.north) {$A_{\text{adr}}$};
  \node[font=\scriptsize, anchor=north] at ([xshift=0.5cm, yshift=-0.05cm]vb.north) {$V_{\text{oku}}$};
  \node[font=\scriptsize, anchor=west] at ([xshift=0.1cm, yshift=-0.85cm]vb.west) {$V_{\text{yaz}}$};

  % AMB sonucu dallanma noktası -> Vb A_adr
  \coordinate (ambsplit) at (11.15, 2.6);
  \node[circle, fill=black, inner sep=1pt] at (ambsplit) {};
  \draw[->] (ambsplit) -- ([xshift=-0.5cm]vb.north);

  % B_veri -> V_yaz teli
  \coordinate (bverijnct) at (7.0, 2.35);
  \node[circle, fill=black, inner sep=1pt] at (bverijnct) {};
  \draw[->] (bverijnct) -- (7.0, -0.35) -- ([yshift=-0.85cm]vb.west) node[bitayar, pos=0.55] {32 bit};

  % BelleğeYaz denetim sinyali (beq için 0)
  \node[font=\scriptsize, text=sinyallacivert, anchor=north] at ([yshift=-0.5cm]vb.south) {BelleğeYaz};
  \draw[->, sinyallacivert, thick] ([yshift=-0.4cm]vb.south) -- (vb.south);

  % --- GeriYazmaKaynağı çoklayıcısı ---
  \node[muxdik] (mux2) at (13.25, 2.35) {};
  \node[font=\tiny] at (mux2.center) {Çok.};
  \node[font=\tiny, anchor=west] at ([xshift=0.02cm, yshift=0.25cm]mux2.west) {0};
  \node[font=\tiny, anchor=west] at ([xshift=0.02cm, yshift=-0.25cm]mux2.west) {1};

  \draw[->] ([xshift=0.5cm]vb.north) |- ([yshift=-0.25cm]mux2.west) node[bitayar, pos=0.6] {32 bit};
  \draw[->] (ab.east) -- ([yshift=0.25cm]mux2.west) node[bitayar, pos=0.15] {32 bit};

  \draw[->] (mux2.east) -- ++(0.5,0)
            |- ++(0,-4.0)
            -| ([xshift=-1.0cm,yshift=-1.0cm]yo.west)
            -- ([yshift=-1.0cm]yo.west);
  \node[bitayar, anchor=west] at ([xshift=0.15cm]mux2.east) {32 bit};

  % --- Anlık değer üreteci ---
  \node[anlikb] (anlik) at (5.4, -0.2) {Anlık\\Üreteci};
  % B türü: 12 bit (buyruktan fiziksel olarak gelen) anlık değer dört parçadan birleşir; üreteç içinde imm[0]=0 eklenip 13-bit etkin uzaklık olarak işaretle genişletilir
  \draw[->] ([xshift=1.125cm]seritsol) -- ++(0,-1.0) -- ++(-0.6,0) |- (anlik.west) node[bitayar, pos=0.75] {12 bit};
  % Anlık çıkışı: mux1'e gider (B türü beq'te kullanılmasa da donanım kalır); "32 bit" etiketi Şekil 6.7'deki gibi anlık.east hemen sağında
  \draw[->] (anlik.east) -- ++(0.75,0) coordinate (anlikbifurc) -- (6.85, 1.85) -- ([yshift=-0.25cm]mux1.west);
  \node[bitayar] at (6.5, -0.2) {32 bit};
  % YENİ: AnlıkBiçim denetim sinyali (üretecin türe göre davranmasını sağlar; alttan giriyor) — kırmızı çerçeve
  \node[font=\scriptsize, text=sinyalkirmizi, anchor=north] at ([yshift=-0.5cm]anlik.south) {AnlıkBiçim};
  \draw[->, sinyalkirmizi, thick] ([yshift=-0.4cm]anlik.south) -- (anlik.south);

  % AMBİşlem denetim sinyali (üstten, 2 bit)
  \node[font=\scriptsize, text=sinyallacivert, anchor=south, align=center] at ([yshift=0.5cm]ab.north) {AMBİşlem\\[-1pt]\tiny 00\,$=$\,Topla,\ 01\,$=$\,Çıkar,\ 10\,$=$\,VEYA};
  \draw[->, sinyallacivert, thick] ([yshift=0.4cm]ab.north) -- (ab.north);

  % AMBKaynak denetim sinyali
  \node[font=\scriptsize, text=sinyallacivert, anchor=north] at ([yshift=-0.4cm]mux1.south) {AMBKaynak};
  \draw[->, sinyallacivert, thick] ([yshift=-0.3cm]mux1.south) -- (mux1.south);

  % GeriYazmaKaynağı denetim sinyali
  \node[font=\scriptsize, text=sinyallacivert, anchor=north, align=center] at ([yshift=-0.4cm]mux2.south) {GeriYazma\\Kaynağı};
  \draw[->, sinyallacivert, thick] ([yshift=-0.3cm]mux2.south) -- (mux2.south);

  % ==========================================
  % EK YENİ ÖĞE (kırmızı): AMB Sıfır çıkışı
  % (B-anlık -> Dal.Top wire yukarıda PS güncelleme bloğunda tanımlandı)
  % ==========================================

  % AMB Sıfır çıkışı
  \node[font=\scriptsize, anchor=west] (sifir) at ([xshift=0.4cm, yshift=0.6cm]ab.east) {Sıfır};
  \draw[->, yeni] ([yshift=0.6cm]ab.east) -- (sifir.west);

  % Uygulanmış buyruklar kutusu (sağ üst köşe; 6/7)
  \node[draw, rounded corners, fill=blue!5, font=\tiny, align=center, inner sep=4pt, anchor=north east] at (13.0, 7.2) {Buyruk durumu (6/7)\\[2pt]\begin{tabular}{@{}l@{\quad}l@{}}
1.~{\color{green!60!black}\checkmark}~\texttt{add} & 5.~{\color{green!60!black}\checkmark}~\texttt{sw}\\
2.~{\color{green!60!black}\checkmark}~\texttt{sub} & 6.~{\color{green!60!black}\checkmark}~\texttt{beq}/\texttt{bne}\\
3.~{\color{green!60!black}\checkmark}~\texttt{ori} & 7.~{\color{gray}$\circ$}~\texttt{jal}\\
4.~{\color{green!60!black}\checkmark}~\texttt{lw} & \\
\end{tabular}};
\end{tikzpicture}
\caption[Dallanma buyruklarını (\texttt{beq}, \texttt{bne}) yürüten veri yolu]{Dallanma buyruklarını (\texttt{beq}, \texttt{bne}) yürüten veri yolu. Şekil~\ref{fig:vy-sw}'ye göre beş yenilik (kırmızı): (1) \textbf{Dal.\ Top.} (dallanma adresi toplayıcısı) PS ile B-türü anlık değeri toplayarak dallanma hedefini üretir; (2) AMB'nin yeni \textbf{Sıfır} çıkışı, \texttt{x5--x6} farkının sıfır olup olmadığını söyler; (3) \textbf{PSKaynak} çoklayıcısı sıradaki PS değerini PS+4 (varsayılan akış) ile dallanma hedefi arasından seçer; (4) \textbf{Dallan} denetim sinyali (dallanma buyruğunda Sıfır bayrağıyla birleştirilir) seçimi yapar; (5) \textbf{AnlıkBiçim} denetim sinyali, anlık değer üretecine buyruğun türünü bildirir (I, S, B parçalarını farklı sırada okur, B türü için alt 0'ı ekler). Buyruk şeridi B türünü göstermektedir; anlık değer dört parçaya dağılmıştır (\texttt{imm[12,10:5]} ve \texttt{imm[4:1,11]}) ve anlık değer üretecinde 12 saklanan bit + örtük 0 ile 13 bitlik etkin uzaklık olarak kurulup 32 bite işaretle genişletilir.}
\label{fig:vy-beq}
\end{figure}

Adımlar:
\begin{enumerate}
    \item \texttt{x5} ve \texttt{x6} değerlerini yazmaç öbeğinden oku.
    \item AMB'de \texttt{x5 - x6} farkını hesapla. Sonuç sıfırsa eşitlik vardır (\texttt{Sıfır} bayrağı 1 olur).
    \item Aynı çevrimde \texttt{Dal.\,Top.} \texttt{PS + işaretle genişletilmiş B-anlık} değerini hesaplar.
    \item \texttt{Dallan} denetim sinyali \texttt{PSKaynak} çoklayıcısına seçim verir: 0\,$=$\,PS+4, 1\,$=$\,dallanma hedefi.
\end{enumerate}

\texttt{beq}'in veri yoluna kattığı dört yeni öğe vardır.

\textbf{Dallanma adresi toplayıcısı (Dal.\ Top.)}, PS ile B-türü işaretle genişletilmiş anlık değeri toplayarak hedef adresi üretir. PS+4 toplayıcısından bağımsız ikinci bir toplayıcıdır ve her çevrimde, dallanılsın dallanılmasın, paralel olarak çalışır. Sonuç \texttt{PSKaynak} çoklayıcısının bir girişine bağlanır. Dallanılmazsa boşa hesaplanmış olur.

\textbf{AMB Sıfır çıkışı}, AMB iç yapısında zaten var olan ``sonucun tüm bitleri 0 mı?'' bilgisini dışarı taşır (Bölüm~\ref{chap:aritmetik}). \texttt{beq} için AMB'ye \texttt{x5 - x6} farkı hesaplatılır (\texttt{AMBİşlem}=01). Sonuç sıfırsa iki yazmaç eşittir. Bu bit, denetim birimine giden tek-bitlik bir bayraktır. Veri yolundaki diğer 32-bitlik yollardan farklı olarak yalnızca koşulu temsil eder.

\textbf{PSKaynak çoklayıcısı}, PS girişini PS+4 ile dallanma hedefi arasında seçer. Seçim bilgisini \texttt{Dallan} denetim sinyalinden alır.

\textbf{\texttt{Dallan} denetim sinyali}, \emph{dallanılıp dallanılmayacağını} söyleyen tek-bitlik bir sinyaldir. \texttt{PSKaynak} çoklayıcısının seçim girişine bağlanır. \texttt{Dallan}=1 olduğunda dallanma hedefi PS'ye yönlenir, \texttt{Dallan}=0 olduğunda PS+4 yönlenir.

\texttt{Dallan} sinyalinin değerine denetim birimi karar verir. İki şeye bakar:
\begin{enumerate}
    \item \emph{Yürütülen buyruk bir dallanma buyruğu mu?} İşlem kodundan (\emph{opcode}) anlaşılır. Değilse (örn. \texttt{add}, \texttt{lw}) \texttt{Dallan}=0 olur, dallanma engellenir.
    \item \emph{Dallanma buyruğu ise, koşul sağlandı mı?} AMB'nin \texttt{Sıfır} bayrağına bakılır. \texttt{beq} (eşitse dallan) için \texttt{Sıfır}=1 ise koşul sağlanmıştır (yazmaçlar eşit); \texttt{bne} (eşit değilse dallan) için ise \texttt{Sıfır}=0 ise koşul sağlanmıştır (yazmaçlar farklı).
\end{enumerate}

Aynı veri yolu donanımı her iki buyruğu da yürütür. Tek fark, denetim biriminin işlem koduna bakarak \texttt{Sıfır} bayrağını hangi yönde yorumladığıdır. Denetim biriminin bu kararı nasıl ürettiği Bölüm~\ref{sec:denetim-birimi}'te ele alınacaktır.

\subsubsection*{B türü anlık değerin yapısı ve 13 bit meselesi}

B türü buyruğun anlık değer alanı dört parçaya dağılmıştır: \texttt{imm[12]} (bit 31), \texttt{imm[10:5]} (bit 30--25), \texttt{imm[4:1]} (bit 11--8), \texttt{imm[11]} (bit 7). Buyruktan toplam 12 bit fiziksel olarak gelir. Oysa bu 12 bit \emph{13-bit'lik bir işaretli uzaklığı} temsil eder. En alt bit (\texttt{imm[0]}) buyrukta saklanmaz, donanım tarafından sabit 0 olarak eklenir.

Bunun nedeni RISC-V'in iki farklı buyruk uzunluğunu birlikte desteklemesidir. Temel buyruk kümesi (RV32I) her buyruğu \emph{4 bayt} olarak kodlar. Bu durumda her buyruğun başlangıç adresi 4'ün katıdır. Ancak RISC-V mimarisi, isteğe bağlı \emph{sıkıştırılmış buyruk uzantısı} (\emph{compressed instruction extension}, kısaca C uzantısı) tanımlar. Bu uzantı sık kullanılan bazı buyrukları daha kısa, \emph{2 baytlık} biçimde kodlar. Aynı program içinde 4 baytlık ve 2 baytlık buyruklar karışık olarak bulunabileceği için RISC-V belirtimi, dallanma hedef adreslerinin, iki olası buyruk uzunluğunun (2 ve 4 bayt) en büyük ortak böleni olan \emph{2'nin katı} olmasını şart koşar. Böylece sıkıştırılmış uzantı kullanılsın ya da kullanılmasın, dallanma hedefi her zaman geçerli bir buyruk başlangıcına denk gelir. İki bayt arasında bir adrese (örneğin tek sayılı bir adrese) düşemez. Sonuç olarak hedef uzaklığın en alt biti (\texttt{imm[0]}) her zaman 0'dır ve bu biti buyrukta saklamak bilgi kaybı olmaksızın atlanır. 12 saklanan bit + 1 örtük sıfır = 13 bitlik işaretli uzaklık, PS'ye göre yaklaşık $\pm 4$ KB (yaklaşık $\pm 1024$ buyruk) menzili 2'nin katı adresler üzerinde ifade eder. Aynı mantık J türü (\texttt{jal}) için de geçerlidir. Orada 20 saklanan bit + 1 örtük sıfır $= 21$ bit uzaklık olur (bir sonraki bölümde).

\subsubsection*{\texttt{AnlıkBiçim} denetim sinyali}

Anlık değer, buyruk türüne göre buyruğun \emph{farklı bit konumlarından} gelir: I türünde bit 31--20 (tek parça, 12 bit), S türünde bit 31--25 ile 11--7 (iki parça, 12 bit), B türünde bit 31, 7, 30--25 ve 11--8 (dört parça, 12 bit). Bu farklı kaynaklardan tek bir 32 bitlik çıkışa erişmek için üretecin içine bir çoklayıcı yerleştirilir. Çoklayıcı, türe göre doğru bit alanlarını doğru sırayla çıkışa aktarır. Sonra ortak bir işaretle genişletme katmanı sonucu 32 bite tamamlar. B türü için ek bir adım daha vardır. En alt bite sabit 0 eklenir (yukarıda anlattığımız 13 bitlik etkin uzaklık).

Bu iç çoklayıcının seçim girişi \texttt{AnlıkBiçim} denetim sinyalidir. Denetim birimi yürütülen buyruğun türüne bakarak \texttt{AnlıkBiçim}'i ayarlar. Üreteç de o türe karşılık gelen bit alanlarını seçip işaretle genişletir. J türü eklendiğinde \texttt{AnlıkBiçim} dördüncü bir değer daha kazanır. Anlık değer üretecinin iç yapısı bu bölümün ileriki kısmında, veri yolu birimlerinin iç yapılarını incelediğimiz Bölüm~\ref{subsec:vy-ic-yapi}'de ele alınacaktır.

PS güncellemesi artık iki olası girdiye sahiptir: PS+4 (normal akış) ve PS+B-anlık (dallanılırsa). Bir sonraki adımda \texttt{jal} üçüncü bir girdi daha ekleyecektir.

\subsection{\texttt{jal} Buyruğunu Eklemek: Koşulsuz Atlama ve Dönüş Adresi}
\label{subsec:vy-jal}

\texttt{jal x1, etiket} (\emph{jump and link}) buyruğu iki iş yapar: koşulsuz olarak \texttt{etiket} adresine atlar ve bir sonraki buyruğun adresini (\texttt{PS+4}) hedef yazmaca yazar. İkinci iş, altyordam (\emph{subroutine}) çağrılarında dönüş adresinin saklanmasını ve altyordam bittikten sonra çağıran koda geri dönülebilmesini sağlar. Adımlar:
\begin{enumerate}
    \item Anlık değer üreteci, J türü 20 bitlik anlık değeri okur, en altına örtük bir 0 ekleyerek 21 bitlik etkin uzaklık olarak kurar ve 32 bite işaretle genişletir.
    \item Dal.\ Top.\ toplayıcısı PS ile bu uzaklığı toplayarak atlama hedef adresini üretir.
    \item Aynı çevrimde Top.\ toplayıcısı \texttt{PS+4}'ü hesaplar. Bu değer dönüş adresi olarak yazmaç öbeğinin yazma kapısına yönlenir.
    \item Denetim birimi \texttt{Dallan}=1 üreterek PSKaynak çoklayıcısının hedef adresi seçmesini ve PS'nin yeni değerle güncellenmesini sağlar.
    \item Aynı zamanda \texttt{YazmacaYaz}=1 yaparak PS+4 değerinin \texttt{rd} ile gösterilen yazmaca yazılmasını etkinleştirir.
\end{enumerate}

Şekil~\ref{fig:vy-jal}'da bu adımdan sonra elde edilen tam veri yolu gösterilmektedir. Aşağıda eklenen tek yapısal yeniliği ve buyruğun denetim sinyalleri üzerindeki etkilerini sırayla ele alıyoruz.

\begin{figure}[htbp]
\centering
\begin{tikzpicture}[>=Stealth, semithick, font=\footnotesize,
    yzm/.style={draw, fill=yellow!25, minimum width=1cm, minimum height=1cm, align=center},
    bllk/.style={draw, fill=green!20, minimum width=1.8cm, minimum height=1.6cm, align=center},
    yobk/.style={draw, fill=blue!10, minimum width=2.6cm, minimum height=3cm},
    amb/.style={draw, fill=red!20, minimum width=1.6cm, minimum height=2cm, align=center},
    artop/.style={draw, fill=gray!15, minimum width=1.1cm, minimum height=0.7cm, align=center},
    anlikb/.style={draw, fill=orange!25, minimum width=1.4cm, minimum height=0.9cm, align=center},
    muxdik/.style={draw, fill=purple!15, minimum width=0.7cm, minimum height=1.0cm, align=center},
    muxyatay/.style={draw, fill=purple!15, minimum width=1.0cm, minimum height=0.7cm, align=center},
    vbll/.style={draw, fill=green!20, minimum width=1.8cm, minimum height=2.2cm, align=center},
    bitayar/.style={fill=white, inner sep=1pt, font=\tiny},
    yeni/.style={draw=red!70!black, very thick}
]
  % --- Sol PS güncelleme bloğu (Şekil 6.8 ile aynı yapı) ---
  \node[bllk] (bb) at (0, 4.6) {Buyruk\\Belleği};
  \node[yzm] (ps) at (0, 2.7) {PS};
  \node[muxyatay] (psmux) at (0, 1.3) {};
  \node[font=\tiny] at (psmux.center) {Çok.};
  \node[font=\tiny, anchor=south] at ([xshift=-0.3cm, yshift=0.02cm]psmux.south) {0};
  \node[font=\tiny, anchor=south] at ([xshift=0.3cm, yshift=0.02cm]psmux.south) {1};

  \node[artop] (top4) at (-1.3, 0) {Top.};
  \node[artop, minimum width=1.2cm, minimum height=0.7cm] (daladd) at (1.3, 0) {Dal.\\Top.};

  \draw[->] (ps.north) -- (bb.south) node[midway, right=1pt, font=\scriptsize] {adres};
  \draw[->] (psmux.north) -- (ps.south);

  \draw[->] (ps.west) -- (-1.5, 2.7) -- ([xshift=-0.2cm]top4.north);
  \draw[->] (ps.east) -- (1.1, 2.7) -- ([xshift=-0.2cm]daladd.north);

  \node[font=\scriptsize] at (-1.1, 0.85) {$4$};
  \draw[->] (-1.1, 0.7) -- ([xshift=0.2cm]top4.north);

  % B/J anlık -> Dal.Top.north_east (dış üst sağ giriş, anlık.east teli üzerinden tap)
  \coordinate (anlikdal) at (6.85, 1.0);
  \node[circle, fill=black, inner sep=1pt] at (anlikdal) {};
  \draw[->] (anlikdal) -- (1.5, 1.0) -- ([xshift=0.2cm]daladd.north);

  % Top4 çıkışı (PS+4): tap ile hem PSKaynak'a hem GeriYazma'ya gider
  \draw[->] (top4.south) -- (-1.3, -0.7) -- (-0.3, -0.7) -- ([xshift=-0.3cm]psmux.south);
  \node[font=\scriptsize, anchor=north] at (-0.8, -0.75) {PS$+$4};
  \draw[->] (daladd.south) -- (1.3, -0.7) -- (0.3, -0.7) -- ([xshift=0.3cm]psmux.south);
  \node[font=\scriptsize, anchor=north] at (0.8, -0.75) {Dal.hedef};

  % Dallan denetim sinyali (psmux'a SOLDAN)
  \draw[->, sinyallacivert, thick] ([xshift=-0.4cm]psmux.west) -- (psmux.west);
  \node[font=\scriptsize, text=sinyallacivert, anchor=south] at ([xshift=-0.5cm, yshift=0.05cm]psmux.west) {Dallan};

  \coordinate (parse) at (2.5, 4.6);
  \draw[->] (bb.east) -- (parse) node[midway, below, font=\scriptsize] {buyruk} node[bitayar, pos=0.55] {32 bit};

  % J türü buyruk şeridi (DEĞİŞTİ: 20 bitlik tek imm alan + rd + opcode; kırmızı çerçeve)
  \coordinate (seritsol) at (3.0, 6.4);
  \draw[gray, dashed, semithick] (parse) -- (seritsol);
  \draw[gray, dashed, semithick] (parse) -- ([xshift=6cm]seritsol);
  \begin{scope}[shift={(seritsol)}]
    \draw[fill=orange!25, yeni] (0,0) rectangle (3.75,0.55) node[pos=.5, font=\tiny] {imm[20,10:1,11,19:12]};
    \draw[fill=blue!18] (3.75,0) rectangle (4.6875,0.55) node[pos=.5, font=\tiny] {rd};
    \draw[fill=gray!10] (4.6875,0) rectangle (6,0.55) node[pos=.5, font=\tiny] {opcode};
    \node[font=\tiny, anchor=north] at (0,0) {31};
    \node[font=\tiny, anchor=north] at (3.75,0) {12};
    \node[font=\tiny, anchor=north] at (4.6875,0) {7};
    \node[font=\tiny, anchor=north] at (6,0) {0};
    \node[font=\scriptsize, anchor=south] at (2.3,0.55) {J türü buyruk};
  \end{scope}

  \node[yobk] (yo) at (5.4, 2.6) {};
  \node[anchor=center, font=\footnotesize\bfseries, align=center] at (yo.center) {Yazmaç\\Öbeği};

  % J türü: rd -> C_adr (tek bağlantı; rs1/rs2 alanları yok)
  \draw[->] ([xshift=4.22cm]seritsol) -- ++(0,-1.2) -| ([xshift=0.6cm]yo.north) node[bitayar, pos=0.75] {5 bit};

  \node[font=\scriptsize, anchor=north] at ([xshift=-0.6cm,yshift=-0.05cm]yo.north) {$A_{\text{adr}}$};
  \node[font=\scriptsize, anchor=north] at ([xshift=0cm,yshift=-0.05cm]yo.north) {$B_{\text{adr}}$};
  \node[font=\scriptsize, anchor=north] at ([xshift=0.6cm,yshift=-0.05cm]yo.north) {$C_{\text{adr}}$};
  \node[font=\scriptsize, anchor=east] at ([xshift=-0.1cm,yshift=0.5cm]yo.east) {$A_{\text{veri}}$};
  \node[font=\scriptsize, anchor=east] at ([xshift=-0.1cm,yshift=-0.25cm]yo.east) {$B_{\text{veri}}$};
  \node[font=\scriptsize, anchor=west] at ([xshift=0.1cm,yshift=-1.0cm]yo.west) {$C_{\text{veri}}$};

  % YazmacaYaz (jal için 1: dönüş adresi yazılacak); etiket okun ÜSTÜNDE, Y harfi okun başlangıcıyla tam hizalı (inner sep=0)
  \node[font=\scriptsize, text=sinyallacivert, anchor=south west, inner sep=0pt] at ([xshift=1.05cm, yshift=0.4cm]yo.north) {YazmacaYaz};
  \draw[->, sinyallacivert, thick] ([xshift=1.05cm, yshift=0.4cm]yo.north) -- ([xshift=1.05cm]yo.north);

  % --- AMBKaynak çoklayıcısı ---
  \node[muxdik] (mux1) at (8.0, 2.1) {};
  \node[font=\tiny] at (mux1.center) {Çok.};
  \node[font=\tiny, anchor=west] at ([xshift=0.02cm, yshift=0.25cm]mux1.west) {0};
  \node[font=\tiny, anchor=west] at ([xshift=0.02cm, yshift=-0.25cm]mux1.west) {1};

  % --- AMB ---
  \node[amb] (ab) at (9.6, 2.6) {AMB};

  \draw[->] ([yshift=0.5cm]yo.east) -- ([yshift=0.5cm]ab.west) node[bitayar, pos=0.5] {32 bit};
  \draw[->] ([yshift=-0.25cm]yo.east) -- ([yshift=0.25cm]mux1.west);
  \draw[->] (mux1.east) -- ([yshift=-0.5cm]ab.west);

  % --- Veri Belleği (aşağı kaydırıldı: PS+4 teli AMB-vb boşluğundan rahat geçsin diye) ---
  \node[vbll] (vb) at (11.65, 0.25) {Veri\\Belleği};
  \node[font=\scriptsize, anchor=north] at ([xshift=-0.5cm, yshift=-0.05cm]vb.north) {$A_{\text{adr}}$};
  \node[font=\scriptsize, anchor=north] at ([xshift=0.5cm, yshift=-0.05cm]vb.north) {$V_{\text{oku}}$};
  \node[font=\scriptsize, anchor=west] at ([xshift=0.1cm, yshift=-0.85cm]vb.west) {$V_{\text{yaz}}$};

  \coordinate (ambsplit) at (11.15, 2.6);
  \node[circle, fill=black, inner sep=1pt] at (ambsplit) {};
  \draw[->] (ambsplit) -- ([xshift=-0.5cm]vb.north);

  \coordinate (bverijnct) at (7.0, 2.35);
  \node[circle, fill=black, inner sep=1pt] at (bverijnct) {};
  \draw[->] (bverijnct) -- (7.0, -0.6) -- ([yshift=-0.85cm]vb.west) node[bitayar, pos=0.55] {32 bit};

  \node[font=\scriptsize, text=sinyallacivert, anchor=north] at ([yshift=-0.5cm]vb.south) {BelleğeYaz};
  \draw[->, sinyallacivert, thick] ([yshift=-0.4cm]vb.south) -- (vb.south);

  % --- GeriYazmaKaynağı çoklayıcısı (3 girdili — kırmızı çerçeve; mux2 aşağı kaydırıldı: AMB çıkışıyla aynı hizada) ---
  \node[muxdik, yeni, minimum height=1.5cm] (mux2) at (13.25, 2.2) {};
  \node[font=\tiny, rotate=90] at ([xshift=0.1cm]mux2.center) {Çoklayıcı};
  \node[font=\tiny, anchor=west] at ([xshift=0.02cm, yshift=0.4cm]mux2.west) {0};
  \node[font=\tiny, anchor=west] at ([xshift=0.02cm, yshift=0cm]mux2.west) {1};
  \node[font=\tiny, anchor=west] at ([xshift=0.02cm, yshift=-0.4cm]mux2.west) {2};

  % AMB çıkışı -> mux2 input 0 (üst giriş, y=2.6) — TAM YATAY, aynı hizada
  \draw[->] (ab.east) -- ([yshift=0.4cm]mux2.west) node[bitayar, pos=0.15] {32 bit};

  % V_oku -> mux2 input 1 (orta giriş, y=2.2)
  \draw[->] ([xshift=0.5cm]vb.north) |- ([yshift=0cm]mux2.west) node[bitayar, pos=0.6] {32 bit};

  % YENİ: PS+4 -> mux2 input 2 (alt giriş, y=1.8) — en aşağıdan dolanır (her şeyin altından, mux2 çıkış teli dahil), AMB-vb boşluğundan yukarı dirsek
  \coordinate (ps4tap) at (-0.3, -0.7);
  \node[circle, fill=black, inner sep=1pt, yeni] at (ps4tap) {};
  \draw[->, yeni] (ps4tap) -- (-0.3, -2.0) -- (10.55, -2.0) -- (10.55, 1.8) -- ([yshift=-0.4cm]mux2.west) node[bitayar, pos=0.55] {32 bit};

  % mux2 çıkışı -> Cveri yolu
  \draw[->] (mux2.east) -- ++(0.5,0)
            |- ++(0,-4.0)
            -| ([xshift=-1.0cm,yshift=-1.0cm]yo.west)
            -- ([yshift=-1.0cm]yo.west);
  \node[bitayar, anchor=west] at ([xshift=0.15cm]mux2.east) {32 bit};

  % --- Anlık değer üreteci ---
  \node[anlikb] (anlik) at (5.4, -0.2) {Anlık\\Üreteci};
  % J türü: 20 bit anlık değer; üreteç içinde imm[0]=0 eklenip 21-bit etkin uzaklık olarak işaretle genişletilir
  \draw[->] ([xshift=1.875cm]seritsol) -- ++(0,-1.0) -- ++(-1.35,0) |- (anlik.west) node[bitayar, pos=0.75] {20 bit};
  \draw[->] (anlik.east) -- ++(0.75,0) coordinate (anlikbifurc) -- (6.85, 1.85) -- ([yshift=-0.25cm]mux1.west);
  \node[bitayar] at (6.5, -0.2) {32 bit};
  % AnlıkBiçim denetim sinyali (J türü için 4. değer kazandı — kırmızı çerçeve)
  \node[font=\scriptsize, text=sinyalkirmizi, anchor=north] at ([yshift=-0.5cm]anlik.south) {AnlıkBiçim};
  \draw[->, sinyalkirmizi, thick] ([yshift=-0.4cm]anlik.south) -- (anlik.south);

  % AMBİşlem denetim sinyali
  \node[font=\scriptsize, text=sinyallacivert, anchor=south, align=center] at ([yshift=0.5cm]ab.north) {AMBİşlem\\[-1pt]\tiny 00\,$=$\,Topla,\ 01\,$=$\,Çıkar,\ 10\,$=$\,VEYA};
  \draw[->, sinyallacivert, thick] ([yshift=0.4cm]ab.north) -- (ab.north);

  % AMBKaynak denetim sinyali
  \node[font=\scriptsize, text=sinyallacivert, anchor=north] at ([yshift=-0.4cm]mux1.south) {AMBKaynak};
  \draw[->, sinyallacivert, thick] ([yshift=-0.3cm]mux1.south) -- (mux1.south);

  % GeriYazmaKaynağı denetim sinyali (genişletildi: artık 2 bit, 3 değer — kırmızı çerçeve)
  \node[font=\scriptsize, text=sinyalkirmizi, anchor=north, align=center] at ([yshift=-0.4cm]mux2.south) {GeriYazma\\Kaynağı};
  \draw[->, sinyalkirmizi, thick] ([yshift=-0.3cm]mux2.south) -- (mux2.south);

  % AMB Sıfır çıkışı (beq'den kalan; jal'da kullanılmaz)
  \node[font=\scriptsize, anchor=west, inner sep=2pt] (sifir) at ([xshift=0.4cm, yshift=0.6cm]ab.east) {Sıfır};
  \draw[->] ([yshift=0.6cm]ab.east) -- (sifir.west);

  % Uygulanmış buyruklar kutusu (sağ üst köşe; 7/7 — TÜMÜ tamamlandı)
  \node[draw, rounded corners, fill=blue!5, font=\tiny, align=center, inner sep=4pt, anchor=north east] at (13.0, 7.2) {Buyruk durumu (7/7)\\[2pt]\begin{tabular}{@{}l@{\quad}l@{}}
1.~{\color{green!60!black}\checkmark}~\texttt{add} & 5.~{\color{green!60!black}\checkmark}~\texttt{sw}\\
2.~{\color{green!60!black}\checkmark}~\texttt{sub} & 6.~{\color{green!60!black}\checkmark}~\texttt{beq}/\texttt{bne}\\
3.~{\color{green!60!black}\checkmark}~\texttt{ori} & 7.~{\color{green!60!black}\checkmark}~\texttt{jal}\\
4.~{\color{green!60!black}\checkmark}~\texttt{lw} & \\
\end{tabular}};
\end{tikzpicture}
\caption[Sekiz RISC-V buyruğunu yürüten tam tek vuruşluk veri yolu]{Sekiz RISC-V buyruğunu yürüten tam tek vuruşluk veri yolu. Şekil~\ref{fig:vy-beq}'e göre tek yapısal yenilik (kırmızı): \textbf{GeriYazmaKaynağı çoklayıcısı} 2 girdiden 3 girdiye genişletildi; üçüncü girdi olarak \texttt{PS+4} (dönüş adresi) eklendi. Böylece \texttt{jal} bir sonraki buyruğun adresini hedef yazmaca yazabilir. Dallanma adresi toplayıcısı (Dal.\ Top.) zaten kuruluydu; J türü 20 bitlik anlık değer üreteç tarafından üretilip Dal.\ Top.'a girdiğinde, jal'ın hedef adresi de aynı toplayıcıyla hesaplanır. Buyruk şeridi J türünü göstermektedir; 20 saklanan bit + örtük 0 ile 21 bitlik etkin uzaklık anlık değer üretecinde kurulup 32 bite işaretle genişletilir. \texttt{jal} koşulsuz olduğundan \texttt{Dallan} denetim sinyali jal için her zaman 1'dir.}
\label{fig:vy-jal}
\end{figure}

\subsubsection*{Yapısal değişiklik: \texttt{GeriYazmaKaynağı} çoklayıcısı genişler}

Önceki tasarımda \texttt{GeriYazmaKaynağı} çoklayıcısı iki kaynaktan birini yazmaç öbeğinin $C_{\text{veri}}$ girişine yönlendiriyordu: aritmetik buyruklarda AMB sonucu, \texttt{lw}'de veri belleğinden okunan değer. \texttt{jal}'ın yazmaç öbeğine yazacağı değer ise PS+4'tür (dönüş adresi); bu, AMB veya bellekten gelmediği için çoklayıcıya \emph{üçüncü} bir girdi gerekir. \texttt{jal}'ın getirdiği tek yapısal yenilik budur.

Çoklayıcı 2 girdiden 3 girdiye genişler. Seçim sinyali \texttt{GeriYazmaKaynağı} de buna bağlı olarak 1 bitten 2 bite çıkar:
\begin{itemize}
    \item \texttt{00}: AMB sonucu (\texttt{add}, \texttt{sub}, \texttt{ori})
    \item \texttt{01}: Veri belleği (\texttt{lw})
    \item \texttt{10}: PS+4 (\texttt{jal})
\end{itemize}

Geri kalan tüm donanım önceki adımlardan miras alınır. AMB'ye yeni bir işlem eklenmez. Dal.\ Top.\ aynı kalır. PSKaynak çoklayıcısı yine 2 girdilidir. \texttt{jal}'ın etkisi büyük ölçüde \emph{denetim katmanındadır}. Var olan donanım üç sinyal aracılığıyla yeni bir biçimde yönlendirilir.

\subsubsection*{Denetim sinyallerinin \texttt{jal} için yorumu}

\textbf{Dal.\ Top.\ \texttt{jal} için yeniden kullanılır.} Toplayıcı her çevrimde PS ile anlık değer üretecinden gelen değeri toplar. Anlık değer üreteci buyruk türüne göre farklı bit alanlarını seçtiği için, B türünde 13 bitlik uzaklık, J türünde 21 bitlik uzaklık çıkar. İkisi de Dal.\ Top.'a aynı 32 bitlik girdi olarak ulaşır. Aynı donanım iki farklı atlama türünü destekler. Davranış farkı tamamen anlık değer üretecine verilen \texttt{AnlıkBiçim} sinyaline bağlıdır.

\textbf{\texttt{Dallan} koşulsuz aktiftir.} \texttt{beq}/\texttt{bne}'de \texttt{Dallan} sinyali \texttt{Sıfır} bayrağına bağlıydı; \texttt{jal}'da koşul yoktur, denetim birimi yalnızca işlem koduna bakarak \texttt{Dallan}=1 üretir. PSKaynak çoklayıcısı atlama hedefini seçer ve PS güncellenir.

\textbf{\texttt{YazmacaYaz} aktiftir.} \texttt{beq}/\texttt{bne}'de bu sinyal pasif tutuluyordu (yazma yok); \texttt{jal} dönüş adresini yazdığı için aktiftir. Yazılacak yazmacın numarası (\texttt{rd}, buyruğun bit 11--7 alanı) doğrudan $C_{\text{adr}}$ kapısına ulaşır. $C_{\text{veri}}$ girişine ise GeriYazmaKaynağı çoklayıcısının seçtiği PS+4 değeri yönlenir.

\textbf{\texttt{AnlıkBiçim} dördüncü değerini kazanır.} Anlık değer üreteci artık dört buyruk türünü destekler: I (\texttt{lw}, \texttt{ori}), S (\texttt{sw}), B (\texttt{beq}, \texttt{bne}) ve J (\texttt{jal}). Üreteç içindeki çoklayıcı, \texttt{AnlıkBiçim} sinyalinin değerine göre buyruğun ilgili bit alanlarını doğru sıraya koyar, gerekirse alt 0'ı ekler ve sonucu işaretle genişletir. Dört durumu ayırmak için sinyalin iki bit olması gerekir. J türünün eklenmesiyle dördüncü değer kullanıma girer.

\subsubsection*{J türü anlık değer biçimi ve 21 bitlik etkin uzaklık}

J türü buyruğun anlık değer alanı, B türündekine benzer biçimde \emph{dört parçaya} dağıtılmıştır: \texttt{imm[20]} (bit 31), \texttt{imm[10:1]} (bit 30--21), \texttt{imm[11]} (bit 20), \texttt{imm[19:12]} (bit 19--12). Buyruktan toplam 20 bit fiziksel olarak gelir, ancak bu \emph{21 bitlik bir işaretli uzaklığı} temsil eder. En alt bit (\texttt{imm[0]}) buyrukta saklanmaz, donanım tarafından sabit 0 olarak eklenir. Sebep B türü ile aynıdır. Tüm atlama hedef adresleri 2'nin katı olmak zorundadır.

21 bitlik işaretli uzaklık, PS'ye göre yaklaşık $\pm 1$ MB ($\pm 2^{20}$ bayt) menzili ifade eder. Bu, B türünün $\pm 4$ KB menzilinden 256 kat daha geniştir. Çoğu altyordam çağrısı için bu menzil yeterlidir. Daha uzak hedeflere atlamak gerektiğinde RISC-V \emph{dolaylı atlama} (\emph{indirect jump}) buyruğu \texttt{jalr}'ı kullanır. Bu buyruk hedef adresi bir yazmaçtan okur ve menzili tüm 32 bitlik adres uzayını kapsar (Ek~\ref{app:risc-v-referans}).

\medskip
Bu adımla birlikte sekiz buyruğu (\texttt{add}, \texttt{sub}, \texttt{ori}, \texttt{lw}, \texttt{sw}, \texttt{beq}, \texttt{bne}, \texttt{jal}) yürüten tek vuruşluk veri yolu tamamlanmıştır. Sonraki alt bölümlerde yazmaç öbeği ve anlık değer üretecinin iç yapısına yakından bakacak, ardından bu sekiz buyruğa karşılık gelen denetim sinyallerini üreten denetim birimini tasarlayacağız.

\begin{karsilastirma}[RISC-V'in Yazmaç Alanları için Sabit Konum Tercihi]
Veri yolu boyunca dikkat çeken bir ayrıntı vardır. Buyruktan yazmaç öbeğine giden üç adres yolu (rs1, rs2, rd) doğrudan buyruğun belirli bit aralıklarından gelir. RISC-V'in tüm buyruk türlerinde \texttt{rd} her zaman bitler \texttt{[11:7]}'de, \texttt{rs1} bitler \texttt{[19:15]}'te ve \texttt{rs2} bitler \texttt{[24:20]}'de yer alır. R, I, S, B ve J türleri arasında bu üç alanın konumu değişmez. Yalnızca buyruğun anlık değer ve işlem kodu alanları yeniden düzenlenir.

Bu tasarım kararı veri yolunu sadeleştirir. MIPS gibi farklı türlerde yazmaç adres alanlarını farklı bit konumlarına yerleştiren mimarilerde, hedef yazmacın hangi alandan okunacağını (\texttt{rt} mi yoksa \texttt{rd} mi) seçen ek bir çoklayıcı gerekir. Bu seçim için de ayrı bir denetim sinyali üretilmek zorundadır. RISC-V'te ise yazmaç adres yolları sabit olduğu için böyle bir çoklayıcıya ve sinyale ihtiyaç kalmaz. Yazmaca yazılıp yazılmayacağı (S ve B türü buyruklarda yazılmaz) yalnızca \texttt{YazmacaYaz} sinyaliyle belirlenir.

Bedel olarak, RISC-V'in anlık değer alanları farklı türlerde dağınık biçimde kodlanır (özellikle B ve J türlerinde bitler dört parçaya ayrılır). Bu dağılım, anlık değer üretecinin biraz daha karmaşık olmasını gerektirir. Ancak yazmaç adres yollarındaki sadeleşme, denetim biriminde kazanılan basitlikle birlikte bu maliyeti fazlasıyla karşılar.
\end{karsilastirma}

\begin{bilginot}[RISC-V'in Dört Anlık Değer Biçimi]
Veri yolundaki anlık değer üreteci, dört buyruk türünden gelen anlık değer bitlerini buyruğun farklı bölgelerinden toplar ve hepsini 32 bite işaretle genişletir. R türü buyruklar anlık değer içermediği için üretecin bu tür için bir görevi yoktur. \texttt{AnlıkBiçim} sinyali (2 bit) üretece I, S, B ve J olmak üzere dört durumdan birini seçtirir.

\begin{center}
\begin{tikzpicture}[>=Stealth, font=\tiny\ttfamily,
    alan/.style={draw, minimum height=0.55cm, inner sep=1pt, anchor=west},
    turadi/.style={font=\small\bfseries\rmfamily, anchor=east},
    okst/.style={->, thick}]

  \def\Iy{0}
  \def\Sy{-1.2}
  \def\By{-2.4}
  \def\Jy{-3.6}

  % Tüm satırlar 8.0 cm (= 32 bit × 0.25 cm/bit). J türünde tek bitlik
  % alanlar etiket okunabilirliği için 1.0 cm; çok bitlik alanlar bunu
  % dengelemek için 0.2 cm/bit ölçeğinde tutulur. Toplam yine 8.0 cm.
  \def\xOK{8.2}
  \def\xCIK{9.0}

  %----- I türü (12 + 5 + 3 + 5 + 7 = 32 bit, 0.25 cm/bit, toplam 8.0 cm) -----
  \node[turadi] at (-0.2, \Iy) {I};
  \node[alan, minimum width=3.0cm, fill=blokmavi!60] at (0, \Iy) {imm[11:0]\ (12)};
  \node[alan, minimum width=1.25cm, fill=bloksari!60] at (3.0, \Iy) {rs1 (5)};
  \node[alan, minimum width=0.75cm, fill=blokyesil!60] at (4.25, \Iy) {f3};
  \node[alan, minimum width=1.25cm, fill=blokkirmizi!40] at (5.0, \Iy) {rd (5)};
  \node[alan, minimum width=1.75cm, fill=bloksari!40] at (6.25, \Iy) {opc (7)};
  \draw[okst] (\xOK, \Iy) -- (\xCIK, \Iy)
        node[right, font=\tiny, align=left] {$\{\text{imm[11]}\}\times 20,$\\imm[11:0]};

  %----- S türü (7 + 5 + 5 + 3 + 5 + 7 = 32 bit, toplam 8.0 cm) -----
  \node[turadi] at (-0.2, \Sy) {S};
  \node[alan, minimum width=1.75cm, fill=blokmavi!60] at (0, \Sy) {imm[11:5]\ (7)};
  \node[alan, minimum width=1.25cm, fill=gray!20] at (1.75, \Sy) {rs2 (5)};
  \node[alan, minimum width=1.25cm, fill=gray!20] at (3.0, \Sy) {rs1 (5)};
  \node[alan, minimum width=0.75cm, fill=blokyesil!60] at (4.25, \Sy) {f3};
  \node[alan, minimum width=1.25cm, fill=blokmavi!60] at (5.0, \Sy) {imm[4:0]\ (5)};
  \node[alan, minimum width=1.75cm, fill=bloksari!40] at (6.25, \Sy) {opc (7)};
  \draw[okst] (\xOK, \Sy) -- (\xCIK, \Sy)
        node[right, font=\tiny, align=left] {$\{\text{imm[11]}\}\times 20,$\\imm[11:5],imm[4:0]};

  %----- B türü (1+6 + 5 + 5 + 3 + 4+1 + 7 = 32 bit, toplam 8.0 cm) -----
  \node[turadi] at (-0.2, \By) {B};
  \node[alan, minimum width=1.75cm, fill=blokmavi!60] at (0, \By) {imm[12,\,10:5]};
  \node[alan, minimum width=1.25cm, fill=gray!20] at (1.75, \By) {rs2 (5)};
  \node[alan, minimum width=1.25cm, fill=gray!20] at (3.0, \By) {rs1 (5)};
  \node[alan, minimum width=0.75cm, fill=blokyesil!60] at (4.25, \By) {f3};
  \node[alan, minimum width=1.25cm, fill=blokmavi!60] at (5.0, \By) {imm[4:1,\,11]};
  \node[alan, minimum width=1.75cm, fill=bloksari!40] at (6.25, \By) {opc (7)};
  \draw[okst] (\xOK, \By) -- (\xCIK, \By)
        node[right, font=\tiny, align=left] {$\{\text{imm[12]}\}\times 19,$\\imm[12,11,10:5,4:1],0};

  %----- J türü (1 + 10 + 1 + 8 + 5 + 7 = 32 bit, toplam 8.0 cm) -----
  % Tek bitlik alanlar 1.0 cm (etiket okunabilirliği için), çok bitlik
  % alanlar 0.2 cm/bit (toplamı 8.0 cm'e tutturmak için).
  \node[turadi] at (-0.2, \Jy) {J};
  \node[alan, minimum width=1.0cm, fill=blokmavi!60] at (0, \Jy) {imm[20]};
  \node[alan, minimum width=2.0cm, fill=blokmavi!60] at (1.0, \Jy) {imm[10:1]\ (10)};
  \node[alan, minimum width=1.0cm, fill=blokmavi!60] at (3.0, \Jy) {imm[11]};
  \node[alan, minimum width=1.6cm, fill=blokmavi!60] at (4.0, \Jy) {imm[19:12]\ (8)};
  \node[alan, minimum width=1.0cm, fill=blokkirmizi!40] at (5.6, \Jy) {rd (5)};
  \node[alan, minimum width=1.4cm, fill=bloksari!40] at (6.6, \Jy) {opc (7)};
  \draw[okst] (\xOK, \Jy) -- (\xCIK, \Jy)
        node[right, font=\tiny, align=left] {$\{\text{imm[20]}\}\times 11,$\\imm[20,19:12,11,10:1],0};

  % Üst başlıklar
  \node[font=\footnotesize\bfseries\rmfamily] at (4.0, 0.8) {Buyruk bitleri (32 bit)};
  \node[font=\footnotesize\bfseries\rmfamily] at (10.5, 0.8) {Genişletilmiş çıkış (32 bit)};

\end{tikzpicture}
\end{center}

Çıkış sütununda $\{\text{imm[X]}\}\times N$ ifadesi imm[X] bitinin $N$ kez yan yana kopyalanmasını (işaretle genişletme) anlatır. Buradaki gösterim bir üst çizgi (tümleyen) değil, aynı bitin yinelenmesidir. Sıfırla genişletme yerine işaretle genişletme tercih edilmesi RISC-V'in tasarım kararıdır. Bu seçim Bölüm~\ref{subsec:vy-ori}'te \texttt{ori} buyruğu eklenirken ele alındı. B ve J türlerinin en alt biti her zaman 0'dır. B türü 12 saklanan bit + örtük 0 ile 13 bitlik, J türü ise 20 saklanan bit + örtük 0 ile 21 bitlik etkin işaretli uzaklığı kodlar (Bölüm~\ref{subsec:vy-beq} ve~\ref{subsec:vy-jal}).
\end{bilginot}

\subsection{Veri Yolu Birimlerinin İç Yapısı}
\label{subsec:vy-ic-yapi}

Önceki alt bölümlerde veri yolunu adım adım inşa ederken birimleri ``kara kutu'' olarak gösterdik. Şimdi bu birimlere daha yakından bakacağız. Bellekleri kısaca işlev ve organizasyonlarıyla, yazmaç öbeğini ise iç yapısıyla ele alacağız. İleride denetim biriminin neyi yönlendirdiğini anlamak için bu ayrıntılar gereklidir.

\subsubsection{Buyruk ve Veri Bellekleri}

Veri yolunda buyruk belleği ile veri belleğini birbirinden ayrı çizmiş olmamız, Bölüm~\ref{subsec:vonneumann-harvard}'de tanıtılan Harvard mimarisi yaklaşımına karşılık gelir. Her iki belleğe aynı çevrimde erişebilmek tek vuruşluk tasarımı sadeleştirir. Von Neumann mimarisinde ise buyruk ve veri tek bir ortak bellekte saklanır ve aynı çevrimde iki erişim yapılamaz. Çağdaş işlemcilerde bu iki yaklaşım iç içedir. Ana bellek tektir (Von Neumann), fakat işlemci içindeki birinci düzey önbellek buyruk önbelleği ve veri önbelleği olarak ikiye ayrılır (Harvard). Böylece üst seviyede esneklik, alt seviyede koşut erişim aynı anda elde edilir. Ayrıntılı tartışma Bölüm~\ref{chap:bellek}'dedir. Buyruk belleğinde çalışan programın makine diline çevrilmiş kodları bulunur ve buyruklar bellekten 32 bit olarak çekilir.

Veri belleği ise işlenecek verileri tutar. Bellekten getirilmesi gereken bir veri buradan alınarak yazmaca yazılır. Kaydedilmesi istenen veriler de buraya yazılır. Ayrıca veri belleği, yazmaçlarda yer kalmadığında geçici saklama alanı olarak da kullanılır.

\subsubsection{Yazmaç Öbeği}

Yazmaçlar, işlemci üzerinde AMB'nin işlem yapmak için veri alabileceği en yakın saklama birimleridir. Bu nedenle veri yolu tasarlanırken bir \emph{yazmaç öbeği} olarak bir araya getirilir. Yazmaçlar\index{yazmaç} kavramsal olarak saat sinyaliyle tetiklenen veri tutuculardır. Anlatım sadeliği için bu bölümde 32 bitlik bir yazmacı 32 kenar tetiklemeli D türü flip-flop\index{flip-flop}'la (Ek~\ref{ch:sayisal-tasarim}) modelleyeceğiz. Her flip-flop bir biti tutar ve saat kenarı geldiğinde girdisini çıktısına geçirir.

Çağdaş işlemcilerde ise yazmaç öbeği fiziksel olarak flip-flop dizisi değil, çok kapılı \emph{SRAM bit hücreleri}\index{SRAM} ile gerçekleştirilir. Her bit için altı transistörlü bir SRAM hücresi temel saklama elemanıdır. İki okuma kapısı ve bir yazma kapısı sağlamak için her hücreye ek erişim transistörleri eklenir (toplam genellikle hücre başına 8--10 transistör). Bu yaklaşım hem daha az alan kaplar hem de aynı çevrim içinde birden çok kapıdan erişimi doğal biçimde destekler. Flip-flop ve SRAM hücresi aynı işlevsel davranışı farklı transistör topolojileriyle sağlar. Bu bölümde flip-flop modelini tercih ediyoruz çünkü saat kenarına duyarlı davranış doğrudan görünür. SRAM hücresinin iç yapısı Bölüm~\ref{chap:bellek}'de ayrıntılı olarak incelenir.

\begin{bilginot}[x0 Yazmacının Özel Davranışı]
RISC-V'te x0 yazmacı donanımda sabit sıfıra bağlıdır (hardwired). Bu bağlantı sayesinde x0 her zaman 0 değerini döndürür ve x0'a yapılan yazma girişimleri donanım tarafından göz ardı edilir. İçeriği değişmez. Bu özellik, koşulsuz dallanmanın \texttt{jal x0, hedef} biçiminde, mutlak sıfırla karşılaştırmaların \texttt{beq rs1, x0, etiket} biçiminde derleyici tarafından ekonomik şekilde kodlanmasını sağlar.
\end{bilginot}

Yazmaç öbeğinin tasarımı Şekil~\ref{fig:yazmac-obegi}'da gösterilmektedir.

\begin{figure}[htbp]
\centering
\begin{tikzpicture}[>=Stealth, font=\footnotesize, scale=0.82, transform shape,
    yzm/.style={draw, fill=bloksari, minimum width=2.5cm, minimum height=0.45cm,
                font=\small, inner sep=2pt},
    cok/.style={draw, fill=blokmavi!40, rounded corners=2pt,
                minimum width=1.6cm, minimum height=1.0cm,
                font=\scriptsize, align=center}]

  % Yazmaç dizisi (orta)
  \node[yzm, pattern=north east lines, pattern color=gray!40] (x0) at (4, 0) {x0 (her zaman 0)};
  \foreach \i in {1,2,3,4,5,6,7} {
    \pgfmathsetmacro{\yval}{-\i * 0.5}
    \node[yzm] at (4, \yval) {x\i};
  }
  \node[font=\small] at (4, -4.2) {\vdots};
  \node[yzm] at (4, -4.9) {x30};
  \node[yzm] at (4, -5.4) {x31};

  % Sağda iki okuma çoklayıcısı
  \node[cok] (mux1) at (8, -1.0) {Çoklayıcı 1\\(32x1)};
  \node[cok] (mux2) at (8, -4.0) {Çoklayıcı 2\\(32x1)};

  % Solda yazma kod çözücüsü (5x32) — yazma izinlerini üretir
  \node[cok] (mux3) at (0, -3.5) {Kod Çözücü\\(5x32)};

  % Yazmaç dizisinden okuma çoklayıcılarına (sembolik 32-yollu kalın gri)
  \draw[line width=2.5pt, gray!50, ->] (5.25, -1.0) -- (mux1.west);
  \node[font=\tiny, text=gray!70, anchor=south] at (6.0, -0.9) {32 bit};
  \draw[line width=2.5pt, gray!50, ->] (5.25, -4.0) -- (mux2.west);
  \node[font=\tiny, text=gray!70, anchor=south] at (6.0, -3.9) {32 bit};

  % OkuAdres seçim girişleri (mux'ların üst/altından)
  \draw[->, thick] (8, 0.8) -- (mux1.north);
  \node[font=\footnotesize, anchor=south] at (8, 0.8) {OkuAdres1};
  \node[font=\tiny, anchor=west] at (8.1, 0.3) {5 bit};

  \draw[->, thick] (8, -5.8) -- (mux2.south);
  \node[font=\footnotesize, anchor=north] at (8, -5.8) {OkuAdres2};
  \node[font=\tiny, anchor=west] at (8.1, -5.3) {5 bit};

  % Okuma çıkışları
  \draw[->, thick] (mux1.east) -- ++(1.7, 0) node[right, font=\footnotesize] {Veri1 (32 bit)};
  \draw[->, thick] (mux2.east) -- ++(1.7, 0) node[right, font=\footnotesize] {Veri2 (32 bit)};

  % Kod çözücü girişleri (solda)
  \draw[->, thick] (-2.5, -3.5) -- (mux3.west);
  \node[font=\footnotesize, anchor=east] at (-2.5, -3.5) {YazAdres};
  \node[font=\tiny, anchor=south] at (-1.5, -3.5) {5 bit};

  \draw[->, thick] (0, -5.3) -- (mux3.south);
  \node[font=\footnotesize, anchor=north] at (0, -5.3) {YazmacaYaz};
  \node[font=\tiny, anchor=south west] at (0.05, -4.05) {Etkin.};

  % Kod çözücü çıkışı → yazmaç dizisi (sembolik 32 yazma izni)
  \draw[line width=2.5pt, gray!50, ->] (mux3.east) -- (2.75, -3.5);
  \node[font=\tiny, text=gray!70, anchor=south] at (1.9, -3.4) {32 izin};

  % YazVeri (32 bit) doğrudan yazmaç dizisine yayım hattı
  \draw[->, thick] (-2.5, -2.0) -- (2.75, -2.0);
  \node[font=\footnotesize, anchor=east] at (-2.5, -2.0) {YazVeri};
  \node[font=\tiny, anchor=south] at (-1.5, -2.0) {32 bit};
  \node[font=\tiny, anchor=south] at (0.6, -2.0) {(tüm yazmaçlara)};

  % Saat — tüm yazmaçların kenar tetikleme girişi
  \draw[->, thick] (4, -6.6) -- (4, -5.75);
  \node[font=\footnotesize, anchor=north] at (4, -6.6) {Saat};
\end{tikzpicture}
\caption[Yazmaç öbeğinin iç yapısı]{Yazmaç öbeğinin iç yapısı. İki okuma çoklayıcısı (32x1) okuma adreslerine göre seçtikleri yazmacın değerini Veri1 ve Veri2 çıkışlarına aktarır. Yazma tarafında YazVeri tüm yazmaçlara yayımlanır; kod çözücü (5x32) YazAdres'i 32 yazma iznine çevirir ve yalnızca seçilen yazmaç, YazmacaYaz=1 iken saat kenarında bu değeri yakalar. \texttt{x0} her zaman 0 değerini döndürür ve yazma girişimleri yoksayılır. Bu yapı, yazmaçlar ister D türü flip-flop ister SRAM bit hücresiyle gerçekleştirilsin aynıdır.}
\label{fig:yazmac-obegi}
\end{figure}

Şekil~\ref{fig:yazmac-obegi}'da belirtilen, yazmaçların önüne bağlanan kod çözücü\index{dekoder}\index{decoder|see{dekoder}} (bkz.\ Ek~\ref{ch:sayisal-tasarim}) ile gelen yazmaç adresine göre verinin kaçıncı yazmaca yazılacağı seçilir. Yazmaçlardan alınan verilerle AMB üzerinde işlem yapılır ve sonuç AMB'nin çıkışından belleğe ve yazmaç öbeğine bağlanır. Denetim biriminden gelecek sinyale göre sonuç, yazmaca ya da belleğe yazılır.

\begin{bilginot}[Gerçek Yazmaç Öbeklerinde Okuma Nasıl Yapılır?]
Şekil~\ref{fig:yazmac-obegi}'daki 32x1 okuma çoklayıcıları kavramsal bir modeldir. Hangi yazmacın okunacağını anlatmak için açıktır, ama çağdaş işlemciler okuma yolunu bu biçimde kurmaz. Gerçek yazmaç öbekleri SRAM bit hücrelerinden yapılır (Bölüm~\ref{chap:bellek}) ve okuma da yazma gibi bir \emph{kod çözücü} ile gerçekleştirilir. Okuma adresinin kod çözücüsü seçilen yazmacın \emph{sözcük hattını} (\emph{word line}) etkinleştirir. O satırdaki bit hücreleri sakladıkları değerleri ortak \emph{bit hatlarına} bırakır.

Bit hatlarındaki gerilim değişimi çok küçük ve yavaş geliştiği için, her bit hattının ucunda \emph{algılama yükselteci}\index{algılama yükselteci} (\emph{sense amplifier}) adı verilen analog bir devre bulunur. Bu devre küçük gerilim farkını hızla algılayıp tam mantık düzeyine (0 veya 1) yükseltir. Böylece 32 girişli kocaman bir çoklayıcının gecikmesine girilmeden okuma birkaç pikosaniyede tamamlanır. İki okuma kapısı için iki ayrı kod çözücü ve iki ayrı bit hattı kümesi kullanılır. Bu sayede iki yazmaç aynı çevrimde okunabilir.
\end{bilginot}

\section{Denetim Birimi}
\label{sec:denetim-birimi}

Veri yolunu bir orkestra olarak düşünürsek, denetim birimi bu orkestranın şefidir. Denetim sinyalleri\index{denetim sinyali}\index{control signal|see{denetim sinyali}} olmadan veri yolu doğru sonuç üretemez. Çoklayıcılar yanlış girişi seçer, yazmaçlara yanlış zamanda yazılır. Denetim birimi, yürütülen buyruğa bakarak veri yolundaki her bileşene ne yapması gerektiğini bildiren sinyalleri üretir.

İşlevsel olarak denetim birimi bir birleşik mantık bloğudur. Girişleri buyruğun işlem kodu (7 bit) ile gerektiğinde funct3 ve funct7 alanları, çıkışları ise önceki bölümde tanımladığımız yedi denetim sinyalidir (Şekil~\ref{fig:denetim-birimi-blok}). Dallanma buyruklarında AMB'nin \texttt{Sıfır} bayrağı da bir giriştir; çünkü dallanılıp dallanılmayacağı, koşulun sağlanıp sağlanmadığına bağlıdır. Birim, her işlem kodu örüntüsünü o buyruğun gerektirdiği sinyal değerlerine \emph{eşler}. Bu eşlemenin tamamı Bölüm~\ref{subsec:denetim-tablosu}'deki denetim tablosunda verilir. Bu birleşik mantık donanımda AND-OR kapı dizileriyle (örneğin \emph{programlanabilir mantık dizisi}, PLA) ya da bir arama belleğiyle gerçekleştirilebilir. İç tasarım ayrıntıları Ek~\ref{ch:sayisal-tasarim}'dedir.

\begin{figure}[htbp]
\centering
\begin{tikzpicture}[>=Stealth, font=\footnotesize]
  \node[draw, fill=blokmavi!30, rounded corners=3pt, minimum width=3cm, minimum height=4.8cm,
        font=\normalsize\bfseries, align=center] (db) at (0,0) {Denetim\\Birimi};

  % Girişler (sol)
  \foreach \y/\lbl/\bit in {1.6/{işlem kodu [6:0]}/7, 0.6/funct3/3, -0.4/{funct7[5]}/1, -1.4/{Sıfır (AMB)}/1} {
    \draw[->, thick] (-4.4, \y) -- (-1.5, \y);
    \node[anchor=east, font=\scriptsize] at (-4.4, \y) {\lbl};
    \node[anchor=south, font=\tiny, text=gray] at (-2.9, \y) {\bit};
  }

  % Çıkışlar (sağ): 7 denetim sinyali
  \foreach \y/\sig/\bit in {%
    1.8/YazmacaYaz/1, 1.2/AMBKaynak/1, 0.6/{GeriYazmaKaynağı}/2, 0.0/AnlıkBiçim/2,
    -0.6/BelleğeYaz/1, -1.2/Dallan/1, -1.8/AMBİşlem/2} {
    \draw[->, thick, sinyallacivert] (1.5, \y) -- (4.1, \y);
    \node[anchor=west, font=\scriptsize, text=sinyallacivert] at (4.15, \y) {\sig};
    \node[anchor=south, font=\tiny, text=gray] at (2.7, \y) {\bit};
  }
\end{tikzpicture}
\caption[Denetim biriminin blok gösterimi]{Denetim biriminin blok gösterimi. Birim, işlem kodu ile funct alanlarından (ve dallanma için AMB'nin \texttt{Sıfır} bayrağından) yedi denetim sinyalini üretir; her okun yanındaki sayı bit genişliğidir. Bu sinyallerin buyruğa göre aldığı değerler Bölüm~\ref{subsec:denetim-tablosu}'deki denetim tablosunda toplanmıştır.}
\label{fig:denetim-birimi-blok}
\end{figure}

Şekil~\ref{fig:denetim-birimi-blok}'deki giriş kümesi, bu bölümdeki sekiz buyruktan oluşan yalın tasarıma özgüdür. Veri yolundan denetim birimine geri dönen tek işaret \texttt{Sıfır} bayrağıdır ve o da yalnızca dallanma kararında kullanılır. Gerçek bir işlemcide denetim birimi çok daha fazla giriş alır. AMB'nin ürettiği diğer durum bayrakları (elde, taşma, işaret), bellek ya da veri yolunun ``hazır'' işaretleri, dış dünyadan gelen \emph{kesme}\index{kesme} (\emph{interrupt}) istekleri ve buyruk yürütümü sırasında ortaya çıkan \emph{kural dışı durum}\index{kural dışı durum} (\emph{exception}) sinyalleri bunlardan birkaçıdır. Bu girişler hem üretilen denetim sinyallerini hem de buyruk akışını değiştirebilir. Örneğin bir kesme geldiğinde PS, çalışan programdan kesme hizmet programına yönlendirilir (Bölüm~\ref{chap:giris-cikis}; ayrıcalık düzeyleri ve CSR yazmaçları Bölüm~\ref{chap:buyruk-kumesi}). Buyruk kümesi ve mimari yetenekler büyüdükçe hem giriş hem çıkış denetim sinyallerinin sayısı artar. Buradaki sadeleştirme, denetim biriminin temel çalışma ilkesini (buyruğu sinyallere eşleme) çıplak biçimde göstermek içindir.

\subsection{Denetim İşaretleri ve Denetim Tablosu}
\label{subsec:denetim-isaretleri}\label{subsec:denetim-tablosu}\index{denetim tablosu}\index{control table|see{denetim tablosu}}

Denetim biriminden çıkan işaretlerin her biri veri yolundaki belirli bir bileşeni denetler. Tablo~\ref{tab:denetim-birimleri}'de yedi denetim sinyali ve işlevleri verilmektedir.

\begin{table}[htbp]
\centering
\caption{Denetim sinyalleri ve işlevleri}
\label{tab:denetim-birimleri}
\begin{tabular}{lp{10cm}}
\toprule
\textbf{Sinyal} & \textbf{İşlev} \\
\midrule
YazmacaYaz & Yazmaç öbeğine yazma izni: 1 = yaz, 0 = yazma. \\
AMBKaynak & AMB'nin B girdisi: 0 = rs2 yazmacı, 1 = anlık değer. \\
GeriYazmaKaynağı & Yazmaç öbeğine yazılacak değerin kaynağı: 00 = AMB sonucu, 01 = veri belleği, 10 = PS+4. \\
AnlıkBiçim & Anlık değer üretecinin buyruk türü seçimi: 00 = I türü, 01 = S türü, 10 = B türü, 11 = J türü. \\
BelleğeYaz & Veri belleğine yazma izni: 1 = yaz, 0 = yazma. Okuma yan etkisiz olduğundan ayrı bir okuma izni gerekmez; veri belleği her çevrim adresindeki değeri çıkışına verir, GeriYazmaKaynağı bu değerin kullanılıp kullanılmayacağını belirler. \\
Dallan & PSKaynak çoklayıcısının seçim girdisi: 0 = PS+4 seç, 1 = atlama hedefi seç. beq/bne için Sıfır bayrağıyla birleştirilir. \\
AMBİşlem & AMB'nin yapacağı işlem: 00 = Topla, 01 = Çıkar, 10 = VEYA. \\
\bottomrule
\end{tabular}
\end{table}

\begin{bilginot}[RISC-V'te Sonuç Yazmacının Sabit Konumu]
RISC-V'te sonuç yazmacı ``rd'' her buyruk türünde aynı bit konumunda (bitler [11:7]) bulunduğu için, MIPS'teki ``SnçYzmç'' (rt/rd seçimi) denetim işaretine gerek yoktur. S ve B türü buyruklarda yazmaca yazma yapılmaz. Bu durum \texttt{YazmacaYaz} işaretiyle sağlanır.
\end{bilginot}

Bu yedi işaretin her buyruk için aldığı değerler, veri yolunu adım adım kurarken (Bölüm~\ref{sec:tek-vurusluk}) zaten belirlenmişti. Tablo~\ref{tab:tum-denetim} bu değerleri tek bir yerde toplar. Tablonun her satırı bir buyruk için yedi çıkış sinyalinin alacağı değeri verir ve böylece denetim biriminin işlevsel tanımını oluşturur.

\begin{table}[htbp]
\centering
\caption{Sekiz buyruk için tam denetim tablosu}
\label{tab:tum-denetim}
\small
\begin{tabular}{l|c|c|c|c|c|c|c}
\toprule
\textbf{Buyruk} & \textbf{Yazmaca} & \textbf{AMB} & \textbf{GeriYazma} & \textbf{Anlık} & \textbf{Belleğe} & \textbf{Dallan} & \textbf{AMB} \\
 & \textbf{Yaz} & \textbf{Kaynak} & \textbf{Kaynağı} & \textbf{Biçim} & \textbf{Yaz} & & \textbf{İşlem} \\
\midrule
add  & 1 & 0 & 00 & X  & 0 & 0      & 00 \\
sub  & 1 & 0 & 00 & X  & 0 & 0      & 01 \\
ori  & 1 & 1 & 00 & 00 & 0 & 0      & 10 \\
lw   & 1 & 1 & 01 & 00 & 0 & 0      & 00 \\
sw   & 0 & 1 & X  & 01 & 1 & 0      & 00 \\
beq  & 0 & 0 & X  & 10 & 0 & Sıfır  & 01 \\
bne  & 0 & 0 & X  & 10 & 0 & \textlnot{}Sıfır & 01 \\
jal  & 1 & X & 10 & 11 & 0 & 1      & X  \\
\bottomrule
\end{tabular}
\footnotesize \\[0.3em]
AMBKaynak: 0=rs2 yazmaç, 1=anlık değer. GeriYazmaKaynakı: 00=AMB sonucu, 01=veri belleği, 10=PS+4. \\
AnlıkBiçim: 00=I türü, 01=S türü, 10=B türü, 11=J türü. AMBİşlem: 00=Topla, 01=Çıkar, 10=VEYA. \\
add/sub aynı işlem kodunu (0110011) paylaşır; funct7{[5]} ayırır. beq/bne aynı işlem kodunu (1100011) paylaşır; funct3 ayırır.
\end{table}

Tablodaki X değeri, o sinyalin ilgili buyruk için \emph{önemsenmeyen} (\emph{don't care}) olduğunu gösterir. Sinyalin aldığı değer o çevrimde sonucu etkilemez. Örneğin \texttt{sw} buyruğunda yazmaç öbeğine yazma yapılmadığından \texttt{GeriYazmaKaynağı} çoklayıcısının seçtiği değer hiçbir yere ulaşmaz. Bu seçim önemsizdir. Benzer biçimde \texttt{jal} AMB'yi kullanmadığından \texttt{AMBİşlem} bu buyruk için X'tir. Önemsenmeyen durumlar, denetim mantığı kurulurken Karno haritalarında sadeleştirme için kullanılabilir.

Bir kural dışı durum vardır. Yazma izni sinyalleri \texttt{YazmacaYaz} ile \texttt{BelleğeYaz} hiçbir buyruk için X olamaz. Her zaman kesin bir değer (0 ya da 1) almak zorundadır. Bu iki sinyal \emph{kalıcı} bir durum değişikliği tetikler. X bırakılırsa donanım o çevrimde yazmaç öbeğine veya belleğe istenmeden yazabilir ve bir yazmacın ya da bellek konumunun içeriği geri döndürülemez biçimde bozulur. Çoklayıcı seçimleri ile okuma yolları yan etkisizdir, yanlış hesaplanan değer kullanılmadan atılabilir; ancak yapılan bir yazma geri alınamaz. Bu nedenle yazmanın gerçekleşmemesi gereken her buyrukta ilgili yazma izni X değil, açıkça 0 olarak belirlenir (tabloda \texttt{sw}, \texttt{beq}, \texttt{bne} için \texttt{YazmacaYaz}=0).

Bu denetim tablosu, işlemcinin tanıdığı buyruk kümesinin tam tanımıdır. Veri yolundaki bileşenler (yazmaç öbeği, AMB, bellekler, çoklayıcılar, anlık değer üreteci) hiç değişmeden, tabloya yeni bir satır eklenerek yeni bir buyruk tanımlanabilir. Yeni satır, var olan donanımı farklı bir denetim sinyali bileşimiyle yönlendirir ve böylece yeni bir işlem elde edilir. İşlemciye buyruk eklemek çoğu zaman donanımı büyütmeyi değil, denetim biriminin ürettiği sinyal örüntülerini genişletmeyi gerektirir.

Bunun bir sonucu da vardır. İki farklı buyruğun tablodaki satırı birebir aynı olamaz. İki satır bütün denetim sinyallerinde aynıysa, o iki buyruk veri yolunu tıpatıp aynı biçimde yönlendiriyor demektir. İşlemci açısından birbirinden ayırt edilemezler ve gerçekte tek bir buyruğun iki ayrı adıdır. Buyrukları birbirinden ayıran şey adları ya da söz dizimleri değil, denetim biriminde ürettikleri sinyal bileşimidir.

Bir sonraki adımda, her bir denetim birimi için çıkarılacak tablo sayesinde denetim birimi mantık kapılarıyla tasarlanacak hale gelir. Denetim biriminin girişlerine buyrukların işlem kodu (``opcode'') değerleri verilir. Bu değerler Tablo~\ref{tab:islem-kodu}'te özetlenmektedir. Denetim birimi bu işlem koduna göre buyruğun gerektirdiği denetim sinyallerini üretir.

\begin{table}[htbp]
\centering
\caption{İşlem kodu (``opcode'') değerleri}
\label{tab:islem-kodu}
\begin{tabular}{lc}
\toprule
\textbf{Buyruk} & \textbf{Opcode (ikilik, 7 bit)} \\
\midrule
Aritmetik--Mantık (R türü: add, sub) & 0110011 \\
Mantıksal VEYA Anlık (I türü: ori) & 0010011 \\
Yükle (I türü: lw) & 0000011 \\
Sakla (S türü: sw) & 0100011 \\
Dallanma (B türü: beq, bne) & 1100011 \\
jal (J türü) & 1101111 \\
\bottomrule
\end{tabular}
\end{table}

\begin{ornek}[6.2]
\textbf{(a)} \texttt{lw x6, 8(x10)} buyruğu tek vuruşluk veri yolunda yürütülürken tüm denetim sinyallerinin aldığı değerleri belirleyin ve buyruğun veri yolundan nasıl aktığını adım adım açıklayın.

\textbf{(b)} Veri yolundaki hiçbir bileşen değiştirilmeden, bir yazmacın değeri ile anlık değeri toplayıp sonucu yazmaca yazan \texttt{addi} buyruğunu denetim tablosuna eklemek istiyoruz. Bu buyruğun denetim satırını yazın ve \texttt{ori} satırından farkını belirtin.

\textbf{Çözüm.}

\emph{(a)} \texttt{lw} (yükle) buyruğu I türüdür; \texttt{rd}=x6, \texttt{rs1}=x10 ve anlık değer 8'dir. Buyruk, $x6 \leftarrow \text{Bel}[x10 + 8]$ işlemini gerçekleştirir. Denetim sinyalleri (Tablo~\ref{tab:tum-denetim}'deki \texttt{lw} satırı) ve her birinin işlevi:

\begin{center}
\begin{tabular}{lcp{7.2cm}}
\toprule
\textbf{Sinyal} & \textbf{Değer} & \textbf{İşlevi} \\
\midrule
YazmacaYaz & 1 & Bellekten okunan sözcük \texttt{x6} yazmacına yazılır \\
AMBKaynak & 1 & AMB'nin B girişine yazmaç yerine anlık değer (8) verilir \\
GeriYazmaKaynağı & 01 & Yazmaca yazılacak değer veri belleğinden gelir \\
AnlıkBiçim & 00 (I) & Anlık üreteç, buyruğun [31:20] bitlerinden 8 değerini üretir \\
BelleğeYaz & 0 & Belleğe yazma yok; veri belleği adresindeki sözcüğü okur \\
Dallan & 0 & Sıradaki buyruk: PS $\leftarrow$ PS+4 \\
AMBİşlem & 00 & Toplama: AMB adresi $x10 + 8$ olarak hesaplar \\
\bottomrule
\end{tabular}
\end{center}

Veri yolundan akış şöyle ilerler. Önce program sayacının gösterdiği adresten buyruk getirilir ve çözülür. Yazmaç öbeği \texttt{x10}'u okuyup değerini AMB'nin A girişine verir. \texttt{AMBKaynak}=1 olduğundan B girişine anlık üretecin ürettiği 8 gelir. \texttt{AMBİşlem}=00 ile AMB bu ikisini toplayarak bellek adresini ($x10+8$) üretir. Bu adres veri belleğine uygulanır. \texttt{BelleğeYaz}=0 olduğu için yazma yapılmaz ve bellek adresindeki sözcüğü çıkışına verir. \texttt{GeriYazmaKaynağı}=01 bu sözcüğü geri yazma çoklayıcısından geçirir, \texttt{YazmacaYaz}=1 ile değer \texttt{x6}'ya yazılır. Bu arada \texttt{Dallan}=0 olduğundan program sayacı PS+4 ile ilerler.

\emph{(b)} \texttt{addi rd, rs1, anlık} buyruğu, bir yazmacın değeriyle anlık değeri toplar ve sonucu yazmaca yazar. Bu işlemin gerektirdiği her bileşen veri yolunda hâlihazırda vardır: anlık üreteç I türü anlık değeri üretebilir, AMB toplama yapabilir (00) ve geri yazma çoklayıcısı AMB sonucunu yazmaca yönlendirebilir (00). Bu nedenle hiçbir bileşen eklemeden, tabloya yalnızca tek bir satır eklemek yeterlidir:

\begin{center}
\small
\begin{tabular}{l|c|c|c|c|c|c|c}
\toprule
\textbf{Buyruk} & \textbf{Yazmaca} & \textbf{AMB} & \textbf{GeriYazma} & \textbf{Anlık} & \textbf{Belleğe} & \textbf{Dallan} & \textbf{AMB} \\
 & \textbf{Yaz} & \textbf{Kaynak} & \textbf{Kaynağı} & \textbf{Biçim} & \textbf{Yaz} & & \textbf{İşlem} \\
\midrule
ori  & 1 & 1 & 00 & 00 & 0 & 0 & 10 \\
addi & 1 & 1 & 00 & 00 & 0 & 0 & 00 \\
\bottomrule
\end{tabular}
\end{center}

\texttt{addi} satırı, \texttt{ori} satırından yalnızca tek bir sinyalde ayrılır. \texttt{AMBİşlem}, \texttt{ori}'de VEYA (10) iken \texttt{addi}'de Toplama'dır (00). Geri kalan altı sinyal aynıdır, çünkü iki buyruk da I türüdür, anlık değer kullanır ve sonucu AMB'den yazmaca yazar. Satırları özdeş olmadığından bunlar işlemci açısından iki ayrı buyruktur. \texttt{addi} ile \texttt{ori} aynı işlem kodunu (0010011) paylaşır. Aralarındaki ayrım, çözme sırasında \texttt{funct3} alanına bakılarak \texttt{AMBİşlem} sinyalinin farklı üretilmesiyle yapılır. Bu durum, \texttt{add} ile \texttt{sub} buyruklarının aynı işlem kodunu paylaşıp \texttt{funct7} ile ayrılmasıyla aynı düzendir.
\end{ornek}

\subsection{AMB Denetimi}
\label{subsec:amb-denetimi}

Denetim tablosu, yedi çıkış sinyalinin her buyruk için alacağı değeri verir. Geriye bu değerleri donanıma, başka bir deyişle kapı devrelerine dönüştürmek kalır. Bu dönüşümü somut görmek için sinyallerden birini, \texttt{AMBİşlem}'i örnek alalım ve kapı düzeyinde nasıl tasarlandığını adım adım izleyelim. AMB'ye işlem kodunu gönderen bu alt birime AMB denetimi\index{AMB denetimi}\index{ALU control|see{AMB denetimi}} denir. Buyruk kümemizdeki sekiz buyruğun AMB'den beklediği işlem Tablo~\ref{tab:amb-islemleri}'te özetlenmiştir.

\begin{table}[htbp]
\centering
\caption{AMB işlemleri}
\label{tab:amb-islemleri}
\begin{tabular}{ll}
\toprule
\textbf{Buyruk} & \textbf{AMB işlemi} \\
\midrule
add & Toplama \\
sub & Çıkarma \\
ori & Mantıksal VEYA \\
lw & Toplama (adres hesabı) \\
sw & Toplama (adres hesabı) \\
beq / bne & Çıkarma (eşitlik testi) \\
jal & Önemsiz (adres hesabı Dal.Top.\ birimi yapar) \\
\bottomrule
\end{tabular}
\end{table}

AMB'nin yapacağı işlem 2 bitlik \textbf{AMBİşlem} sinyaliyle belirlenir. Denetim birimi bu sinyali doğrudan işlem kodu, funct3 ve funct7[5] bitlerinden üretir. Üç farklı işlem yeterlidir: 00 = Topla, 01 = Çıkar, 10 = VEYA. add ve sub aynı işlem kodunu (0110011) paylaşır. Aralarındaki fark funct7 alanının 5.\ biti (bit[30]) ile anlaşılır: add için bu bit 0, sub için 1'dir. beq ve bne aynı işlem kodunu (1100011) paylaşır. Aralarındaki fark funct3 alanıyla anlaşılır: beq için 000, bne için 001. ori ise I türü işlem kodu (0010011) altında funct3 = 110 ile tanımlanır.

\begin{figure}[htbp]
\centering
\begin{tikzpicture}[>=Stealth, node distance=1.5cm]
  % Giriş portları (sol)
  \node[font=\small] (opcode) at (0, 1.5) {opcode (7 bit)};
  \node[font=\small] (f3) at (0, 0) {funct3 (3 bit)};
  \node[font=\small] (f7) at (0, -1.5) {funct7{[5]} (1 bit)};
  % Birleşik mantık kutusu
  \node[blok, minimum width=3cm, minimum height=3cm, right=2.5cm of f3] (logik) {Birleşik\\Mantık};
  % Çıkış portu (sağ)
  \node[font=\small, right=2.5cm of logik] (cikis) {AMBİşlem (2 bit)};
  % Bağlantılar
  \draw[veriyolu] (opcode.east) -- ++(1.2,0) |- ([yshift=0.8cm]logik.west);
  \draw[veriyolu] (f3.east)     -- ++(0.5,0) -- ([yshift=0cm]logik.west);
  \draw[veriyolu] (f7.east)     -- ++(1.2,0) |- ([yshift=-0.8cm]logik.west);
  \draw[veriyolu] (logik.east)  -- (cikis.west);
  % Açıklama notları
  \node[font=\tiny, text=gray, right=0cm of cikis, yshift=-0.5cm]
       {00=Topla, 01=Çıkar, 10=VEYA};
\end{tikzpicture}
\caption[AMB Denetimi bloğu]{AMB Denetimi bloğu: işlem kodu, funct3 ve funct7{[5]} girişlerinden 2 bitlik AMBİşlem kodunu üretir.}
\label{fig:amb-denetimi-alt-blok}
\end{figure}

\subsubsection*{Birleşik mantığın iç tasarımı}

Şekil~\ref{fig:amb-denetimi-alt-blok}'deki ``Birleşik Mantık'' kutusunun içini açalım. Gerçek RISC-V'te girişler 7 bitlik işlem kodu ile funct alanlarıdır. Tam çözüm bunların hepsini birden ele alır ve çok sayıda kapı gerektirir. Tasarım yöntemini açıkça görmek için varsayımsal, basitleştirilmiş bir kodlama kullanalım. Sekiz buyruğa 3 bitlik birer kod ($b_2 b_1 b_0$) atadığımızı varsayalım. (Uygulamada bu kod, ana çözücünün işlem kodundan ürettiği bir ara gösterim olarak düşünülebilir.) Tablo~\ref{tab:amb-denetim-dogruluk} bu kodları ve her birinin gerektirdiği \texttt{AMBİşlem} değerini listeler.

\begin{table}[htbp]
\centering
\caption{Basitleştirilmiş kodlamayla AMB Denetimi doğruluk tablosu}
\label{tab:amb-denetim-dogruluk}
\begin{tabular}{llcc}
\toprule
\textbf{Buyruk} & \textbf{Kod} ($b_2 b_1 b_0$) & \textbf{AMBİşlem} & \textbf{İşlem} \\
\midrule
add & 000 & 00 & Topla \\
sub & 001 & 01 & Çıkar \\
ori & 010 & 10 & VEYA \\
lw  & 011 & 00 & Topla \\
sw  & 100 & 00 & Topla \\
beq & 101 & 01 & Çıkar \\
bne & 110 & 01 & Çıkar \\
jal & 111 & XX & (önemsiz) \\
\bottomrule
\end{tabular}
\end{table}

\texttt{jal} buyruğu AMB'yi kullanmadığından çıkışı önemsizdir (x). Bu serbestlik Karno haritasında sadeleştirmeye yardımcı olur. Tablodaki değerleri, her çıkış biti için birer Karno haritasına yerleştirip komşu 1 değerlerini (uygun olduğunda önemsiz hücreyi de) gruplarız (Şekil~\ref{fig:amb-karno}; Karno haritası yöntemi Ek~\ref{ch:sayisal-tasarim}'de ayrıntılı anlatılmıştır).

\begin{figure}[htbp]
\centering
\begin{tikzpicture}[scale=0.82, font=\footnotesize]
  % ===== Harita 1: AMBİşlem[1] =====
  \begin{scope}
    \node[font=\small\bfseries] at (2, 3.65) {\texttt{AMBİşlem}[1]};
    \draw[thick] (0,0) grid (4,2);
    \node[above, font=\scriptsize] at (0.5,2) {00};
    \node[above, font=\scriptsize] at (1.5,2) {01};
    \node[above, font=\scriptsize] at (2.5,2) {11};
    \node[above, font=\scriptsize] at (3.5,2) {10};
    \node[left, font=\scriptsize] at (0,1.5) {0};
    \node[left, font=\scriptsize] at (0,0.5) {1};
    \node[above left, font=\scriptsize] at (0,2) {$b_2$\textbackslash$b_1b_0$};
    \node at (0.5,1.5) {0}; \node at (1.5,1.5) {0}; \node at (2.5,1.5) {0}; \node at (3.5,1.5) {1};
    \node at (0.5,0.5) {0}; \node at (1.5,0.5) {0}; \node at (2.5,0.5) {x}; \node at (3.5,0.5) {0};
    % değişken bölgeleri (Ek B ile aynı biçim)
    \draw[decorate, decoration={brace, amplitude=4pt}, thick] (-0.55,0) -- (-0.55,1) node[midway, left=4pt, font=\scriptsize] {$b_2$};
    \draw[decorate, decoration={brace, amplitude=4pt}, thick] (2,2.5) -- (4,2.5) node[midway, above=3pt, font=\scriptsize] {$b_1$};
    \draw[decorate, decoration={brace, amplitude=4pt, mirror}, thick] (1,-0.2) -- (3,-0.2) node[midway, below=3pt, font=\scriptsize] {$b_0$};
    \draw[rounded corners=3pt, very thick, sinyalkirmizi] (3.06,1.06) rectangle (3.94,1.94);
    \node[sinyalkirmizi, font=\scriptsize] at (2,-1.05) {$\overline{b_2}\,b_1\,\overline{b_0}$ \ (\texttt{ori})};
  \end{scope}
  % ===== Harita 2: AMBİşlem[0] =====
  \begin{scope}[xshift=7.5cm]
    \node[font=\small\bfseries] at (2, 3.65) {\texttt{AMBİşlem}[0]};
    \draw[thick] (0,0) grid (4,2);
    \node[above, font=\scriptsize] at (0.5,2) {00};
    \node[above, font=\scriptsize] at (1.5,2) {01};
    \node[above, font=\scriptsize] at (2.5,2) {11};
    \node[above, font=\scriptsize] at (3.5,2) {10};
    \node[left, font=\scriptsize] at (0,1.5) {0};
    \node[left, font=\scriptsize] at (0,0.5) {1};
    \node[above left, font=\scriptsize] at (0,2) {$b_2$\textbackslash$b_1b_0$};
    \node at (0.5,1.5) {0}; \node at (1.5,1.5) {1}; \node at (2.5,1.5) {0}; \node at (3.5,1.5) {0};
    \node at (0.5,0.5) {0}; \node at (1.5,0.5) {1}; \node at (2.5,0.5) {x}; \node at (3.5,0.5) {1};
    % değişken bölgeleri (Ek B ile aynı biçim)
    \draw[decorate, decoration={brace, amplitude=4pt}, thick] (-0.55,0) -- (-0.55,1) node[midway, left=4pt, font=\scriptsize] {$b_2$};
    \draw[decorate, decoration={brace, amplitude=4pt}, thick] (2,2.5) -- (4,2.5) node[midway, above=3pt, font=\scriptsize] {$b_1$};
    \draw[decorate, decoration={brace, amplitude=4pt, mirror}, thick] (1,-0.2) -- (3,-0.2) node[midway, below=3pt, font=\scriptsize] {$b_0$};
    % grup A: sütun 01 (dikey) → i1'·i0
    \draw[rounded corners=3pt, very thick, sinyalmavi] (1.06,0.06) rectangle (1.94,1.94);
    \node[sinyalmavi, font=\scriptsize] at (0.7,-1.05) {$\overline{b_1}\,b_0$};
    % grup B: alt satır 11+10 (önemsiz hücre dahil) → i2·i1
    \draw[rounded corners=3pt, very thick, blokyesil!70!black] (2.06,0.06) rectangle (3.94,0.94);
    \node[blokyesil!70!black, font=\scriptsize] at (3.2,-1.05) {$b_2\,b_1$};
  \end{scope}
\end{tikzpicture}
\caption[AMB Denetimi çıkış bitleri için Karno haritaları]{AMB Denetimi çıkış bitleri için Karno haritaları (satır $b_2$, sütun $b_1b_0$ Gray kodu sırasıyla). \texttt{jal} (kod 111) önemsiz hücredir (x); AMBİşlem[0] haritasında bu hücre alttaki grupla birlikte alınarak $b_2 b_1$ terimini sadeleştirir, AMBİşlem[1] haritasında ise gruba katılmaz.}
\label{fig:amb-karno}
\end{figure}

Gruplama sonucunda iki çıkış biti şu yalın ifadelere iner:
\begin{align*}
\text{AMBİşlem}[1] &= \overline{b_2}\cdot b_1 \cdot \overline{b_0} \\
\text{AMBİşlem}[0] &= \overline{b_1}\cdot b_0 \;+\; b_2 \cdot b_1
\end{align*}
İlk ifade yalnızca \texttt{ori} için (kod 010), ikincisi ise çıkarma yapan üç buyruk (\texttt{sub}, \texttt{beq}, \texttt{bne}) için 1 olur.

Gerçek RISC-V tasarımlarında bu mantık çoğunlukla iki aşamalı kurulur. Ana çözücü, işlem kodundan buyruğun hangi sınıfa ait olduğunu (aritmetik--mantık, bellek erişimi ya da dallanma) belirten birkaç bitlik bir ara sinyal üretir. AMB denetimi ise bu ara sinyali \texttt{funct3} ve \texttt{funct7} alanlarıyla birleştirerek nihai \texttt{AMBİşlem} kodunu seçer. Bu iki aşamalı bölüştürme iki yarar sağlar. Ana çözücü buyruk türü ayrıntılarından arınmış kalır ve AMB'ye sonradan yeni bir işlem eklendiğinde yalnızca ikinci aşamanın güncellenmesi yeter. Yukarıdaki üç bitlik basitleştirilmiş kod, işte bu ara sinyalin yerini tutmakta ve yöntemin özünü tek bir küçük örnekte göstermektedir.

\begin{figure}[htbp]
\centering
\begin{tikzpicture}[kapiboy, large circuit symbols, scale=1.0, transform shape]
    % Tüm VE kapıları 3 girişli çizilir (eşit boy); 2 girişli işlevler
    % üst (giriş 1) ile alt (giriş 3) girişi kullanır, orta giriş boş kalır.
    % Doğrudan sinyaller DEĞİL kapılarının ARASINDAKİ boşluklardan saptırılır;
    % tüm dikey yollar DEĞİL kapılarının SAĞINDA toplanır (kapı içinden tel geçmez).

    \coordinate (xin) at (-0.4, 0);   % giriş etiketleri x
    \coordinate (xnot) at (1.6, 0);   % DEĞİL kapıları x

    % --- VE kapıları (hepsi 3 girişli; eşit boy; sola çekildi -> dar şekil) ---
    \node[and gate US, draw, thick, fill=blokgri, inputs={nnn}] (and3) at (5.0, 3.6) {};
    \node[and gate US, draw, thick, fill=blokgri, inputs={nnn}] (anda) at (5.0, 1.9) {};
    \node[and gate US, draw, thick, fill=blokgri, inputs={nnn}] (andb) at (5.0, 0.2) {};

    % --- VEYA kapısı (3 girişli; VE kapılarıyla aynı boy, orta giriş boş; sola çekildi) ---
    \node[or gate US, draw, thick, fill=blokgri, inputs={nnn}] (or1) at (7.2, 1.05) {};

    % --- DEĞİL kapıları: ni2 ve ni1 ilgili VE giriş hizasına çekildi (düz tel) ---
    \node[not gate US, draw, thick, fill=blokgri] (ni2) at (xnot |- and3.input 1) {};
    \node[not gate US, draw, thick, fill=blokgri] (ni1) at (xnot |- anda.input 1) {};
    \node[not gate US, draw, thick, fill=blokgri] (ni0) at (1.6, 0.4) {};

    % --- Giriş düğümleri (DEĞİL kapılarıyla hizalı) ---
    \coordinate (ini2) at (xin |- ni2);
    \coordinate (ini1) at (xin |- ni1);
    \coordinate (ini0) at (xin |- ni0);
    \node[left] at (ini2) {$b_2$};
    \node[left] at (ini1) {$b_1$};
    \node[left] at (ini0) {$b_0$};

    % --- Girişlerden DEĞİL kapılarına; doğrudan dal için fan-out noktası ---
    % (taplar farklı x'te: dikey hatlar üst üste binmesin; i0 dikeyi i1 yatayının solunda)
    \draw (ini2) -- ++(0.3,0) coordinate (t2); \filldraw (t2) circle (1.6pt);
    \draw (t2) -- (ni2.input);
    \draw (ini1) -- ++(0.9,0) coordinate (t1); \filldraw (t1) circle (1.6pt);
    \draw (t1) -- (ni1.input);
    \draw (ini0) -- ++(0.5,0) coordinate (t0); \filldraw (t0) circle (1.6pt);
    \draw (t0) -- (ni0.input);

    % === i2' -> and3 giriş 1 : DÜZ (hizalı, dirseksiz) ===
    \draw (ni2.output) -- (and3.input 1);

    % === i1' -> anda giriş 1 : DÜZ (hizalı, dirseksiz) ===
    \draw (ni1.output) -- (anda.input 1);

    % === i0' -> and3 giriş 3 : x=2.4 dikey yolu (DEĞİL sağında) ===
    \draw (ni0.output) -- (2.4,0.4) coordinate (b0);
    \draw (b0) -- (b0 |- and3.input 3) -- (and3.input 3);

    % === i1 (doğrudan) -> and3 giriş 2 + andb giriş 3 : x=2.7 yolu, alt boşluktan ===
    \draw (t1) |- (2.7,1.2) coordinate (i1y);
    \filldraw (i1y) circle (1.6pt);
    \draw (i1y) -- (i1y |- and3.input 2) -- (and3.input 2);
    \draw (i1y) -- (i1y |- andb.input 3) -- (andb.input 3);

    % === i2 (doğrudan) -> andb giriş 1 : noktadan AŞAĞI, ni2-ni1 arasından x=3.0 yoluna ===
    \draw (t2) -- (-0.1,3.0) -- (3.0,3.0) coordinate (b2);
    \draw (b2) -- (b2 |- andb.input 1) -- (andb.input 1);

    % === i0 (doğrudan) -> anda giriş 3 : noktadan YUKARI tam giriş hizasına, sonra DÜZ ===
    % (dikey kısım t0.x'te; dirsek anda.input 3 hizasında -> kapıya dirseksiz girer.
    %  i1 yatay hattı x=0.9'dan başladığı için bu dikey hattı kesmez.)
    \draw (t0) -- (t0 |- anda.input 3) -- (anda.input 3);

    % === VE çıkışları -> VEYA (giriş 1 ve 3; orta boş) ===
    \draw (anda.output) -- ++(0.4,0) |- (or1.input 1);
    \draw (andb.output) -- ++(0.4,0) |- (or1.input 3);

    % === Çıkış etiketleri ===
    \draw[okstil] (and3.output) -- ++(1.3,0) node[right, font=\footnotesize] {\texttt{AMBİşlem}[1]};
    \draw[okstil] (or1.output)  -- ++(1.3,0) node[right, font=\footnotesize] {\texttt{AMBİşlem}[0]};
\end{tikzpicture}
\caption[AMB Denetimi birleşik mantığının kapı devresi]{AMB Denetimi birleşik mantığının kapı devresi. Üç bitlik buyruk kodundan ($b_2 b_1 b_0$) iki bitlik AMBİşlem üretilir; tümleyenler DEĞİL kapılarıyla elde edilir. \texttt{AMBİşlem}[1]$=\overline{b_2}\,b_1\,\overline{b_0}$, \texttt{AMBİşlem}[0]$=\overline{b_1}\,b_0+b_2\,b_1$.}
\label{fig:amb-denetimi-devre}
\end{figure}

Şekil~\ref{fig:amb-denetimi-devre} bu iki ifadeyi gerçekleştiren kapı devresini gösterir. $b_2$, $b_1$ ve $b_0$ girişlerinin tümleyenleri DEĞİL kapılarıyla üretilir. \texttt{AMBİşlem}[1] tek bir üç girişli VE kapısıyla, \texttt{AMBİşlem}[0] ise iki VE kapısı ile bir VEYA kapısıyla elde edilir. Gerçek bir işlemcide aynı yöntem (doğruluk tablosu $\rightarrow$ Karno haritası $\rightarrow$ kapılar) tam işlem kodu ve funct alanlarına uygulanır. Yalnızca giriş sayısı arttığından devre büyür, ilke değişmez.

\subsection{Diğer Denetim Sinyallerinin Tasarımı}
\label{subsec:ana-denetim}

\texttt{AMBİşlem} için izlediğimiz yol (doğruluk tablosu $\rightarrow$ Karno haritası $\rightarrow$ kapı devresi) geri kalan altı denetim sinyali için de aynen geçerlidir. Bu sinyallerin tümü, AMB denetiminde kullandığımız aynı üç bitlik buyruk kodundan ($b_2 b_1 b_0$; Tablo~\ref{tab:amb-denetim-dogruluk}) türetilir. Tablo~\ref{tab:ana-denetim-dogruluk} bu sinyallerden beşinin her buyruk için aldığı değerleri verir. Altıncı sinyal olan \texttt{Dallan}'ı çalışma anına bağlı olduğundan ayrıca ele alacağız.

\begin{table}[htbp]
\centering
\caption{Ana denetim sinyalleri için doğruluk tablosu}
\label{tab:ana-denetim-dogruluk}
\small
\begin{tabular}{cl|c|c|c|c|c}
\toprule
\textbf{Kod} & \textbf{Buyruk} & \textbf{Yazmaca} & \textbf{AMB} & \textbf{GeriYazma} & \textbf{Anlık} & \textbf{Belleğe} \\
$b_2b_1b_0$ & & \textbf{Yaz} & \textbf{Kaynak} & \textbf{Kaynağı} & \textbf{Biçim} & \textbf{Yaz} \\
\midrule
000 & add  & 1 & 0 & 00 & X  & 0 \\
001 & sub  & 1 & 0 & 00 & X  & 0 \\
010 & ori  & 1 & 1 & 00 & 00 & 0 \\
011 & lw   & 1 & 1 & 01 & 00 & 0 \\
100 & sw   & 0 & 1 & X  & 01 & 1 \\
101 & beq  & 0 & 0 & X  & 10 & 0 \\
110 & bne  & 0 & 0 & X  & 10 & 0 \\
111 & jal  & 1 & X & 10 & 11 & 0 \\
\bottomrule
\end{tabular}
\end{table}

Her çıkış biti, bu tablodaki değerlerin Karno haritasıyla sadeleştirilmesiyle bir Boole ifadesine iner. Örnek olarak \texttt{YazmacaYaz} sinyalini ele alalım. Haritası Şekil~\ref{fig:yazmacayaz-karno}'te verilmiştir.

\begin{figure}[htbp]
\centering
\begin{tikzpicture}[scale=0.82, font=\footnotesize]
    \node[font=\small\bfseries] at (2, 3.65) {\texttt{YazmacaYaz}};
    \draw[thick] (0,0) grid (4,2);
    \node[above, font=\scriptsize] at (0.5,2) {00};
    \node[above, font=\scriptsize] at (1.5,2) {01};
    \node[above, font=\scriptsize] at (2.5,2) {11};
    \node[above, font=\scriptsize] at (3.5,2) {10};
    \node[left, font=\scriptsize] at (0,1.5) {0};
    \node[left, font=\scriptsize] at (0,0.5) {1};
    \node[above left, font=\scriptsize] at (0,2) {$b_2$\textbackslash$b_1b_0$};
    \node at (0.5,1.5) {1}; \node at (1.5,1.5) {1}; \node at (2.5,1.5) {1}; \node at (3.5,1.5) {1};
    \node at (0.5,0.5) {0}; \node at (1.5,0.5) {0}; \node at (2.5,0.5) {1}; \node at (3.5,0.5) {0};
    \draw[decorate, decoration={brace, amplitude=4pt}, thick] (-0.55,0) -- (-0.55,1) node[midway, left=4pt, font=\scriptsize] {$b_2$};
    \draw[decorate, decoration={brace, amplitude=4pt}, thick] (2,2.5) -- (4,2.5) node[midway, above=3pt, font=\scriptsize] {$b_1$};
    \draw[decorate, decoration={brace, amplitude=4pt, mirror}, thick] (1,-0.2) -- (3,-0.2) node[midway, below=3pt, font=\scriptsize] {$b_0$};
    % grup 1: üst satır tamamı -> i2'
    \draw[rounded corners=3pt, very thick, sinyalmavi] (0.06,1.06) rectangle (3.94,1.94);
    \node[sinyalmavi, font=\scriptsize] at (0.9,-1.05) {$\overline{b_2}$};
    % grup 2: 11 sütunu -> i1 i0
    \draw[rounded corners=3pt, very thick, blokyesil!70!black] (2.06,0.06) rectangle (2.94,1.94);
    \node[blokyesil!70!black, font=\scriptsize] at (3.0,-1.05) {$b_1 b_0$};
\end{tikzpicture}
\caption[\texttt{YazmacaYaz} sinyali için Karno haritası]{\texttt{YazmacaYaz} sinyali için Karno haritası. Üst satırın tamamı $\overline{b_2}$ terimini, 11 sütunu ise $b_1 b_0$ terimini verir.}
\label{fig:yazmacayaz-karno}
\end{figure}

Haritadaki iki grup \texttt{YazmacaYaz}$=\overline{b_2}+b_1 b_0$ ifadesini verir: üst satırın tamamı (\texttt{add}, \texttt{sub}, \texttt{ori}, \texttt{lw}) $\overline{b_2}$ ile, alttaki tek 1 değeri (\texttt{jal}) ise $b_1 b_0$ ile kapsanır. Aynı yöntem kalan dört sinyale (ve iki bitlik sinyallerin her bir bitine) uygulandığında şu ifadeler elde edilir:
\begin{align*}
\text{YazmacaYaz} &= \overline{b_2} + b_1 b_0 \\
\text{AMBKaynak} &= \overline{b_2}\,b_1 + b_2\,\overline{b_1}\,\overline{b_0} \\
\text{GeriYazmaKaynağı}[1] &= b_2 \\
\text{GeriYazmaKaynağı}[0] &= \overline{b_2}\,b_1\,b_0 \\
\text{AnlıkBiçim}[1] &= b_2\,(b_1 + b_0) \\
\text{AnlıkBiçim}[0] &= \overline{b_1}\,\overline{b_0} + b_2\,b_1\,b_0 \\
\text{BelleğeYaz} &= b_2\,\overline{b_1}\,\overline{b_0}
\end{align*}
Bu ifadelerin sadeliği, sekiz buyruktan oluşan küçük kümede her sinyalin yalnızca birkaç buyruk için etkin olmasından kaynaklanır. \texttt{add} ile \texttt{sub} buyruklarında anlık değer kullanılmadığından \texttt{AnlıkBiçim} bu iki buyruk için önemsizdir (X); bu serbestlik yukarıdaki ifadelerin sadeleşmesinde kullanılmıştır.

\texttt{Dallan} sinyali bu altısından ayrılır. Değeri yalnızca buyruğa değil, çalışma anında AMB'nin ürettiği \texttt{Sıfır} bayrağına da bağlıdır. Tablo~\ref{tab:tum-denetim}'de \texttt{beq} için \texttt{Sıfır}, \texttt{bne} için $\lnot$\texttt{Sıfır} yazılmasının nedeni budur. Ana çözücü bu sinyali iki ara işaret üzerinden üretir: koşulsuz atlama yapan \texttt{jal} için \emph{Atla} ve koşullu dallanan \texttt{beq}/\texttt{bne} için \emph{Dal}:
\begin{align*}
\text{Atla} &= b_2\,b_1\,b_0 \\
\text{Dal} &= b_2\,(b_1 \oplus b_0)
\end{align*}
Program sayacı kaynağını seçen \texttt{Dallan} sinyali, bu iki işareti dallanma koşuluyla birleştirir:
\[
\text{Dallan} = \text{Atla} + \text{Dal}\cdot(\text{Sıfır} \oplus b_1).
\]
Buradaki $\text{Sıfır} \oplus b_1$ çarpanı, dallanmanın \texttt{beq} ($b_1=0$) için \texttt{Sıfır}$=1$ olduğunda, \texttt{bne} ($b_1=1$) için \texttt{Sıfır}$=0$ olduğunda gerçekleşmesini sağlar. Gerçek RISC-V'te \texttt{beq} ile \texttt{bne} arasındaki bu yön farkı \texttt{funct3} alanından okunur.

Bu bölümde ifadeleri verilen denetim sinyallerinin her birini üreten kapı devresi, \texttt{AMBİşlem} için Şekil~\ref{fig:amb-denetimi-devre}'te izlenen yöntemle doğrudan çizilebilir. Bu çizimler bölüm sonu alıştırmalarına bırakılmıştır.

\subsection{Denetim Biriminin Birleştirilmesi}
\label{subsec:denetim-birlestirme}

Böylece denetim biriminin her iki parçasını da tasarlamış olduk. Ana çözücü, buyruğun işlem kodundan altı denetim sinyalini üretir (Bölüm~\ref{subsec:ana-denetim}). AMB denetimi ise işlem kodunu \texttt{funct3} ve \texttt{funct7}[5] alanlarıyla birleştirerek \texttt{AMBİşlem} sinyalini belirler (Bölüm~\ref{subsec:amb-denetimi}). Tam denetim birimi bu iki birleşik mantık öbeğinin yan yana getirilmesinden oluşur. Girişleri buyruğun işlem kodu ile funct alanları, çıkışları ise veri yolundaki her bileşeni yönlendiren yedi denetim sinyalidir. Her sinyal, bu bölümde türetilen Boole ifadesini gerçekleştiren küçük bir kapı öbeğiyle üretilir.

Denetim sinyallerinin kapı düzeyindeki bu ayrıntısının üzerine çıkıldığında, tasarımı yalnızca verinin yazmaçlar ile birimler arasında nasıl aktığı ve hangi işlemlerden geçtiği düzeyinde betimlemek de olanaklıdır. Bu soyutlama seviyesine \emph{Yazmaç Aktarım Seviyesi}\index{yazmaç aktarım seviyesi}\index{register transfer level|see{yazmaç aktarım seviyesi}} (YAS) denir.

Denetim birimi ve veri yoluyla birlikte tam olarak tamamlanmış tek vuruşluk işlemci Şekil~\ref{fig:tek-vurusluk-tam}'da verilmektedir.

\begin{figure}[htbp]
\centering
\begin{tikzpicture}[>=Stealth, semithick, font=\footnotesize,
    yzm/.style={draw, fill=yellow!25, minimum width=1cm, minimum height=1cm, align=center},
    bllk/.style={draw, fill=green!20, minimum width=1.8cm, minimum height=1.6cm, align=center},
    yobk/.style={draw, fill=blue!10, minimum width=2.6cm, minimum height=3cm},
    amb/.style={draw, fill=red!20, minimum width=1.6cm, minimum height=2cm, align=center},
    artop/.style={draw, fill=gray!15, minimum width=1.1cm, minimum height=0.7cm, align=center},
    anlikb/.style={draw, fill=orange!25, minimum width=1.4cm, minimum height=0.9cm, align=center},
    muxdik/.style={draw, fill=purple!15, minimum width=0.7cm, minimum height=1.0cm, align=center},
    muxyatay/.style={draw, fill=purple!15, minimum width=1.0cm, minimum height=0.7cm, align=center},
    vbll/.style={draw, fill=green!20, minimum width=1.8cm, minimum height=2.2cm, align=center},
    bitayar/.style={fill=white, inner sep=1pt, font=\tiny},
    yeni/.style={}
]
  % =================================================================
  % VERİ YOLU (fig:vy-jal ile birebir; yeni/.style={} -- kırmızı vurgu yok)
  % =================================================================

  % --- Sol PS güncelleme bloğu ---
  \node[bllk] (bb) at (0, 4.6) {Buyruk\\Belleği};
  \node[yzm] (ps) at (0, 2.7) {PS};
  \node[muxyatay] (psmux) at (0, 1.3) {};
  \node[font=\tiny] at (psmux.center) {Çok.};
  \node[font=\tiny, anchor=south] at ([xshift=-0.3cm, yshift=0.02cm]psmux.south) {0};
  \node[font=\tiny, anchor=south] at ([xshift=0.3cm, yshift=0.02cm]psmux.south) {1};

  \node[artop] (top4) at (-1.3, 0) {Top.};
  \node[artop, minimum width=1.2cm, minimum height=0.7cm] (daladd) at (1.3, 0) {Dal.\\Top.};

  \draw[->] (ps.north) -- (bb.south) node[midway, right=1pt, font=\scriptsize] {adres};
  \draw[->] (psmux.north) -- (ps.south);

  \draw[->] (ps.west) -- (-1.5, 2.7) -- ([xshift=-0.2cm]top4.north);
  \draw[->] (ps.east) -- (1.1, 2.7) -- ([xshift=-0.2cm]daladd.north);

  \node[font=\scriptsize] at (-1.1, 0.85) {$4$};
  \draw[->] (-1.1, 0.7) -- ([xshift=0.2cm]top4.north);

  % B/J anlık -> Dal.Top.north_east (anlık.east teli üzerinden tap)
  \coordinate (anlikdal) at (6.85, 1.0);
  \node[circle, fill=black, inner sep=1pt] at (anlikdal) {};
  \draw[->] (anlikdal) -- (1.5, 1.0) -- ([xshift=0.2cm]daladd.north);

  % Top4 çıkışı (PS+4): tap ile hem PSKaynak'a hem GeriYazma'ya gider
  \draw[->] (top4.south) -- (-1.3, -0.7) -- (-0.3, -0.7) -- ([xshift=-0.3cm]psmux.south);
  \node[font=\scriptsize, anchor=north] at (-0.8, -0.75) {PS$+$4};
  \draw[->] (daladd.south) -- (1.3, -0.7) -- (0.3, -0.7) -- ([xshift=0.3cm]psmux.south);
  \node[font=\scriptsize, anchor=north] at (0.8, -0.75) {Dal.hedef};

  % Dallan stub -- MAVİ (Denetim Birimi çıkışı; bileşende etiketli)
  \draw[->, sinyallacivert, thick] ([xshift=-0.4cm]psmux.west) -- (psmux.west);
  \node[font=\scriptsize, text=sinyallacivert, anchor=south] at ([xshift=-0.5cm, yshift=0.05cm]psmux.west) {Dallan};

  \coordinate (parse) at (2.5, 4.6);
  \draw[->] (bb.east) -- (parse) node[midway, below, font=\scriptsize] {buyruk} node[bitayar, pos=0.55] {32 bit};

  % R türü buyruk şeridi
  \coordinate (seritsol) at (3.0, 6.4);
  \draw[gray, dashed, semithick] (parse) -- (seritsol);
  \draw[gray, dashed, semithick] (parse) -- ([xshift=6cm]seritsol);
  \begin{scope}[shift={(seritsol)}]
    \draw[fill=gray!10]  (0,0)      rectangle (1.3125,0.55) node[pos=.5, font=\tiny] {funct7};
    \draw[fill=green!18] (1.3125,0) rectangle (2.25,0.55)   node[pos=.5, font=\tiny] {rs2};
    \draw[fill=green!18] (2.25,0)   rectangle (3.1875,0.55) node[pos=.5, font=\tiny] {rs1};
    \draw[fill=gray!10]  (3.1875,0) rectangle (3.75,0.55)   node[pos=.5, font=\tiny] {f3};
    \draw[fill=blue!18]  (3.75,0)   rectangle (4.6875,0.55) node[pos=.5, font=\tiny] {rd};
    \draw[fill=gray!10]  (4.6875,0) rectangle (6,0.55)      node[pos=.5, font=\tiny] {opcode};
    \node[font=\tiny, anchor=north] at (0,0)      {31};
    \node[font=\tiny, anchor=north] at (1.3125,0) {25};
    \node[font=\tiny, anchor=north] at (2.25,0)   {20};
    \node[font=\tiny, anchor=north] at (3.1875,0) {15};
    \node[font=\tiny, anchor=north] at (3.75,0)   {12};
    \node[font=\tiny, anchor=north] at (4.6875,0) {7};
    \node[font=\tiny, anchor=north] at (6,0)      {0};
  \end{scope}
  \node[font=\scriptsize, anchor=south] at ([xshift=2.06cm, yshift=0.6cm]seritsol) {R türü buyruk};

  \node[yobk] (yo) at (5.4, 2.6) {};
  \node[anchor=center, font=\footnotesize\bfseries, align=center] at (yo.center) {Yazmaç\\Öbeği};

  % R türü: rs1 -> A_adr, rs2 -> B_adr, rd -> C_adr
  \draw[->] ([xshift=2.71875cm]seritsol) -- ++(0,-1.0) -| ([xshift=-0.6cm]yo.north) node[bitayar, pos=0.2]{5 bit};
  \draw[->] ([xshift=1.78125cm]seritsol) -- ++(0,-1.4) -| ([xshift=0cm]yo.north);
  \draw[->] ([xshift=4.21875cm]seritsol) -- ++(0,-0.6) -| ([xshift=0.6cm]yo.north);

  \node[font=\scriptsize, anchor=north] at ([xshift=-0.6cm,yshift=-0.05cm]yo.north) {$A_{\text{adr}}$};
  \node[font=\scriptsize, anchor=north] at ([xshift=0cm,yshift=-0.05cm]yo.north) {$B_{\text{adr}}$};
  \node[font=\scriptsize, anchor=north] at ([xshift=0.6cm,yshift=-0.05cm]yo.north) {$C_{\text{adr}}$};
  \node[font=\scriptsize, anchor=east] at ([xshift=-0.1cm,yshift=0.5cm]yo.east) {$A_{\text{veri}}$};
  \node[font=\scriptsize, anchor=east] at ([xshift=-0.1cm,yshift=-0.25cm]yo.east) {$B_{\text{veri}}$};
  \node[font=\scriptsize, anchor=west] at ([xshift=0.1cm,yshift=-1.0cm]yo.west) {$C_{\text{veri}}$};

  % YazmacaYaz stub -- MAVİ (Denetim Birimi çıkışı; bileşende etiketli)
  \node[font=\scriptsize, text=sinyallacivert, anchor=south west, inner sep=0pt] at ([xshift=1.05cm, yshift=0.4cm]yo.north) {YazmacaYaz};
  \draw[->, sinyallacivert, thick] ([xshift=1.05cm, yshift=0.4cm]yo.north) -- ([xshift=1.05cm]yo.north);

  % --- AMBKaynak çoklayıcısı ---
  \node[muxdik] (mux1) at (8.0, 2.1) {};
  \node[font=\tiny] at (mux1.center) {Çok.};
  \node[font=\tiny, anchor=west] at ([xshift=0.02cm, yshift=0.25cm]mux1.west) {0};
  \node[font=\tiny, anchor=west] at ([xshift=0.02cm, yshift=-0.25cm]mux1.west) {1};

  % --- AMB ---
  \node[amb] (ab) at (9.6, 2.6) {AMB};

  \draw[->] ([yshift=0.5cm]yo.east) -- ([yshift=0.5cm]ab.west) node[bitayar, pos=0.5] {32 bit};
  \draw[->] ([yshift=-0.25cm]yo.east) -- ([yshift=0.25cm]mux1.west);
  \draw[->] (mux1.east) -- ([yshift=-0.5cm]ab.west);

  % --- Veri Belleği ---
  \node[vbll] (vb) at (11.65, 0.25) {Veri\\Belleği};
  \node[font=\scriptsize, anchor=north] at ([xshift=-0.5cm, yshift=-0.05cm]vb.north) {$A_{\text{adr}}$};
  \node[font=\scriptsize, anchor=north] at ([xshift=0.5cm, yshift=-0.05cm]vb.north) {$V_{\text{oku}}$};
  \node[font=\scriptsize, anchor=west] at ([xshift=0.1cm, yshift=-0.85cm]vb.west) {$V_{\text{yaz}}$};

  \coordinate (ambsplit) at (11.15, 2.6);
  \node[circle, fill=black, inner sep=1pt] at (ambsplit) {};
  \draw[->] (ambsplit) -- ([xshift=-0.5cm]vb.north);

  \coordinate (bverijnct) at (7.0, 2.35);
  \node[circle, fill=black, inner sep=1pt] at (bverijnct) {};
  \draw[->] (bverijnct) -- (7.0, -0.6) -- ([yshift=-0.85cm]vb.west) node[bitayar, pos=0.55] {32 bit};

  % BelleğeYaz stub -- MAVİ (Denetim Birimi çıkışı; bileşende etiketli)
  \node[font=\scriptsize, text=sinyallacivert, anchor=north] at ([yshift=-0.5cm]vb.south) {BelleğeYaz};
  \draw[->, sinyallacivert, thick] ([yshift=-0.4cm]vb.south) -- (vb.south);

  % --- GeriYazmaKaynağı çoklayıcısı (3 girdili) ---
  \node[muxdik, yeni, minimum height=1.5cm] (mux2) at (13.25, 2.2) {};
  \node[font=\tiny, rotate=90] at ([xshift=0.1cm]mux2.center) {Çoklayıcı};
  \node[font=\tiny, anchor=west] at ([xshift=0.02cm, yshift=0.4cm]mux2.west) {0};
  \node[font=\tiny, anchor=west] at ([xshift=0.02cm, yshift=0cm]mux2.west) {1};
  \node[font=\tiny, anchor=west] at ([xshift=0.02cm, yshift=-0.4cm]mux2.west) {2};

  % AMB çıkışı -> mux2 input 0
  \draw[->] (ab.east) -- ([yshift=0.4cm]mux2.west) node[bitayar, pos=0.15] {32 bit};

  % V_oku -> mux2 input 1
  \draw[->] ([xshift=0.5cm]vb.north) |- ([yshift=0cm]mux2.west) node[bitayar, pos=0.6] {32 bit};

  % PS+4 -> mux2 input 2
  \coordinate (ps4tap) at (-0.3, -0.7);
  \node[circle, fill=black, inner sep=1pt, yeni] at (ps4tap) {};
  \draw[->, yeni] (ps4tap) -- (-0.3, -2.0) -- (10.55, -2.0) -- (10.55, 1.8) -- ([yshift=-0.4cm]mux2.west) node[bitayar, pos=0.55] {32 bit};

  % mux2 çıkışı -> Cveri yolu
  \draw[->] (mux2.east) -- ++(0.5,0)
            |- ++(0,-4.0)
            -| ([xshift=-1.0cm,yshift=-1.0cm]yo.west)
            -- ([yshift=-1.0cm]yo.west);
  \node[bitayar, anchor=west] at ([xshift=0.15cm]mux2.east) {32 bit};

  % --- Anlık değer üreteci ---
  \node[anlikb] (anlik) at (5.4, -0.2) {Anlık\\Üreteci};
  % R türü: anlık yok; buyruk hattından (parse) besle
  \node[circle, fill=black, inner sep=1pt] at (parse) {};
  \draw[->] (parse) -- (2.5,-0.2) -- (anlik.west) node[bitayar, pos=0.5]{buyruk};
  \draw[->] (anlik.east) -- ++(0.75,0) coordinate (anlikbifurc) -- (6.85, 1.85) -- ([yshift=-0.25cm]mux1.west);
  \node[bitayar] at (6.5, -0.2) {32 bit};

  % AnlıkBiçim stub -- MAVİ (Denetim Birimi çıkışı; bileşende etiketli)
  \node[font=\scriptsize, text=sinyallacivert, anchor=north] at ([yshift=-0.5cm]anlik.south) {AnlıkBiçim};
  \draw[->, sinyallacivert, thick] ([yshift=-0.4cm]anlik.south) -- (anlik.south);

  % AMBİşlem stub -- MAVİ (Denetim Birimi çıkışı; bileşende etiketli)
  \node[font=\scriptsize, text=sinyallacivert, anchor=south, align=center] at ([yshift=0.5cm]ab.north) {AMBİşlem\\[-1pt]\tiny 00\,$=$\,Topla,\ 01\,$=$\,Çıkar,\ 10\,$=$\,VEYA};
  \draw[->, sinyallacivert, thick] ([yshift=0.4cm]ab.north) -- (ab.north);

  % AMBKaynak stub -- MAVİ (Denetim Birimi çıkışı; bileşende etiketli)
  \node[font=\scriptsize, text=sinyallacivert, anchor=north] at ([yshift=-0.4cm]mux1.south) {AMBKaynak};
  \draw[->, sinyallacivert, thick] ([yshift=-0.3cm]mux1.south) -- (mux1.south);

  % GeriYazmaKaynağı stub -- MAVİ (Denetim Birimi çıkışı; bileşende etiketli)
  \node[font=\scriptsize, text=sinyallacivert, anchor=north, align=center] at ([yshift=-0.4cm]mux2.south) {GeriYazma\\Kaynağı};
  \draw[->, sinyallacivert, thick] ([yshift=-0.3cm]mux2.south) -- (mux2.south);

  % AMB Sıfır çıkışı
  \node[font=\scriptsize, anchor=west, inner sep=2pt] (sifir) at ([xshift=0.4cm, yshift=0.6cm]ab.east) {Sıfır};
  \draw[->] ([yshift=0.6cm]ab.east) -- (sifir.west);

  % =================================================================
  % DENETİM BİRİMİ -- ÜSTTE (buyruk şeridinin üzerinde, y=9.3)
  % Buyruk şeridi y=6.4..6.95; Denetim Birimi kutusu merkez y=9.3
  % denetim.south ~ (6, 8.8); denetim.north ~ (6, 9.8)
  % =================================================================

  % Denetim Birimi kutusu (biraz aşağıda: merkez y=8.7; south=8.2, north=9.2)
  \node[draw, fill=blokmavi!25, rounded corners=3pt,
        minimum width=7cm, minimum height=1.0cm, align=center]
    (denetim) at (6, 8.7) {Denetim Birimi};

  % -----------------------------------------------------------------
  % GİRİŞ 1: opcode [6:0] -- R türü opcode alanı merkezi abs x=3.0+5.34375=8.34375
  % (8.34375, 6.95) yukarı dikey -> denetim.south
  % -----------------------------------------------------------------
  \draw[->, sinyallacivert, thick] (8.34375, 6.95) -- (8.34375, 8.2);
  \node[font=\scriptsize, text=sinyallacivert, anchor=west] at (8.38, 7.7) {opcode [6:0]};

  % -----------------------------------------------------------------
  % GİRİŞ 2: funct -- iki ayrı alan: funct7 (abs x=3.66, sol) + funct3 (abs x=6.47)
  % İkisi de şerit üst kenarından (y=6.95) yukarı Denetim Birimi'ne girer
  % -----------------------------------------------------------------
  \draw[->, sinyallacivert, thick] (3.66, 6.95) -- (3.66, 8.2);
  \node[font=\scriptsize, text=sinyallacivert, anchor=east] at (3.62, 7.7) {funct7};
  \draw[->, sinyallacivert, thick] (6.47, 6.95) -- (6.47, 8.2);
  \node[font=\scriptsize, text=sinyallacivert, anchor=west] at (6.51, 7.7) {funct3};

  % -----------------------------------------------------------------
  % GİRİŞ 3: Sıfır -- AMB Sıfır çıkışından sağ dış kanal x=11.6, Denetim Birimi'ne SAĞDAN
  % sifir.east -> sağ (11.6, 3.2) -> yukarı (11.6, 8.7) -> sola denetim.east
  % -----------------------------------------------------------------
  \draw[->, sinyallacivert, thick] (sifir.east) -- (11.6, 3.2) -- (11.6, 8.7) -- (denetim.east);
  \node[font=\scriptsize, text=sinyallacivert, anchor=west] at (11.64, 6.0) {Sıfır};

  % -----------------------------------------------------------------
  % ÇIKIŞLAR -- Denetim Birimi'nin ÜST kenarından 7 kısa ok (yukarı)
  % x = 3.0, 3.9, 4.8, 5.7, 6.6, 7.5, 8.4
  % Sinyaller: YazmacaYaz, AMBKaynak, GeriYazmaKaynağı, AnlıkBiçim,
  %            BelleğeYaz, Dallan, AMBİşlem
  % -----------------------------------------------------------------
  % 7 çıkış: merkez x=6 etrafında simetrik (x=3,4,5,6,7,8,9), 1.0 aralık
  % Etiketler anchor=west + rotate=90 -> hepsi alttan (y=9.8) hizalı, yukarı okunur (kutuya sarkmaz)
  % (1) YazmacaYaz
  \draw[->, sinyallacivert, thick] (3.0, 9.2) -- (3.0, 9.75);
  \node[font=\tiny, text=sinyallacivert, anchor=west, rotate=45] at (3.0, 9.8) {YazmacaYaz};

  % (2) AMBKaynak
  \draw[->, sinyallacivert, thick] (4.0, 9.2) -- (4.0, 9.75);
  \node[font=\tiny, text=sinyallacivert, anchor=west, rotate=45] at (4.0, 9.8) {AMBKaynak};

  % (3) GeriYazmaKaynağı
  \draw[->, sinyallacivert, thick] (5.0, 9.2) -- (5.0, 9.75);
  \node[font=\tiny, text=sinyallacivert, anchor=west, rotate=45] at (5.0, 9.8) {GeriYazmaKaynağı};

  % (4) AnlıkBiçim
  \draw[->, sinyallacivert, thick] (6.0, 9.2) -- (6.0, 9.75);
  \node[font=\tiny, text=sinyallacivert, anchor=west, rotate=45] at (6.0, 9.8) {AnlıkBiçim};

  % (5) BelleğeYaz
  \draw[->, sinyallacivert, thick] (7.0, 9.2) -- (7.0, 9.75);
  \node[font=\tiny, text=sinyallacivert, anchor=west, rotate=45] at (7.0, 9.8) {BelleğeYaz};

  % (6) Dallan
  \draw[->, sinyallacivert, thick] (8.0, 9.2) -- (8.0, 9.75);
  \node[font=\tiny, text=sinyallacivert, anchor=west, rotate=45] at (8.0, 9.8) {Dallan};

  % (7) AMBİşlem
  \draw[->, sinyallacivert, thick] (9.0, 9.2) -- (9.0, 9.75);
  \node[font=\tiny, text=sinyallacivert, anchor=west, rotate=45] at (9.0, 9.8) {AMBİşlem};

\end{tikzpicture}
\caption[Tek vuruşluk veri yolu ve denetim birimi]{Tek vuruşluk veri yolu ve denetim birimi. Denetim birimi, işlem kodu ve funct alanları (ve dallanma için \texttt{Sıfır} bayrağı) girişlerinden yedi denetim sinyalini üretir. Bu sinyaller veri yolunda üretildikleri bileşende mavi olarak etiketlenmiştir; okunabilirlik için denetim biriminden her bileşene giden tek tek teller çizilmemiş, sinyaller ad eşleşmesiyle gösterilmiştir.}
\label{fig:tek-vurusluk-tam}
\end{figure}


\begin{ornek}[6.3]
Tek vuruşluk işlemcimize, \emph{dizinli yükleme} adını vereceğimiz yeni bir \texttt{dyük} buyruğu eklemek istiyoruz:
\[
\texttt{dyük rd, rs1, rs2}\colon\quad R[\text{rd}] \leftarrow \text{Bel}[\,R[\text{rs1}] + R[\text{rs2}]\,]
\]
Bu buyruk, bellek adresini bir yazmaç ile bir anlık değerden değil, iki yazmacın toplamından üretir (bir dizi erişiminde taban adres ile dizin yazmacının toplanması gibi). Veri yolundaki bileşenlerden hiçbirini değiştirmeden bu buyruk eklenebilir mi? Eklenebiliyorsa denetim tablosundaki satırını yazın.

\textbf{Çözüm.}

Buyruğun veri yolundan akışını izleyelim: yazmaç öbeği \texttt{rs1} ile \texttt{rs2}'yi okur (öbekte zaten iki okuma kapısı vardır, R türü buyruklar bu ikisini kullanır). AMB bu iki değeri toplayarak bellek adresini üretir. Veri belleği bu adresteki sözcüğü okur. Okunan sözcük geri yazma çoklayıcısından geçip \texttt{rd}'ye yazılır. Bu yolun gerektirdiği her bileşen ve her bağlantı veri yolunda hâlihazırda bulunur. Yeni hiçbir parça gerekmez. Tek yapılması gereken, denetim tablosuna karşılık gelen satırı eklemektir:

\begin{center}
\small
\begin{tabular}{l|c|c|c|c|c|c|c}
\toprule
\textbf{Buyruk} & \textbf{Yazmaca} & \textbf{AMB} & \textbf{GeriYazma} & \textbf{Anlık} & \textbf{Belleğe} & \textbf{Dallan} & \textbf{AMB} \\
 & \textbf{Yaz} & \textbf{Kaynak} & \textbf{Kaynağı} & \textbf{Biçim} & \textbf{Yaz} & & \textbf{İşlem} \\
\midrule
add  & 1 & 0 & 00 & X  & 0 & 0 & 00 \\
lw   & 1 & 1 & 01 & 00 & 0 & 0 & 00 \\
\texttt{dyük} & 1 & 0 & 01 & X  & 0 & 0 & 00 \\
\bottomrule
\end{tabular}
\end{center}

Bu satır, var olan iki satırın bir bileşimidir. Adres hesabını \texttt{add} gibi yapar (\texttt{AMBKaynak}=0, AMB'nin B girişine \texttt{rs2} gelir), geri yazılacak değeri ise \texttt{lw} gibi bellekten alır (\texttt{GeriYazmaKaynağı}=01). \texttt{lw}'den tek farkı \texttt{AMBKaynak} sinyalidir. \texttt{lw} adres için anlık değer kullanırken (\texttt{AMBKaynak}=1), \texttt{dyük} ikinci yazmacı kullanır (\texttt{AMBKaynak}=0). Anlık değer kullanılmadığından \texttt{AnlıkBiçim} bu buyruk için önemsizdir (X).

Bu, Bölüm~\ref{subsec:denetim-tablosu}'de belirtilen ilkenin somut bir örneğidir. Mimarideki bileşenler hiç değişmeden, denetim tablosuna yeni bir satır eklenerek yeni bir buyruk tanımlanabilir. \texttt{dyük}'ün satırı tablodaki hiçbir satırla aynı olmadığından, işlemci açısından kendine özgü, yeni bir buyruktur. Buyruk R türü biçiminde kodlanabilir (\texttt{rd}, \texttt{rs1}, \texttt{rs2}; anlık değer yok). Denetim birimine yalnızca bu yeni biçimi tanıyıp yukarıdaki satırı üretecek mantık eklenir, veri yolu olduğu gibi kalır.
\end{ornek}

%% ================================================================
%% ÇOK VURUŞLUK İŞLEMCİ
%% ================================================================

\section{Çok Vuruşluk İşlemci}\index{çok vuruşluk}\index{multi-cycle|see{çok vuruşluk}}
\label{sec:cok-vurusluk}

Tek vuruşluk tasarımın temel kısıtı, saat çevriminin her zaman en yavaş buyruğa göre ayarlanmak zorunda olmasıdır. Çevrim süresi yükle buyruğunun gecikmesine sabitlendiğinden, yalnızca birkaç adım gerektiren dallanma gibi buyruklar işlerini erken bitirse bile bir sonraki buyruk ancak çevrim dolduğunda başlayabilir. Aradaki boş süre geri kazanılamaz.

Çok vuruşluk işlemci bu kısıtı, her buyruğu birkaç küçük adıma bölerek aşar. Her adım tek bir saat çevriminde tamamlanır ve çevrim süresi artık en yavaş \emph{buyruğa} değil, en yavaş \emph{adıma} göre belirlenir. Çevrim böylece belirgin biçimde kısalır. Buna karşılık her buyruk birden çok çevrim sürer ve her buyruk yalnızca gereksinim duyduğu adım sayısı kadar çevrim kullanır (yükle beş, dallanma üç adım gibi). Bir adımda üretilen değerin sonraki adımda kullanılabilmesi için veri yoluna adımlar arası ara yazmaçlar eklenir. Bu bölümde tek vuruşluk veri yolunu adım adım çok vuruşluk bir tasarıma dönüştüreceğiz.

\subsection{Tek Vuruşluk İşlemcide Karşılaşılan Sorunlar}
\label{subsec:tek-vurusluk-sorunlar}

Tek vuruşluk işlemcinin başarımını, Bölüm~\ref{chap:basarim}'de tanımladığımız iki çarpan belirler: buyruk başına çevrim sayısı (BBÇ\index{BBÇ}\index{CPI|see{BBÇ}}) ile çevrim süresi\index{çevrim zamanı}\index{cycle time|see{çevrim zamanı}}. Her buyruk tek bir çevrimde tamamlandığından BBÇ değeri 1'dir. Bu, ilk bakışta ideal görünür. Ne var ki başarımın asıl belirleyicisi çevrim süresidir ve bu süre, az önce gördüğümüz gibi, en yavaş buyruğun (yükle) gecikmesine sabitlenir.

Kısa buyrukların, işlerini bitirdikten sonra çevrimin dolmasını beklemesi boşa geçen zamandır. Bu israfı somut görmek için, buyrukların veri yolundaki ``getir--çöz--yürüt--bellek--sonuç yaz'' aşamalarından hangilerinden geçtiğine bakalım (Tablo~\ref{tab:buyruk-asamalar}). Bu aşamaların tek bir çevrim içindeki zaman dağılımını ve kısa buyruklarda boşa geçen süreyi Şekil~\ref{fig:tek-vurusluk-zaman}'de görmüştük.

\begin{table}[htbp]
\centering
\caption{Buyrukların veri yolunda geçtikleri aşamalar (en çok çevrimden en aza)}
\label{tab:buyruk-asamalar}
\begin{tabular}{lcccccr}
\toprule
\textbf{Buyruk} & \textbf{Getir} & \textbf{Çöz} & \textbf{Yürüt} & \textbf{Bellek} & \textbf{Sonucu Yaz} & \textbf{Toplam} \\
\midrule
Yükle & \checkmark & \checkmark & \checkmark & \checkmark & \checkmark & 5 \\
Sakla & \checkmark & \checkmark & \checkmark & \checkmark & & 4 \\
Aritmetik--Mantık (R) & \checkmark & \checkmark & \checkmark & & \checkmark & 4 \\
Mantıksal VEYA Anlık (ori) & \checkmark & \checkmark & \checkmark & & \checkmark & 4 \\
jal & \checkmark & \checkmark & \checkmark & & & 3 \\
Dallanma (beq/bne) & \checkmark & \checkmark & \checkmark & & & 3 \\
\bottomrule
\end{tabular}
\\[0.3em]
{\footnotesize jal, dönüş adresini (PS+4) ayrı bir Sonucu Yaz çevriminde değil, Yürüt çevriminde rd yazmacına yazar. Bu nedenle çok vuruşluk gerçekleştiriminde üç çevrimde tamamlanır (Bölüm~\ref{subsec:sinirli-durum-makinesi}'teki durum makinesiyle uyumludur).}
\end{table}

Yükle buyruğunun veri yolundaki aşamalardan geçişine göre çevrim süresinin\index{bellek erişim zamanı}\index{memory access time|see{bellek erişim zamanı}} hesaplanması:
\begin{align*}
\text{Yükle} &= \text{PS'nin saat gecikmesi} + \text{Buyruk belleğinin erişim zamanı} \\
&\quad + \text{Yazmaç öbeğinin erişim zamanı} + \text{AMB'nin 32 bit toplama yapması} \\
&\quad + \text{Veri belleği erişim zamanı} + \text{Yazmaç öbeğinin yazma için hazırlanması} \\
&\quad + \text{Saatin gecikmesi}
\end{align*}

Tek vuruşluk veri yolu gerçek hayattaki işlemcilerde kullanılmaz. RISC-V'in alt kümesi için işlemler ortalama bir sürede bitecek gibi görünse de, kayan nokta işlemleri çok daha uzun çevrimlere neden olur. Ayrıca bellek hızı da çevrimin süresini etkileyen bir diğer etkendir.

\begin{bilginot}[Açık Kaynak RISC-V Çekirdekleri]
Açık kaynak RISC-V çekirdekleri arasında PicoRV32 basit, çok vuruşluk bir tasarımdır (buyruk başına ortalama yaklaşık dört çevrim, BBÇ$\approx$4 harcar). Western Digital'in SweRV EH1 çekirdeği ise boru hatlı, çift yollu (dual-issue) bir mimariye sahiptir. Bu çekirdekler eğitim ve endüstriyel uygulamalarda yaygın olarak kullanılmaktadır. İşlemci tasarımının temelleri bu bölümde anlatılan ilkelerle aynıdır. Fark, gerçek dünya tasarımlarının çok daha fazla buyruk ve karmaşık denetim gerektirmesidir.
\end{bilginot}

\begin{ornek}[6.4]
Tek vuruşluk bir RISC-V işlemcide birimlerin gecikme süreleri aşağıdaki gibidir. (Örnek~6.1'deki kaba modelden farklı olarak burada çoklayıcı ve anlık değer üreteci gecikmeleri de ayrıca sayılır. Birim değerleri de farklı seçilmiştir.)

\begin{center}
\begin{tabular}{lc}
\toprule
\textbf{Birim} & \textbf{Gecikme (ps)} \\
\midrule
Buyruk belleği / Veri belleği & 200 \\
Yazmaç öbeği (okuma veya yazma) & 50 \\
AMB & 100 \\
Çoklayıcı & 25 \\
Toplayıcı (PS+4) & 50 \\
Anlık değer üreteci & 25 \\
\bottomrule
\end{tabular}
\end{center}

(a)~Her buyruk türü için \textbf{kritik yol}\index{kritik yol}\index{critical path|see{kritik yol}} gecikmesini hesaplayın. (b)~Tek vuruşluk işlemcinin çevrim zamanını ve saat sıklığını belirleyin.

\textbf{Çözüm.}

\textbf{(a)} Her buyruk türü için gecikme (ps cinsinden):

\begin{center}
\small
\begin{tabular}{@{}lp{7cm}r@{}}
\toprule
\textbf{Buyruk} & \textbf{Kritik yol} & \textbf{Toplam} \\
\midrule
R türü & ByrBel+YazÖb(oku)+Mux+AMB+Mux+YazÖb(yaz) & 450\,ps \\
ori & ByrBel+YazÖb(oku)+AMB+Mux+YazÖb(yaz) & 425\,ps \\
Yükle & ByrBel+YazÖb(oku)+AMB+VerBel+Mux+YazÖb(yaz) & 625\,ps \\
Sakla & ByrBel+YazÖb(oku)+AMB+VerBel & 550\,ps \\
beq/bne & ByrBel+YazÖb(oku)+Mux+AMB & 375\,ps \\
\bottomrule
\end{tabular}
\end{center}

\textbf{(b)} Çevrim zamanı en uzun kritik yol tarafından belirlenir: $t_{\text{çevrim}} = 625$~ps (yükle buyruğu). R türü ile beq buyruklarında ikinci yazmaç (rs2) AMBKaynak çoklayıcısından geçtiği için kritik yolları 25~ps daha uzundur. Anlık değer kullanan ori, lw ve sw buyruklarında ise rs1 doğrudan AMB'ye girer ve anlık değer paralel daldan geldiğinden çoklayıcı bu buyrukların kritik yolunda yer almaz. Program sayacı güncelleme yolundaki toplayıcı, Dallan kapıları ve PS çoklayıcısı da ayrıca izlenmemiştir. Bu yol veri yoluyla koşut çalışır ve toplamı en uzun yolu aşmaz. En uzun kritik yol yine yükle buyruğunundur (625~ps).

\[
\text{Saat sıklığı} = \frac{1}{625 \times 10^{-12}} \approx 1{,}60 \text{~GHz}
\]

Yükle buyruğu programın yalnızca \%12'sini oluştursa bile tüm buyruklar 625~ps beklemek zorundadır. Bu, tek vuruşluk tasarımın temel verimsizliğidir. Bu kritik yol Şekil~\ref{fig:kritik-yol-lw}'de görsel olarak gösterilmektedir.
\end{ornek}

\begin{figure}[htbp]
\centering
\begin{tikzpicture}[>=Stealth, semithick, font=\footnotesize,
    yzm/.style={draw, fill=yellow!25, minimum width=1cm, minimum height=1cm, align=center},
    bllk/.style={draw, fill=green!20, minimum width=1.8cm, minimum height=1.6cm, align=center},
    yobk/.style={draw, fill=blue!10, minimum width=2.6cm, minimum height=3cm},
    amb/.style={draw, fill=red!20, minimum width=1.6cm, minimum height=2cm, align=center},
    artop/.style={draw, fill=gray!15, minimum width=1.1cm, minimum height=0.7cm, align=center},
    anlikb/.style={draw, fill=orange!25, minimum width=1.4cm, minimum height=0.9cm, align=center},
    muxdik/.style={draw, fill=purple!15, minimum width=0.7cm, minimum height=1.0cm, align=center},
    muxyatay/.style={draw, fill=purple!15, minimum width=1.0cm, minimum height=0.7cm, align=center},
    vbll/.style={draw, fill=green!20, minimum width=1.8cm, minimum height=2.2cm, align=center},
    bitayar/.style={fill=white, inner sep=1pt, font=\tiny},
    yeni/.style={},
    % Kritik yol bindirmesi için stiller
    kritik/.style={line width=1.6pt, sinyalkirmizi, ->,
                   -{Stealth[length=4pt,width=4pt]}},
    psure/.style={font=\tiny\bfseries, text=sinyalkirmizi,
                  fill=white, inner sep=1pt, rounded corners=1pt,
                  draw=sinyalkirmizi, line width=0.4pt}
]
  % --- Sol PS güncelleme bloğu ---
  \node[bllk] (bb) at (0, 4.6) {Buyruk\\Belleği};
  \node[yzm] (ps) at (0, 2.7) {PS};
  \node[muxyatay] (psmux) at (0, 1.3) {};
  \node[font=\tiny] at (psmux.center) {Çok.};
  \node[font=\tiny, anchor=south] at ([xshift=-0.3cm, yshift=0.02cm]psmux.south) {0};
  \node[font=\tiny, anchor=south] at ([xshift=0.3cm, yshift=0.02cm]psmux.south) {1};

  \node[artop] (top4) at (-1.3, 0) {Top.};
  \node[artop, minimum width=1.2cm, minimum height=0.7cm] (daladd) at (1.3, 0) {Dal.\\Top.};

  \draw[->] (ps.north) -- (bb.south) node[midway, right=1pt, font=\scriptsize] {adres};
  \draw[->] (psmux.north) -- (ps.south);

  \draw[->] (ps.west) -- (-1.5, 2.7) -- ([xshift=-0.2cm]top4.north);
  \draw[->] (ps.east) -- (1.1, 2.7) -- ([xshift=-0.2cm]daladd.north);

  \node[font=\scriptsize] at (-1.1, 0.85) {$4$};
  \draw[->] (-1.1, 0.7) -- ([xshift=0.2cm]top4.north);

  % Anlık -> Dal.Top.
  \coordinate (anlikdal) at (6.85, 1.0);
  \node[circle, fill=black, inner sep=1pt] at (anlikdal) {};
  \draw[->] (anlikdal) -- (1.5, 1.0) -- ([xshift=0.2cm]daladd.north);

  % Top4 çıkışı (PS+4)
  \draw[->] (top4.south) -- (-1.3, -0.7) -- (-0.3, -0.7) -- ([xshift=-0.3cm]psmux.south);
  \node[font=\scriptsize, anchor=north] at (-0.8, -0.75) {PS$+$4};
  \draw[->] (daladd.south) -- (1.3, -0.7) -- (0.3, -0.7) -- ([xshift=0.3cm]psmux.south);
  \node[font=\scriptsize, anchor=north] at (0.8, -0.75) {Dal.hedef};

  % Dallan denetim sinyali
  \draw[->, sinyallacivert, thick] ([xshift=-0.4cm]psmux.west) -- (psmux.west);
  \node[font=\scriptsize, text=sinyallacivert, anchor=south] at ([xshift=-0.5cm, yshift=0.05cm]psmux.west) {Dallan};

  \coordinate (parse) at (2.5, 4.6);
  \draw[->] (bb.east) -- (parse) node[midway, below, font=\scriptsize] {buyruk} node[bitayar, pos=0.55] {32 bit};

  % ── lw (I türü) buyruk şeridi ──────────────────────────────────────
  % seritsol = (3.0, 6.4) — bağlantı dikey: parse (2.5,4.6) → (2.5,6.4) → (3.0,6.4)
  \coordinate (seritsol) at (3.0, 6.4);
  \draw[gray, dashed, semithick] (parse) -- (2.5, 6.4) -- (seritsol);
  \draw[gray, dashed, semithick] (2.5, 6.4) -- (2.5, 5.3);  % dikey bağlayıcı alt ucu

  \begin{scope}[shift={(seritsol)}]
    % imm[11:0]  0–2.25 cm
    \draw[fill=orange!25] (0,0) rectangle (2.25,0.55) node[pos=.5, font=\tiny] {imm[11:0]};
    % rs1  2.25–3.1875 cm
    \draw[fill=cyan!20]   (2.25,0) rectangle (3.1875,0.55) node[pos=.5, font=\tiny] {rs1};
    % f3  3.1875–3.75 cm
    \draw[fill=gray!15]   (3.1875,0) rectangle (3.75,0.55) node[pos=.5, font=\tiny] {f3};
    % rd  3.75–4.6875 cm
    \draw[fill=blue!18]   (3.75,0) rectangle (4.6875,0.55) node[pos=.5, font=\tiny] {rd};
    % opcode  4.6875–6 cm
    \draw[fill=gray!10]   (4.6875,0) rectangle (6,0.55) node[pos=.5, font=\tiny] {opcode};
    % bit sınır etiketleri
    \node[font=\tiny, anchor=north] at (0,0)      {31};
    \node[font=\tiny, anchor=north] at (2.25,0)   {20};
    \node[font=\tiny, anchor=north] at (3.1875,0) {15};
    \node[font=\tiny, anchor=north] at (3.75,0)   {12};
    \node[font=\tiny, anchor=north] at (4.6875,0) {7};
    \node[font=\tiny, anchor=north] at (6,0)      {0};
    \node[font=\scriptsize, anchor=south] at (3.0, 0.55) {lw buyruğu (I türü)};
  \end{scope}

  % rs1 (x≈2.72 cm içinde seritsol → rs1 merkezi) → A_adr (dikey+dikey)
  \draw[->] ([xshift=2.7188cm]seritsol) -- ++(0,-0.6) -- ++(0,-0.7)
            -| ([xshift=-0.6cm]yo.north) node[bitayar, pos=0.6] {5 bit};

  % rd (x≈4.22 cm içinde seritsol) → C_adr (dikey)
  \draw[->] ([xshift=4.22cm]seritsol) -- ++(0,-0.6) -- ++(0,-0.7)
            -| ([xshift=0.6cm]yo.north) node[bitayar, pos=0.75] {5 bit};

  % imm → Anlık Üreteci (dikey sonra yatay)
  \draw[->] ([xshift=1.125cm]seritsol) -- ++(0,-1.5) -- (4.7, 4.9) |- (anlik.west)
            node[bitayar, pos=0.8] {12 bit};

  \node[yobk] (yo) at (5.4, 2.6) {};
  \node[anchor=center, font=\footnotesize\bfseries, align=center] at (yo.center) {Yazmaç\\Öbeği};

  \node[font=\scriptsize, anchor=north] at ([xshift=-0.6cm,yshift=-0.05cm]yo.north) {$A_{\text{adr}}$};
  \node[font=\scriptsize, anchor=north] at ([xshift=0cm,yshift=-0.05cm]yo.north) {$B_{\text{adr}}$};
  \node[font=\scriptsize, anchor=north] at ([xshift=0.6cm,yshift=-0.05cm]yo.north) {$C_{\text{adr}}$};
  \node[font=\scriptsize, anchor=east] at ([xshift=-0.1cm,yshift=0.5cm]yo.east) {$A_{\text{veri}}$};
  \node[font=\scriptsize, anchor=east] at ([xshift=-0.1cm,yshift=-0.25cm]yo.east) {$B_{\text{veri}}$};
  \node[font=\scriptsize, anchor=west] at ([xshift=0.1cm,yshift=-1.0cm]yo.west) {$C_{\text{veri}}$};

  % YazmacaYaz
  \node[font=\scriptsize, text=sinyallacivert, anchor=south west, inner sep=0pt] at ([xshift=1.05cm, yshift=0.4cm]yo.north) {YazmacaYaz};
  \draw[->, sinyallacivert, thick] ([xshift=1.05cm, yshift=0.4cm]yo.north) -- ([xshift=1.05cm]yo.north);

  % --- AMBKaynak çoklayıcısı ---
  \node[muxdik] (mux1) at (8.0, 2.1) {};
  \node[font=\tiny] at (mux1.center) {Çok.};
  \node[font=\tiny, anchor=west] at ([xshift=0.02cm, yshift=0.25cm]mux1.west) {0};
  \node[font=\tiny, anchor=west] at ([xshift=0.02cm, yshift=-0.25cm]mux1.west) {1};

  % --- AMB ---
  \node[amb] (ab) at (9.6, 2.6) {AMB};

  % A_veri → AMB doğrudan (rs2 için B_veri → Mux → AMB)
  \draw[->] ([yshift=0.5cm]yo.east) -- ([yshift=0.5cm]ab.west) node[bitayar, pos=0.5] {32 bit};
  \draw[->] ([yshift=-0.25cm]yo.east) -- ([yshift=0.25cm]mux1.west);
  \draw[->] (mux1.east) -- ([yshift=-0.5cm]ab.west);

  % --- Veri Belleği ---
  \node[vbll] (vb) at (11.65, 0.25) {Veri\\Belleği};
  \node[font=\scriptsize, anchor=north] at ([xshift=-0.5cm, yshift=-0.05cm]vb.north) {$A_{\text{adr}}$};
  \node[font=\scriptsize, anchor=north] at ([xshift=0.5cm, yshift=-0.05cm]vb.north) {$V_{\text{oku}}$};
  \node[font=\scriptsize, anchor=west] at ([xshift=0.1cm, yshift=-0.85cm]vb.west) {$V_{\text{yaz}}$};

  \coordinate (ambsplit) at (11.15, 2.6);
  \node[circle, fill=black, inner sep=1pt] at (ambsplit) {};
  \draw[->] (ambsplit) -- ([xshift=-0.5cm]vb.north);

  \coordinate (bverijnct) at (7.0, 2.35);
  \node[circle, fill=black, inner sep=1pt] at (bverijnct) {};
  \draw[->] (bverijnct) -- (7.0, -0.6) -- ([yshift=-0.85cm]vb.west) node[bitayar, pos=0.55] {32 bit};

  \node[font=\scriptsize, text=sinyallacivert, anchor=north] at ([yshift=-0.5cm]vb.south) {BelleğeYaz};
  \draw[->, sinyallacivert, thick] ([yshift=-0.4cm]vb.south) -- (vb.south);

  % --- GeriYazmaKaynağı çoklayıcısı (3 girdili) ---
  \node[muxdik, minimum height=1.5cm] (mux2) at (13.25, 2.2) {};
  \node[font=\tiny, rotate=90] at ([xshift=0.1cm]mux2.center) {Çoklayıcı};
  \node[font=\tiny, anchor=west] at ([xshift=0.02cm, yshift=0.4cm]mux2.west) {0};
  \node[font=\tiny, anchor=west] at ([xshift=0.02cm, yshift=0cm]mux2.west) {1};
  \node[font=\tiny, anchor=west] at ([xshift=0.02cm, yshift=-0.4cm]mux2.west) {2};

  % AMB çıkışı -> mux2 giriş 0
  \draw[->] (ab.east) -- ([yshift=0.4cm]mux2.west) node[bitayar, pos=0.15] {32 bit};

  % V_oku -> mux2 giriş 1
  \draw[->] ([xshift=0.5cm]vb.north) |- ([yshift=0cm]mux2.west) node[bitayar, pos=0.6] {32 bit};

  % PS+4 -> mux2 giriş 2
  \coordinate (ps4tap) at (-0.3, -0.7);
  \node[circle, fill=black, inner sep=1pt] at (ps4tap) {};
  \draw[->] (ps4tap) -- (-0.3, -2.0) -- (10.55, -2.0) -- (10.55, 1.8) -- ([yshift=-0.4cm]mux2.west) node[bitayar, pos=0.55] {32 bit};

  % mux2 çıkışı -> C_veri yolu
  \draw[->] (mux2.east) -- ++(0.5,0) -- ++(0,-4.0)
            -| ([xshift=-1.0cm,yshift=-1.0cm]yo.west)
            -- ([yshift=-1.0cm]yo.west);
  \node[bitayar, anchor=west] at ([xshift=0.15cm]mux2.east) {32 bit};

  % --- Anlık değer üreteci ---
  \node[anlikb] (anlik) at (5.4, -0.2) {Anlık\\Üreteci};
  \draw[->] (anlik.east) -- ++(0.75,0) coordinate (anlikbifurc) -- (6.85, 1.85) -- ([yshift=-0.25cm]mux1.west);
  \node[bitayar] at (6.5, -0.2) {32 bit};

  % AnlıkBiçim denetim sinyali
  \node[font=\scriptsize, text=sinyallacivert, anchor=north] at ([yshift=-0.5cm]anlik.south) {AnlıkBiçim};
  \draw[->, sinyallacivert, thick] ([yshift=-0.4cm]anlik.south) -- (anlik.south);

  % AMBİşlem denetim sinyali
  \node[font=\scriptsize, text=sinyallacivert, anchor=south, align=center] at ([yshift=0.5cm]ab.north) {AMBİşlem\\[-1pt]\tiny 00\,$=$\,Topla,\ 01\,$=$\,Çıkar,\ 10\,$=$\,VEYA};
  \draw[->, sinyallacivert, thick] ([yshift=0.4cm]ab.north) -- (ab.north);

  % AMBKaynak denetim sinyali
  \node[font=\scriptsize, text=sinyallacivert, anchor=north] at ([yshift=-0.4cm]mux1.south) {AMBKaynak};
  \draw[->, sinyallacivert, thick] ([yshift=-0.3cm]mux1.south) -- (mux1.south);

  % GeriYazmaKaynağı denetim sinyali
  \node[font=\scriptsize, text=sinyallacivert, anchor=north, align=center] at ([yshift=-0.4cm]mux2.south) {GeriYazma\\Kaynağı};
  \draw[->, sinyallacivert, thick] ([yshift=-0.3cm]mux2.south) -- (mux2.south);

  % AMB Sıfır çıkışı
  \node[font=\scriptsize, anchor=west, inner sep=2pt] (sifir) at ([xshift=0.4cm, yshift=0.6cm]ab.east) {Sıfır};
  \draw[->] ([yshift=0.6cm]ab.east) -- (sifir.west);

  % =================================================================
  % lw KRİTİK YOL BİNDİRMESİ (625 ps)
  % rs1 dalı: PS→ByrBel(200)→YazÖb oku(50)→AMB doğrudan→AMB(100)→VerBel(200)→GeriYazma Çok.(25)→YazÖb yaz(50)
  % Anlık dalı: Anlık Üreteci(25)→AMBKaynak Çok.(25) → AMB'de birleşir (paralel, toplamda 50 ps)
  % =================================================================

  % 1) PS → Buyruk Belleği (200 ps)
  \draw[kritik] (ps.north) -- node[psure, right=2pt]{200 ps} (bb.south);

  % 2) Buyruk Belleği → Yazmaç Öbeği rs1 giriş (50 ps) — yatay+dikey+yatay+dikey, dik açılı
  \draw[kritik] (bb.east) -- (2.5, 4.6) -- (2.5, 5.3) -- (4.8, 5.3) -- ([xshift=-0.6cm]yo.north);
  \node[psure] at (3.6, 5.5) {50 ps};

  % 3) A_veri → AMB doğrudan (rs1 yolu — çoklayıcı YOK)
  \draw[kritik] ([yshift=0.5cm]yo.east) -- node[psure, above]{} ([yshift=0.5cm]ab.west);
  \node[psure] at ([yshift=0.4cm]ab.center) {100 ps};

  % 3b) Anlık dalı (paralel): Anlık Üreteci → AMBKaynak Çoklayıcısı → AMB B girişi
  %     25 ps (AnlıkÜrt) + 25 ps (Çoklayıcı) = toplam 50 ps, paralel rs1 ile
  \draw[kritik] (anlik.east) -- (6.85, -0.2) -- (6.85, 1.85)
                node[psure, right=2pt, pos=0.3]{25 ps} -- ([yshift=-0.25cm]mux1.west);
  \draw[kritik] (mux1.east) -- node[psure, above, pos=0.5]{25 ps} ([yshift=-0.5cm]ab.west);

  % 4) AMB → Veri Belleği (200 ps)
  \draw[kritik] (ab.east) -- (ambsplit) -- node[psure, right=2pt]{200 ps} ([xshift=-0.5cm]vb.north);

  % 5) Veri Belleği V_oku → GeriYazmaKaynağı çoklayıcısı (25 ps)
  \draw[kritik] ([xshift=0.5cm]vb.north) -- (12.15, 1.35) -- (12.15, 2.2)
                node[psure, right=2pt, pos=0.5]{25 ps} -- ([yshift=0cm]mux2.west);

  % 6) GeriYazmaKaynağı çoklayıcısı → Yazmaç Öbeği geri yazma (50 ps)
  %    yo.west=(4.1,2.6), yshift=-1.0cm → (4.1,1.6); tüm segmentler yatay veya dikey
  \draw[kritik] (mux2.east) -- (13.75, 2.2) -- (13.75, -1.8)
                -- (3.1, -1.8) -- (3.1, 1.6) -- ([yshift=-1.0cm]yo.west);
  \node[psure] at (8.0, -1.55) {50 ps};

  % ─── Toplam kritik yol kutusu ───
  \node[font=\small\bfseries, text=sinyalkirmizi,
        draw=sinyalkirmizi, rounded corners=3pt, fill=sinyalkirmizi!10,
        inner sep=4pt, align=center] at (6.8, -3.3)
       {Toplam kritik yol: $200 + 50 + 100 + 200 + 25 + 50 = 625$~ps\\
        \normalfont\tiny\textit{(rs1 ve anlık dalları paraleldir; ikisi de AMB'ye 50~ps'de varır)}};

\end{tikzpicture}
\caption{lw buyruğunun kritik yolu (kırmızı yolla işaretlenmiş)}
\label{fig:kritik-yol-lw}
\end{figure}

\subsection{Çok Vuruşluk Veri Yolu}
\label{subsec:cok-vurusluk-veriyolu}

Çok vuruşluk tasarımda her buyruk, Tablo~\ref{tab:buyruk-asamalar}'deki aşamalara (getir, çöz, yürüt, bellek, sonuç yaz) bölünür ve her aşama bir saat çevriminde tamamlanır. Böylece her buyruk yalnızca gereksinim duyduğu aşama sayısı kadar çevrim alır: yükle beş, dallanma üç. Tek vuruşluk tasarımda olduğu gibi, çok vuruşluk işlemcide de her an \emph{tek bir buyruk} bulunur. Sıradaki buyruk, içerideki buyruk bütün aşamalarını bitirip tamamlanmadan başlamaz. Birden çok buyruğun aşamalarını örtüştürerek bunları aynı anda yürütmek ise ayrı bir yöntemdir: boru hattı\index{boru hattı} (Bölüm~\ref{chap:boruhatti}).

Bu bölünmenin neden başarım kazandırabileceğini önce basit bir kestirimle görelim. Beş aşamanın hepsinin eşit gecikmede ($T$) olduğunu varsayalım. Tek vuruşlukta saat çevrimi en uzun buyruğa (yükle, beş aşama) göre sabitlendiğinden her buyruk $5T$ sürer. Çok vuruşlukta ise çevrim tek bir aşama kadardır ($T$) ve her buyruk yalnızca kendi aşamalarını yürütür. Örnek bir buyruk dağılımıyla (\%25 yükle, \%10 sakla, \%45 R türü, \%12 ori, \%8 dallanma) buyruk başına ortalama çevrim sayısı
\[
\text{BBÇ} = 0{,}25\times5 + 0{,}10\times4 + 0{,}45\times4 + 0{,}12\times4 + 0{,}08\times3 = 4{,}17
\]
olur. Buna göre çok vuruşluk tasarımın tek vuruşluğa hızlanması
\[
\text{Hızlanma} = \frac{5T}{4{,}17\,T} \approx 1{,}20
\]
çıkar. Bu da yaklaşık \%20'lik bir kazançtır. Kazanç, kısa buyrukların en yavaş buyruğun bütün aşamalarını beklemek zorunda kalmamasından doğar.

Bu kazancı elde etmek için veri yolu aşamalara bölünür ve her aşamada üretilen ara değerlerin bir sonraki çevrime taşınabilmesi için birimlerin çıkışlarına \emph{ara yazmaçlar} eklenir. Aksi halde bu değerler bir sonraki saat vuruşunda kaybolur (Şekil~\ref{fig:ara-yazmac}). Yukarıdaki kestirim, aşamaların eşit gecikmede olduğu varsayımına dayanır. Gerçekte aşamalar eşit değildir (bellek erişimi, yazmaç okumadan çok daha uzundur) ve çevrim en yavaş aşamaya göre ayarlandığından kazanç azalır, hatta tersine dönebilir. Bu gerçekçi karşılaştırmayı bu bölümün ilerleyen örneğinde yapacağız.

\begin{figure}[htbp]
\centering
\begin{tikzpicture}[>=Stealth, node distance=2cm]
  % Giriş
  \node[font=\small] (giris) {Giriş};
  % Mantık birimi
  \node[blok, right=1.5cm of giris, minimum width=2.5cm] (mantik) {Mantık Birimi};
  % Ara yazmaç
  \node[yazmac, right=1.5cm of mantik, minimum height=1.5cm, minimum width=0.4cm] (yzm) {};
  \node[font=\tiny, rotate=90] at (yzm) {Yazmaç};
  % Çıkış
  \node[font=\small, right=1.5cm of yzm] (cikis) {Çıkış};
  % Bağlantılar
  \draw[veriyolu] (giris) -- (mantik);
  \draw[veriyolu] (mantik) -- (yzm);
  \draw[veriyolu] (yzm) -- (cikis);
  % Saat
  \draw[denetimstil] (yzm.south) -- ++(0,-1) node[below, font=\footnotesize, text=sinyalkirmizi]{Saat};
  % Açıklama
  \node[font=\footnotesize, text=gray, align=center, below=2cm of mantik] {Yazmaç, saat vuruşu arasında\\veriyi tutarak bir sonraki\\aşamaya aktarır.};
\end{tikzpicture}
\caption{Mantık birimine ara yazmaç konulması}
\label{fig:ara-yazmac}
\end{figure}

Veri yolu parçalara bölünürken tüm aşamaların tuttuğu sürenin birbirine hemen hemen eşit olması gerekir; çünkü çevrim süresi en yavaş aşamaya göre belirlenir. Tablo~\ref{tab:veriyolu-bolme} aynı veri yolunun iki ayrı bölünmesini karşılaştırır.

\begin{table}[htbp]
\centering
\caption{Tek vuruşluk veri yolunun birimlere ayrılması}
\label{tab:veriyolu-bolme}
\begin{tabular}{lccc}
\toprule
 & \textbf{Çok vuruşluk 1} & \textbf{Çok vuruşluk 2} & \textbf{Tek vuruşluk} \\
\midrule
1.\ birim & 5 ns & 17 ns & \multirow{5}{*}{90 ns} \\
2.\ birim & 15 ns & 19 ns & \\
3.\ birim & 10 ns & 18 ns & \\
4.\ birim & 10 ns & 19 ns & \\
5.\ birim & 50 ns & 17 ns & \\
\midrule
\textbf{Çevrim süresi} & \textbf{50 ns} & \textbf{19 ns} & \textbf{90 ns} \\
\bottomrule
\end{tabular}
\end{table}

Tablodaki ``Çok vuruşluk 2'' bölünmesinde aşama süreleri birbirine yakındır ve çevrim süresi 19 ns'ye iner. Dengesiz ``Çok vuruşluk 1'' bölünmesinde ise çevrim en yavaş aşamaya (50 ns) takılır. Bu bölümde tasarlanacak veri yolu beş aşamadan oluşur (getir--çöz--yürüt--bellek--sonuç yaz). Aşamaların arasındaki bölünme noktalarına ara yazmaçlar yerleştirilir (Şekil~\ref{fig:bolunme-noktalari}).

\begin{figure}[htbp]
\centering
\begin{tikzpicture}[>=Stealth]
  % Aşama kutuları
  \node[blok, minimum width=2.2cm, minimum height=1.2cm] (s1) at (0,0) {Getir};
  \node[blok, minimum width=2.2cm, minimum height=1.2cm] (s2) at (3,0) {Çöz};
  \node[blok, minimum width=2.2cm, minimum height=1.2cm] (s3) at (6,0) {Yürüt};
  \node[blok, minimum width=2.2cm, minimum height=1.2cm] (s4) at (9,0) {Bellek};
  \node[blok, minimum width=2.2cm, minimum height=1.2cm] (s5) at (12,0) {Sonuç\\Yaz};
  % Oklar
  \draw[veriyolu] (s1) -- (s2);
  \draw[veriyolu] (s2) -- (s3);
  \draw[veriyolu] (s3) -- (s4);
  \draw[veriyolu] (s4) -- (s5);
  % Bölünme noktaları (dikey kesikli çizgiler)
  \foreach \x in {1.5, 4.5, 7.5, 10.5} {
    \draw[dashed, thick, sinyalkirmizi] (\x, -1.2) -- (\x, 1.2);
  }
  % Yazmaç etiketleri
  \node[font=\tiny, text=sinyalkirmizi, rotate=90] at (1.5, -1.6) {Ara yazmaç};
  \node[font=\tiny, text=sinyalkirmizi, rotate=90] at (4.5, -1.6) {Ara yazmaç};
  \node[font=\tiny, text=sinyalkirmizi, rotate=90] at (7.5, -1.6) {Ara yazmaç};
  \node[font=\tiny, text=sinyalkirmizi, rotate=90] at (10.5, -1.6) {Ara yazmaç};
\end{tikzpicture}
\caption{İşlemcinin bölüneceği noktalar}
\label{fig:bolunme-noktalari}
\end{figure}

Her bölünme noktasına, o aşamada üretilen değeri bir sonraki çevrime taşıyan bir \emph{ara yazmaç} yerleştirilir (Şekil~\ref{fig:ara-yazmaclar}). Bu yazmaçların her biri belirli bir değeri tutar. \emph{Buyruk yazmacı} bellekten getirilen buyruğu, \emph{A} ve \emph{B} yazmaçları yazmaç öbeğinden okunan iki yazmaç değerini, \emph{AMB yazmacı} AMB'nin ürettiği sonucu, \emph{bellek yazmacı} ise bellekten okunan veriyi bir sonraki aşamada kullanılmak üzere saklar.

\begin{figure}[htbp]
\centering
\begin{tikzpicture}[>=Stealth, font=\footnotesize,
    rol/.style={font=\scriptsize, text=gray!55!black, align=center, anchor=north, text width=2.5cm}]
  % Çevrim etiketleri (üst)
  \node[font=\small\bfseries] at (0,2.5)  {Getir};
  \node[font=\small\bfseries] at (3,2.5)  {Çöz};
  \node[font=\small\bfseries] at (6,2.5)  {Yürüt};
  \node[font=\small\bfseries] at (9,2.5)  {Bellek};
  \node[font=\small\bfseries] at (12,2.5) {Sonucu Yaz};
  % Çevrimler arası ayraçlar
  \foreach \x in {1.5,4.5,7.5,10.5} { \draw[gray!45, dashed] (\x,2.9) -- (\x,-1.5); }
  % Her çevrimin çıktısını tutan ara yazmaç
  \node[yazmac, minimum height=1.7cm, minimum width=0.55cm] (byz) at (0,0.7) {};
  \node[font=\scriptsize, rotate=90] at (byz) {Buyruk Yzm.};
  \node[yazmac, minimum height=0.7cm, minimum width=0.55cm] (areg) at (3,1.2) {};
  \node[font=\scriptsize] at (areg) {A};
  \node[yazmac, minimum height=0.7cm, minimum width=0.55cm] (breg) at (3,0.2) {};
  \node[font=\scriptsize] at (breg) {B};
  \node[yazmac, minimum height=1.7cm, minimum width=0.55cm] (ambyz) at (6,0.7) {};
  \node[font=\scriptsize, rotate=90] at (ambyz) {AMB Yzm.};
  \node[yazmac, minimum height=1.7cm, minimum width=0.55cm] (belyz) at (9,0.7) {};
  \node[font=\scriptsize, rotate=90] at (belyz) {Bellek Yzm.};
  \node[font=\scriptsize, text=gray!60!black, align=center, text width=2cm] at (12,0.7) {(ara yazmaç\\gerekmez)};
  % Rol notları (alt)
  \node[rol] at (0,-0.4)  {getirilen buyruk};
  \node[rol] at (3,-0.4)  {okunan iki yazmaç değeri};
  \node[rol] at (6,-0.4)  {AMB sonucu};
  \node[rol] at (9,-0.4)  {bellekten okunan veri};
  \node[rol] at (12,-0.4) {sonuç doğrudan yazmaç öbeğine yazılır};
  % Çevrim yönü
  \draw[->, thick] (-0.9,-2.0) -- (12.9,-2.0) node[right, font=\scriptsize] {çevrim sırası};
\end{tikzpicture}
\caption[Çok vuruşluk tasarıma eklenecek ara yazmaçlar]{Eklenecek ara yazmaçlar: her çevrimin ürettiği değeri bir sonraki çevrimde kullanılmak üzere tutarlar. Her buyruk bütün çevrimleri kullanmaz; örneğin R türü Bellek çevrimini, dallanma ise hem Bellek hem Sonucu Yaz çevrimini atlar.}
\label{fig:ara-yazmaclar}
\end{figure}

Ara yazmaç eklemek hız kazandırır, ancak veri yoluna fazladan donanım yükler. Çok vuruşluk tasarımın asıl üstünlüğü, aynı birimin farklı çevrimlerde yeniden kullanılabilmesidir. Buyruk ile veri bellekleri tek bir belleğe birleştirilir ve ayrı program sayacı toplayıcılarına gerek kalmadığından tek bir AMB hem adres hesaplarını hem aritmetik işlemleri üstlenir.

Belleğin bu biçimde paylaşılması, tek vuruşluk tasarımda bulunmayan yeni bir denetim gerektirir: \texttt{BuyruğaYaz}. Tek vuruşlukta buyruk belleği ayrıydı ve sürekli okunduğundan buyruk kendiliğinden veri yolunda dururdu. Çok vuruşlukta ise aynı bellek, Bellek aşamasında \texttt{lw} buyruğunun verisini de okur. Buyruk yazmacı bu yüzden getirilen buyruğu yalnızca Getir çevriminde yakalamalı ve buyruğun geri kalan çevrimleri boyunca değişmeden tutmalıdır. \texttt{BuyruğaYaz} yalnızca Getir çevriminde 1 yapılır, sonraki çevrimlerde 0 kalır. Bu denetim olmasaydı bir \texttt{lw} buyruğu Bellek aşamasında kendi verisini okurken bu değer buyruğun üzerine yazılır ve buyruktan türeyen bütün denetim sinyalleri işin ortasında kaybolurdu.

Program sayacının Getir çevriminde PS$+$4'e ilerletilmesi de ikinci bir ekleme gerektirir: \texttt{EskiPS} yazmacı. RISC-V'te dallanma ve \texttt{jal} hedefi, dallanma buyruğunun kendi adresine göre hesaplanır (adres $+$ anlık değer). Oysa hedef Çöz çevriminde hesaplandığında PS çoktan PS$+$4 olmuştur, bu yüzden PS $+$ anlık dört bayt kaymış bir adres verirdi. \texttt{EskiPS}, Getir çevriminde buyruğun asıl adresini saklar ve yazma iznini \texttt{BuyruğaYaz} ile paylaşır. Böylece hedef, EskiPS $+$ anlık ile tam doğru hesaplanır. Bu düzen \texttt{jal} buyruğunu da doğal biçimde çözer. PS zaten dönüş adresi olan PS$+$4 değerini tutarken, EskiPS $+$ anlık atlama hedefini verir.

Bütün bu bileşenler bir araya geldiğinde Şekil~\ref{fig:cok-vurusluk-veriyolu}'deki çok vuruşluk veri yolu elde edilir.

\begin{figure}[htbp]
\centering
% Cok vurusluk veri yolu -- TAM PAYLASIMLI tasarim (Secenek A):
% Top4 ve Dal.Top YOK; PS+4 Getir'de, dal/jal hedefi Coz'de AMB'de hesaplanir.
% Eski cizim: fig621-eski-yedek.tex.bak
% Tum \draw segmentleri eksenel (yalnizca yatay/dikey); |- ve -| KULLANILMAZ.
% Kesisen ama BAGLANMAYAN teller noktasiz kesisir; baglantilar siyah nokta ile.
\begin{tikzpicture}[>={Stealth[length=2mm, width=1.6mm]}, semithick, font=\footnotesize, scale=0.68, transform shape,
    yzm/.style={draw, fill=yellow!25, minimum width=1cm, minimum height=1cm, align=center},
    bllk/.style={draw, fill=green!20, minimum width=2.0cm, minimum height=2.4cm, align=center},
    yobk/.style={draw, fill=blue!10, minimum width=2.2cm, minimum height=3.0cm, align=center},
    amb/.style={draw, fill=red!20, minimum width=1.6cm, minimum height=2.0cm, align=center},
    anlikb/.style={draw, fill=orange!25, minimum width=1.4cm, minimum height=0.9cm, align=center},
    muxdik/.style={draw, fill=purple!15, minimum width=0.7cm, minimum height=1.0cm, align=center},
    kapi/.style={draw, fill=gray!12, minimum width=0.55cm, minimum height=0.45cm, align=center, font=\tiny},
    bitayar/.style={fill=white, inner sep=1pt, font=\tiny, text=black!70},
    sabit/.style={draw, fill=gray!10, minimum width=0.6cm, minimum height=0.5cm, align=center},
    geri/.style={->, draw=black!45, thin},
]

  % ============================================================
  % DUGUMLER -- sol->sag akis
  % ============================================================

  % --- PSKaynak coklaycisi (2 girisli: 0=AMB cikisi/PS+4, 1=AMB Yzm./dal-jal hedefi) ---
  \node[muxdik, minimum height=1.2cm] (psmux) at (0.3, 0.0) {};
  \node[font=\tiny] at (psmux.center) {Çok.};
  \node[font=\tiny, anchor=west] at (-0.03, 0.27) {0};
  \node[font=\tiny, anchor=west] at (-0.03,-0.27) {1};

  % --- PS ---
  \node[yzm] (ps) at (2.0, 0.0) {PS};

  % --- PS yazma izni kapilari: PSYaz VEYA (PSKoşul VE (Sıfır XOR i1)) ---
  \node[kapi] (vekapi)   at (1.1, 2.0) {VE};
  \node[kapi, minimum width=0.6cm] (veyakapi) at (2.0, 1.2) {VEYA};
  % Sifir XOR i1 kapisi: beq'te (i1=0) Sifir'i, bne'de (i1=1) degilini gecirir
  \node[kapi, minimum width=0.5cm, minimum height=0.4cm] (xorkapi) at (1.1, 2.52) {$=1$};

  % --- AdresKaynagi coklaycisi (0=PS, 1=AMB Yzm.) ---
  \node[muxdik, minimum height=1.2cm] (adrmux) at (3.8, 0.0) {};
  \node[font=\tiny] at (adrmux.center) {Çok.};
  \node[font=\tiny, anchor=west] at (3.47, 0.27) {0};
  \node[font=\tiny, anchor=west] at (3.47,-0.27) {1};

  % --- Birlesik Bellek (Buyruk+Veri) ---
  \node[bllk] (bel) at (6.5, 0.4) {Bellek\\(Buyruk\\+Veri)};
  \node[font=\scriptsize, anchor=west] at (5.55, 1.2) {$A_{\text{adr}}$};
  \node[font=\scriptsize, anchor=west] at (5.55,-0.55) {$V_{\text{yaz}}$};

  % --- Buyruk Yazmaci (IR, ara yazm.) ---
  \node[yzm, minimum width=0.55cm, minimum height=2.0cm] (iryz) at (8.2, 1.2) {};
  \node[font=\scriptsize, rotate=90] at (8.2, 1.2) {Buyruk Yzm.};

  % --- Bellek Veri Yazmaci (MDR, ara yazm.) ---
  \node[yzm, minimum width=0.55cm, minimum height=1.9cm] (mdryz) at (8.2,-0.9) {};
  \node[font=\scriptsize, rotate=90] at (8.2,-0.9) {Bellek Yzm.};

  % --- Yazmac Obegi ---
  \node[yobk] (yob) at (11.0, 0.0) {};
  \node[anchor=center, font=\footnotesize\bfseries, align=center] at (yob.center) {Yazmaç\\Öbeği};
  \node[font=\scriptsize, anchor=north] at (10.5, 1.55) {$A_{\text{adr}}$};
  \node[font=\scriptsize, anchor=north] at (11.0, 1.55) {$B_{\text{adr}}$};
  \node[font=\scriptsize, anchor=north] at (11.5, 1.55) {$C_{\text{adr}}$};
  \node[font=\scriptsize, anchor=east]  at (12.05, 0.5)  {$A_{\text{veri}}$};
  \node[font=\scriptsize, anchor=east]  at (12.05,-0.5)  {$B_{\text{veri}}$};
  \node[font=\tiny, anchor=south] at (11.5,-1.45) {$C_{\text{veri}}$};

  % --- Anlik Ureteci ---
  \node[anlikb] (anlik) at (11.0,-3.5) {Anlık\\Üreteci};

  % --- A ve B Yazmaclari (ara yazm.) ---
  \node[yzm, minimum width=0.55cm, minimum height=1.2cm] (ayz) at (13.0, 0.9) {};
  \node[font=\scriptsize, rotate=90] at (13.0, 0.9) {A Yzm.};
  \node[yzm, minimum width=0.55cm, minimum height=1.2cm] (byz) at (13.0,-0.9) {};
  \node[font=\scriptsize, rotate=90] at (13.0,-0.9) {B Yzm.};

  % --- AMBKaynakA coklaycisi (3 girisli: 10=EskiPS ust, 01=PS orta, 00=A Yzm. alt) ---
  \node[muxdik, minimum height=1.6cm] (amuxa) at (14.7, 0.9) {};
  \node[font=\tiny, rotate=90] at ([xshift=0.18cm]amuxa.center) {Çok.};
  \node[font=\tiny, anchor=west] at (14.37, 1.35) {10};
  \node[font=\tiny, anchor=west] at (14.37, 0.90) {01};
  \node[font=\tiny, anchor=west] at (14.37, 0.45) {00};

  % --- EskiPS yazmaci (ara yazm.): Getir'de PS'yi saklar; yazma izni BuyrugaYaz ile ortak ---
  \node[yzm, minimum width=0.55cm, minimum height=0.85cm] (eskips) at (13.0, 2.35) {};
  \node[font=\scriptsize, rotate=90] at (13.0, 2.35) {EskiPS};

  % --- AMBKaynakB coklaycisi (00=B Yzm., 01=4, 10=anlik) ---
  \node[muxdik, minimum height=1.5cm] (amuxb) at (14.7,-1.5) {};
  \node[font=\tiny, rotate=90] at ([xshift=0.18cm]amuxb.center) {Çok.};
  \node[font=\tiny, anchor=west] at (14.37,-1.1) {00};
  \node[font=\tiny, anchor=west] at (14.37,-1.5) {01};
  \node[font=\tiny, anchor=west] at (14.37,-1.9) {10};

  % --- 4 sabiti ---
  \node[sabit] (dort) at (13.2,-2.5) {$4$};

  % --- AMB ---
  \node[amb] (amb) at (16.5, 0.0) {AMB};

  % --- AMB Yazmaci (ara yazm.) ---
  \node[yzm, minimum width=0.55cm, minimum height=1.6cm] (ambyz) at (18.2, 0.0) {};
  \node[font=\scriptsize, rotate=90] at (18.2, 0.0) {AMB Yzm.};

  % --- GeriYazmaKaynagi coklaycisi (00=AMB Yzm., 01=Bellek Yzm., 10=PS) ---
  \node[muxdik, minimum height=1.5cm] (gerimux) at (20.0,-1.5) {};
  \node[font=\tiny, rotate=90] at ([xshift=0.18cm]gerimux.center) {Çok.};
  \node[font=\tiny, anchor=west] at (19.67,-1.1) {00};
  \node[font=\tiny, anchor=west] at (19.67,-1.5) {01};
  \node[font=\tiny, anchor=west] at (19.67,-1.9) {10};

  % ============================================================
  % VERI TELLERI -- her ardisik koordinat cifti tek eksende degisir
  % ============================================================

  % (Y1) psmux cikisi -> PS girisi
  \draw[->] (0.65, 0.0) -- (1.5, 0.0);

  % (Y2) PS dogu -> AdresKaynagi giris 0 (Getir adresi)
  \draw[->] (2.5, 0.0) -- (3.0, 0.0) -- (3.0, 0.27) -- (3.43, 0.27);
  \node[circle, fill=black, inner sep=1pt] at (3.0, 0.0) {};

  % (Y3) PS -> AMBKaynakA giris 01 (ust kanal y=3.0)
  \draw[geri] (3.0, 0.0) -- (3.0, 3.0) -- (14.0, 3.0) -- (14.0, 0.9) -- (14.33, 0.9);
  \node[circle, fill=black, inner sep=1pt] at (3.0, 3.0) {};
  \node[bitayar, anchor=south] at (8.5, 3.02) {PS};

  % (Y3b) PS -> EskiPS girisi (kanal y=2.6, PS telinden ayrilir)
  \draw[geri] (3.0, 2.6) -- (12.5, 2.6) -- (12.5, 2.35) -- (12.725, 2.35);
  \node[circle, fill=black, inner sep=1pt] at (3.0, 2.6) {};
  % (Y3c) EskiPS -> AMBKaynakA giris 10 (PS telini noktasiz keser x=14.0)
  \draw[->] (13.275, 2.35) -- (14.15, 2.35) -- (14.15, 1.35) -- (14.33, 1.35);

  % (Y4) PS -> GeriYazma giris 10 (jal donus adresi; ust kanal y=3.8)
  \draw[geri] (3.0, 3.0) -- (3.0, 3.8) -- (20.5, 3.8) -- (20.5,-2.55) -- (19.4,-2.55) -- (19.4,-1.9) -- (19.63,-1.9);
  \node[bitayar, anchor=south] at (10.5, 3.82) {PS $\to$ jal dönüş adresi};

  % (Y5) AdresKaynagi -> Bellek adres girisi
  \draw[->] (4.15, 0.0) -- (4.8, 0.0) -- (4.8, 1.2) -- (5.5, 1.2);

  % (Y6) Bellek -> Buyruk Yzm. ve Bellek Yzm.
  \draw[->] (7.5, 1.6) -- (7.925, 1.6);
  \draw[->] (7.5,-0.4) -- (7.925,-0.4);

  % (Y7) Buyruk Yzm. -> Yazmac Obegi adresleri (A, B, C)
  \node[circle, fill=black, inner sep=1pt] at (9.2, 1.2) {};
  \draw[-]  (8.475, 1.2) -- (9.2, 1.2) -- (9.2, 2.0);
  \node[circle, fill=black, inner sep=1pt] at (9.2, 2.0) {};
  \draw[->] (9.2, 2.0) -- (10.5, 2.0) -- (10.5, 1.55);
  \node[circle, fill=black, inner sep=1pt] at (10.5, 2.0) {};
  \draw[->] (10.5, 2.0) -- (11.0, 2.0) -- (11.0, 1.55);
  \node[circle, fill=black, inner sep=1pt] at (11.0, 2.0) {};
  \draw[->] (11.0, 2.0) -- (11.5, 2.0) -- (11.5, 1.55);

  % (Y8) Buyruk Yzm. -> Anlik Ureteci
  \draw[->] (9.2, 1.2) -- (9.2,-3.5) -- (10.3,-3.5);

  % (Y9) Yazmac Obegi -> A ve B yazmaclari
  \draw[->] (12.1, 0.5) -- (12.5, 0.5) -- (12.5, 0.9) -- (12.725, 0.9);
  \draw[->] (12.1,-0.5) -- (12.5,-0.5) -- (12.5,-0.9) -- (12.725,-0.9);

  % (Y10) A Yzm. -> AMBKaynakA giris 00 (erken asagi kirilir; PS teliyle ayni hizada kalmasin)
  \draw[->] (13.275, 0.9) -- (13.5, 0.9) -- (13.5, 0.45) -- (14.33, 0.45);

  % (Y11) B Yzm. -> AMBKaynakB giris 00 (tap: belleğe yazma için)
  \draw[->] (13.275,-0.9) -- (13.7,-0.9) -- (13.7,-1.1) -- (14.33,-1.1);
  \node[circle, fill=black, inner sep=1pt] at (13.55,-0.9) {};

  % (Y12) B Yzm. -> Bellek V_yaz (sw verisi; alt kanal y=-5.8)
  \draw[geri] (13.55,-0.9) -- (13.55,-5.8) -- (5.2,-5.8) -- (5.2,-0.55) -- (5.5,-0.55);
  \node[bitayar, anchor=north] at (9.0,-5.82) {B Yzm. $\to$ belleğe yazılacak veri (sw)};

  % (Y13) 4 sabiti -> AMBKaynakB giris 01
  \draw[->] (13.5,-2.5) -- (13.95,-2.5) -- (13.95,-1.5) -- (14.33,-1.5);

  % (Y14) Anlik -> AMBKaynakB giris 10
  \draw[->] (11.7,-3.5) -- (12.5,-3.5) -- (12.5,-1.9) -- (14.33,-1.9);

  % (Y15) AMBKaynakA/B -> AMB girisleri
  \draw[->] (15.05, 0.9) -- (15.4, 0.9) -- (15.4, 0.5) -- (15.8, 0.5);
  \draw[->] (15.05,-1.5) -- (15.4,-1.5) -- (15.4,-0.5) -- (15.8,-0.5);

  % (Y16) AMB -> AMB Yzm. (tap: PS+4 geri beslemesi icin)
  \draw[->] (17.3, 0.0) -- (17.925, 0.0);
  \node[circle, fill=black, inner sep=1pt] at (17.6, 0.0) {};

  % (Y17) AMB cikisi -> PSKaynak giris 00 (Getir'de PS+4; ust kanal y=3.4)
  \draw[geri] (17.6, 0.0) -- (17.6, 3.4) -- (-0.9, 3.4) -- (-0.9, 0.27) -- (-0.05, 0.27);
  \node[bitayar, anchor=south] at (8.5, 3.42) {PS$+$4 (AMB çıkışı)};

  % (Y18) AMB Yzm. dogu tap noktasi
  \node[circle, fill=black, inner sep=1pt] at (19.0, 0.0) {};

  % (Y19) AMB Yzm. -> PSKaynak giris 01 ve 10 (dal/jal hedefi; en ust kanal y=4.2)
  \draw[geri] (18.475, 0.0) -- (19.0, 0.0) -- (19.0, 4.2) -- (-1.2, 4.2) -- (-1.2,-0.27) -- (-0.05,-0.27);
  \node[bitayar, anchor=south] at (8.5, 4.22) {dal/jal hedefi (AMB Yzm.)};

  % (Y20) AMB Yzm. -> GeriYazma giris 00
  \draw[->] (19.0, 0.0) -- (19.0,-1.1) -- (19.63,-1.1);
  \node[circle, fill=black, inner sep=1pt] at (19.0,-1.1) {};

  % (Y21) AMB Yzm. -> AdresKaynagi giris 1 (veri erisim adresi; alt kanal y=-4.75)
  \draw[geri] (19.0,-1.1) -- (19.0,-4.75) -- (2.85,-4.75) -- (2.85,-0.27) -- (3.43,-0.27);
  \node[bitayar, anchor=south] at (5.8,-4.73) {veri erişim adresi (AMB Yzm.)};

  % (Y22) Bellek Yzm. -> GeriYazma giris 01 (alt kanal y=-5.2)
  \draw[geri] (8.475,-0.8) -- (9.5,-0.8) -- (9.5,-5.2) -- (19.3,-5.2) -- (19.3,-1.5) -- (19.63,-1.5);
  \node[bitayar, anchor=north] at (14.0,-5.22) {Bellek Yzm. $\to$ geri yazma};

  % (Y23) GeriYazma cikisi -> Yazmac Obegi C_veri (alt kanal y=-6.4)
  \draw[geri] (20.35,-1.5) -- (21.0,-1.5) -- (21.0,-6.4) -- (12.1,-6.4) -- (12.1,-2.0) -- (11.5,-2.0) -- (11.5,-1.5);
  \node[bitayar, anchor=north] at (15.5,-6.42) {geri yazma ($C_{\text{veri}}$)};

  % ============================================================
  % SIFIR CIKISI ve PS YAZMA IZNI
  % ============================================================

  % AMB Sifir cikisi -> XOR kapisi (kanal y=2.85; kapiya bitisik girdigi icin ok ucu yok)
  \draw[draw=black!45, thin] (17.0, 1.0) -- (17.0, 2.85) -- (1.1, 2.85) -- (1.1, 2.72);
  \node[bitayar, anchor=south] at (15.8, 2.87) {Sıfır};

  % i1 (buyruk[12]) -> XOR sol girdi; XOR cikisi -> VE ust girdi
  \node[font=\scriptsize, text=sinyallacivert, anchor=east] at (0.5, 2.52) {$b_1$};
  \draw[->, sinyallacivert, thick] (0.55, 2.52) -- (0.85, 2.52);
  \draw[draw=black!45, thin] (1.1, 2.32) -- (1.1, 2.225);

  % PSKosul -> VE, PSYaz -> VEYA, VE -> VEYA, VEYA -> PS yazma izni
  \node[font=\scriptsize, text=sinyallacivert, anchor=east] at (0.45, 2.0) {PSKoşul};
  \draw[->, sinyallacivert, thick] (0.5, 2.0) -- (0.825, 2.0);
  \node[font=\scriptsize, text=sinyallacivert, anchor=east] at (1.0, 1.08) {PSYaz};
  \draw[->, sinyallacivert, thick] (1.05, 1.08) -- (1.7, 1.08);
  \draw[->, sinyallacivert, thick] (1.375, 2.0) -- (1.55, 2.0) -- (1.55, 1.32) -- (1.7, 1.32);
  \draw[->, sinyallacivert, thick] (2.0, 0.95) -- (2.0, 0.5);

  % ============================================================
  % DENETIM SINYALLERI (sinyallacivert)
  % ============================================================

  % PSKaynak -> psmux guney
  \node[font=\scriptsize, text=sinyallacivert, anchor=north] at (0.3,-1.35) {PSKaynak};
  \draw[->, sinyallacivert, thick] (0.3,-1.3) -- (0.3,-0.6);

  % AdresKaynagi -> adrmux guney
  \node[font=\scriptsize, text=sinyallacivert, anchor=north, align=center] at (3.8,-1.15) {AdresKaynağı};
  \draw[->, sinyallacivert, thick] (3.8,-1.1) -- (3.8,-0.6);

  % BellegeYaz -> bellek guney
  \node[font=\scriptsize, text=sinyallacivert, anchor=north] at (6.5,-1.35) {BelleğeYaz};
  \draw[->, sinyallacivert, thick] (6.5,-1.3) -- (6.5,-0.8);

  % BuyrugaYaz -> Buyruk Yzm. kuzey (etiket okun sagina; y=2.6 kanaliyla cakismasin)
  \node[font=\scriptsize, text=sinyallacivert, anchor=west] at (8.32, 2.38) {BuyruğaYaz};
  \draw[->, sinyallacivert, thick] (8.2, 2.52) -- (8.2, 2.2);

  % YazmacaYaz -> yob kuzey (etiket okun soluna; y=2.6 kanaliyla cakismasin)
  \node[font=\scriptsize, text=sinyallacivert, anchor=east] at (11.78, 2.38) {YazmacaYaz};
  \draw[->, sinyallacivert, thick] (11.9, 2.52) -- (11.9, 1.5);

  % AnlikBicim -> anlik guney
  \node[font=\scriptsize, text=sinyallacivert, anchor=north] at (11.0,-4.35) {AnlıkBiçim};
  \draw[->, sinyallacivert, thick] (11.0,-4.3) -- (11.0,-3.95);

  % AMBKaynakA -> amuxa guney (2 bit)
  \node[font=\scriptsize, text=sinyallacivert, anchor=north] at (14.7,-0.1) {AMBKaynakA};
  \draw[->, sinyallacivert, thick] (14.7, 0.0) -- (14.7, 0.1);

  % AMBKaynakB -> amuxb guney
  \node[font=\scriptsize, text=sinyallacivert, anchor=north] at (14.7,-2.75) {AMBKaynakB};
  \draw[->, sinyallacivert, thick] (14.7,-2.7) -- (14.7,-2.25);

  % AMBIslem -> amb kuzey
  \node[font=\scriptsize, text=sinyallacivert, anchor=south] at (16.1, 1.55) {AMBİşlem};
  \draw[->, sinyallacivert, thick] (16.1, 1.5) -- (16.1, 1.0);

  % GeriYazmaKaynagi -> gerimux guney
  \node[font=\scriptsize, text=sinyallacivert, anchor=north, align=center] at (20.3,-2.75) {GeriYazma\\Kaynağı};
  \draw[->, sinyallacivert, thick] (20.3,-2.7) -- (20.3,-2.25);

\end{tikzpicture}
\caption[Çok vuruşluk veri yolu: tek birleşik bellek ve tek AMB]{Çok vuruşluk veri yolu: tek birleşik bellek ve tek AMB ile tam paylaşımlı tasarım.}
\label{fig:cok-vurusluk-veriyolu}
\end{figure}

Bu veri yolu kaynak paylaşımı üzerine kuruludur. Tek birleşik bellek hem buyruk getirme (Getir) hem de veri erişimi (Bellek) çevrimlerinde kullanılır. Tek AMB ise PS$+$4 hesabını (Getir), dal ve jal hedef adresini (Çöz) ve asıl aritmetik işlemi (Yürüt) sırayla üstlenir. Çöz aşamasında hedef adres EskiPS $+$ anlık olarak hesaplanır. EskiPS, Getir çevriminde buyrukla birlikte (BuyruğaYaz izniyle) saklanan buyruk adresidir. Böylece hedef, RISC-V tanımındaki gibi buyruğun kendi adresine göre çıkar. AMBKaynakA çoklayıcısı AMB'nin A girişini üç kaynak arasından seçer: A yazmacı (00), PS (01) ve EskiPS (10). Sarı ara yazmaçlar bir çevrimin ürettiği değeri sonraki çevrimlere taşır.

PSKaynak çoklayıcısının iki girişi çevrim zamanlamasına göre ayrışır. \textbf{0} girişi AMB'nin \emph{o çevrimdeki} çıkışını, Getir'de hesaplanan PS$+$4 değerini doğrudan alır. \textbf{1} girişi ise bir önceki çevrimde (Çöz) hesaplanıp AMB yazmacında bekleyen dal ya da jal hedefini alır. Dal ile jal ayrımı veri kaynağında değil yazma izninde yapılır. Koşullu dallanmada PS'ye yazma, PSKoşul VE (Sıfır XOR $b_1$) koşuluyla, jal'da ise PSYaz ile etkinleşir. Buradaki $b_1 = $ buyruk[12] değeri beq'te 0, bne'de 1 olduğundan aynı donanım her iki dallanmayı da yürütür. Şekilde ``PS'' teli PS'yi AMB'nin A girişine taşır (girdi), ``PS$+$4 (AMB çıkışı)'' teli ise AMB'nin sonucunu PSKaynak'a geri getirir (çıktı). Bu iki tel kısa devre değil, program sayacını her çevrimde artıran geri besleme halkasının iki ucudur.

\begin{ornek}[6.5 -- Çok vuruşluk işlemcide aşama gecikmeleri ve çevrim zamanı]
Örnek~6.4'teki birim gecikmelerini kullanarak Şekil~\ref{fig:cok-vurusluk-veriyolu}'deki çok vuruşluk veri yolu için (a)~her aşamanın kritik yol gecikmesini hesaplayın, (b)~çevrim zamanını ve saat sıklığını belirleyin, (c)~her buyruk türünün toplam yürütme süresini bulup tek vuruşluk tasarımla karşılaştırın. (Ara yazmaçların kurulum ve yayılım süreleri ile kapı gecikmeleri yalınlık için ihmal edilsin.)

\tcblower
\textbf{Çözüm.}

\textbf{(a)} Tek vuruşlukta kritik yol buyruk boyunca uzanıyordu. Çok vuruşlukta ise her aşama kendi içinde bir kritik yola sahiptir ve aşamanın sonunda değer bir ara yazmaca kilitlenir:

\begin{center}
\small
\begin{tabular}{@{}lp{8.2cm}r@{}}
\toprule
\textbf{Aşama} & \textbf{Kritik yol} & \textbf{Gecikme} \\
\midrule
Getir & Çok.(AdresKaynağı) + Bellek; paralel PS$+$4 yolu (Çok. + AMB + Çok.) 150~ps'te kalır & 225\,ps \\
Çöz & Üreteç + Çok.(AMBKaynakB) + AMB (dal/jal hedef hesabı); yazmaç okuma (50~ps) paralel yürür & 150\,ps \\
Yürüt & Üreteç + Çok. + AMB (en kötü durum: ori); R türü ve dallanmada 125~ps & 150\,ps \\
Bellek & Çok.(AdresKaynağı) + Bellek & 225\,ps \\
Sonuç Yaz & Çok.(GeriYazma) + Yazmaç öbeği (yazma) & 75\,ps \\
\bottomrule
\end{tabular}
\end{center}

\textbf{(b)} Çevrim zamanını en yavaş \emph{aşama} belirler: $t_{\text{çevrim}} = 225$~ps (belleğe erişen Getir ve Bellek aşamaları).

\[
\text{Saat sıklığı} = \frac{1}{225 \times 10^{-12}} \approx 4{,}44 \text{~GHz}
\]

Tek vuruşluğun 625~ps'lik çevrimine karşılık çevrim 2{,}8 kat kısalmıştır. Saat artık en yavaş buyruğa değil en yavaş aşamaya ayarlanır.

\textbf{(c)} Her buyruğun toplam süresi, kullandığı çevrim sayısı (Tablo~\ref{tab:buyruk-asamalar}) ile çevrim zamanının çarpımıdır:

\begin{center}
\small
\begin{tabular}{@{}lccc@{}}
\toprule
\textbf{Buyruk} & \textbf{Çevrim sayısı} & \textbf{Çok vuruşluk süre} & \textbf{Tek vuruşluk süre} \\
\midrule
Yükle (lw) & 5 & 1125\,ps & 625\,ps \\
Sakla (sw) & 4 & 900\,ps & 625\,ps \\
R türü & 4 & 900\,ps & 625\,ps \\
ori & 4 & 900\,ps & 625\,ps \\
beq/bne & 3 & 675\,ps & 625\,ps \\
jal & 3 & 675\,ps & 625\,ps \\
\bottomrule
\end{tabular}
\end{center}

Sonuç çarpıcıdır. Çevrim 2{,}8 kat kısalmasına karşın \emph{tek başına her buyruk} çok vuruşlukta daha uzun sürer, çünkü aşama sınırlarındaki bekleme her çevrimi en yavaş aşamaya (225~ps) yuvarlar. Kısa buyruklar (675~ps) ile uzun buyruklar (1125~ps) arasındaki fark ise artık donanıma yansımaktadır. Ortalamada hangi tasarımın öne geçtiğini buyruk karışımı belirler. Bu hesabı Örnek~6.7'de yapacağız. Çok vuruşluğun asıl meyvesini ise aşamaları örtüştüren boru hattı toplayacaktır (Bölüm~\ref{chap:boruhatti}).
\end{ornek}

\subsubsection{Aşama Sayısı Bir Tasarım Kararıdır}

Veri yolunu beş aşamaya böldük; ancak beş sayısı zorunlu değil, bir tasarım kararıdır. Sınırlar doğal birim sınırlarına (bellek erişimi, yazmaç öbeği, AMB) denk geldiği ve RISC işlemcilerin klasik bölünmesi olduğu için bu sayı seçilmiştir. Kararın ölçütü ise bellidir. Çok vuruşluk veri yolunun varlık nedeni, çevrimi kısaltarak saat vuruş sıklığını yükseltmektir. Beş aşamalı bölme bunu 625~ps ve 1{,}60~GHz'den 225~ps ve 4{,}44~GHz'e taşıdı. İlkece aşama sayısı arttıkça aşamalar incelir ve saat daha sık vurdurulabilir. Aşama sayısı azaldıkça tersi olur. Örnek~6.5'in sayılarıyla iki yönün de bedelini görebiliriz.

\emph{Aşama sayısını azaltmak} çevrimi uzatır, çünkü aynı aşamaya sıkışan işler artık ardışık yürür. Çöz ile Yürüt tek aşamada birleştirilse (dört aşamalı tasarım) bu aşamanın gecikmesi $150 + 150 = 300$~ps olur ve çevrim 300~ps'e çıkar. Yükle buyruğu $4 \times 300 = 1200$~ps sürer, beş aşamalı tasarımın 1125~ps'inden yavaştır. Buyruk karışımı üzerinden ortalama hesabı da benzer sonucu verir. Kazanç tarafında ise daha az ara yazmaç, daha az durum ve daha yalın bir denetim birimi vardır. Bu yolun ucunda, bütün işlerin tek aşamaya sıkıştığı tasarım zaten tanıdıktır: 625~ps çevrimli tek vuruşluk işlemci.

\emph{Aşama sayısını artırmak} ise ancak \emph{en yavaş} aşama bölünebiliyorsa işe yarar. Örneğin Yürüt aşamasını AMB'yi parçalayarak ikiye bölmek çevrimi değiştirmez. Taban hâlâ belleğe erişen aşamalardır (225~ps). Bu durumda aşama eklemek yalnızca çevrim sayısını artırır. Yükle $6 \times 225 = 1350$~ps'e \emph{yavaşlar}. Asıl hedef olan bellek aşamasını bölmek ise kolay değildir, çünkü 200~ps'lik erişim süresi bellek dizisinin iç yapısından gelir ve bir ara yazmaç sınırıyla ortadan ikiye kesilemez. Üstelik her yeni sınırın bir bedeli vardır. Eklenen ara yazmacın kurulum ve yayılım süreleri (bu örnekte ihmal ettik) her çevrime biner ve çevrim süresi gerçekte
\[
t_{\text{çevrim}} = t_{\text{yazmaç}} + t_{\text{aşama}}
\]
biçimini alır. Aşamalar inceldikçe $t_{\text{yazmaç}}$'ın payı büyür ve saat vuruş sıklığı, aşama gecikmesi sıfıra inse bile $1/t_{\text{yazmaç}}$ tavanını aşamaz. Örneğin ara yazmaç vergisi 25~ps alınsaydı beş aşamalı tasarımın çevrimi $225+25 = 250$~ps'e, saati 4{,}0~GHz'e gerilerdi. Bellek engeli bir yana, bütün aşamalar bir biçimde yarıya inceltilebilseydi bile çevrim $112{,}5 + 25 \approx 138$~ps olurdu. Saat, beklenen 8{,}0~GHz'e değil 7{,}3~GHz'e çıkar ve her yeni bölmede getiri biraz daha azalır. Aşama sayısını artırmak saati daha sık vurdurmanın yoludur, ama bu yolun sınırını ara yazmaç gecikmeleri çizer. Denetim birimi de büyür. Durum sayısı ve geçişler artar. Çok vuruşlukta aşamalar örtüşmediğinden bu bedellere karşılık alınan tek şey dengesizliğin giderilmesidir. Bu yüzden aşama sayısını büyütmenin getirisi burada sınırlıdır. Aynı soru boru hattında yeniden ve çok daha keskin biçimde karşımıza çıkacak. Orada aşamalar örtüştüğü için derinlik doğrudan işlem hacmini artırır, tasarımcılar bu uğurda belleği bile çok çevrimli erişime bölerler ve ara yazmaç vergisi ile dengesizlik, derinliğin sınırını çizen ana etkenler olur (Bölüm~\ref{chap:boruhatti}).

\subsection{Çok Vuruşluk Veri Yolu Denetim Birimleri}
\label{subsec:cok-vurusluk-denetim}

Veri yolu kurulduğuna göre artık onu yönetecek sinyalleri sayabiliriz. Listenin çoğu tek vuruşluk tasarımdan tanıdıktır. YazmacaYaz, BelleğeYaz, AMBİşlem, AnlıkBiçim ve GeriYazmaKaynağı işlevlerini olduğu gibi korur. Yeni ve genişleyen sinyallerin tümü ise önceki kesimde yapılan değişikliklerin doğrudan sonucudur. Kaynak paylaşımı üç sinyal getirir: tek belleğin adresini seçen AdresKaynağı, buyruğu Getir çevriminde yakalayıp buyruk boyunca tutan BuyruğaYaz ve üç ayrı hesabı sırayla üstlenen AMB'nin girişlerini seçen AMBKaynakA ile AMBKaynakB (EskiPS gibi yeni kaynaklar eklendiğinden ikisi de iki bite genişlemiştir). Program sayacının güncellenmesi de artık tek bir noktada değil, buyruğa göre farklı çevrimlerde gerçekleşir. Bu zamanlamayı PSKaynak, PSYaz ve PSKoşul üçlüsü yönetir ve tek vuruşluktaki Dallan sinyalinin yerini alır. Tam liste Tablo~\ref{tab:cok-vurusluk-denetim}'da verilmektedir.

\begin{table}[H]
\centering
\caption{Çok vuruşluk veri yolu denetimleri}
\label{tab:cok-vurusluk-denetim}
\small
\begin{tabular}{lp{10cm}}
\toprule
\textbf{Denetim} & \textbf{İşlev} \\
\midrule
YazmacaYaz & Sonuç yazmacına gelen değerin yazılıp yazılmayacağını bildirir. \\
AMBKaynakA & AMB'nin A girişini seçer (2 bit): 00 = A yazmacı, 01 = PS, 10 = EskiPS (Getir'de saklanan buyruk adresi). \\
AMBKaynakB & B yolu için B yazmacı, 4 sabiti ve anlık değer arasından seçim yapar (2 bit genişliğinde). \\
AMBİşlem & AMB'nin yapacağı işlemi seçer (00 = Topla, 01 = Çıkar, 10 = VEYA). \\
AnlıkBiçim & Anlık değer üretecinin buyruk biçimi seçimi (I, S, B, J). \\
PSKaynak & Program sayacına gidecek değeri seçer (0 = PS+4, 1 = dal/jal hedefi). \\
BuyruğaYaz & Buyruk yazmacına bellekten gelen değerin yazılıp yazılmayacağını denetler. \\
GeriYazmaKaynağı & Sonuç yazmacına yazılacak olan değeri seçer. \\
BelleğeYaz & Belleğe değer yazılmasını denetler. \\
AdresKaynağı & Belleğe gelen adresin PS mi yoksa veri adresi mi olduğunu seçer. \\
PSKoşul & Koşullu dallanmayı denetler. PS yazma izni $=$ PSYaz $\lor$ (PSKoşul $\land$ (Sıfır $\oplus\, b_1$)); $b_1 = $ buyruk[12] beq'te 0, bne'de 1 olur. \\
PSYaz & Program sayacına gelen değerin yazılıp yazılmayacağını denetler. \\
\bottomrule
\end{tabular}
\end{table}

Tabloda bir okuma izni sinyalinin bulunmadığına dikkat edin. Okuma yan etkisiz olduğundan bellek de yazmaç öbeği de her çevrim serbestçe okunabilir, yalnızca yazma işlemleri izne bağlanır. Tek vuruşluktan asıl ayrılık ise sinyallerin \emph{zamana} bağlanmasıdır. Aynı buyruk için aynı sinyal çevrimden çevrime farklı değerler alır. Örneğin bir yükle buyruğunda AdresKaynağı, Getir çevriminde 0 (buyruk adresi) iken Bellek çevriminde 1 (veri adresi) olmalı. BuyruğaYaz yalnızca Getir'de 1 yapılmalıdır. Bu nedenle liste tek başına yeterli değildir. Her aşamada hangi sinyalin hangi değeri alacağını Bölüm~\ref{subsec:cok-vurusluk-asamalar}'te buyruk türlerini tek tek ele alarak çıkaracağız. Bu değerleri doğru sırayla üretecek denetim birimini ise Bölüm~\ref{subsec:sinirli-durum-makinesi}'te sonlu durum makinesi olarak tasarlayacağız.

\subsection{Çok Vuruşluk Veri Yolu Aşamaları}
\label{subsec:cok-vurusluk-asamalar}

\subsubsection{Getir Aşaması}

Buyruğun bellekten program sayacındaki değer yardımıyla getirildiği aşamadır. Program sayacındaki değerle belleğe gidilir ve buyruk okunur. PS'nin geçebilmesi için AdresKaynağı = 0 olmalıdır. Bellekten okunan buyruğun buyruk yazmacına yazılabilmesi için BuyruğaYaz denetimi etkinleştirilmelidir. Aynı çevrimde program sayacındaki buyruk adresi EskiPS yazmacına da saklanır; EskiPS'in yazma izni BuyruğaYaz ile ortaktır, böylece buyruk ile buyruğun adresi birlikte yakalanır ve sonraki çevrimlerde dallanma hedefi bu adrese göre hesaplanabilir. Ayrıca bir sonraki buyruğun yerinin hesaplanması için program sayacına 4 eklenir. AMB'de toplama işlemi yapılır. AMBKaynakA = 01, AMBKaynakB = 01 seçilerek PS ve 4 değeri AMB'ye verilir. AMB'nin çıkışındaki bu PS+4 değeri, PSKaynak = 0 seçiliyken doğrudan (AMB yazmacına uğramadan) program sayacına yönlenir. PSYaz etkinleştirilince getir çevrimi biterken PS bir sonraki buyruğa ilerlemiş olur. Dallanma ve jal hedefleri ise sonraki çevrimlerde AMB yazmacından (PSKaynak = 1) alınır. Dal ile jal ayrımı veri kaynağında değil, yazma izninde (PSKoşul / PSYaz) yapılır.

\subsubsection{Çöz Aşaması}

Buyruk yazmacındaki buyruğun üzerindeki alanlarda bulunan verilerin görevlerine göre veri yoluna dağıtıldığı aşamadır. Tüm buyruklar için çöz aşaması aynı şekilde gerçekleştirilir. Tüm buyrukların aynı yolu takip ediyor olması \emph{yalınlık düzenden gelir} tasarım ilkesini de destekler.

Çöz aşamasında A ve B yazmaçlarına yazmaç öbeğinden okunan değerler yazılır. İşaretle genişletilmiş anlık değer de üretilir. Bu çevrimde AMB boşta olduğundan, dallanma ve jal buyruklarında gerekebilecek hedef adres de (EskiPS + anlık değer; AMBKaynakA = 10, AMBKaynakB = 10, AMBİşlem = 00) AMB'de hesaplanıp AMB yazmacına konur. Buyruk dallanmazsa bu değer kullanılmadan atılır. Burada EskiPS, Getir çevriminde saklanan buyruk adresidir. Hedef bu adrese göre hesaplandığından RISC-V tanımıyla birebir aynı (buyruk adresi + anlık) olur. Program sayacı bu çevrime kadar zaten PS+4'e ilerlediği için, hedefin PS yerine EskiPS üzerinden hesaplanması dört baytlık kaymayı önler.

\subsubsection{Yürüt Aşaması}

Getir ve çöz aşamalarında tüm buyrukların yaptığı işlemler aynı iken yürüt aşamasında yapılan işlemler buyruk türüne göre farklılık gösterir.

\begin{itemize}
    \item \textbf{Dallanma (beq/bne):} AMB'de A ile B yazmaçlarındaki değerler karşılaştırılır (AMBKaynakA = 00, AMBKaynakB = 00, AMBİşlem = 01). Çöz aşamasında hesaplanıp AMB yazmacında bekleyen dallanma hedefi program sayacına yazılır ya da yazılmaz. PS yazma izni $=$ PSYaz $\lor$ (PSKoşul $\land$ (Sıfır $\oplus\, b_1$)) biçiminde üretilir (PSKoşul etkin, PSKaynak = 1). Buradaki $b_1 = $ buyruk[12], funct3'ün ilgili bitidir. Bu bit beq'te 0 olduğundan koşul Sıfır$=1$'de, bne'de 1 olduğundan koşul Sıfır$=0$'da sağlanır. Böylece aynı donanım her iki dallanmayı da yürütür.
    \item \textbf{Yükle:} Bellekten okunacak değerin adresi AMB'de hesaplanır (AMBİşlem = 00). AMBKaynakA = 00, AMBKaynakB = 10 seçilir.
    \item \textbf{Sakla:} Adres hesabı yükle buyruğuyla aynıdır.
    \item \textbf{Aritmetik--mantık (R türü):} A ve B yazmaçlarındaki değerler AMB'ye gönderilir (AMBKaynakA = 00, AMBKaynakB = 00, AMBİşlem = 00).
    \item \textbf{Mantıksal VEYA Anlık (ori):} A yazmacındaki değer ve işaretle genişletilmiş anlık değer AMB'ye gönderilir (AMBKaynakA = 00, AMBKaynakB = 10, AMBİşlem = 10).
    \item \textbf{jal:} Çöz aşamasında hesaplanıp AMB yazmacında bekleyen atlama hedefi program sayacına yazılır (PSYaz etkin, PSKaynak = 1). Aynı çevrimde PS'de duran PS$+$4 değeri dönüş adresi olarak rd yazmacına yazılır (YazmacaYaz etkin, GeriYazmaKaynağı = 10).
\end{itemize}

\subsubsection{Bellek Aşaması}

Bellek aşaması sadece bellekle veri alış verişinde bulunulan buyrukların geçtiği bir aşamadır.
\begin{itemize}
    \item \textbf{Yükle:} AMB yazmacındaki adres değeri belleğe verilir (AdresKaynağı = 1). Bellekten okunan veri bellek yazmacına (Bellek Yzm.) yazılır.
    \item \textbf{Sakla:} AMB yazmacından gelen adresle belleğe gidilir (AdresKaynağı = 1), B yazmacındaki değer belleğe yazılır (BelleğeYaz etkin).
\end{itemize}

\subsubsection{Sonuç Yaz Aşaması}

Sonuç yazma aşamasını gerçekleştiren buyruklar yükle buyruğu, aritmetik--mantık buyrukları ve ori buyruğudur; jal dönüş adresini yürüt çevriminde yazdığından bu aşamaya gelmez.
\begin{itemize}
    \item \textbf{Yükle:} GeriYazmaKaynağı = 01, YazmacaYaz etkin (rd yazmacına yazılır).
    \item \textbf{Aritmetik--mantık (R türü):} GeriYazmaKaynağı = 00, YazmacaYaz etkin (rd yazmacına yazılır).
    \item \textbf{Mantıksal VEYA Anlık (ori):} GeriYazmaKaynağı = 00, YazmacaYaz etkin (rd yazmacına yazılır).
\end{itemize}

\subsection{Çok Vuruşluk Veri Yolu İçin Denetim Birimi: Sonlu Durum Makinesi}
\label{subsec:sinirli-durum-makinesi}

Tek vuruşluk işlemcide denetim birimi yalnızca bir birleşik devreydi. Buyruk belliyken üretilecek sinyaller de belliydi ve buyruk boyunca hiç değişmiyordu. Çok vuruşlukta bu yaklaşım tek başına yetmez, çünkü denetim artık yalnızca \emph{hangi buyruğun} yürütüldüğüne değil, o buyruğun \emph{hangi çevrimde} olduğuna da bağlıdır. Aynı buyruk için aynı sinyal çevrimden çevrime farklı değerler alır. Örneğin AMBKaynakA, Getir'de 01 (PS), Çöz'de 10 (EskiPS), Yürüt'te ise 00 (A yazmacı) olmalıdır. Üstelik buyruklar farklı sayıda çevrim sürer. Yükle beş çevrimde biterken dallanma üçüncü çevrimin sonunda yeni buyruğa geçer. Denetim biriminin doğru sinyalleri doğru anda üretebilmesi için ``şu anda hangi buyruğun hangi aşamasındayız'' bilgisini bir çevrimden diğerine \emph{hatırlaması} gerekir. Salt birleşik mantığın belleği yoktur. Belleği olan denetim düzeneği ise \textbf{sonlu durum makinesi}\index{sonlu durum makinesi}\index{finite state machine|see{sonlu durum makinesi}} (kısaca SDM) olarak bilinir. Bulunulan aşama bir \emph{durum yazmacında}\index{durum yazmacı}\index{state register|see{durum yazmacı}} tutulur, birleşik mantık bu duruma ve buyruğa bakarak hem o çevrimin denetim sinyallerini hem de bir sonraki durumu üretir (Şekil~\ref{fig:durum-makinesi-genel}).

\begin{figure}[htbp]
\centering
\begin{tikzpicture}[>=Stealth, node distance=2cm]
  % Durum yazmacı
  \node[yazmacblok, minimum width=2.5cm, minimum height=1.2cm] (durum) {Durum\\Yazmacı};
  % Birleşik mantık
  \node[blok, right=3cm of durum, minimum width=3cm, minimum height=2cm] (mantik) {Birleşik\\Mantık};
  % Girişler
  \draw[<-, thick] (mantik.north) -- ++(0,0.8) node[above, font=\footnotesize, align=center]{opcode\\funct3, funct7};
  % Çıkışlar
  \draw[->, thick] (mantik.east) -- ++(2,0) node[right, font=\footnotesize, align=center]{Denetim\\İşaretleri\\(Çıkışlar)};
  % Sonraki durum
  \draw[->, thick] (mantik.south) -- ++(0,-1) -| (durum.south);
  \node[font=\footnotesize, above] at ($(mantik.south)!0.5!(durum.south)+(0,-1)$) {Sonraki Durum};
  % Durum -> Mantık
  \draw[->, thick] (durum.east) -- node[above, font=\footnotesize]{Şimdiki Durum} (mantik.west);
  % Saat
  \draw[denetimstil] (durum.west) -- ++(-1.5,0) node[left, font=\footnotesize, text=sinyalkirmizi]{Saat};
\end{tikzpicture}
\caption{Durum makinesi ile tasarlanacak denetim birimi tasarımı}
\label{fig:durum-makinesi-genel}
\end{figure}

Şekil~\ref{fig:durum-makinesi-genel}'deki yapıda her durum, veri yolunun geçirdiği bir saat çevrimine karşılık gelir. Makine bir durumda kaldığı sürece birleşik mantık o çevrimin denetim sinyallerini üretir, saat kenarı geldiğinde ise durum yazmacı sonraki duruma geçer. Çizeceğimiz durum çizgelerinde şu gösterimi kullanacağız. Her daire bir durumdur ve yanındaki yazı hangi aşamayı yürüttüğünü söyler. Bir durumda üretilen denetim değerleri, o durumdan çıkan okun üzerinde listelenir. Listelenmeyen bütün sinyaller etkisizdir (yazma izinleri 0, çoklayıcı seçimleri önemsizdir). Çöz durumundan ayrılan okların üzerindeki adlar ise üretilen sinyal değil \emph{geçiş koşuludur}. İşlem kodu hangi buyruğu gösteriyorsa makine o buyruğun yoluna sapar. Buyruğun son durumu tamamlandığında makine, yeni buyruğu getirmek üzere başlangıç durumuna döner.

Önce en kısa makineyi, dallanma buyruklarınınkini izleyelim (Şekil~\ref{fig:durum-beq}). S0, Getir çevrimidir. BuyruğaYaz = 1 ile buyruk (ve onunla birlikte EskiPS) yakalanır, PS $\leftarrow$ PS$+$4 ile program sayacı ilerletilir. S1, Çöz çevrimidir. Yazmaçlar okunur, AMB dal/jal olasılığına karşı hedef adresi hesaplayıp AMB yazmacına koyar ve işlem kodu çözülür. Makine beq/bne görürse S8'e geçer. S8'de AMB bu kez A ile B yazmaçlarını karşılaştırır (AMBİşlem = 01) ve PSKoşul = 1, PSKaynak = 1 ile bekleyen hedef, Sıfır $\oplus\, b_1$ koşulu sağlanırsa PS'ye yazılır. Üçüncü çevrimin sonunda makine S0'a döner. S0 ile S1'in içeriğinin buyruktan bağımsız olduğuna dikkat edin. Getir ve Çöz bütün buyruklar için özdeştir, bu yüzden bu iki durum bütün makinelerin ortak başlangıcıdır ve yol ayrımı ancak işlem kodunun çözüldüğü S1'in çıkışında yapılabilir.

\begin{figure}[htbp]
\centering
\begin{tikzpicture}[>=Stealth, node distance=3cm]
  \node[durum] (s0) {S0};
  \node[font=\scriptsize, below=0.1cm of s0] {Getir};
  \node[durum, right=3cm of s0] (s1) {S1};
  \node[font=\scriptsize, below=0.1cm of s1] {Çöz};
  \node[durum, right=3cm of s1] (s8) {S8};
  \node[font=\scriptsize, below=0.1cm of s8, align=center] {beq/bne\\Yürüt};
  % Geçişler
  \draw[durumok] (s0) -- node[above, font=\scriptsize, align=center]{BuyruğaYaz=1\\PS$\leftarrow$PS+4} (s1);
  \draw[durumok] (s1) -- node[above, font=\scriptsize]{beq/bne} (s8);
  % S8'den S0'a dönüş
  \draw[durumok] (s8) to[bend left=40] node[below, font=\scriptsize, align=center]{AMBİşlem=01\\PSKoşul=1, PSKaynak=1\\(PS koşula bağlı yazılır)} (s0);
\end{tikzpicture}
\caption{Dallanma buyruklarının durum makinesi}
\label{fig:durum-beq}
\end{figure}

Yükle ve sakla buyruklarının makinesi bir tasarım inceliği taşır (Şekil~\ref{fig:durum-yukle-sakla}). İkisi de aynı adres hesabını yaptığından S2 (Adres Hesapla) durumunu \emph{paylaşırlar}. Yol ancak bellek çevriminde ayrılır. \emph{lw} bellekten okur (S3) ve okuduğu değeri rd yazmacına yazmak için bir çevrim daha harcar (S4); \emph{sw} ise B yazmacındaki değeri belleğe yazar (S5) ve orada biter, çünkü yazmaç öbeğine yazacağı bir sonucu yoktur. Durum paylaşımı yalnızca çizimi kısaltmaz. Birleşik makinede durum sayısını, dolayısıyla denetim donanımını da küçültür.

\begin{figure}[htbp]
\centering
\begin{tikzpicture}[>=Stealth]
  \node[durum] (s0) at (0,0) {S0};
  \node[font=\scriptsize, below=0.1cm of s0] {Getir};
  \node[durum] (s1) at (2.4,0) {S1};
  \node[font=\scriptsize, below=0.1cm of s1] {Çöz};
  \node[durum] (s2) at (4.8,0) {S2};
  \node[font=\scriptsize, below=0.1cm of s2, align=center] {Adres\\Hesapla};
  % lw yolu (üst)
  \node[durum] (s3) at (7.0,1.7) {S3};
  \node[font=\scriptsize, below=0.1cm of s3, align=center] {Belleği Oku\\(AdresKaynağı=1)};
  \node[durum] (s4) at (9.2,1.7) {S4};
  \node[font=\scriptsize, below=0.1cm of s4, align=center] {Yaz\\Yazmaç};
  % sw yolu (alt)
  \node[durum] (s5) at (7.0,-1.3) {S5};
  \node[font=\scriptsize, above=0.1cm of s5] {Belleğe Yaz};
  % Geçişler
  \draw[durumok] (s0) -- (s1);
  \draw[durumok] (s1) -- node[above, font=\scriptsize]{lw/sw} (s2);
  \draw[durumok] (s2) -- node[above left, font=\scriptsize]{lw} (s3);
  \draw[durumok] (s3) -- (s4);
  \draw[durumok] (s2) -- node[below left, font=\scriptsize]{sw} (s5);
  % Geri dönüşler -- dik açılı dar koridorlar (üst y=2.9, alt y=-2.6)
  \draw[durumok] (s4.north) -- (9.2,2.9) -- (0,2.9) -- (s0.north);
  \node[font=\scriptsize, above] at (4.6,2.9) {YazmacaYaz=1, GeriYazmaKaynağı=01};
  \draw[durumok] (s5.south) -- (7.0,-2.6) -- (-0.95,-2.6) -- (-0.95,0) -- (s0.west);
  \node[font=\scriptsize, below] at (3.0,-2.6) {AdresKaynağı=1, BelleğeYaz=1};
\end{tikzpicture}
\caption{Yükle -- Sakla buyruğunun durum makinesi}
\label{fig:durum-yukle-sakla}
\end{figure}

R türü aritmetik--mantık buyrukları belleğe uğramaz. Dört durumda tamamlanırlar: getir, çöz, AMB yürüt ve sonuç yaz (Şekil~\ref{fig:durum-aritmetik}). S6'daki AMB işlemi tek bir koda bağlanamaz, çünkü add, sub ve ori dışındaki mantık buyrukları aynı yoldan akar. İşlemi buyruğun funct3/funct7 alanları belirler ve tek vuruşluk tasarımdaki AMB denetim mantığı burada da aynen iş başındadır. S7'de sonuç, GeriYazmaKaynağı = 00 seçilerek rd yazmacına yazılır.

\begin{figure}[htbp]
\centering
\begin{tikzpicture}[>=Stealth, node distance=3cm]
  \node[durum] (s0) {S0};
  \node[font=\scriptsize, below=0.1cm of s0] {Getir};
  \node[durum, right=3cm of s0] (s1) {S1};
  \node[font=\scriptsize, below=0.1cm of s1] {Çöz};
  \node[durum, right=3cm of s1] (s6) {S6};
  \node[font=\scriptsize, below=0.1cm of s6, align=center] {AMB\\Yürüt};
  \node[durum, right=3cm of s6] (s7) {S7};
  \node[font=\scriptsize, below=0.1cm of s7, align=center] {Sonuç\\Yaz};
  % Geçişler
  \draw[durumok] (s0) -- (s1);
  \draw[durumok] (s1) -- node[above, font=\scriptsize]{R türü} (s6);
  \draw[durumok] (s6) -- node[above, font=\scriptsize]{AMBİşlem: funct'a göre} (s7);
  \draw[durumok] (s7) to[bend left=40] node[below, font=\scriptsize]{YazmacaYaz=1, GeriYazmaKaynağı=00} (s0);
\end{tikzpicture}
\caption{Aritmetik -- Mantık buyruğunun durum makinesi}
\label{fig:durum-aritmetik}
\end{figure}

Mantıksal VEYA anlık (\emph{ori}) buyruğunun makinesi aynı dört aşamayı izler. Tek fark S9 durumunda AMB'nin B girişine yazmaç yerine anlık değerin verilmesidir (AMBKaynakB = 10) ve işlem sabittir (AMBİşlem = 10, VEYA) (Şekil~\ref{fig:durum-ori}).

\begin{figure}[htbp]
\centering
\begin{tikzpicture}[>=Stealth, node distance=3cm]
  \node[durum] (s0) {S0};
  \node[font=\scriptsize, below=0.1cm of s0] {Getir};
  \node[durum, right=3cm of s0] (s1) {S1};
  \node[font=\scriptsize, below=0.1cm of s1] {Çöz};
  \node[durum, right=3cm of s1] (s9) {S9};
  \node[font=\scriptsize, below=0.1cm of s9, align=center] {AMB+\\anlık};
  \node[durum, right=3cm of s9] (s10) {S10};
  \node[font=\scriptsize, below=0.1cm of s10, align=center] {Sonuç\\Yaz};
  % Geçişler
  \draw[durumok] (s0) -- (s1);
  \draw[durumok] (s1) -- node[above, font=\scriptsize]{ori} (s9);
  \draw[durumok] (s9) -- node[above, font=\scriptsize]{AMBİşlem=10} (s10);
  \draw[durumok] (s10) to[bend left=40] node[below, font=\scriptsize]{YazmacaYaz=1, GeriYazmaKaynağı=00} (s0);
\end{tikzpicture}
\caption{Mantıksal VEYA Anlık (ori) buyruğunun durum makinesi}
\label{fig:durum-ori}
\end{figure}

\emph{jal} buyruğu en kısa makinelerden birine sahiptir, çünkü işini Çöz'de hazırlanan malzemeyle bitirir. Atlama hedefi zaten S1'de hesaplanıp AMB yazmacına konmuştur. S11'de bu değer PS'ye yazılırken (PSKaynak = 1, PSYaz = 1), PS'de duran PS$+$4 değeri de dönüş adresi olarak rd yazmacına yazılır (YazmacaYaz = 1, GeriYazmaKaynağı = 10); tek durumda iki yazma işlemi birden yapılır (Şekil~\ref{fig:durum-jal}).

\begin{figure}[htbp]
\centering
\begin{tikzpicture}[>=Stealth, node distance=3cm]
  \node[durum] (s0) {S0};
  \node[font=\scriptsize, below=0.1cm of s0] {Getir};
  \node[durum, right=3cm of s0] (s1) {S1};
  \node[font=\scriptsize, below=0.1cm of s1] {Çöz};
  \node[durum, right=3cm of s1] (s11) {S11};
  \node[font=\scriptsize, below=0.1cm of s11, align=center] {PS güncelle\\rd$\leftarrow$PS+4};
  % Geçişler
  \draw[durumok] (s0) -- (s1);
  \draw[durumok] (s1) -- node[above, font=\scriptsize]{jal} (s11);
  \draw[durumok] (s11) to[bend left=40] node[below, font=\scriptsize, align=center]{PSKaynak=1, PSYaz=1\\YazmacaYaz=1, GeriYazmaKaynağı=10} (s0);
\end{tikzpicture}
\caption{jal buyruğunun durum makinesi}
\label{fig:durum-jal}
\end{figure}

Beş makine tek çatı altında toplandığında Şekil~\ref{fig:durum-birlesik}'deki on iki durumlu birleşik makine ortaya çıkar. Ortak S0--S1 omurgası bir kez çizilir. S1'in çıkışı, işlem koduna göre beş yola ayrılan tek karar noktasıdır ve her yolun son durumu yeni buyruk için S0'a döner. Buyruk kümesine yeni bir buyruk eklemenin bedeli de bu çizimden okunabilir. Makineye yalnızca o buyruğun Yürüt ve sonrası için gereken durumlar eklenir, S0--S1 ve varsa paylaşılabilen ara durumlar (S2 gibi) olduğu gibi kalır.

\begin{figure}[htbp]
\centering
\resizebox{\textwidth}{!}{%
\begin{tikzpicture}[>=Stealth, node distance=2cm]
  % Ortak durumlar
  \node[durum] (s0) at (0,0) {S0};
  \node[font=\scriptsize, below=0.05cm of s0] {Getir};
  \node[durum] (s1) at (3,0) {S1};
  \node[font=\scriptsize, below=0.05cm of s1] {Çöz};
  % Dallanma yolu (üst)
  \node[durum] (s8) at (7,4) {S8};
  \node[font=\scriptsize, below=0.05cm of s8] {beq/bne};
  % lw/sw yolu
  \node[durum] (s2) at (7,1.5) {S2};
  \node[font=\scriptsize, below=0.05cm of s2] {Adres H.};
  \node[durum] (s3) at (10,2.5) {S3};
  \node[font=\scriptsize, below=0.05cm of s3] {Belleği Oku};
  \node[durum] (s4) at (13,2.5) {S4};
  \node[font=\scriptsize, below=0.05cm of s4] {Yzm.Yaz};
  \node[durum] (s5) at (10,0.5) {S5};
  \node[font=\scriptsize, below=0.05cm of s5] {Belleğe Yaz};
  % R türü yolu
  \node[durum] (s6) at (7,-1.5) {S6};
  \node[font=\scriptsize, below=0.05cm of s6] {AMB Y.};
  \node[durum] (s7) at (10,-1.5) {S7};
  \node[font=\scriptsize, below=0.05cm of s7] {Sonucu Yaz};
  % ori yolu
  \node[durum] (s9) at (7,-3.5) {S9};
  \node[font=\scriptsize, below=0.05cm of s9] {AMB+anl.};
  \node[durum] (s10) at (10,-3.5) {S10};
  \node[font=\scriptsize, below=0.05cm of s10] {Sonucu Yaz};
  % jal yolu
  \node[durum] (s11) at (7,-5.5) {S11};
  \node[font=\scriptsize, below=0.05cm of s11] {jal Yürüt};
  % Geçişler
  \draw[->, >=Stealth, thick] (s0) -- (s1);
  \draw[->, >=Stealth, thick] (s1) -- node[above left, font=\scriptsize]{beq/bne} (s8);
  \draw[->, >=Stealth, thick] (s1) -- node[above, font=\scriptsize]{lw/sw} (s2);
  \draw[->, >=Stealth, thick] (s1) -- node[below left, font=\scriptsize]{R türü} (s6);
  \draw[->, >=Stealth, thick] (s1) -- node[below, pos=0.6, font=\scriptsize]{ori} (s9);
  \draw[->, >=Stealth, thick] (s1) -- node[below left, pos=0.7, font=\scriptsize]{jal} (s11);
  \draw[->, >=Stealth, thick] (s2) -- node[above left, font=\scriptsize]{lw} (s3);
  \draw[->, >=Stealth, thick] (s2) -- node[below left, font=\scriptsize]{sw} (s5);
  \draw[->, >=Stealth, thick] (s3) -- (s4);
  \draw[->, >=Stealth, thick] (s6) -- (s7);
  \draw[->, >=Stealth, thick] (s9) -- (s10);
  % === Geri dönüşler: dik açılı, iç içe (laminar) koridorlar, çakışmasız ===
  % --- Üst koridorlar: S8 (y=5.0), S4 (y=5.8), S5 (y=6.6) ---
  % S8 -> S0 : en iç koridor, sol iniş x=-0.25 -> s0.north
  \draw[->, >=Stealth, thick, dashed, sinyalkirmizi]
    (s8.north) -- (7,5.0) -- (-0.25,5.0) -- ([xshift=-0.25cm]s0.north);
  % S4 -> S0 : sağ çıkış x=14.5, sol iniş x=-1.0 -> s0.north west
  \draw[->, >=Stealth, thick, dashed, sinyalkirmizi]
    (s4.east) -- (14.5,2.5) -- (14.5,5.8) -- (-1.0,5.8) -- (-1.0,0.42) -- (s0.north west);
  % S5 -> S0 : sağ çıkış x=15.2, sol iniş x=-1.8 -> s0.west
  \draw[->, >=Stealth, thick, dashed, sinyalkirmizi]
    (s5.east) -- (15.2,0.5) -- (15.2,6.6) -- (-1.8,6.6) -- (-1.8,0.15) -- ([yshift=0.15cm]s0.west);
  % --- Alt koridorlar: S11 en yakın (y=-7.0), S10 (y=-7.8), S7 en dış (y=-9.0) ---
  % S11 -> S0 : sağ iniş x=14.5, sol çıkış x=-0.25 -> s0.south
  \draw[->, >=Stealth, thick, dashed, sinyalkirmizi]
    (s11.east) -- (14.5,-5.5) -- (14.5,-7.0) -- (-0.25,-7.0) -- ([xshift=-0.25cm]s0.south);
  % S10 -> S0 : sağ iniş x=15.2, sol çıkış x=-1.0 -> s0.south west
  \draw[->, >=Stealth, thick, dashed, sinyalkirmizi]
    (s10.east) -- (15.2,-3.5) -- (15.2,-7.8) -- (-1.0,-7.8) -- (-1.0,-0.42) -- (s0.south west);
  % S7 -> S0 : sağ iniş x=15.9, sol çıkış x=-1.8 -> s0.west
  \draw[->, >=Stealth, thick, dashed, sinyalkirmizi]
    (s7.east) -- (15.9,-1.5) -- (15.9,-9.0) -- (-1.8,-9.0) -- (-1.8,-0.15) -- ([yshift=-0.15cm]s0.west);
  % --- Koridor etiketleri: boş orta alanda, kendi koridor çizgisinin üstünde ---
  \node[font=\scriptsize, text=sinyalkirmizi, align=center] at (3.5,5.35) {AMBİşlem=01, PSKaynak=1, PSKoşul=1 (dal)};
  \node[font=\scriptsize, text=sinyalkirmizi, align=center] at (3.5,6.15) {YazmacaYaz=1, GeriYazmaKaynağı=01};
  \node[font=\scriptsize, text=sinyalkirmizi, align=center] at (3.5,6.95) {AdresKaynağı=1, BelleğeYaz=1};
  \node[font=\scriptsize, text=sinyalkirmizi, align=center] at (3.5,-6.55) {PSKaynak=1, PSYaz=1, YazmacaYaz=1,\\GeriYazmaKaynağı=10 (jal)};
  \node[font=\scriptsize, text=sinyalkirmizi, align=center] at (3.5,-7.5) {YazmacaYaz=1, GeriYazmaKaynağı=00 (ori)};
  \node[font=\scriptsize, text=sinyalkirmizi, align=center] at (3.5,-8.7) {YazmacaYaz=1, GeriYazmaKaynağı=00};
\end{tikzpicture}%
}
\caption{Tüm denetimlerin birleştiği durum makinesi}
\label{fig:durum-birlesik}
\end{figure}

Bu makineyi donanıma dökmek, Ek~\ref{ch:sayisal-tasarim}'deki bildik yöntemlerin işidir. On iki durum dört bitlik bir durum koduyla adlandırılır. Durum yazmacı dört flip-floptan kurulur. Birleşik mantık, şimdiki durum kodu ile işlem kodundan hem o çevrimin denetim sinyallerini hem de sonraki durum kodunu üretir. Tek vuruşluğun ana çözücüsünde doğruluk tablosundan Karno haritalarıyla kapılara indiğimiz sürecin aynısı, burada durum geçiş tablosuna uygulanır. Denetim birimi veri yoluna bağlandığında çok vuruşluk işlemci tamamlanmış olur:
\begin{itemize}
    \item Dallanma buyrukları (beq/bne) için getir, çöz, yürüt olarak \textbf{3 adım},
    \item Yükle buyruğu için getir, çöz, yürüt, bellek, sonuç yaz olarak \textbf{5 adım},
    \item Sakla buyruğu için getir, çöz, yürüt, bellek olarak \textbf{4 adım},
    \item Aritmetik--mantık buyrukları (R türü) için getir, çöz, yürüt, sonuç yaz olarak \textbf{4 adım},
    \item Mantıksal VEYA anlık (ori) buyruğu için getir, çöz, yürüt, sonuç yaz olarak \textbf{4 adım},
    \item jal buyruğu için getir, çöz, yürüt olarak \textbf{3 adım} (dönüş adresi yürüt çevriminde yazmaca yazılır)
\end{itemize}
gerçekleştirilir.

Böylece tasarım tamamlanmıştır. Çok vuruşluk veri yolunun tüm bileşenleri ve sonlu durum makinesi tabanlı denetim birimi, Şekil~\ref{fig:cok-vurusluk-tam}'da topluca gösterilmektedir.

\begin{figure}[htbp]
\centering
\resizebox{\textwidth}{!}{%
\begin{tikzpicture}[>=Stealth, font=\small,
    vy/.style={->, semithick}]
  % === Ana bloklar (soldan sağa tek ana hat) ===
  \node[blok, minimum width=1.2cm, minimum height=1cm] (ps) at (0,0) {PS};
  \node[blok, fill=green!20, minimum width=1.8cm, minimum height=2cm] (bel) at (2.6,0) {Bellek};
  \node[blok, minimum width=1.4cm, minimum height=1.6cm] (byrk) at (5.2,0) {Byrk.\\Yzm.};
  \node[blok, minimum width=2cm, minimum height=2.4cm] (yob) at (8.4,0) {Yazmaç\\Öbeği};
  \node[blok, minimum width=1.2cm, minimum height=1.8cm] (abyz) at (10.9,0) {A, B\\Yzm.};
  \node[blok, fill=blokkirmizi!30, minimum width=1.2cm, minimum height=1.8cm] (amb) at (13.1,0) {AMB};
  \node[blok, minimum width=1.2cm, minimum height=1.8cm] (ambyzm) at (15.1,0) {AMB\\Yzm.};
  % Bellek yazmacı (Bellek altında) ve Anlık Üreteci (alt koridor)
  \node[blok, minimum width=1.4cm, minimum height=0.9cm] (belyz) at (2.6,-2.4) {Bellek\\Yzm.};
  \node[blok, minimum width=1.6cm, minimum height=0.9cm] (anlik) at (10.4,-1.9) {Anlık\\Üreteci};
  % Denetim birimi (üstte) ve AMB Denetimi
  \node[denetim, minimum width=3.6cm, minimum height=1.4cm] (denetim) at (7,3.2) {Denetim Birimi\\(SDM)};
  \node[blok, minimum width=1.4cm, minimum height=0.8cm, fill=blokkirmizi!30] (ambden) at (13.1,2.6) {AMB\\Den.};
  % === Ana veri hattı ===
  \draw[vy] (ps.east) -- (bel.west);
  \draw[vy] (bel.east) -- (byrk.west);
  \draw[vy] (byrk.east) -- (yob.west);
  \draw[vy] (yob.east) -- (abyz.west);
  \draw[vy] (abyz.east) -- (amb.west);
  \draw[vy] (amb.east) -- (ambyzm.west);
  % Bellek Yzm.: yan kenardan çıkıp kendi düzeyinden Yazmaç Öbeği'ne
  \draw[vy] (bel.south) -- (belyz.north);
  \draw[vy] (belyz.east) -- (7.9,-2.4) -- ([xshift=-0.5cm]yob.south);
  % Anlık Üreteci: Byrk.Yzm.'den sol kenarına, çıkışı AMB'nin batı kenarına dik
  \draw[vy] (byrk.south) -- (5.2,-1.9) -- (anlik.west);
  \draw[vy] (anlik.east) -- (12.1,-1.9) -- (12.1,-0.5) -- ([yshift=-0.5cm]amb.west);
  % === Geri beslemeler (alt koridorlar, dik açılı, ayrı seviyeler) ===
  % Sonuç -> Yazmaç Öbeği (güneyden)
  \draw[vy] (ambyzm.south) -- (15.1,-3.2) -- (9.0,-3.2) -- ([xshift=0.6cm]yob.south);
  % Adres -> PS/Bellek (doğudan çıkar, en alttan döner)
  \draw[vy] (ambyzm.east) -- (16.3,0) -- (16.3,-4.0) -- (0,-4.0) -- (ps.south);
  % === Denetim: tek gövde çizgisi + her bloğa kısa iniş ===
  \coordinate (cbus) at (0,1.9);
  \draw[denetimstil, -] (0,1.9) -- (15.1,1.9);
  \draw[denetimstil, -] (denetim.south) -- (7,1.9);
  \foreach \b in {ps,bel,byrk,yob,abyz,ambyzm}
    \draw[denetimstil] (\b.north |- cbus) -- (\b.north);
  % AMB denetim zinciri (yesil): Denetim -> AMB Den. -> AMB
  \draw[->, sinyalyesil, dashed, semithick] (denetim.east) -- (12.0,3.2) -- (12.0,2.6) -- (ambden.west);
  \draw[->, sinyalyesil, dashed, semithick] (ambden.south) -- (amb.north);
  % opcode -> Denetim (lacivert, soldan tek bukumle bati ucuna girer)
  \draw[->, sinyallacivert, dashed, semithick] ([xshift=-0.45cm]byrk.north) -- (4.75,3.2) -- (denetim.west);
  \node[font=\scriptsize, text=sinyallacivert, anchor=east] at (4.6,2.35) {opcode};
\end{tikzpicture}%
}
\caption[Çok vuruşluk veri yolu ve denetim birimi (genel bakış)]{Çok vuruşluk veri yolu ve denetim birimi (genel bakış). Kırmızı kesikli teller denetim sinyallerini, lacivert kesikli tel denetime giren işlem kodunu, yeşil kesikli teller AMB denetim zincirini gösterir. Çoklayıcılar, 4 sabiti girişi, EskiPS yazmacı ve kimi bağlantılar yalınlık için gösterilmemiştir; tam veri yolu Şekil~\ref{fig:cok-vurusluk-veriyolu}'dedir.}
\label{fig:cok-vurusluk-tam}
\end{figure}

%% ================================================================
%% ÖRNEKLER 6.6-6.8
%% ================================================================

\begin{ornek}[6.6 -- Tek Vuruşluk ve Çok Vuruşluk İşlemci Karşılaştırması]
Aşağıdaki yüzdelerle buyrukların kullanıldığı bir program, 1 vuruşu 90 ns olan tek vuruşluk bir işlemci ile 1 vuruşu 20 ns olan çok vuruşluk bir işlemcide çalıştırılacaktır. Program işlemcilerin her birinde kaç nanosaniyede tamamlanır?

\begin{center}
\%12 yükle, \%30 sakla, \%43 aritmetik--mantık, \%7 ori, \%8 dallanma
\end{center}

\tcblower
\textbf{Çözüm:} Buyruk sayısı 100 alınırsa programda 12 tane yükle, 30 tane sakla, 43 tane aritmetik--mantık, 7 tane ori ve 8 tane dallanma buyruğu vardır.

\textbf{Tek vuruşluk işlemci için yürütme zamanı:}
\[
\text{Yürütme zamanı} = \text{buyruk sayısı} \times \text{çevrim zamanı} \times \text{BBÇ} = 100 \times 90 \times 1 = 9000 \text{ ns}
\]

\textbf{Çok vuruşluk işlemci için yürütme zamanı:}

\begin{center}
\begin{tabular}{lcccc}
\toprule
\textbf{Buyruk} & \textbf{Sayı} & \textbf{Çevrim (ns)} & \textbf{BBÇ} & \textbf{Süre (ns)} \\
\midrule
Yükle & 12 & 20 & 5 & 1200 \\
Sakla & 30 & 20 & 4 & 2400 \\
Aritmetik--Mantık & 43 & 20 & 4 & 3440 \\
Mantıksal VEYA Anlık (ori) & 7 & 20 & 4 & 560 \\
Dallanma & 8 & 20 & 3 & 480 \\
\midrule
\textbf{Toplam} & & & & \textbf{8080} \\
\bottomrule
\end{tabular}
\end{center}

Çok vuruşluk işlemcide bazı buyruklar tek vuruşluk tasarımdakinden daha uzun sürer. Fakat bazı buyrukların tüm aşamalardan geçmiyor olması çok vuruşluk işlemcinin yürütme zamanının daha kısa olmasını sağlar.
\end{ornek}

\begin{ornek}[6.7]
Örnek~6.5'te bulunan 225~ps'lik çevrim zamanını kullanarak çok vuruşluk işlemcinin ortalama BBÇ'sini ve tek vuruşluk işlemciye göre hızlanmasını hesaplayın. Buyruk dağılımı: \%25 yükle, \%10 sakla, \%45 R türü, \%12 ori, \%8 dallanma.

\textbf{Çözüm.}

Her buyruk türünün çevrim sayısı (Tablo~\ref{tab:buyruk-asamalar}'den): yükle = 5, sakla = 4, R türü = 4, ori = 4, dallanma = 3.

Ortalama BBÇ:
\[
\text{BBÇ}_{\text{ort}} = 0{,}25 \times 5 + 0{,}10 \times 4 + 0{,}45 \times 4 + 0{,}12 \times 4 + 0{,}08 \times 3 = 1{,}25 + 0{,}40 + 1{,}80 + 0{,}48 + 0{,}24 = 4{,}17
\]

Yürütme zamanları (100 buyruk için):
\begin{itemize}
    \item Tek vuruşluk: $100 \times 1 \times 625 = 62\,500$~ps
    \item Çok vuruşluk: $100 \times 4{,}17 \times 225 = 93\,825$~ps
\end{itemize}

\[
\text{Hızlanma} = \frac{62\,500}{93\,825} \approx 0{,}67
\]

Bu örnekte çok vuruşluk işlemci tek vuruşluktan \emph{daha yavaştır}! Bunun nedeni BBÇ'nin 1'den çok yüksek olmasıdır ($4{,}17$). Çok vuruşluk işlemcinin avantajı, bu aşamaların boru hattı ile örtüştürülmesiyle ortaya çıkar (Bölüm~\ref{chap:boruhatti}).
\end{ornek}

İki tasarım yaklaşımı arasındaki temel farklar Tablo~\ref{tab:tek-cok-karsilastirma}'da özetlenmektedir.

\begin{table}[htbp]
\centering
\caption{Tek vuruşluk ve çok vuruşluk işlemcinin karşılaştırılması}
\label{tab:tek-cok-karsilastirma}
\small
\setlength{\tabcolsep}{4pt}
\begin{tabular}{lcc}
\toprule
\textbf{Ölçüt} & \textbf{Tek Vuruşluk} & \textbf{Çok Vuruşluk} \\
\midrule
Buyruk başına çevrim (BBÇ) & 1 & 3--5 (buyruğa göre değişir) \\
Çevrim süresi (Örnek 6.4 ve 6.5) & 625 ps (en yavaş buyruk) & 225 ps (en yavaş aşama) \\
Ortalama yürütme süresi (lw) & $1 \times 625 = 625$ ps & $5 \times 225 = 1125$ ps \\
Donanım maliyeti & Tek AMB, basit denetim & Tek AMB + ara yazmaçlar + SDM \\
Bellek erişimi & 1 saat çevriminde 2 erişim & Aşamalar arasında 1 erişim/çevrim \\
En uygun olduğu yer & Eğitim, basit gömülü & Karmaşık ISA, mikrokodlu CISC \\
\bottomrule
\end{tabular}
\end{table}

\begin{ornek}[6.8 -- Kod parçası üzerinde tek ve çok vuruşluk karşılaştırması]
Aşağıdaki program, \texttt{a0} adresinden başlayan 4 elemanlı bir diziyi toplayıp sonucu \texttt{a1} adresine yazmaktadır (yalnızca bu bölümün buyruk alt kümesi kullanılmıştır):

\begin{lstlisting}[style=riscv]
       ori  t0, x0, 0       # toplam = 0
       ori  t1, x0, 4       # sayac = 4 (eleman sayisi)
       ori  t3, x0, 1       # sabit 1 (sayac adimi)
       ori  t4, x0, 4       # sabit 4 (adres adimi)
dongu: lw   t2, 0(a0)       # siradaki elemani yukle
       add  t0, t0, t2      # toplama ekle
       add  a0, a0, t4      # adresi 4 ilerlet
       sub  t1, t1, t3      # sayaci 1 azalt
       bne  t1, x0, dongu   # sayac 0 degilse basa don
       sw   t0, 0(a1)       # toplami bellege yaz
\end{lstlisting}

(a)~Program tek vuruşluk işlemcide kaç çevrimde biter; Örnek~6.4'teki 625~ps'lik çevrimle süresi nedir? (b)~Aynı program çok vuruşluk işlemcide kaç çevrim sürer; Örnek~6.5'teki 225~ps'lik çevrimle süresi nedir? (c)~Çok vuruşluk işlemcinin bu programda tek vuruşluğu geçebilmesi için çevrim zamanı en çok, saat sıklığı en az ne olmalıdır? Bu değer ulaşılabilir midir?

\tcblower
\textbf{Çözüm.}

\textbf{(a)} Önce yürütülen (devingen) buyruk sayısını sayalım: girişteki 4 \texttt{ori}, döngü gövdesindeki 5 buyruğun (lw, add, add, sub, bne) 4 yinelemesi ve son \texttt{sw}:
\[
4 + 5 \times 4 + 1 = 25 \text{ buyruk}
\]
Tek vuruşlukta her buyruk bir çevrimdir: \textbf{25 çevrim}, süre $25 \times 625 = 15\,625$~ps.

\textbf{(b)} Çok vuruşlukta her buyruk türünün çevrim sayısıyla (Tablo~\ref{tab:buyruk-asamalar}) ağırlıklandırırız:

\begin{center}
\small
\begin{tabular}{@{}lccc@{}}
\toprule
\textbf{Buyruk grubu} & \textbf{Adet} & \textbf{Çevrim} & \textbf{Toplam} \\
\midrule
ori (giriş) & 4 & 4 & 16 \\
lw (döngü) & 4 & 5 & 20 \\
add, sub (döngü) & 12 & 4 & 48 \\
bne (döngü) & 4 & 3 & 12 \\
sw (çıkış) & 1 & 4 & 4 \\
\midrule
\textbf{Toplam} & 25 & & \textbf{100} \\
\bottomrule
\end{tabular}
\end{center}

\textbf{100 çevrim}; süre $100 \times 225 = 22\,500$~ps. Bu programda BBÇ $= 100/25 = 4{,}0$'dır ve çok vuruşluk işlemci $22\,500 / 15\,625 = 1{,}44$ kat \emph{daha yavaştır}.

\textbf{(c)} Eşitlik koşulundan çevrim üst sınırı bulunur:
\[
100 \times t_{\text{çevrim}} < 25 \times 625 \text{ ps} \quad\Rightarrow\quad t_{\text{çevrim}} < 156{,}25 \text{ ps} \quad\Rightarrow\quad f > 6{,}4 \text{ GHz}
\]
Bu değer ulaşılabilir değildir. Tek başına bellek erişimi 200~ps sürer ve bir aşama en az bir bellek erişimi içerdiğinden çevrim 200~ps'nin altına inemez (aşama sayısı tartışmasında gördüğümüz bölünemezlik engeli). Sonuç, Örnek~6.7'deki dersin kod üzerinde doğrulanmasıdır. Çok vuruşluk tasarım donanımı sadeleştirir ama BBÇ 1'in çok üzerine çıktığından bu iş yükünde tek vuruşluğu geçemez. Aradaki farkı ancak aşamaları örtüştüren boru hattı kapatacaktır (Bölüm~\ref{chap:boruhatti}).
\end{ornek}

Buraya kadarki örneklerde çok vuruşluk tasarım hep geride kaldı. Bunun nedeni, kümemizdeki buyrukların maliyetçe birbirine yakın olmasıdır. Kümeye pahalı bir buyruk girdiğinde durum tersine döner.

\begin{ornek}[6.9 -- Buyruk kümesine bölme eklenirse]
Buyruk kümesine \texttt{div} (bölme) eklensin ve buyrukların \%2'si bölme olsun (kalan dağılım: \%25 yükle, \%10 sakla, \%43 R türü, \%12 ori, \%8 dallanma). Tek vuruşluk tasarımda bölmenin tek çevrimde bitmesi için 1800~ps gecikmeli birleşik bir dizi bölücü kullanılacak. Çok vuruşluk tasarımda ise bölme, Bölüm~\ref{chap:aritmetik}'teki yinelemeli bölücüyle Yürüt aşamasında 32 çevrim sürecektir. İki tasarımın başarımını karşılaştırın.

\tcblower
\textbf{Çözüm.} Tek vuruşlukta çevrim süresi en yavaş buyruğa göre belirlenir ve en yavaş buyruk artık bölmedir:
\[
t_{\text{çevrim}} = 200 + 50 + 25 + 1800 + 25 + 50 = 2150 \text{ ps}
\]
Bu bedeli yalnızca bölme değil, \emph{bütün buyruklar} öder. 100 buyruk $100 \times 2150 = 215\,000$~ps sürer.

Çok vuruşlukta çevrim 225~ps'te kalır, çünkü bölücü kendi aşamasında yinelenir ve hiçbir aşamayı uzatmaz. Bölme buyruğu $1 + 1 + 32 + 1 = 35$ çevrim sürer (Getir, Çöz, 32 Yürüt yinelemesi, Sonucu Yaz); bedeli yalnızca onu kullanan öder:
\[
\text{BBÇ}_{\text{ort}} = 0{,}25 \times 5 + 0{,}10 \times 4 + 0{,}43 \times 4 + 0{,}12 \times 4 + 0{,}08 \times 3 + 0{,}02 \times 35 = 4{,}79
\]
\[
100 \times 4{,}79 \times 225 = 107\,775 \text{ ps} \qquad \text{Hızlanma} = \frac{215\,000}{107\,775} \approx 2{,}0
\]
Çok vuruşluk işlemci bu kez tam iki kat hızlıdır. Aradaki farkın kaynağı bellidir. Tek vuruşlukta en pahalı buyruk herkesin çevrimini belirler. Çok vuruşlukta ise her buyruk yalnızca kendine gereken çevrimleri harcar. Çarpma, bölme ve kayan nokta gibi pahalı işlemler gerçek buyruk kümelerinin hemen hepsinde bulunduğundan çok vuruşluk yaklaşım kaçınılmazdı. İlk mikroişlemcilerin yaklaşık yirmi yıl boyunca hep çok vuruşluk olması bundandır.
\end{ornek}

%% ================================================================
%% MİKROPROGRAMLAMA
%% ================================================================

\section{Mikroprogramlama}
\label{sec:mikroprogramlama}

Çok vuruşluk veri yolundaki sonlu durum makinesi, her durumda veri yoluna gönderilecek denetim sinyalleri kümesini ve bir sonraki duruma geçişi belirler. Buyruk kümesi büyüdükçe ya da denetim mantığı karmaşıklaştıkça, durumları ve geçişleri sabit donanımla kodlamak güçleşir. Bu noktada, durumların her birini birer bellek sözcüğünde saklamak ve donanımın bu sözcükleri sırayla okuyup işlemesini sağlamak akla gelir. İşte bu yaklaşıma \textbf{mikroprogramlama}\index{mikroprogramlama}\index{microprogramming|see{mikroprogramlama}} denir. Her mikrobuyruk, durum makinesindeki bir durumun denetim sinyallerini taşıyan bir ``sözcük''tür.

İki yaklaşım arasındaki fark Şekil~\ref{fig:donanim-vs-mikro}'da özetlenmiştir. Donanım tabanlı denetimde sinyaller her çevrimde kapı ağında yeniden \emph{hesaplanır}. Durum ile işlem kodu kapılardan geçer ve sinyal değerleri çıkışta belirir. Mikroprogramlamada ise aynı değerler bellek satırlarında hazır bekler. Mikro-PS hangi satırı gösteriyorsa o satır olduğu gibi \emph{okunur}. Şekildeki bellek içeriği rastgele değildir. Satırlar, ileride Tablo~\ref{tab:r-turu-mikrobuyruk}'de ayrıntısıyla vereceğimiz R türü mikrobuyruk dizisinin ta kendisidir ve iki yolun aynı örnekte (R türü buyruğun Yürüt çevrimi) aynı sinyalleri ürettiği görülmektedir.

\begin{figure}[htbp]
\centering
\begin{tikzpicture}[>=Stealth, font=\small, semithick]
  % ================= (a) Donanim tabanli =================
  \node[anchor=west, font=\small\bfseries] at (-0.2,4.55) {(a) Donanım tabanlı denetim: sinyaller kapı ağında \emph{hesaplanır}};
  \node[blok, minimum width=2.1cm, minimum height=1.0cm, font=\scriptsize] (gir1) at (0.9,3.4) {durum: S6 (Yürüt)\\opcode: R türü};
  \node[blok, fill=blokkirmizi!20, minimum width=2.2cm, minimum height=1.3cm, font=\scriptsize] (kapi) at (4.0,3.4) {kapı ağı\\(VE, VEYA, DEĞİL)};
  \node[blok, fill=bloksari, minimum width=2.4cm, minimum height=1.0cm, font=\scriptsize] (cik1) at (7.3,3.4) {AMBKaynakA = 00\\AMBKaynakB = 00};
  \draw[->] (gir1.east) -- (kapi.west);
  \draw[->] (kapi.east) -- (cik1.west);
  % ================= (b) Mikroprogramlama =================
  \node[anchor=west, font=\small\bfseries] at (-0.2,2.0) {(b) Mikroprogramlama: sinyaller bellekten \emph{okunur}};
  \node[blok, minimum width=1.7cm, minimum height=1.0cm, font=\scriptsize] (mps) at (0.9,0.2) {Mikro-PS\\\texttt{= 2}};
  % ROM satirlari (icerik = Tablo 6.12, R turu dizisi)
  \node[font=\scriptsize\bfseries, anchor=west] at (2.7,1.45) {Mikrobuyruk Belleği (ROM)};
  \node[draw, minimum width=4.6cm, minimum height=0.48cm, anchor=west, font=\tiny\ttfamily] (r0) at (2.7,0.95) {0: 1 0 0 1 01 01 xx 0 \ Sıra};
  \node[draw, minimum width=4.6cm, minimum height=0.48cm, anchor=west, font=\tiny\ttfamily] (r1) at (2.7,0.47) {1: 0 x x 0 10 10 xx 0 \ Dağıtım};
  \node[draw, fill=bloksari, minimum width=4.6cm, minimum height=0.48cm, anchor=west, font=\tiny\ttfamily] (r2) at (2.7,-0.01) {2: 0 x x 0 00 00 xx 0 \ Sıra};
  \node[draw, minimum width=4.6cm, minimum height=0.48cm, anchor=west, font=\tiny\ttfamily] (r3) at (2.7,-0.49) {3: 0 x x 0 xx xx 00 1 \ Yeni};
  \node[blok, fill=bloksari, minimum width=2.4cm, minimum height=1.0cm, font=\scriptsize] (cik2) at (9.5,0.2) {AMBKaynakA = 00\\AMBKaynakB = 00};
  \draw[->] (mps.east) -- (2.2,0.2) -- (2.2,-0.01) -- (r2.west);
  \draw[->] (r2.east) -- (7.7,-0.01) -- (7.7,0.2) -- (cik2.west);
\end{tikzpicture}
\caption[Donanım tabanlı denetim ile mikroprogramlamanın karşılaştırması]{Aynı denetim sorusuna iki yanıt: donanım tabanlı denetim sinyalleri kapılarla hesaplar, mikroprogramlama hazır satırdan okur. Bellek satırları Tablo~\ref{tab:r-turu-mikrobuyruk}'deki R türü dizisidir (sütun sırası da aynıdır); Mikro-PS $=$ 2, R türü buyruğun Yürüt çevrimine karşılık gelen satırı seçer.}
\label{fig:donanim-vs-mikro}
\end{figure}

Sonlu durum makineleriyle denetim birimi tasarlamak, küçük buyruk kümeleri için işe yarıyordu. Ancak buyruk sayısı arttıkça bu yaklaşım sürdürülemez hale gelir. RISC-V'in temel buyruk kümesinde bile yaklaşık 40 buyruk vardır. Uzantılarla (M, A, F, D vb.) bu sayı önemli ölçüde artar. Her buyruk için ortalama 4--5 aşama düşünüldüğünde yüzlerce durum, binlerce mantık kapısı gerekir. Daha pratik bir çözüm yok mu?

Mikroprogramlama, donanım tabanlı denetim biriminin bu sorunlarını ortadan kaldırmak için geliştirilmiştir. Mikroprogramlar \textbf{mikrobuyruk}\index{mikrobuyruk}\index{microinstruction|see{mikrobuyruk}}lardan oluşur. Mikrobuyruklar buyrukların daha küçük parçalara bölünmüş halleridir. Buyruklar için gerekli denetimler mikrobuyruklar tarafından sağlanır. Örneğin, bir toplama buyruğu getir, çöz, yürüt, sonuç yaz mikrobuyrukları tarafından gerçekleştirilir.

Bir bellek biriminde\index{mikro program belleği}\index{control store|see{mikro program belleği}} tutulan mikrobuyruklar yapılacak işleme göre bellekten alınarak veri yolunda kullanılır. Bir ROM'a ya da EPROM'a yerleştirilen mikrobuyruklar o anda bulunulan duruma göre sırayla çekilir. EPROM kullanılan denetim birimlerinde yeni eklenecek buyruklar için mikrobuyruklar ekleme kolaylıkla yapılır. Ayrıca bazı sonradan fark edilen tasarım hataları da mikrobuyrukların tekrar düzenlenmesi sayesinde düzeltilebilir.

Mikroprogramlama işlemciye kendisini başka işlemcilere benzetme olanağı verir. Eğer veri yolu bir başka işlemcinin gerçekleştirdiği buyrukların sonuçlarını verdirecek şekilde kullanılırsa işlemci iki farklı işlemci olarak da kullanılabilir.

\subsection{Dikey ve Yatay Mikroprogramlama}
\label{subsec:dikey-yatay}

Mikroprogramlama için iki yöntem vardır:

\begin{tanim}[Dikey ve Yatay Mikroprogramlama]
\textbf{Yatay mikroprogramlama:}\index{yatay mikroprogramlama}\index{horizontal microprogramming|see{yatay mikroprogramlama}} Mikrobuyrukların doğrudan veri yolundaki denetim girişlerini etkilemesidir. Daha hızlı sonuçlar verir.

\textbf{Dikey mikroprogramlama:}\index{dikey mikroprogramlama}\index{vertical microprogramming|see{dikey mikroprogramlama}} Veri yolundaki denetim girişlerine verilecek değerlerin mikrobuyruğun çözümlenmesinden sonra oluşturulmasıdır. Denetim birimlerine değerlerin ulaşması daha uzun sürer.
\end{tanim}

Aradaki farkı küçük bir örnekle görelim. \texttt{sw} buyruğunun Bellek çevriminde yatay mikrobuyruk tüm denetim bitlerini açık biçimde taşır (AdresKaynağı = 1, BelleğeYaz = 1, öbür alanlar etkisiz). Dikey kodlamada ise aynı çevrim, mikrobuyrukta yalnızca kısa bir işlem koduyla (örneğin \texttt{BELYAZ}) saklanır. Bu kod, mikrobuyruk belleğinden okunduktan sonra bir kod çözücüde yukarıdaki sinyallere açılır. Sözcük kısalır. Karşılığında her çevrime bir kod çözme gecikmesi eklenir.

Yatay mikroprogramlamada mikrobuyruk sözcüğünün her alanı doğrudan bir denetim noktasına bağlanır ve kod çözücüye gerek yoktur. Bu yapı Şekil~\ref{fig:yatay-mikro}'de gösterilmektedir.

\begin{figure}[htbp]
\centering
\begin{tikzpicture}[>=Stealth, node distance=2cm]
  % Mikrobuyruk belleği
  \node[bellekblok, minimum width=3cm, minimum height=1.5cm] (mbel) {Mikrobuyruk\\Belleği};
  % Doğrudan bağlantı
  \node[blok, right=3cm of mbel, minimum width=3cm, minimum height=1.5cm] (denetim) {Denetim\\Noktaları};
  % Veri yolu
  \node[blok, right=3cm of denetim, minimum width=2.5cm, minimum height=1.5cm, fill=bloksari] (veriyolu) {Veri yolu};
  % Bağlantılar
  \draw[->, semithick] (mbel.east) -- node[above, font=\footnotesize, align=center]{Doğrudan\\bağlantı} (denetim.west);
  \draw[denetimstil] (denetim.east) -- node[above, font=\footnotesize]{Denetim sinyalleri} (veriyolu.west);
  % Mikro-PS
  \node[yazmacblok, below=1.5cm of mbel, minimum width=2cm] (mps) {Mikro-PS};
  \draw[->, semithick] (mps) -- (mbel);
  % Açıklama
  \node[font=\footnotesize, text=gray, below=0.5cm of denetim] {Kod çözücü yok $\rightarrow$ Hızlı};
\end{tikzpicture}
\caption{Yatay mikroprogramlama}
\label{fig:yatay-mikro}
\end{figure}

Dikey mikroprogramlamada ise mikrobuyruk alanları önce bir kod çözücüden geçer. Bu ek gecikme karşılığında daha kısa sözcük boyutu elde edilir (Şekil~\ref{fig:dikey-mikro}).

\begin{figure}[htbp]
\centering
\begin{tikzpicture}[>=Stealth, node distance=2cm]
  % Mikrobuyruk belleği
  \node[bellekblok, minimum width=3cm, minimum height=1.5cm] (mbel) {Mikrobuyruk\\Belleği};
  % Kod çözücü
  \node[blok, right=2cm of mbel, minimum width=2cm, minimum height=1.5cm, fill=blokkirmizi-acik] (kodcoz) {Kod\\Çözücü};
  % Denetim noktaları
  \node[blok, right=2cm of kodcoz, minimum width=3cm, minimum height=1.5cm] (denetim) {Denetim\\Noktaları};
  % Veri yolu
  \node[blok, right=2.5cm of denetim, minimum width=2.5cm, minimum height=1.5cm, fill=bloksari] (veriyolu) {Veri yolu};
  % Bağlantılar
  \draw[->, semithick] (mbel.east) -- node[above, font=\footnotesize]{Kodlanmış} (kodcoz.west);
  \draw[->, semithick] (kodcoz.east) -- node[above, font=\footnotesize]{Çözülmüş} (denetim.west);
  \draw[denetimstil] (denetim.east) -- node[above, font=\footnotesize]{Denetim} (veriyolu.west);
  % Mikro-PS
  \node[yazmacblok, below=1.5cm of mbel, minimum width=2cm] (mps) {Mikro-PS};
  \draw[->, semithick] (mps) -- (mbel);
  % Açıklama
  \node[font=\footnotesize, text=gray, below=0.5cm of kodcoz] {Kod çözücü gecikme ekler};
\end{tikzpicture}
\caption{Dikey mikroprogramlama}
\label{fig:dikey-mikro}
\end{figure}

\subsection{Mikrobuyruk Yapısı}
\label{subsec:mikrobuyruk-yapisi}

Veri yolunda işlem yapan her bir birim için mikrobuyruk üzerinde bir alan ayrılır. Mikrobuyruk dokuz alandan oluşur: PS denetimi, yazmaç öbeği denetimi, A ve B giriş seçimleri, AMB denetimi, bellek denetimi, geri yazma kaynağı seçimi, anlık biçim seçimi ve bir sonraki mikrobuyruğun nereden alınacağını söyleyen \emph{Sonraki} alanı. O anda işlenecek buyruğun her aşamada gerek duyduğu denetim değerleri, o aşamanın mikrobuyruğu üzerinde taşınır.

Bir sonraki mikrobuyruğun seçilmesi için sıralama düzeni:
\begin{itemize}
    \item \textbf{Sıra:} Sıralı şekilde devam edecek olan mikrobuyruk seçilir.
    \item \textbf{Yeni:} Buyruğun işlenmesi bitmiştir ve yeni buyruğun ilk mikrobuyruğu çekilecektir.
    \item \textbf{Dağıtım (tablo\_no):} Buyruk görev dağıtım tablosundan çekilecektir.
\end{itemize}

Her alanın alabileceği değerlerin tam listesi Tablo~\ref{tab:mikrobuyruk-degerler}'de verilmektedir.

\begin{table}[htbp]
\centering
\caption{Mikrobuyruk üzerindeki alanların alabileceği değerler}
\label{tab:mikrobuyruk-degerler}
\small
\begin{tabular}{llp{8cm}}
\toprule
\textbf{Alan} & \textbf{Değer} & \textbf{Açıklama} \\
\midrule
\multirow{4}{*}{PS Denetim}
 & AMB çıkışı & Getir çevriminde hesaplanan PS+4 değeri PS'ye yazılır. \\
 & AMB Yzm. VE PSKoşul & Çöz'de hesaplanmış dal hedefi, koşul sağlanırsa PS'ye yazılır. \\
 & AMB Yzm. (jal) & Çöz'de hesaplanmış atlama hedefi PS'ye koşulsuz yazılır. \\
 & Yazma yok & PS bu çevrimde değişmez (PSYaz $=$ 0, PSKoşul $=$ 0). \\
\midrule
\multirow{3}{*}{Yazmaç Öbeği Den.}
 & Oku & rs1 ve rs2 yazmaçları okunur. \\
 & rd'ye yaz & Sonuç rd yazmacına yazılır. \\
 & Etkisiz & Bu çevrimde yazmaç öbeğine yazılmaz (YazmacaYaz $=$ 0). \\
\midrule
\multirow{3}{*}{A yazmacı Den.}
 & A & A yazmacının değeri verilir (Yürüt). \\
 & PS & AMB A girişine güncel PS değeri verilir (Getir'de PS+4). \\
 & EskiPS & Getir'de saklanan buyruk adresi verilir (Çöz'de dal/jal hedefi). \\
\midrule
\multirow{3}{*}{B yazmacı Den.}
 & B & B yazmacındaki değer verilir. \\
 & 4 & PS'nin 4 artırılmasını sağlar. \\
 & Anlık & İşaretle genişletilmiş anlık değer (biçimini AnlıkBiçim alanı seçer). \\
\midrule
AnlıkBiçim Den. & I / S / B / J & Anlık değer üretecinin buyruk biçimi seçimi. \\
\midrule
\multirow{3}{*}{AMB Denetim}
 & Topla & Toplama işlemi yapılır. \\
 & Çıkar & Çıkarma işlemi yapılır. \\
 & R türü & funct3 ve funct7 değerine göre işlem yapılır. \\
\midrule
\multirow{3}{*}{Bellek Denetim}
 & Buyruk oku & PS adresindeki buyruk okunur. \\
 & Veri oku & Yükle buyruğu için veri okunur. \\
 & Veri yaz & Sakla buyruğu için veri yazılır. \\
\midrule
\multirow{3}{*}{GeriYazma Kaynağı}
 & AMB sonucu & AMB çıkışı rd yazmacına yazılır (R türü, ori). \\
 & Veri belleği & Bellekten okunan değer rd yazmacına yazılır (lw). \\
 & PS$+$4 & Dönüş adresi rd yazmacına yazılır (jal). \\
\midrule
\multirow{3}{*}{Sonraki}

 & Sıra & Sıradaki mikrobuyruk alınır. \\
 & Yeni & Yeni bir mikrobuyruk alınır. \\
 & Dağıtım tablo\_no & Belirtilen tablodan çekilir. \\
\bottomrule
\end{tabular}
\\[0.3em]
{\scriptsize Bellek Denetim alanı donanımda AdresKaynağı, BuyruğaYaz ve BelleğeYaz sinyallerine açılır. Bu tasarımda okuma yan etkisiz olduğundan ayrı bir okuma izni sinyali yoktur, ``Buyruk oku'' ve ``Veri oku'' değerleri yalnızca adres kaynağının seçimini anlatır.}
\end{table}

Alan genişlikleri toplandığında mikrobuyruk sözcüğünün boyutu da ortaya çıkar: PS denetimi 3, yazmaç öbeği 1, A kaynağı 2, B kaynağı 2, AMB 2, bellek denetimi 3, GeriYazmaKaynağı 2, AnlıkBiçim 2 ve Sonraki 2 bit olmak üzere toplam 19 bit. On iki satırlık mikroprogram böylece $12 \times 19 = 228$ bitlik, birkaç on baytlık minicik bir ROM'a sığar. Gerçek bir buyruk kümesinde tablo yüzlerce satıra, sözcük onlarca bite büyür. Mikroprogramlamanın gücü de tam burada kendini gösterir. Buyruk eklemek kapı eklemek değil, satır eklemektir.

Örneğin R türü buyruklar için mikrobuyruk dizisi Tablo~\ref{tab:r-turu-mikrobuyruk}'de verilmektedir. R türü dört çevrimde tamamlanır (Getir, Çöz, Yürüt, Yaz) ve dizinin sütun düzeni lw dizisiyle (Tablo~\ref{tab:lw-mikrobuyruk}) aynıdır.

\begin{table}[htbp]
\centering
\caption{R türü buyruklar için mikrobuyruk dizisi}
\label{tab:r-turu-mikrobuyruk}
\footnotesize
\resizebox{\textwidth}{!}{%
\begin{tabular}{l|c|c|c|c|c|c|c|c|c}
\toprule
\textbf{Aşama} & \textbf{PSYaz} & \textbf{PSKaynak} & \textbf{Adres-} & \textbf{Buyruğa-} & \textbf{AMB-} & \textbf{AMB-} & \textbf{GeriYazma-} & \textbf{Yazmaca-} & \textbf{Sonraki} \\
 & & & \textbf{Kaynağı} & \textbf{Yaz} & \textbf{KaynakA} & \textbf{KaynakB} & \textbf{Kaynağı} & \textbf{Yaz} & \\
\midrule
Getir  & 1 & 0 & 0 & 1 & 01 & 01 & xx & 0 & Sıra \\
Çöz    & 0 & x & x & 0 & 10 & 10 & xx & 0 & Dağıtım \\
Yürüt  & 0 & x & x & 0 & 00 & 00 & xx & 0 & Sıra \\
Yaz    & 0 & x & x & 0 & xx & xx & 00 & 1 & Yeni \\
\bottomrule
\end{tabular}%
}
\\[0.3em]
{\scriptsize BelleğeYaz tüm aşamalarda 0'dır. PSKoşul kullanılmaz. AMBKaynakA: 00 = A yazmacı, 01 = PS, 10 = EskiPS. AMBKaynakB: 00 = B yazmacı, 01 = 4, 10 = anlık. Getir'de PS$+$4, Çöz'de EskiPS $+$ anlık (dal hedefi; R türü kullanmaz) hesaplanır. Yürüt aşamasında AMB işlemi funct3/funct7 alanlarından çözülür (Tablo~\ref{tab:mikrobuyruk-degerler}'deki ``R türü'' değeri). Sonraki: Sıra, Dağıtım (işlem koduna göre), Yeni.}
\end{table}

Yükle (lw) buyruğu beş aşamadan geçtiği için mikroprogramı beş satırlık bir mikrobuyruk dizisidir. Tablo~\ref{tab:lw-mikrobuyruk}, her aşamada veri yoluna gönderilen denetim sinyallerini gösterir. Son sütun bir sonraki mikrobuyruğun adresini taşır.

\begin{table}[htbp]
\centering
\caption{lw buyruğu için mikrobuyruk dizisi}
\label{tab:lw-mikrobuyruk}
\footnotesize
\resizebox{\textwidth}{!}{%
\begin{tabular}{l|c|c|c|c|c|c|c|c|c}
\toprule
\textbf{Aşama} & \textbf{PSYaz} & \textbf{PSKaynak} & \textbf{Adres-} & \textbf{Buyruğa-} & \textbf{AMB-} & \textbf{AMB-} & \textbf{GeriYazma-} & \textbf{Yazmaca-} & \textbf{Sonraki} \\
 & & & \textbf{Kaynağı} & \textbf{Yaz} & \textbf{KaynakA} & \textbf{KaynakB} & \textbf{Kaynağı} & \textbf{Yaz} & \\
\midrule
Getir  & 1 & 0 & 0 & 1 & 01 & 01 & xx & 0 & Sıra \\
Çöz    & 0 & x & x & 0 & 10 & 10 & xx & 0 & Dağıtım \\
Yürüt  & 0 & x & x & 0 & 00 & 10 & xx & 0 & Sıra \\
Bellek & 0 & x & 1 & 0 & xx & xx & xx & 0 & Sıra \\
Yaz    & 0 & x & x & 0 & xx & xx & 01 & 1 & Yeni \\
\bottomrule
\end{tabular}%
}
\\[0.3em]
{\scriptsize BelleğeYaz tüm aşamalarda 0'dır. PSKoşul kullanılmaz. AMBKaynakA: 00 = A yazmacı, 01 = PS, 10 = EskiPS. Çöz aşamasında AMB, dal/jal olasılığına karşı hedef adresi (EskiPS + anlık) hesaplar; lw bu değeri kullanmaz. Sonraki alanı: Sıra = sıradaki mikrobuyruk, Dağıtım = işlem koduna göre dağıtım, Yeni = yeni buyruğun getirilmesi.}
\end{table}

Bu mikrobuyruk dizilerini bellekten sırayla çeken ve sonraki adresi hesaplayan donanım, Şekil~\ref{fig:mikroprogramlama-donanim}'te gösterilmektedir.

\begin{figure}[htbp]
\centering
\begin{tikzpicture}[>=Stealth]
  % Mikro-PS ve Mikrobuyruk Belleği (açık koordinatlar)
  \node[yazmacblok, minimum width=2cm, minimum height=1cm] (mps) at (0,0) {Mikro-PS};
  \node[bellekblok, minimum width=4cm, minimum height=1.5cm] (mbel) at (6,0) {Mikrobuyruk Belleği};
  % Mikrobuyruk sözcüğü alanları (belleğe yakın, soldan sağa zincir)
  \node[bitalani, minimum width=1.3cm, minimum height=0.6cm, font=\scriptsize] (a1) at (0.65,-2.0) {PS Den.};
  \node[bitalani, minimum width=1.3cm, minimum height=0.6cm, anchor=west, font=\scriptsize] (a2) at (a1.east) {Yzm.Ö.};
  \node[bitalani, minimum width=1cm, minimum height=0.6cm, anchor=west, font=\scriptsize] (a3) at (a2.east) {A den.};
  \node[bitalani, minimum width=1cm, minimum height=0.6cm, anchor=west, font=\scriptsize] (a4) at (a3.east) {B den.};
  \node[bitalani, minimum width=1.3cm, minimum height=0.6cm, anchor=west, font=\scriptsize] (a5) at (a4.east) {AMB D.};
  \node[bitalani, minimum width=1.2cm, minimum height=0.6cm, anchor=west, font=\scriptsize] (a6) at (a5.east) {Bel. D.};
  \node[bitalani, minimum width=1.2cm, minimum height=0.6cm, anchor=west, font=\scriptsize, fill=bloksari] (a7) at (a6.east) {Sonraki};
  % Mikro-PS -> Bellek
  \draw[->, semithick] (mps.east) -- (mbel.west);
  % Bellek -> alanların 0.4cm üstünde yatay dağıtım hattı -> her alana kısa ok
  \draw[->, semithick] (mbel.south) -- (6,-1.3);
  \draw[semithick] ($(a1.north)+(0,0.4)$) -- ($(a7.north)+(0,0.4)$);
  \foreach \a in {a1,a2,a3,a4,a5,a6,a7} {
    \draw[->, semithick] ($(\a.north)+(0,0.4)$) -- (\a.north);
  }
  % Denetim sinyalleri aşağı (a1..a6)
  \foreach \a in {a1,a2,a3,a4,a5,a6} {
    \draw[denetimstil] (\a.south) -- ++(0,-1);
  }
  \node[font=\footnotesize, text=sinyalkirmizi] at (4.15,-3.75) {Veri yoluna denetim sinyalleri};
  % Sonraki -> Mikro-PS geri besleme: alanların ve denetim sinyallerinin ALTINDAN dik açılı dolaş
  \draw[->, semithick] (a7.south) -- (a7.south |- 0,-4.3) -- (-1.5,-4.3) -- (-1.5,0) -- (mps.west);
  \node[font=\scriptsize, anchor=north] at (3,-4.4) {Sonraki adres};
\end{tikzpicture}
\caption[Mikroprogramlama donanımı]{Mikroprogramlama donanımı. Mikrobuyruk sekiz alandan oluşur (Tablo~\ref{tab:mikrobuyruk-degerler}); AnlıkBiçim alanı ile sonraki adres seçim mantığı (işlem kodu/Dağıtım girişi) yalınlık için gösterilmemiştir.}
\label{fig:mikroprogramlama-donanim}
\end{figure}

Bu noktada mikroprogramın aslında ne olduğu görünür hâle gelir. Mikroprogram, Şekil~\ref{fig:durum-birlesik}'deki birleşik durum makinesinin bellekte saklanan biçimidir. Her durum bir mikrobuyruk satırına karşılık gelir. On iki durumlu makinemiz on iki satırlık bir mikroprograma dönüşür ve ilk iki satır (Getir ile Çöz) bütün buyrukların ortak girişidir. Sonraki alanının üç değeri de donanımda Mikro-PS'nin girişindeki küçük bir seçiciyle karşılanır. \emph{Sıra} için Mikro-PS bir artırılır; \emph{Yeni} için adres sıfıra, Getir satırına döndürülür; \emph{Dağıtım} için ise işlem kodu, \emph{dağıtım tablosu} denen küçük bir ROM'dan geçirilerek ilgili buyruğun ilk mikrobuyruğunun adresine çevrilir. Durum makinesindeki S1 çıkışının beş yola ayrılması, mikroprogram dünyasında bu tablo araması olarak yaşar.

Yatay mikroprogramlamada mikrobuyruk sözcüğünün her alanı doğrudan ilgili denetim noktasına bağlanır. Bu donanım düzenlemesi Şekil~\ref{fig:yatay-mikro-donanim}'te görülmektedir.

\begin{figure}[H]
\centering
\begin{tikzpicture}[>=Stealth, scale=1.1, transform shape, node distance=1.2cm]
  % Mikrobuyruk sözcüğü alanları
  \node[bitalani, minimum width=2cm, minimum height=0.7cm, font=\scriptsize] (f1) at (0,1.5) {PS Den. (3 bit)};
  \node[bitalani, minimum width=2cm, minimum height=0.7cm, anchor=west, font=\scriptsize] (f2) at (f1.east) {Yzm. Den. (1 bit)};
  \node[bitalani, minimum width=1.5cm, minimum height=0.7cm, anchor=west, font=\scriptsize] (f3) at (f2.east) {A Den. (2 bit)};
  \node[bitalani, minimum width=2cm, minimum height=0.7cm, anchor=west, font=\scriptsize] (f4) at (f3.east) {B Den. (2 bit)};
  \node[bitalani, minimum width=1.5cm, minimum height=0.7cm, anchor=west, font=\scriptsize] (f5) at (f4.east) {AMB (2 bit)};
  \node[bitalani, minimum width=1.5cm, minimum height=0.7cm, anchor=west, font=\scriptsize] (f6) at (f5.east) {Bel. (3 bit)};
  \node[bitalani, minimum width=1.7cm, minimum height=0.7cm, anchor=west, font=\scriptsize] (f7) at (f6.east) {GeriYzm. (2 bit)};
  % Başlık: kutu dizisinin ortasına
  \node[font=\small\bfseries, anchor=south] at ($(f1.north)!0.5!(f7.north) + (0,0.12)$) {Mikrobuyruk Sözcüğü};
  % Doğrudan denetim noktalarına (kısa oklar)
  \foreach \i/\lbl in {f1/PS, f2/Yzm.Ö., f3/AMB-A, f4/AMB-B, f5/AMB, f6/Bellek, f7/GeriYazma} {
    \draw[->, semithick, sinyalmavi] (\i.south) -- ++(0,-0.9) node[below, font=\scriptsize, align=center]{\lbl\\denetim};
  }
  \node[font=\footnotesize, text=gray] at ($(f1.south)!0.5!(f7.south) + (0,-2.15)$) {Her alan doğrudan ilgili denetim noktasına bağlanır};
\end{tikzpicture}
\caption[Yatay mikroprogramlama donanımı]{Yatay Mikroprogramlama donanımı. Yalınlık için yalnızca veri yoluna giden alanlar çizilmiştir; AnlıkBiçim ile Sonraki alanları gösterilmemiştir (tam liste Tablo~\ref{tab:mikrobuyruk-degerler}).}
\label{fig:yatay-mikro-donanim}
\end{figure}

Mikroprogramlamanın kötü yönü donanım tabanlı denetim biriminden daha yavaş çalışmasıdır. Çünkü bellekten veri alınması, mantık devrelerinden daha çok vakit alır. Mikroprogramlama, bir dönem karmaşık denetimi düzenli biçimde tasarlamanın kolay yolu olduğu için yaygınlaşmıştı; ancak bellek teknolojisi ilerleyip erişim gecikmeleri kısaldıkça, sade buyrukları doğrudan donanımla işleyen RISC yaklaşımı ve donanım tabanlı denetim öne geçti, mikroprogramlama görece gözden düştü.

Donanım tabanlı denetim birimine göre daha yavaş olan mikrokodlamanın tercih edilme sebebi tasarımının daha yalın ve esnek olmasıdır.

Bu esnekliğin asıl değeri, değişikliğin tasarımın ne kadar \emph{geç} bir aşamasında yapılabildiğinde ortaya çıkar. Karmaşık bir işlemcinin mimari tasarımı, doğrulanması ve üretime hazırlanması yıllar, büyük çekirdeklerde çoğu zaman üç ilâ beş yıl sürer. Donanım tabanlı denetimde bir buyruğun davranışını değiştirmek ya da yeni bir buyruk eklemek, denetim mantığının kapı ağını yeniden tasarlayıp yongayı baştan üretmek demektir. Bu da yıllar süren döngünün büyük bölümünü yeniden çalıştırmak kadar pahalıdır. Mikroprogramlamada aynı değişiklik, denetim belleğindeki satırları düzenlemekten ibarettir. Denetim belleği alanda güncellenebilir bir teknolojiyle gerçekleştirildiğinde (eskiden EPROM, çağdaş işlemcilerde yüklenebilir mikrokod), davranış üretimden sonra bile değiştirilebilir. Sonradan ortaya çıkan tasarım hataları donanım değiştirmeden, bir yazılım güncellemesi gibi yamalanır. Yazılabilir denetim belleği kullanan tasarımlarda bu esneklik, üretilmiş bir işlemciye yeni mikrobuyruk dizileri ekleyerek buyruk kümesini genişletmeye kadar varır. Donanım tabanlı bir tasarımda aynı düzeltme, yeni bir silikon dönüşü ve aylarca bekleme gerektirirdi.

Bu esnekliği bir örnekle somutlaştıralım.

\begin{ornek}[6.10 -- Mikroprograma yeni buyruk eklemek]
Buyruk kümesine, kaynak yazmaçla anlık değeri bit düzeyinde VE'leyen \texttt{andi} buyruğu eklenecektir. AMB'nin VE işlemi yapabildiğini ve AMBİşlem kodlamasında 11 değerinin boş olduğunu varsayarak (a)~mikroprogramlanmış denetimde nelerin değişeceğini gösterin, (b)~aynı eklemenin donanım tabanlı denetimde neler gerektireceğini tartışın.

\tcblower
\textbf{Çözüm.}

\textbf{(a)} Getir ile Çöz satırları bütün buyrukların ortak girişi olduğundan onlara dokunulmaz. Yapılacak yalnızca iki şey vardır: mikrobuyruk belleğine \texttt{andi}'nin Yürüt ve Yaz satırlarını eklemek, dağıtım tablosuna da \texttt{andi}'nin işlem kodunu bu yeni dizinin ilk satırına eşleyen bir giriş yazmak.

\begin{center}
\footnotesize
\resizebox{\textwidth}{!}{%
\begin{tabular}{l|c|c|c|c|c|c|c|c|c}
\toprule
\textbf{Aşama} & \textbf{PSYaz} & \textbf{PSKaynak} & \textbf{Adres-} & \textbf{Buyruğa-} & \textbf{AMB-} & \textbf{AMB-} & \textbf{GeriYazma-} & \textbf{Yazmaca-} & \textbf{Sonraki} \\
 & & & \textbf{Kaynağı} & \textbf{Yaz} & \textbf{KaynakA} & \textbf{KaynakB} & \textbf{Kaynağı} & \textbf{Yaz} & \\
\midrule
Yürüt (andi) & 0 & x & x & 0 & 00 & 10 & xx & 0 & Sıra \\
Yaz (andi)   & 0 & x & x & 0 & xx & xx & 00 & 1 & Yeni \\
\bottomrule
\end{tabular}%
}
\end{center}

Satırlar ori dizisinin neredeyse kopyasıdır. Tek anlam farkı AMB Denetim alanının VE işlemini (AMBİşlem $=$ 11) seçmesidir. Toplam değişiklik: iki bellek satırı ve bir dağıtım girişi. Veri yolunda tek bir tel bile değişmez.

\textbf{(b)} Donanım tabanlı denetimde aynı ekleme, denetim biriminin \emph{devresini} değiştirmek demektir. Buyruk kodu tablosuna yeni satır eklenir, ana çözücünün ve AMB denetiminin Karno haritaları yeniden çizilir, sadeleşmiş Boole ifadeleri değiştiği için kapı ağı yeniden kurulur ve tüm buyrukların hâlâ doğru çalıştığı yeniden doğrulanır. Tasarım aşamasında bu yalnızca emek demektir. Yonga üretildikten sonra ise olanaksızdır. Mikroprogramlanmış tasarımda denetim belleği yazılabilirse aynı iki satır, sahadaki işlemciye bile eklenebilir.
\end{ornek}

Bu iki yaklaşımın karşılaştırması Tablo~\ref{tab:donanim-mikro-karsilastirma}'te verilmektedir.

\begin{table}[htbp]
\centering
\caption{Donanım tabanlı denetim ile mikroprogramlama karşılaştırması}
\label{tab:donanim-mikro-karsilastirma}
\begin{tabular}{lcc}
\toprule
\textbf{Özellik} & \textbf{Donanım tabanlı} & \textbf{Mikroprogramlama} \\
\midrule
Hız & Yüksek & Düşük \\
Esneklik & Düşük & Yüksek \\
Tasarım karmaşıklığı & Yüksek & Düşük \\
Hata düzeltme & Zor & Kolay \\
Büyük buyruk kümeleri & Uygun değil & Uygun \\
Donanım maliyeti & Mantık kapıları & ROM/RAM bellek \\
\bottomrule
\end{tabular}
\end{table}

\begin{bilginot}[Mikroprogramlamanın Tarihçesi]
Mikroprogramlama kavramı Maurice Wilkes tarafından 1951'de önerilmiştir. IBM System/360 (1964) ailesi mikroprogramlamayı yaygınlaştırmıştır. Aynı buyruk kümesini destekleyen farklı fiyat/başarım seviyesinde modeller, yalnızca mikro program belleğinin içeriği değiştirilerek üretilmiştir. Günümüzde x86 işlemcileri karmaşık CISC buyrukları iç mikro işlemlere (micro-ops) çevirmek için hâlâ mikrokod kullanmaktadır.
\end{bilginot}

%% ================================================================
%% Gerçek Dünya: İşlemci Tasarım Kararları
%% ================================================================

\section{Gerçek Dünya: İşlemci Tasarım Kararları}
\label{sec:islemci-gercek-dunya}

İşlemci tasarımı, ders kitaplarındaki temiz veri yolu şemalarından çok daha karmaşık mühendislik kararlarını içerir.

\subsection{Bu bölümde anlatılanlar gerçekte nerede kullanılıyor?}
Bu bölümde incelediğimiz teknikler yalnızca öğretim araçları değildir. İlk mikroişlemciler yaklaşık yirmi yıl boyunca hep çok vuruşluk tasarımlardı. Intel'in ilk mikroişlemcisi 4004 (1971) buyruklarını 8 saat vuruşluk çevrimlerle işliyordu (uzun buyruklar iki çevrim sürüyordu). 8086'dan 386'ya uzanan kuşaklar da mikrokodlu çok vuruşluk tasarımlardı ve bir buyruk, türüne göre birkaç çevrimden yüzlerce çevrime kadar sürebiliyordu. Bu tercihin sayısal gerekçesini Örnek~6.9'da görmüştük. Çarpma ve bölme gibi pahalı buyruklar barındıran bir kümede tek vuruşluk yaklaşım herkesi en yavaş buyruğa mahkûm eder. Intel ancak 486 (1989) ile boru hattına geçti. O günden beri hiçbir Intel işlemcisi bu bölümdeki anlamıyla tek ya da çok vuruşluk değildir. Saf \emph{tek vuruşluk} işlemci ise ticari dünyada yok denecek kadar azdır, çünkü Örnek~6.4'te gördüğümüz gibi bellek erişimli buyruklar çevrimi herkes için uzatır. Bu ideale en çok yaklaşanlar, çoğu buyruğu tek çevrimde bitiren AVR gibi Harvard mimarili mikrodenetleyicilerdir (onlar da getirmeyi bir önceki çevrimle örtüştürür). Çok vuruşluk yaklaşım bugün de üretimdedir. Küçük mikrodenetleyicilerde, çevre birim denetleyicilerinde ve düşük güç hedefleyen gömülü çekirdeklerde, buyrukları birkaç çevrimde işleyen ve sonlu durum makinesiyle denetlenen çekirdekler (Bölüm~\ref{subsec:tek-vurusluk-sorunlar}'de andığımız PicoRV32 gibi) alanın ve gücün başarımdan değerli olduğu her yerde tercih edilir. Yüksek başarım tarafında ise merdiven yukarı doğru devam eder. İlk basamak, bir sonraki bölümün konusu olan boru hattıdır. Çağdaş Intel, AMD ve Apple işlemcileri ise tek saat çevriminde birden fazla buyruğu aynı anda başlatabilen \emph{çoklu buyruk başlatımlı} (\emph{superscalar}) mimarilere ulaşır. Apple'ın M serisi işlemcileri tek çevrimde 8'den fazla buyruk başlatabilir ve sırasız yürütüm (out-of-order execution) ile buyrukları en verimli sırada işler. Kaç kat merdiven çıkılırsa çıkılsın temel aynıdır: bir veri yolu, onu yöneten denetim ve buyruk başına düşen çevrim sayısını düşürme mücadelesi.

\subsection{Gerçek sistemlerde mikroprogramlama}
Mikroprogramlama 1970'lerde işlemci tasarımının temel yöntemiydi. RISC devrimiyle donanım tabanlı denetim öne çıktı. Ancak mikrokod ölmedi, iş bölümü yaptı. Çağdaş bir x86 işlemcisi iki dünyayı birlikte kullanır. Sık kullanılan yalın buyruklar donanım çözücülerden geçip doğrudan birkaç mikro-işleme (µop) çevrilirken, karmaşık ve seyrek buyruklar (dizgi kopyalama, sistem geçişleri gibi) yonga içindeki mikrokod ROM'undan okunan uzun mikro-işlem dizileri olarak yürütülür. Bölüm~\ref{subsec:mikrobuyruk-yapisi}'deki dağıtım tablosunun endüstriyel karşılığı tam da budur. İşlem kodu, ilgili mikro dizinin ROM'daki başlangıç adresine yönlendirilir.

Mikrokodun bugünkü en değerli yüzü ise güncellenebilirliğidir. İşlemcilerin mikrokod yamaları, her açılışta BIOS ya da işletim sistemi tarafından yüklenip yonga içindeki küçük bir yama belleğine yazılır. İmzalıdır ve kalıcı değildir. 2018'de ortaya çıkan Spectre ve Meltdown güvenlik açıklarına karşı ilk savunmaların bir bölümü, sahadaki milyonlarca işlemciye işte bu yolla, donanım değişmeden dağıtılmıştır. IBM'in z ana bilgisayar ailesi de karmaşık buyruklarını ve sistem işlevlerini \emph{milikod} adı verilen benzer bir katmanla yürütür. Buna karşılık Arm ve RISC-V gibi RISC çekirdeklerinde mikrokod ya hiç yoktur ya da yok denecek kadar azdır. Buyruklar zaten yalın olduğundan denetim doğrudan donanımla kurulur. Bu da mikroprogramlamanın tarihsel dersini doğrular, karmaşıklık nereye taşınırsa esneklik ve maliyet oraya taşınır.

\subsection{RISC-V'in basitliğinin değeri}
RISC-V'in sabit 32 bit buyruk biçimi ve düzenli yapısı, tek vuruşluk işlemcinin birkaç bin satır Verilog ile gerçekleştirilebilmesini sağlar. Bu basitlik, üniversite projelerinden ticari gömülü işlemcilere kadar geniş bir yelpazede avantaj sunar. Ayrıca denetim mantığı küçük olduğundan tasarımın biçimsel yöntemlerle doğrulanması da kolaylaşır. Buna karşılık, x86 kod çözücüsünün transistör maliyeti tek başına küçük bir RISC-V çekirdeğinden fazladır.


%% ================================================================
%% Yanılgılar ve Tuzaklar
%% ================================================================

\section{Yanılgılar ve Tuzaklar}
\label{sec:islemci-yanilgilar}

İşlemci tasarımında sezgisel kararlar çoğu zaman yanıltıcıdır. Kâğıt üzerinde doğru görünen bir tasarımın gerçek donanımda nasıl davranacağı, zamanlama, mantık derinliği ve denetim mantığı gibi birden çok etkenin etkileşimine bağlıdır. Bu bölümde işlemci tasarımında en sık karşılaşılan yanılgıları ve tasarım tuzaklarını somut örneklerle inceliyoruz.

\yanilgi{Tek vuruşluk işlemci en hızlısıdır çünkü her buyruk bir çevrimde biter}
Her buyruk bir çevrimde tamamlansa da o çevrimin süresi en yavaş buyruğa göre belirlenir. Yükle buyruğu beş aşama gerektirirken, dallanma buyruğu üç aşamada biter; ama ikisi de aynı uzun çevrimi bekler. Üstelik tek vuruşluk tasarımda bellek ve toplayıcı gibi birimlerin ayrı kopyaları gerekir (iki bellek, üç toplama birimi); çünkü aynı birim tek çevrim içinde iki farklı görevde kullanılamaz. Bu nedenle çok vuruşluk tasarım aynı işi daha az donanımla, boru hatlı tasarım ise benzer donanımla çok daha yüksek başarımla yapar. Tek vuruşluk her iki ölçütte de geride kalır.

\yanilgi{Çok vuruşluk tasarım her zaman tek vuruşluk tasarımdan daha verimlidir}
Çok vuruşluk tasarımda donanım yeniden kullanılır ve saat çevrim zamanı en uzun aşamayla sınırlanır; ancak aşamalar arasına yerleştirilen ara yazmaçlar (Buyruk Yzm., EskiPS, A/B, AMB Yzm., Bellek Yzm.) ek kurulum ve yayılım süreleri ekler. Aynı zamanda buyruk başına çevrim sayısı (BBÇ) artar. En yalın buyruklar bile birden çok çevrim alır. Toplam yürütme zamanı $\text{BBÇ} \times \text{çevrim zamanı} \times \text{buyruk sayısı}$ olduğu için, bu iki etken birbirini kısmen götürür. Hangi tasarımın hızlı olduğu hedef saat sıklığına, iş yüküne ve buyruk türü dağılımına bağlıdır. Bağımsız bir ``her zaman'' yargısı yapılamaz. Bölümdeki örnekler iki ucu birlikte gösterir. Yakın maliyetli buyruklarda tek vuruşluk öndedir (Örnek~6.8), kümeye pahalı bir buyruk girdiğinde çok vuruşluk iki kat kazanır (Örnek~6.9).

\yanilgi{Daha fazla denetim işareti daha iyi denetim demektir}
Denetim birimindeki işaret sayısının artması, olası hata noktalarını da artırır. Yatay mikroprogramlama çok esnek olsa da mikrokod ROM'unun boyutu hızla büyür. Dikey mikroprogramlama ya da donanım tabanlı denetim, daha sınırlı ama daha hızlı ve güvenilir olabilir.

\yanilgi{Mikroprogramlama her zaman donanım tabanlı denetim biriminden yavaştır}
Mikroprogramlama denetim sinyallerini her buyruk için mikrokod belleğinden okuduğundan basit görünüşte ek bir bellek erişimi gerektirir. Ancak çağdaş işlemcilerin mikrokod ROM'u önbellek hızında çalışır ve sıkça yürütülen buyrukların mikrokodu sürekli erişilebilir konumda tutulur. Daha da önemlisi, x86 gibi karmaşık buyruk kümelerinde donanım tabanlı denetim biriminin tasarlanması ve doğrulanması pratik değildir. Mikrokod, denetim mantığını yazılım gibi güncellenebilir kılarak hata düzeltmesini ve yeni buyruk eklemesini olanaklı kılar. Intel'in 1994 Pentium FDIV hatası ise bu esnekliğin sınırını gösterir. Hata, bölme biriminin donanıma gömülü arama tablosundaydı ve alanda mikrokodla düzeltilemediğinden Intel yongaları geri çağırıp değiştirmek zorunda kaldı. Alanda güncellenebilir mikrokod, bir sonraki kuşak olan P6 (Pentium Pro) ile yaygınlaştı ve sonraki hata düzeltmelerinin yazılım gibi dağıtılmasını olanaklı kıldı.

\tuzak{Veri yolu şemasını doğrudan donanıma çevirmek}
Ders kitaplarındaki veri yolu şemaları mantıksal bağlantıları gösterir. Fiziksel gerçekleştirmede zamanlama, yol gecikmeleri, dağıtım yükü ve güç dağıtımı gibi sorunlar ortaya çıkar. Bir tasarımın kâğıt üzerinde doğru çalışması, silikon üzerinde de çalışacağı anlamına gelmez.

\tuzak{Kritik yolu hesaplarken flip-flop kurulum ve yayılım sürelerini ihmal etmek}
Saat çevriminin asgari süresi yalnızca birleşik mantıkta geçen yol gecikmesiyle değil, aynı zamanda flip-flop kurulum süresi ($t_{su}$) ve yayılım süresi ($t_{cq}$) ile de belirlenir: $T_{\text{çevrim}} \geq t_{cq} + t_{\text{mantık}} + t_{su}$. Bu üç bileşenden birini atlamak gerçeğe uymayan iyimser bir saat sıklığı öngörüsüne yol açar. Gerçek donanımda ise yazmaç kararsız bir değer yakalayabilir ve sistem öngörülemez biçimde davranır. Aynı vergi tasarım düzeyinde de karşımıza çıkmıştı. Bölüm~\ref{subsec:cok-vurusluk-veriyolu}'deki aşama sayısı tartışmasında saat sıklığının $1/t_{\text{yazmaç}}$ tavanını aşamamasının nedeni tam olarak budur.

\tuzak{Denetim biriminin doğruluk tablosunu Karno haritası kullanmadan elle sadeleştirmeye çalışmak}
RISC-V'in 32 bitlik buyruğunda işlem kodu, fonk3 ve fonk7 alanları toplamda 13 bite ulaşır. Bu kadar girişli bir denetim mantığını sezgisel olarak sadeleştirmek hata yapmaya çok açıktır. Bir kapı eksik ya da yanlış bağlandığında belirsiz bir buyruk türünde işlemci yanlış sonuç üretir ve sebebini bulmak çok zordur. Karno haritası (Ek~\ref{ch:sayisal-tasarim}) ya da Quine-McCluskey ile sistematik sadeleştirme bu tür hataları en aza indirir. Üstelik önemsenmeyen durumları kullanarak daha küçük denetim devresi elde etmeyi sağlar.

\tuzak{Dallanma hedefini, program sayacı ilerletildikten sonra PS üzerinden hesaplamak}
RISC-V'te dallanma ve jal hedefi buyruğun \emph{kendi} adresine göre tanımlıdır. Çok vuruşluk tasarımda ise PS, bir sonraki buyruğu getirebilmek için daha Getir çevriminde PS$+$4'e ilerletilir. Hedef Çöz çevriminde ``PS $+$ anlık'' olarak hesaplanırsa dört bayt kaymış bir adres çıkar ve her dallanma bir buyruk ileriye atlar. Bu bölümde EskiPS yazmacını tam da bu tuzağa karşı ekledik. Hedef, Getir'de saklanan buyruk adresi üzerinden hesaplanır. Tuzak özellikle MIPS kaynaklarından uyarlama yapanları bekler, çünkü MIPS'te dallanma ofseti tanım gereği PC$+$4'e görecelidir ve aynı veri yolu orada doğru çalışır. RISC-V'e taşındığında sessizce bozulur.

\tuzak{Dallanma buyruğunda yeni adresi geç yazıp sonraki buyruğun yanlış adresten getirilmesine izin vermek}
Tek vuruşluk tasarımda dallanma sonucu çevrim sonunda program sayacına yazılır. Çok vuruşluk ve özellikle Bölüm~\ref{chap:boruhatti}'de göreceğimiz boru hatlı tasarımlarda ise PS güncellemesi gecikirse bir sonraki buyruk yanlış adresten getirilir. Dallanma kararını mümkün olduğu kadar erken üretmek ve sonraki bölümde göreceğimiz dallanma öngörüsü\index{dallanma öngörüsü}\index{branch prediction|see{dallanma öngörüsü}} (\emph{branch prediction}) bu tuzaktan kurtulmanın yollarıdır.


%% ================================================================
%% ÖZET
%% ================================================================

\section{Özet}
\label{sec:islemci-ozet}

Bu bölümde bir işlemciyi sıfırdan tasarlamanın adımlarını öğrendik:
\begin{enumerate}
    \item Gerçekleştirilecek buyruk kümesini belirledik.
    \item Gereken donanım birimlerini ve saat stratejisini (tek/çok vuruşluk) seçtik.
    \item Veri yolunu birleştirdik.
    \item Her buyruğun veri yolundaki yolculuğunu izleyerek denetim işaretlerini belirledik.
    \item Denetim birimini tasarladık.
    \item Denetim birimi ile veri yolunu birleştirdik.
\end{enumerate}

İlk adımda tek vuruşluk veri yolunu buyruk buyruk inşa ettik. Her buyruk bir saat vuruşunda tamamlanır, denetim sinyalleri üç bitlik buyruk kodundan ($b_2 b_1 b_0$) Karno haritalarıyla sadeleştirilen birleşik mantıkla üretilir. Ancak bu tasarımda çevrim süresi en yavaş buyruğun (yükleme) \emph{kritik yoluna} göre belirlenir ve bütün buyruklar onu bekler. Üstelik bellek ile toplayıcıların ayrı kopyaları gerekir.

Bu verimsizliği çözmek için veri yolunu aşamalara bölerek çok vuruşluk tasarıma geçtik. Kaynaklar paylaşıldı. Tek birleşik bellek hem buyruk hem veri erişimini, tek AMB üç ayrı hesabı üstlendi. Paylaşım iki yeni öğe doğurdu: buyruğu Getir çevriminde yakalayıp tutan \texttt{BuyruğaYaz} ile dallanma hedefinin dört bayt kaymasını önleyen \texttt{EskiPS} yazmacı; beq ile bne ayrımını ise tek bir Sıfır $\oplus\, b_1$ kapısı çözdü. Aşama sayısının bir tasarım kararı olduğunu, çevrimin $t_{\text{çevrim}} = t_{\text{yazmaç}} + t_{\text{aşama}}$ biçiminde ara yazmaç vergisi taşıdığını ve saat sıklığının bu vergiyle sınırlandığını gördük. Örneklerimizden iki yönlü bir sonuç çıktı. Çevrim 625 ps'den 225 ps'ye indi, ancak BBÇ 1'den 4 dolayına çıktığından, maliyetçe birbirine yakın buyruklardan oluşan kümede çok vuruşluk işlemci tek vuruşluğu geçemedi. Kümeye bölme gibi pahalı bir buyruk girdiğinde ise üstünlük iki kat farkla çok vuruşluğa geçti, çünkü tek vuruşlukta en pahalı buyruk herkesin çevrimini belirler. Kalan kazancı boru hattı toplayacak.

Denetimi, buyrukların farklı sayıda çevrim sürmesi nedeniyle belleği olan bir düzenekle, sonlu durum makinesiyle (SDM) kurduk: on iki durum, ortak Getir--Çöz omurgası ve buyruk yollarına dağılan geçişler. Mikroprogramlamanın bu makinenin bellekte saklanan biçimi olduğunu gördük. Her durum bir mikrobuyruk satırıdır, akışı Sıra/Dağıtım/Yeni değerli \emph{Sonraki} alanı yönetir. Sinyalleri kapılarla hesaplamak yerine hazır satırdan okuyan bu yaklaşım biraz yavaştır. Karşılığında tasarım yalınlaşır ve davranış, tasarımın çok geç aşamalarında, hatta üretimden sonra bile değiştirilebilir.

Bir sonraki bölümde, çok vuruşluk tasarımın aşamalarını örtüştürerek her çevrimde bir buyruk bitirmeyi hedefleyen \emph{boru hattı} tekniğini işleyecek ve bu bölümde ertelenen kazancı toplayacağız.
%% ================================================================
%% ALIŞTIRMALAR
%% ================================================================

\section*{Alıştırmalar}
\addcontentsline{toc}{section}{Alıştırmalar}

\begin{alistirma}[\zor]
RISC-V'te J türü buyruklar için anlık değer üretecini tasarlayın. Anlık değer üretecine 20 bitlik anlık değer (buyruktan) ve program sayacı değeri girecektir. Üretecin 21 bitlik etkin uzaklığı (en düşük anlamlı bit daima 0) işaretle genişletip 32 bite çıkararak program sayacıyla toplanmasını sağlayın.
\end{alistirma}

\begin{alistirma}[\kolay]
Bir işlemcide çok sayıda mimari yazmacının bulunmasının artıları ve eksileri nelerdir? Yazmaç sayısı az olursa ne olur? Açıklayın.
\end{alistirma}

\begin{alistirma}[\kolay]
Mikroprogramlama nedir? Neden kullanılır? Artıları ve eksileri nelerdir?
\end{alistirma}

\begin{alistirma}[\orta]
Tablo~\ref{tab:r-turu-mikrobuyruk} ve Tablo~\ref{tab:lw-mikrobuyruk}'ü örnek alarak \texttt{sw}, \texttt{beq}/\texttt{bne} ve \texttt{jal} buyrukları için mikrobuyruk dizilerini yazın. Her satırda tüm alanların değerlerini belirtin; \texttt{beq}/\texttt{bne} satırında PSKoşul'un, \texttt{jal} satırında PSYaz ile YazmacaYaz'ın nasıl kullanıldığına dikkat edin.
\end{alistirma}

\begin{alistirma}[\orta]
Şekil~\ref{fig:cok-vurusluk-veriyolu}'deki tasarımdan EskiPS yazmacı çıkarılsın ve dallanma hedefi, karar çevriminde ``PS $+$ anlık değer'' biçiminde hesaplansın.
\begin{enumerate}[(a)]
    \item Atlayan bir dallanma buyruğu amaçlanan hedef yerine hangi adrese sapar? Kaymanın kaç bayt olduğunu, PS'nin o anda hangi değeri tuttuğunu izleyerek gösterin.
    \item \texttt{dongu: beq x0, x0, dongu} biçimindeki sonsuz döngü bu hatalı tasarımda nasıl davranır? Buyruğun yürütülüşünü çevrim çevrim izleyin.
    \item Derleyici, ürettiği her dallanma uzaklığından 4 çıkararak bu donanım hatasını görünmez kılabilir mi? Bu çözümün neden kötü bir fikir olduğunu tartışın.
\end{enumerate}
\end{alistirma}

\begin{alistirma}[\orta]
Buyruk alt kümesine \texttt{jalr rd, anlık(rs1)} buyruğu eklenecektir. Yeni program sayacı değeri rs1 $+$ anlık değerdir ve rd yazmacına dönüş adresi yazılır.
\begin{enumerate}[(a)]
    \item Birleşik durum makinesine (Şekil~\ref{fig:durum-birlesik}) hangi durum ya da durumlar eklenir? Geçiş koşulunu ve yeni durumda etkin olan denetim sinyallerini yazın.
    \item Hedef adresin hesabında AMBKaynakA ve AMBKaynakB hangi değerleri almalıdır? Bu buyruk için EskiPS yazmacına gerek var mıdır? Gerekçelendirin.
    \item Hedefi program sayacına taşımak için PSKaynak çoklayıcısının hangi girişi kullanılabilir? (İpucu: hedef, AMB çıkışında hesaplandığı çevrimde hazırdır. AMB Yzm'e yazılmasını beklemek şart mıdır?)
    \item \texttt{jalr} kaç çevrimde tamamlanır?
\end{enumerate}
\end{alistirma}

\begin{alistirma}[\orta]
Örnek~6.5'teki aşama gecikmelerini kullanın ve aşamalar arasına giren her yazmacın çevrim süresine $t_{\text{yazmaç}} = 20$~ps eklediğini varsayın.
\begin{enumerate}[(a)]
    \item Beş aşamalı çok vuruşluk tasarımın gerçek çevrim süresini ve saat vuruşu sıklığını hesaplayın.
    \item 225~ps'lik aşamaların her biri 115'er ps'lik ikişer alt aşamaya bölünebilseydi yeni çevrim süresi ne olurdu? Bu tasarımda \texttt{lw} kaç çevrim ve toplam kaç pikosaniye sürer? Beş aşamalı tasarıma göre kazanç var mıdır?
    \item Aşamalar ne kadar inceltilirse inceltilsin saat vuruşu sıklığının aşamayacağı üst sınır kaç GHz'dir?
\end{enumerate}
\end{alistirma}

\begin{alistirma}[\zor]
Tasarlanacak bir bilgisayar için buyruk biçimi ve buyruk kümesi aşağıda verilmektedir.

\textbf{Buyruk Biçimi:}
\begin{center}
\small
\begin{tabular}{|c|c|c|c|c|}
\hline
İşlem (3 bit) & Y1 (3 bit) & Y2 (3 bit) & S (3 bit) & Anlık Değer -- A (4 bit) \\
\hline
\multicolumn{1}{|c}{\scriptsize İşlem kodu} & \multicolumn{1}{c}{\scriptsize Yazmaç1} & \multicolumn{1}{c}{\scriptsize Yazmaç2} & \multicolumn{1}{c}{\scriptsize Yazmaç$_{\text{sonuç}}$} & \multicolumn{1}{c|}{\scriptsize Anlık değer} \\
\hline
\end{tabular}
\end{center}

\textbf{Buyruk Kümesi:}

\begin{center}
\small
\begin{tabular}{@{}llp{6.2cm}@{}}
\toprule
\textbf{Buyruk} & \textbf{İşlem} & \textbf{Açıklama} \\
\midrule
TOPLA S, Y1, Y2 & S $\leftarrow$ Y1 + Y2 & Y1 ile Y2'yi topla, S'ye yaz \\
ÇIKAR S, Y1, Y2 & S $\leftarrow$ Y1 $-$ Y2 & Y1'den Y2'yi çıkar, S'ye yaz \\
VE S, Y1, Y2 & S $\leftarrow$ Y1 \& Y2 & VE işlemi, S'ye yaz \\
VEYA S, Y1, Y2 & S $\leftarrow$ Y1 $|$ Y2 & VEYA işlemi, S'ye yaz \\
YÜKLE S, Y1, A & S $\leftarrow$ Bellek[Y1+SfGen(A)] & Bellekten veri oku, S'ye yaz \\
SAKLA Y2, Y1, A & Bellek[Y1+SfGen(A)] $\leftarrow$ Y2 & Y2 değerini belleğe yaz \\
ATLA Y1, Y2, A & Koşullu dallanma & Y1$=$Y2 ise PS+4, değilse PS+İşGen(A) \\
ZIPLA Y2, A & PS $\leftarrow$ PS+Y2+İşGen(A) & Program sayacını artır \\
\bottomrule
\end{tabular}
\end{center}

Bu bilgisayarı aşağıda giriş ve çıkışları gösterilen parçaları kullanarak tasarlayın:
\begin{itemize}
    \item Buyruk Belleği: girişler [adres], çıkışlar [buyruk]
    \item Veri Belleği: girişler [adres, yazma/okuma denetimi, veri girişi], çıkışlar [veri çıkışı]
    \item Yazmaç Öbeği: girişler [Y1 adresi, Y2 adresi, S adresi, S Verisi, Yazma denetimi, saat], çıkışlar [Y1 verisi, Y2 verisi]
    \item Program sayacı yazmacı
    \item İşaret Genişletici: Hem sıfırla hem de birle genişletebilir
    \item AMB: Toplama, çıkarma, VE, VEYA işlemlerini yapabilir
    \item Toplayıcı, değişik boyutta çoklayıcılar
\end{itemize}
Tasarım mümkün olduğunca verimli olmalı ve az sayıda donanım kullanmalıdır. Denetim işaretleri açıkça tanımlanmalı, tabloları ve devreleri gösterilmelidir.
\end{alistirma}

\begin{alistirma}[\zor]
Tasarlanacak bir bilgisayar için buyruk biçimi aşağıdaki gibidir:

\begin{center}
\begin{tabular}{|c|c|c|c|}
\hline
İşlem (4 bit) & Yazmaç1 -- Y1 (3 bit) & Yazmaç2 -- Y2 (3 bit) & Anlık Değer -- A (6 bit) \\
\hline
\end{tabular}
\end{center}

16 buyruktan oluşan bir buyruk kümesi tanımlanmıştır (TOPLA, ÇIKAR, VE, ÇARP, VEYA, ÖZELVEYA, DEĞİL, ARTIR, AKTAR, YÜKLE (anlık değerli), YÜKLE (yazmaçla), SAKLA (yazmaçla), SAKLA (anlık değerli), AED, AEŞ, ZIPLA). Bu bilgisayarı verilen parçaları kullanarak tasarlayın. Denetim işaretleri açıkça tanımlanmalı ve tabloları gösterilmelidir.
\end{alistirma}

\begin{alistirma}[\orta]
Aşağıda bir döngünün derleyici tarafından çevirici diline dönüştürülmüş hâli verilmiştir:

\begin{lstlisting}[style=riscv, caption={Kabarcık sıralaması -- çevirici dili (RISC-V)}]
D2:
B1:  addi x3, x7, 0     # x3->a[i]  (mv x3, x7)
B2:  lw   x8, 0(x3)     # a[i] degerini yukle
B3:  addi x3, x3, 4     # x3->a[i+1]
B4:  lw   x9, 0(x3)     # a[i+1] degerini yukle
B5:  bge  x9, x8, D3    # a[i+1] >= a[i] ise dallan
B6:  addi x3, x7, 0     # x3->a[i]
B7:  addi x3, x3, 4     # x3->a[i+1]
B8:  sw   x8, 0(x3)     # a[i+1] konumuna sakla
B9:  sw   x9, -4(x3)    # a[i] konumuna sakla
B10: addi x5, x5, 1     # degisim++
D3:
B11: addi x6, x6, 1     # i++
B12: addi x7, x7, 4     # x7->a[i]
B13: blt  x6, x4, D2    # i < son ise dallan
\end{lstlisting}

Döngünün bir kere döndüğü (son = 1) varsayılırsa:

\begin{enumerate}[(a)]
    \item Saat vuruşu sıklığı 100 MHz olan tek vuruşluk bir RISC-V işlemci kullanılırsa:
    \begin{enumerate}[I.]
        \item if satırındaki koşul doğru ise bu kod parçasının işlemesi için kaç çevrim gerekir?
        \item Koşul yanlış olsaydı kaç çevrim gerekirdi?
        \item BBÇ nedir?
        \item Toplam yürütme zamanı ne kadardır?
    \end{enumerate}
    \item Saat vuruşu sıklığı 400 MHz olan çok vuruşluk bir RISC-V işlemci kullanılırsa (yazmaç/saklama: 4 vuruş, dallanma: 3 vuruş, yükleme: 5 vuruş):
    \begin{enumerate}[I.]
        \item Koşul doğru ise kaç çevrim gerekir?
        \item Koşul yanlış olsaydı kaç çevrim gerekirdi?
        \item BBÇ nedir?
        \item Toplam yürütme zamanı ne kadardır?
        \item Tek vuruşluk işlemciye göre hızlanma ne kadardır?
    \end{enumerate}
\end{enumerate}
\end{alistirma}

\begin{alistirma}[\orta]
Kasırga-1 tek vuruşluk bir işlemcidir. Aşağıdaki buyruk kümesini gerçekleştirecek tek vuruşluk veri yolunu tasarlayın ve denetim tablosunu oluşturun. (Buyruk kümesi: add, addi, sub, subi, mul, muli, not, sll, srl, mov, movi, beq, ba, bgt, blt)
\end{alistirma}

%% ===================================================================
%% Aşağıdaki sorular BİL 361 dersinin geçmiş yıllardan sınav
%% sorularından uyarlanmıştır.
%% ===================================================================

\begin{alistirma}[\zor]
Aşağıda bir buyruk kümesi verilmiştir. Bu buyruk kümesini destekleyen tek vuruşluk bir işlemci tasarlayın.

\begin{center}
\begin{tabular}{|l|l|}
\hline
\textbf{Buyruk} & \textbf{Açıklama} \\ \hline
\texttt{çarp rx, ry, \#anlık} & \texttt{R[rx] = R[ry] * İşaretleGenişlet(anlık)} \\ \hline
\texttt{yükle rx, ry, rz, rt} & \texttt{R[rx] = Bellek[R[ry]] + (R[rz] \% R[rt])} \\ \hline
\texttt{sakla rx, ry} & eğer R[rx] > R[ry] ise Bellek[R[rx]] = R[rx], değilse Bellek[R[rx]] = R[ry] \\ \hline
\texttt{emre rx, ry} & \texttt{R[rx] = R[ry] * R[ry]} \\ \hline
\texttt{dallan rx, ry, \#anlık} & eğer R[rx] > R[ry] ise PS += ?, değilse PS = İşaretleGenişlet(anlık) \\ \hline
\end{tabular}
\end{center}

\begin{enumerate}[(a)]
    \item Buyruklar 16 bit genişliğindedir ve işlemcideki veriler 32 bittir. Yazmaç öbeğinde 8 yazmaç vardır. Buyruk tiplerini ve bit vektörlerini tasarlayın.
    \item İşlemci için gerekli veri yolunu çizin. Gerekli çoklayıcıları ve denetim sinyallerini belirtin.
    \item Tasarladığınız işlemci için denetim tablosunu oluşturun.
\end{enumerate}

\emph{Not:} \texttt{dallan} buyruğundaki \texttt{PS += ?} artış miktarını buyruk genişliğine göre siz belirlemelisiniz.
\end{alistirma}

\begin{alistirma}[\orta]
Tek vuruşluk bir işlemcide aşağıdaki bileşenlerin gecikmeleri verilmiştir:

\begin{center}
\begin{tabular}{|l|c|}
\hline
\textbf{Bileşen} & \textbf{Gecikme} \\ \hline
Buyruk belleği & 200~ps \\ \hline
Yazmaç okuma & 100~ps \\ \hline
AMB & 150~ps \\ \hline
Veri belleği & 250~ps \\ \hline
Yazmaç yazma & 100~ps \\ \hline
Çoklayıcılar & 25~ps \\ \hline
\end{tabular}
\end{center}

\begin{enumerate}[(a)]
    \item \texttt{add} buyruğunun kritik yol gecikmesi ne kadardır?
    \item \texttt{lw} (yükleme) buyruğunun kritik yol gecikmesi ne kadardır?
    \item \texttt{beq} (dallanma) buyruğunun kritik yol gecikmesi ne kadardır?
    \item Bu işlemcinin en yüksek saat sıklığı ne olabilir?
    \item Çok vuruşluk bir işlemcide aynı buyruklar kaçar çevrimde tamamlanır?
\end{enumerate}
\end{alistirma}

%% ===================================================================
%% Aşağıdaki sorular BİL 361 dersinin geçmiş yıllardan sınav
%% sorularından uyarlanmıştır.
%% ===================================================================

\begin{alistirma}[\zor]
Aşağıdaki buyruk kümesini uygulayan bir işlemci tasarlanacaktır. İşlemcide 16 adet 16~bit genişliğinde genel amaçlı yazmaç, bir adet AMB (Aritmetik Mantık Birimi) ve ayrı buyruk/veri bellekleri (Harvard mimarisi) bulunmaktadır. Bellekte 4~bitlik adresleme yapılmaktadır (her adres 4~bitlik bir değer gösterir).

Buyruk kümesi:
\begin{center}
\small
\begin{tabular}{ll}
\toprule
\textbf{Buyruk} & \textbf{Açıklama} \\
\midrule
\texttt{add rd, rs1, rs2} & rd $\leftarrow$ rs1 + rs2 \\
\texttt{sub rd, rs1, rs2} & rd $\leftarrow$ rs1 $-$ rs2 \\
\texttt{lw rd, offset(rs1)} & rd $\leftarrow$ Bellek[rs1 + offset] \\
\texttt{sw rs2, offset(rs1)} & Bellek[rs1 + offset] $\leftarrow$ rs2 \\
\texttt{beq rs1, rs2, etiket} & rs1 = rs2 ise etiket'e atla \\
\texttt{sll rd, rs1, mik} & rd $\leftarrow$ rs1 $\ll$ mik \\
\bottomrule
\end{tabular}
\end{center}

\begin{enumerate}[(a)]
    \item Her buyruğun bit vektörünün yapısını tasarlayın. Kaç tür buyruk biçimi gereklidir? Buyruk genişliğini olabilecek en küçük değer olarak belirleyin.
    \item Bu buyruk kümesine göre çalışan tek vuruşluk işlemciyi tasarlayın (veri yolu şemasını çizin).
    \item Tasarladığınız işlemcinin denetim tablosunu oluşturun.
\end{enumerate}
\end{alistirma}

\begin{alistirma}[\orta]
Aşağıdaki üst düzey dilde yazılmış programı, yukarıdaki alıştırmada verilen buyruk kümesini kullanarak çevirici dile (assembly) çevirin. \texttt{A} dizisinin başlangıç adresi \texttt{x3}, \texttt{B} dizisinin başlangıç adresi \texttt{x4} yazmacında hazır bulunmaktadır. Dizilerde 4~baytlık tam sayılar saklanmaktadır.
\begin{lstlisting}[language=C, style=alistirma-kod]
    int toplam = 0;
    for (int i = 0; i < N; i++)
        toplam += A[i] * B[i];
\end{lstlisting}
\end{alistirma}

\begin{alistirma}[\orta]
\textit{Not: Bu alıştırma, bölümün ana tasarımından farklı olarak ayrı okuma izin sinyalleri (yazmaç öbeği için YazmacaOku, bellek için BelleğiOku) İÇEREN alternatif bir tasarımı inceler. Ana tasarımda okuma yan etkisiz sayıldığından böyle sinyaller kullanılmamıştı. Bu tür okuma izinleri ileride ayrı bir kesimde ele alınacaktır.}

Bir işlemcinin denetim tablosu aşağıda verilmiştir:

\begin{center}
\small
\begin{tabular}{lcccccc}
\toprule
\textbf{Buyruk} & \textbf{YazmacaOku} & \textbf{AMBİşlem} & \textbf{BelleğiOku} & \textbf{BelleğeYaz} & \textbf{YazmacaYaz} & \textbf{Dallan} \\
\midrule
? & 1 & Topla & 0 & 0 & 1 & 0 \\
? & 1 & Topla & 1 & 0 & 1 & 0 \\
? & 1 & Topla & 0 & 1 & 0 & 0 \\
? & 1 & Çıkar & 0 & 0 & 0 & 1 \\
\bottomrule
\end{tabular}
\end{center}

\begin{enumerate}[(a)]
    \item Bu denetim tablosuna göre işlemcinin kaç buyruğu vardır? Her buyruğun ne iş yaptığını, işlenenlerini ve işlevlerini belirleyin.
    \item Belirlediğiniz buyrukları kullanarak aşağıdaki C kodunu çevirici dilde yazın:
\begin{lstlisting}[language=C, style=alistirma-kod]
    if (x == y)
        z = x + y;
    else
        z = A[x];
\end{lstlisting}
\end{enumerate}
\end{alistirma}

\begin{alistirma}[\orta]
Çok vuruşluk bir işlemcide buyrukların aşama gecikmeleri aşağıda verilmiştir:

\begin{center}
\small
\begin{tabular}{lc}
\toprule
\textbf{Buyruk türü} & \textbf{Çevrim sayısı} \\
\midrule
Aritmetik/mantık & 4 \\
Bellek (yükleme) & 6 \\
Bellek (saklama) & 5 \\
Dallanma & 3 \\
Çarpma & 8 \\
\bottomrule
\end{tabular}
\end{center}

Bir döngü gövdesinde 2 bellek yükleme, 1 çarpma, 3 aritmetik ve 1 dallanma buyruğu bulunmaktadır. Döngü $N$ kez yineleniyorsa:
\begin{enumerate}[(a)]
    \item Program toplam kaç çevrimde tamamlanır ($N$ cinsinden)?
    \item Programın ortalama BBÇ'si nedir?
    \item Tek vuruşluk işlemci ile aynı saat sıklığında çalıştırılsaydı program kaç çevrimde biterdi? Çok vuruşluk işlemci daha hızlı mıdır?
    \item Çok vuruşluk işlemcinin tek vuruşluk işlemciden hızlı olabilmesi için çevrim zamanının en az ne kadar olması gerekir?
\end{enumerate}
\end{alistirma}

\begin{alistirma}[\orta]
Aşağıdaki tabloda TeknolojiMerkezi-217 buyruk kümesindeki buyrukların bir kısmı listelenmiştir.

\begin{center}
\begin{tabular}{ll}
\toprule
\textbf{Buyruk} & \textbf{Açıklama} \\
\midrule
topla hy, ky1, ky2 & hy = ky1 + ky2 \\
çıkar hy, ky1, ky2 & hy = ky1 $-$ ky2 \\
böl hy, ky1, ky3 & hy = ky1 / ky3 \\
belböl ky1, ky2, ky3 & Bellek[ky1] = ky2 / ky3 \\
dallan ky1, ky2 & ky1 $=$ ky2 ise PS = ky1 + ky2, \\
 & değilse PS = PS + 4 \\
dallanma & PS = PS + 4 \\
kaykay ky1, anlık & Bellek[ky1] = anlık \\
\bottomrule
\end{tabular}
\end{center}

Bu buyruk kümesi için:
\begin{enumerate}[(a)]
    \item Buyruk kümesini gerçekleyen tek vuruşluk veri yolunu, gerekli çoklayıcıları ve denetim sinyallerini belirterek çizin.
    \item Her buyruk için çoklayıcı seçim bitlerini, YAZ (yazmaç yazma) ve AMB işlem kodu sinyallerini içeren denetim tablosunu oluşturun.
    \item \texttt{belböl} ve \texttt{kaykay} gibi belleğe yazan buyruklar için hangi denetim sinyalinin zorunlu olduğunu bir cümleyle açıklayın.
\end{enumerate}
\end{alistirma}

\begin{alistirma}[\zor]
NEDEN buyruk kümesi mimarisi aşağıdaki tabloda verilen buyrukları içermektedir:

\begin{center}
\begin{tabular}{ll}
\toprule
\textbf{Buyruk} & \textbf{Açıklama} \\
\midrule
\texttt{topla hy ky1 ky2} & \texttt{hy = ky1 + ky2} \\
\texttt{atopla hy ky1 anlık1} & \texttt{hy = ky1 + anlık} \\
\texttt{hopla ky1 ky2 anlık1} & \texttt{ky1 < ky2 ? PS = PS + anlık : PS = PS + X} \\
\texttt{değiş ky1 ky2} & ky1 ve ky2'deki değerleri değiştirir \\
\texttt{zıpla ky2 anlık2} & \texttt{PS = ky2 + anlık} \\
\texttt{kaydet ky1 ky2 anlık1} & \texttt{Bellek[ky1 + anlık] = ky2} \\
\texttt{yükle hy ky1 anlık1} & \texttt{hy = Bellek[ky1 + anlık]} \\
\texttt{ismail ky1 ky2 anlık1} & \texttt{Bellek[ky1 + anlık] = Bellek[ky1 + anlık] + ky2} \\
\bottomrule
\end{tabular}
\end{center}

NEDEN mimarisinde:
\begin{itemize}
    \item Buyruklar 32 bit genişliğindedir.
    \item Yazmaçlar 32 bit genişliğindedir.
    \item 64 tane genel amaçlı yazmaç vardır.
    \item Belleğe erişim iki-bayt adresleme ile gerçekleştirilmektedir (her bellek adresi bellekte ardışık 16 bite işaret etmektedir).
    \item Tüm anlık değerler 32 bite işaretle genişletilmektedir.
    \item Bellek adresleri 32 bit genişliğindedir.
\end{itemize}

\begin{enumerate}[(a)]
    \item \texttt{anlık1} ve \texttt{anlık2} değerleri en fazla kaç bit olabilir?
    \item NEDEN mimarisindeki buyrukları en az sayıda buyruk tipi kullanarak gruplayın.
    \item NEDEN mimarisinin adresleyebildiği bellek uzayı kaç bayt boyutundadır? Açıklayın.
    \item Buyruk belleğine yerleştirilen buyrukların dört bayta hizalandığı varsayıldığında aşağıdaki adreslerden hangileriyle buyruk belleğine erişildiğinde kural dışı durum (\emph{exception}) oluşmasını beklersiniz? Neden?
    \begin{enumerate}[(i)]
        \item \texttt{0x00000000}
        \item \texttt{0x00000001}
        \item \texttt{0x00001010}
    \end{enumerate}
    \item Yukarıdaki buyruk kümesi mimarisini gerçekleyen tek vuruşluk bir işlemci (NEDEN-1) tasarlayın. Denetim tablosunu çizin.
    \item NEDEN-1 tasarımında hangi buyruğun kritik yol üzerinde olmasını (saat sıklığını belirlemesini) beklersiniz? Açıklayın.
\end{enumerate}
\end{alistirma}

\begin{alistirma}[\zor]
Bir şifreleme işlemcisi tasarlamak istiyorsunuz. Aşağıda bu işlemcinin buyruk kümesindeki buyruklar ve işlevleri listelenmiştir:

\begin{center}
\begin{tabular}{ll}
\toprule
\textbf{Buyruk} & \textbf{Açıklama} \\
\midrule
\texttt{taşı hy, \#anlık} & hy $\leftarrow$ anlık \\
\texttt{artıartı hy} & hy $\leftarrow$ hy + 1 \\
\texttt{cozyukle hy, ky1, \#anlık} & hy $\leftarrow$ xnor(Bellek[ky1 + anlık], anlık) \\
\texttt{sifrelesakla ky1, ky2, \#anlık} & Bellek[ky2 + anlık] $\leftarrow$ xor(ky1, anlık) \\
\texttt{cozatla ky1, ky2, \#anlık} & xor(ky1, anlık) $=$ xnor(ky2, anlık) ise \\
 & \quad PS $=$ PS + anlık, değilse PS $=$ PS + 4 \\
\bottomrule
\end{tabular}
\end{center}

İşlemcinin özellikleri şunlardır:
\begin{itemize}
    \item İşlemcide 12 adet 32 bitlik genel amaçlı yazmaç bulunmaktadır.
    \item Buyruk belleği ve veri belleği için bayt adresleme kullanılmaktadır.
    \item Her buyruk için buyruk genişliği sabittir ve tek bir buyruk tipi vardır. Anlık değerler her zaman 12 bit genişliğindedir.
    \item İşlemci Harvard mimarisine göre tasarlanmalıdır.
    \item Buyruk belleği, veri belleği ve yazmaç öbeği dışında istediğiniz kadar bileşen kullanabilirsiniz (AMB, genişletici, toplayıcı vb.).
\end{itemize}

\begin{enumerate}[(a)]
    \item Yukarıdaki özelliklere göre işlemcinin buyruk genişliği en az kaç bit olmalıdır? Buyruk kümesindeki buyrukların bit vektörlerinin yapısını gösterin.
    \item Buyruk kümesi mimarisi verilen işlemciyi belirtilen kısıtlara göre en az sayıda bileşen kullanarak tasarlayın. Veri yolu genişliklerini tasarımınızda belirtin.
    \item Tasarladığınız işlemcinin denetim tablosunu oluşturun.
\end{enumerate}
\end{alistirma}

%% ===================================================================
%% Aşağıdaki sorular BİL 361 dersinin geçmiş yıllardan sınav
%% sorularından uyarlanmıştır.
%% ===================================================================

\begin{alistirma}[\zor]
Özel tasarım 16~bitlik bir işlemci için aşağıdaki buyruk kümesi mimarisi tanımlanmıştır:

\begin{center}
\small
\begin{tabular}{ll}
\toprule
\textbf{Buyruk} & \textbf{Açıklama} \\
\midrule
\texttt{topla hy, ky1, ky2} & \texttt{x[hy] = x[ky1] + x[ky2]} \\
\texttt{ayükle hy, \#anlık} & \texttt{x[hy] = İşaretleGenişlet(anlık)} \\
\texttt{kaydır hy, ky, \#anlık} & \texttt{x[hy] = x[ky] <\null< anlık} \\
\texttt{çarp hy, ky1, ky2} & \texttt{x[hy] = x[ky1] * x[ky2]} \\
\texttt{eriş hy, ky1} & \texttt{x[hy] = Bellek[x[ky1]]} \\
\texttt{sakla ky1, ky2} & \texttt{Bellek[x[ky1]] = x[ky2]} \\
\texttt{uçur \#anlık} & \texttt{PS = PS + İşaretleGenişlet(anlık)} \\
\texttt{budaklan ky1, ky2, \#anlık} & eğer \texttt{x[ky1] > x[ky2]} ise \texttt{PS = PS + anlık}, \\
 & değilse \texttt{PS = PS + ?} \\
\bottomrule
\end{tabular}
\end{center}

İşlemcinin özellikleri şunlardır:
\begin{itemize}
    \item Veriler 16 bit genişliğindedir.
    \item İşlemcide 16 adet genel amaçlı yazmaç (\texttt{x0}--\texttt{x15}) ve 1 adet program sayacı (PS) bulunmaktadır.
    \item \texttt{x0} yazmacı her zaman 0 değerini içerir (sıfır yazmacı).
    \item Bellekte bayt adresleme kullanılmaktadır. Bellek okuma ve yazma işlemleri 16~bit genişliğindedir.
    \item İşlemci Harvard mimarisine göre tasarlanmalıdır (ayrı buyruk belleği ve veri belleği).
    \item Buyruk belleği, veri belleği ve yazmaç öbeği dışında 1 adet AMB kullanılabilir.
    \item \texttt{ayükle} buyruğundaki anlık değer 9~bit, \texttt{kaydır} ve \texttt{budaklan} buyruklarındaki anlık değer 5~bit genişliğindedir.
\end{itemize}

\begin{enumerate}[(a)]
    \item \texttt{topla} buyruğunun bit vektörünü gösterin. İşlem kodu (\emph{opcode}), \texttt{hy}, \texttt{ky1}, \texttt{ky2} alanlarını ve her alanın bit genişliğini belirtin.
    \item Yukarıdaki buyruk kümesini gerçekleyen tek vuruşluk bir işlemci tasarlayın. Veri yolu şemasını çizin. Veri yolu genişliklerini ve denetim sinyallerini şemada belirtin.
    \item Tasarladığınız işlemcinin denetim tablosunu oluşturun.
\end{enumerate}

\emph{Not:} \texttt{budaklan} buyruğundaki \texttt{?} değerini belirlemeniz gerekmektedir. PS artış miktarı buyruk genişliğine bağlıdır.
\end{alistirma}

\begin{alistirma}[\orta]
Aşağıda bir işlemcinin buyruk kümesi mimarisi verilmiştir:

\begin{center}
\small
\begin{tabular}{ll}
\toprule
\textbf{Buyruk} & \textbf{Açıklama} \\
\midrule
\texttt{topla hy, ky1, ky2} & \texttt{hy = ky1 + ky2} \\
\texttt{katla hy, ky1} & \texttt{hy = ky1 * ky1} \\
\texttt{çıkla ky1, ky2} & \texttt{Bellek[ky2] = ky2 $-$ ky1} \\
\texttt{kutla ky1, \#anlık} & eğer \texttt{ky1 < anlık} ise \texttt{PS = ky1 * anlık}, \\
 & değilse \texttt{PS = PS + ?} \\
\texttt{boş.geç} & \texttt{PS = PS + ?} \\
\texttt{tıkla hy, ky1, \#anlık} & \texttt{hy = Bellek[ky1 + anlık]} \\
\bottomrule
\end{tabular}
\end{center}

Buyruk genişliği 64~bittir. Bu işlemci için kısmi denetim tablosu ve kısmi veri yolu şeması verilmiştir.

\begin{enumerate}[(a)]
    \item Denetim tablosundaki tüm eksik alanları doldurun. Her buyruk için buyruk açıklamalarını, denetim sinyallerini ve veri yolu etiketlerini belirleyin.
    \item 64~bitlik buyruk genişliği göz önüne alındığında bu buyruk kümesinde kullanılan adresleme kipini belirleyin. Buyruk tiplerini ve bit vektörlerinin yapısını gösterin.
\end{enumerate}

\emph{Not:} \texttt{kutla} ve \texttt{boş.geç} buyruklarındaki \texttt{?} değerlerini belirlemeniz gerekmektedir. PS artış miktarı buyruk genişliğine bağlıdır.
\end{alistirma}

%% Aşağıdaki sorular BİL 361 dersinin 2024--2025 öğretim yılı
%% sınav sorularından uyarlanmıştır.

\begin{alistirma}[\zor]
(\emph{HALİS BKM -- Güz 2024 Ara Sınav})
HALİS buyruk kümesi mimarisi aşağıda verilmiştir. Bu kümeyi gerçekleyen tek vuruşluk bir işlemci tasarlanacaktır.
\begin{itemize}
    \item Toplamda bir adet yazmaç öbeği, bir adet aritmetik mantık birimi, bir adet veri belleği, bir adet buyruk belleği, bir adet genişletici, bir adet mod alma birimi ve bir adet gizemli birim bulunmaktadır.
    \item Buyruklar 32~bit genişliğindedir ve 32 adet yazmaç bulunmaktadır.
    \item Bayt adresleme kullanılmaktadır.
\end{itemize}

\begin{center}
\small
\begin{tabular}{ll}
\toprule
\textbf{Buyruk} & \textbf{Açıklama} \\
\midrule
\texttt{genislet\_yaz hy2, \#anlık} & \texttt{x[hy2] = sıfırla\_genislet(anlık)} \\
\texttt{carp hy1, ky1, ky2} & \texttt{x[hy1] = x[ky1] * x[ky2]} \\
\texttt{mod\_al hy1, ky1, \#anlık} & \texttt{x[hy1] = x[ky1] \% isaretle\_genislet(anlık)} \\
\texttt{carp\_topla hy2, ky1, ky2, \#anlık} & \texttt{x[hy2] = x[ky1] * x[ky2] + sıfırla\_genislet(anlık)} \\
\texttt{neva hy2, ky1, ky2, \#anlık} & \texttt{x[hy2] = (x[ky1] \% x[ky2]) + (x[ky1] * isaretle\_genislet(anlık))} \\
\texttt{dallan ky1, ky2, ky3, \#anlık} & eğer \texttt{(x[ky1] \% x[ky2]) < x[ky3]} ise \texttt{PS = isaretle\_genislet(anlık)}, \\
 & değilse \texttt{PS += ?} \\
\texttt{yükle hy2, ky1, ky3} & \texttt{x[hy2] = Bellek[x[ky1]] / x[ky3]} \\
\texttt{sakla ky1, ky2} & \texttt{Bellek[x[ky1]] = (x[ky1] * x[ky2]) $-$ x[ky2]} \\
\texttt{yusuf hy1, hy2, ky1, ky2} & eğer \texttt{x[ky1] == x[ky2]} ise \texttt{x[hy2] = x[ky1] * x[ky2]}, \\
 & değilse \texttt{x[hy1] = x[ky1] \% x[ky2]} \\
\bottomrule
\end{tabular}
\end{center}

Bu buyruk kümesini gerçekleyen tek vuruşluk veri yolunu tasarlayın. Her buyruk için gereken denetim sinyallerini bir denetim tablosunda gösterin. Ayrıca \texttt{dallan} buyruğundaki \texttt{PS += ?} değerini belirleyin.

\emph{Not:} İşlem birimlerinde \texttt{değer1} ve \texttt{değer2} sırası önemli değildir.
\end{alistirma}

\begin{alistirma}[\zor]
(\emph{CISC-35 BKM -- Bahar 2025 Ara Sınav})
CISC-35 buyruk kümesi mimarisi aşağıda verilmiştir. Bu kümeyi gerçekleyen tek vuruşluk bir işlemci tasarlanacaktır. Buyruklar 16~bit genişliğindedir. İşlemcideki veriler ise 32~bittir. Toplamda 1~yazmaç öbeği ve yazmaç öbeğinde 8~adet yazmaç bulunmaktadır. Birer adet AMB, genişletici, mod alma birimi, veri belleği ve buyruk belleği bulunmaktadır. İki bayt adresleme kullanılmaktadır.

\begin{center}
\small
\begin{tabular}{ll}
\toprule
\textbf{Buyruk} & \textbf{Açıklama} \\
\midrule
\texttt{çarp rx, ry, \#anlık} & \texttt{R[rx] = R[ry] * işaretle\_genişlet(anlık)} \\
\texttt{umut rx, ry, rz, \#anlık} & \texttt{R[rx] = (R[ry] + R[rz]) \% işaretle\_genişlet(anlık)} \\
\texttt{orhun rx, ry, rz, rt} & eğer \texttt{(R[rz] / R[rt])} tekse \texttt{R[rx] = R[rz] / R[rt]}, \\
 & çiftse \texttt{R[ry] = R[rz] / R[rt]} \\
\texttt{yükle rx, ry, rz, rt} & \texttt{R[rx] = Bellek[R[ry]] + (R[rz] \% R[rt])} \\
\texttt{sakla rx, ry} & eğer \texttt{R[rx] > R[ry]} ise \texttt{Bellek[R[rx]] = R[rx]}, \\
 & değilse \texttt{Bellek[R[rx]] = R[ry]} \\
\texttt{kerem rx, ry, \#anlık} & \texttt{Bellek[R[rx]] = R[ry] $-$ sıfırla\_genişlet(anlık)} \\
\texttt{emre rx, ry} & \texttt{R[rx] = R[ry] * R[ry]} \\
\texttt{dallan rx, ry, \#anlık} & eğer \texttt{R[rx] > R[ry]} ise \texttt{PS += ?}, \\
 & değilse \texttt{PS = işaretle\_genişlet(anlık)} \\
\bottomrule
\end{tabular}
\end{center}

\begin{enumerate}[(a)]
    \item Verilen buyruk kümesindeki buyrukların bit vektörlerinin yapısını ayrıntılı olarak tasarlayarak kaç tip buyruk olduğunu açıkça gösterin. Hangi buyruğun hangi tipi kullandığını açıklayın. (\emph{Not:} En az sayıda buyruk tipi kullanmak zorundasınız.)
    \item Anlıklar en fazla kaç bit olabilir? (\emph{Not:} Farklı buyruklardaki anlıkların genişlikleri farklı olabilir.)
    \item Yukarıdaki buyruk kümesini destekleyen tek vuruşluk işlemciyi tasarlayın. Veri yolunu çizin, denetim sinyallerini ve alanların bit genişliklerini belirtin. Program sayacını kaç artırmanız gerektiğini açıklayın.
\end{enumerate}
\end{alistirma}

\begin{alistirma}[\orta]
Bölüm~\ref{subsec:ana-denetim}'te ana denetim sinyalleri için türetilen Boole ifadelerini kullanarak, aşağıdaki sinyallerin her birini üreten kapı devresini çizin: \texttt{YazmacaYaz}, \texttt{AMBKaynak}, \texttt{GeriYazmaKaynağı}[1], \texttt{GeriYazmaKaynağı}[0], \texttt{AnlıkBiçim}[1], \texttt{AnlıkBiçim}[0] ve \texttt{BelleğeYaz}. Devrelerin girişleri üç bitlik buyruk kodunun bitleri ($b_2$, $b_1$, $b_0$) ile gereken yerlerde bunların tümleyenleridir. \texttt{AMBİşlem} için Şekil~\ref{fig:amb-denetimi-devre}'te verilen kapı devresini örnek alın.
\end{alistirma}

\begin{alistirma}[\zor]
\texttt{Dallan} sinyali, ana çözücünün ürettiği \emph{Atla} ve \emph{Dal} ara işaretleri ile AMB'den gelen \texttt{Sıfır} bayrağından üretilir:
\[
\text{Atla} = b_2 b_1 b_0, \quad \text{Dal} = b_2\,(b_1 \oplus b_0), \quad \text{Dallan} = \text{Atla} + \text{Dal}\cdot(\text{Sıfır} \oplus b_1).
\]
Bu üç ifadeyi gerçekleştiren kapı devresini çizin. Girişler üç bitlik buyruk kodunun bitleri ($b_2 b_1 b_0$) ile \texttt{Sıfır} bayrağı, çıkış ise \texttt{Dallan} sinyalidir. Tasarımda ÖZELVEYA (XOR) kapısını kullanabilirsiniz.
\end{alistirma}
